\Opensolutionfile{ans}[ans/CD19/Muc_9_10]
\setcounter{ex}{0}
\setcounter{dang}{0}
\section{Mức độ 9,10 điểm}
\begin{dang}
	{Phương trình logarit chứa tham số}
\end{dang}
\begin{ex}
	[Đề Minh Họa 2020 Lần 1]%Câu 1.
	Cho phương trình $\log_2^2(2x)-(m+2)\log_2x+m-2=0$ ($m$ là tham số thực). Tập hợp tất cả các giá trị của $m$ để phương trình đã cho có hai nghiệm phân biệt thuộc đoạn $[1;2]$ là
	\choice
	{$(1;2)$}
	{$[1;2]$}
	{\True $[1;2)$}
	{$[2;+\infty)$}
	\loigiai{
		$\log_2^2(2x)-(m+2)\log_2x+m-2=0\Leftrightarrow[1+\log(x)]^2-(m+2)\log_2x+m-2=0 (*)$.\\
		Đặt $t=\log_2x=g(x)\Rightarrow 0\leq t\leq 1$ và mỗi giá trị của $x$ sẽ cho một giá trị của $t$.\\
		$(*)$ trở thành $(1+t)^2-(m+2)t+m-2=0$ \\
		$ \Leftrightarrow t^2+2t+1-mt-2t+m-2=0 $ \\
		$ \Leftrightarrow t^2-1=m(t-1) $ \\
		$ \Leftrightarrow(t-1)(t+1-m)=0 $ \\
		$ \Leftrightarrow\hoac{&t=m-1(1)\\&t=1(2).} $ \\
		Với $t=1$ thì phương trình có một nghiệm $x=2$.\\
		Vậy để phương trình ban đầu có hai nghiệm phân biệt thì phương trình $(1)$ phải có một nghiệm $t\neq 1$.\\
		$0\leq m-1<1\Leftrightarrow 1\leq m<2$.\\
		Vậy $m\in[1;2)$ để thoả mãn yêu cầu bài toán.
	}
\end{ex}
\begin{ex}
	[Chuyên Lam Sơn Thanh Hóa 2019]%Câu 2.
	Cho phương trình sau $3\log_{27}\left[2x^2-(m+3)x+1-m\right]+\log_{\tfrac{1}{3}}\left(x^2-x+1-3m\right)=0$. Số các giá trị nguyên của $m$ để phương trình đã cho có hai nghiệm phân biệt $x_1,x_2$ thỏa mãn $|x_1-x_2|<15$ là 
	\choice
	{$14$}
	{$11$}
	{$12$}
	{\True $13$}
	\loigiai{
		Ta có $3\log_{27}\left[2x^2-(m+3)x+1-m\right]+\log_{\tfrac{1}{3}}\left(x^2-x+1-3m\right)=0$ \\
		$\Leftrightarrow\log_3\left[2x^2-(m+3)x+1-m\right]=\log_3\left(x^2-x+1-3m\right) $ \\
		$\Leftrightarrow\heva{&x^2-x+1-3m>0\\&2x^2-(m+3)x+1-m=x^2-x+1-3m} $ \\
		$ \Leftrightarrow\heva{&x^2-x+1-3m>0(*)\\&x^2-(m+2)x+2m=0&(1)}\Leftrightarrow\heva{&x^2-x+1-3m>0(*)\\&\hoac{&x=m\\&x=2}.} $ \\
		Phương trình đã cho có hai nghiệm phân biệt khi và chỉ khi phương trình (1) có hai nghiệm phân biệt thỏa mãn (*) $\Leftrightarrow\heva{&m^2-m+1-3m>0\\&2^2-1+1-3m>0\\&m\neq 2}\Leftrightarrow\heva{&m^2-4m+1>0\\&4-3m>0}\Leftrightarrow m<2-\sqrt{3.}$ \\
		Theo giả thiết $|x_1-x_2|<15\Leftrightarrow(x_1+x_2)^2-4x_1x_2<225\Leftrightarrow m^2-4m-221<0\Leftrightarrow-13<m<17$ Do đó $-13<m<2-\sqrt{3}$. Vậy số các giá trị nguyên của $m$ thỏa mãn là 13.
	}
\end{ex}
\begin{ex}
	[THPT Yên Phong Số 1 Bắc Ninh 2019]%Câu 3.
	Gọi $S$ là tập tất cả các giá trị nguyên của tham số m với $m<64$ để phương trình $\log_{\tfrac{1}{5}}(x+m)+\log_5(2-x)=0$ có nghiệm. Tính tổng tất cả các phần tử của $S$. 
	\choice
	{$2018$}
	{$2016$}
	{\True $2015$}
	{$2013$}
	\loigiai{
		Ta có $\log_{\tfrac{1}{5}}(x+m)+\log_5(2-x)=0\Leftrightarrow\log_5(x+m)=\log_5(2-x)\Leftrightarrow\heva{&x<2\\&x=\dfrac{2-m}{2}.}$ \\
		Vì $x<2$ nên $\dfrac{2-m}{2}<2\Leftrightarrow m >-2$.\\
		Kết hợp với $m<64$. Khi đó $-2<m<64$.\\
		Vì $m\in\mathbb{Z}$ nên $m=\{-1;0;1\cdots 63\}$ có 65 giá trị.\\
		Vậy tổng $S$ các giá trị của $m$ để phương trình có nghiệm là $S=\dfrac{(-1+63)\cdot 65}{2}=2015$.
	}
\end{ex}
\begin{ex}
	[Mã 102 2019]%Câu 4.
	Cho phương trình $\log_9x^2-\log_3(6x-1)=-\log_3m$ ($m$ là tham số thực). Có tất cả bao nhiêu giá trị nguyên của $m$ để phương trình đã cho có nghiệm?
	\choice
	{$7$}
	{$6$}
	{\True $5$}
	{Vô số}
	\loigiai{
		Xét phương trình $\log_9x^2-\log_3(6x-1)=-\log_3m$.\\
		Điều kiện: $\heva{&x>\dfrac{1}{6}\\&m>0.}$ \\
		Khi đó\\
		$\begin{aligned}&\log_9x^2-\log_3(6x-1)=-\log_3m\Leftrightarrow\log_3x+\log_3m=\log_3(6x-1)\\&\Leftrightarrow mx=6x-1\Leftrightarrow x(6-m)=1(1).\end{aligned}$\\
		Với $m=6$, phương trình (1) trở thành $0=1$ (vô lý).\\
		Với $m\neq 6$, phương trình (1) có nghiệm $x=\dfrac{1}{6-m}$ \\
		$\Rightarrow\dfrac{1}{6-m}>\dfrac{1}{6}\Leftrightarrow\dfrac{1}{6-m}-\dfrac{1}{6}>0\Leftrightarrow\dfrac{m}{6-m}>0\Leftrightarrow 0<m<6 $.\\
		Vậy $0<m<6$. Mà $m\in\mathbb{Z}\Rightarrow m\in\left\{1; 2; 3; 4; 5\right\}$. Vậy có 5 giá trị nguyên của $m$ thỏa mãn.
	}
\end{ex}
\begin{ex}
	[Mã 103 2019]%Câu 5.
	Cho phương trình $\log_9x^2-\log_3(5x-1)=-\log_3m$ ($m$ là tham số thực). Có tất cả bao nhiêu giá trị nguyên của $m$ để phương trình đã cho có nghiệm?
	\choice
	{\True $4$}
	{$6$}
	{Vô số}
	{$5$}
	\loigiai{
		Điều kiện: $\heva{&x>\dfrac{1}{5}\\&m>0.}$ \\
		Xét phương trình: $\log_9x^2-\log_3(5x-1)=-\log_3m (1)$.\\
		Cách 1.\\
		$(1)\Leftrightarrow\log_3x-\log_3(5x-1)=-\log_3m\Leftrightarrow\log_3\dfrac{5x-1}{x}=\log_3m\Leftrightarrow\dfrac{5x-1}{x}=m\Leftrightarrow 5-\dfrac{1}{x}=m\quad(2)$.\\
		Xét $f(x)=5-\dfrac{1}{x}$ trên khoảng $\left(\dfrac{1}{5};+\infty\right)$.\\
		Có $f'(x)=\dfrac{1}{x^2}>0,\forall x\in\left(\dfrac{1}{5};+\infty\right)$ và $\lim\limits_{x\to+\infty} f(x)=\lim\limits_{x\to+\infty}\left(5-\dfrac{1}{x}\right)=5$.\\
		Ta có bảng biến thiên của hàm số $f(x)$:\\
		{\color{red} HINH O DAY}\\
		Phương trình $(1)$ có nghiệm khi và chỉ phương trình $(2)$ có nghiệm $x>\dfrac{1}{5}$.\\
		Từ bảng biến thiên suy ra phương trình $(1)$ có nghiệm khi và chỉ khi $0<m<5$.\\
		Mà $m\in\mathbb{Z}$ và $m>0$ nên $m\in\{1; 2; 3; 4\}$.\\
		Vậy có 4 giá trị nguyên của $m$ để phương trình đã cho có nghiệm.\\
		Cách 2.\\
		Với $\heva{&x>\dfrac{1}{5}\\&m>0}$, ta có\\
		$(1)\Leftrightarrow\log_3x-\log_3(5x-1)=-\log_3m\Leftrightarrow\log_3\dfrac{5x-1}{x}=\log_3m\Leftrightarrow\dfrac{5x-1}{x}=m\Leftrightarrow(5-m)x=1\quad(2)$ Với $m=5$, phương trình $(2)$ thành $0.x=1$ (vô nghiệm).\\
		Với $m\neq 5$, $(2)\Leftrightarrow x=\dfrac{1}{5-m}$.\\
		Xét $x>\dfrac{1}{5}\Leftrightarrow\dfrac{1}{5-m}>\dfrac{1}{5}\Leftrightarrow\dfrac{m}{5\cdot (5-m)}>0\Leftrightarrow 0<m<5$.\\
		Mà $m\in\mathbb{Z}$ và $m>0$ nên $m\in\{1; 2; 3; 4\}$.\\
		Vậy có 4 giá trị nguyên của $m$ để phương trình đã cho có nghiệm.
	}
\end{ex}
\begin{ex}
	[Mã 101 - 2019]%Câu 6.
	Cho phương trình $\log_9x^2-\log_3(3x-1)=-\log_3m$ ($m$ là tham số thực). Có tất cả bao nhiêu giá trị nguyên của tham số $m$ để phương trình đã cho có nghiệm?
	\choice
	{\True $2$}
	{$4$}
	{$3$}
	{Vô số}
	\loigiai{
		Điều kiện: $x>\dfrac{1}{3}$ và $m>0$.\\
		Phương trình đã cho tương đương: $\log_3x-\log_3(3x-1)=\log_3\dfrac{1}{m}\Leftrightarrow\dfrac{x}{3x-1}=\dfrac{1}{m}$.\\
		Xét hàm số $f(x)=\dfrac{x}{3x-1}$ với $x>\dfrac{1}{3}$.\\
		Có $f'(x)=-\dfrac{1}{(3x-1)^2}<0,\forall x>\dfrac{1}{3}$\\
		{\color{red} HINH O DAY}\\
		Dựa vào bảng biến thiên, phương trình có nghiệm khi $\dfrac{1}{m}>\dfrac{1}{3}\Leftrightarrow 0<m<3$.\\
		Do $m\in\mathbb{Z}\Rightarrow m\in\{1,2\}$.
	}
\end{ex}
\begin{ex}
	[Mã 104 2019]%Câu 7.
	Cho phương trình $\log_9x^2-4\log_3(4x-1)=-\log_3m$ ($m$ là tham số thực). Có tất cả bao nhiêu giá trị nguyên của $m$ để phương trình đã cho có nghiệm?
	\choice
	{$5$}
	{$3$}
	{\True Vô số}
	{$4$}
	\loigiai{
		Điều kiện: $x>\dfrac{1}{4}$. Phương trình đã cho $\Leftrightarrow\log_3x-4\log_3(4x-1)=-\log_3m$ \\
		$\Leftrightarrow\log_3x-\log_3(4x-1)^4=\log_3\dfrac{1}{m}\Leftrightarrow\log_3\dfrac{x}{(4x-1)^4}=\log_3\dfrac{1}{m}\Leftrightarrow m=\dfrac{(4x-1)^4}{x}=f(x) $.\\
		Xét hàm số $f(x)=\dfrac{(4x-1)^4}{x}$ có $f'(x)=\dfrac{16x(4x-1)^3-(4x-1)^4}{x^2}=\dfrac{(4x-1)^3(12x+1)}{x^2}>0,\forall x>\dfrac{1}{4}$.\\
		Suy ra bảng biến thiên:\\
		{\color{red} HINH O DAY}\\
		Do đó phương trình có nghiệm khi $m>0$. Vậy có vô số giá trị nguyên của $m$.
	}
\end{ex}
\begin{ex}
	[THPT Lương Thế Vinh Hà Nội 2019]%Câu 8.
	Cho phương trình $\log_{mx-5}\left(x^2-6x+12\right)=\log_{\sqrt{mx-5}}\sqrt{x+2}$, gọi $S$ là tập hợp tất cả các giá trị của tham số $m\in\mathbb{Z}$ để phương trình đã cho có nghiệm duy nhất. Tìm số phần tử của $S$. 
	\choice
	{\True $2$}
	{$0$}
	{$3$}
	{$1$}
	\loigiai{
		Điều kiện $\heva{&x+2>0\\&0<mx-5\neq 1}\Leftrightarrow\heva{&x >-2\\&5<mx\neq 6.}$ \\
		Với điều kiện trên, phương trình $\log_{mx-5}\left(x^2-6x+12\right)=\log_{\sqrt{mx-5}}\sqrt{x+2}\quad(*)$ \\
		$\Leftrightarrow\log_{mx-5}\left(x^2-6x+12\right)=\log_{mx-5}(x+2) $ \\
		$ \Leftrightarrow x^2-6x+12=x+2\Leftrightarrow\hoac{&x=2\\&x=5.} $ \\
		$x=2$ là nghiệm phương trình $(*)$ khi $5<2m\neq 6\Leftrightarrow\dfrac{5}{2}<m\neq 3$, vì $m\in \mathbb{Z}\Rightarrow\heva{&m\geq 4\\&m\in \mathbb{Z}.}$ \\
		$x=5$ là nghiệm phương trình $(*)$ khi $5<5m\neq 6\Leftrightarrow 1<m\neq\dfrac{6}{5}$, vì $m\in \mathbb{Z}\Rightarrow\heva{&m\geq 2\\&m\in \mathbb{Z}.}$ \\
		Phương trình $\log_{mx-5}\left(x^2-6x+12\right)=\log_{\sqrt{mx-5}}\sqrt{x+2}$ có nghiệm duy nhất khi $m=2$ hoặc $m=3$.\\
		Thử lại.\\
		$m=2$: $\log_{2x-5}\left(x^2-6x+12\right)=\log_{\sqrt{2x-5}}\sqrt{x+2}\Leftrightarrow\log_{2x-5}\left(x^2-6x+12\right)=\log_{2x-5}(x+2)$ \\
		$ \Leftrightarrow\heva{&x^2-6x+12=x+2\\&x+2>0\\&0<2x-5\neq 1}\Leftrightarrow x=5 $.\\
		$m=3$: $\log_{3x-5}\left(x^2-6x+12\right)=\log_{\sqrt{3x-5}}\sqrt{x+2}\Leftrightarrow\log_{3x-5}\left(x^2-6x+12\right)=\log_{3x-5}(x+2)$ \\
		$ \Leftrightarrow\heva{&x^2-6x+12=x+2\\&x+2>0\\&0<4x-5\neq 1}\Leftrightarrow x=5 $.\\
		Vậy có hai giá trị $m\in \mathbb{Z}$ thỏa mãn yêu cầu bài toán.
	}
\end{ex}
\begin{ex}
	Cho phương trình $\log_{2+\sqrt{5}}\left(2x^2-x-4m^2+2m\right)+\log_{\sqrt{\sqrt{5}-2}}\sqrt{x^2+mx-2m^2}=0$. Hỏi có bao nhiêu giá trị nguyên của tham số $m$ để phương trình đã cho có hai nghiệm $x_1^2+x_2^2=3$?
	\choice
	{$1$}
	{\True $0$}
	{$3$}
	{$4$}
	\loigiai{
		Phương trình đã cho tương đương với phương trình:\\
		$\log_{2+\sqrt{5}}\left(2x^2-x-4m^2+2m\right)+\log_{\sqrt{5}-2}\left(x^2+mx-2m^2\right)=0$ \\
		$\Leftrightarrow\log_{\sqrt{5}+2}\left(2x^2-x-4m^2+2m\right)-\log_{\sqrt{5}+2}\left(x^2+mx-2m^2\right)=0 $ \\
		$\Leftrightarrow\heva{&x^2+2mx-2m^2>0\\&2x^2-x+2m-4m^2=x^2+mx-2m^2}\Leftrightarrow\heva{&x^2+2mx-2m^2>0\\&x^2-(m+1)x+2m-2m^2=0} $ \\
		$\Leftrightarrow\heva{&x^2+mx-2m^2>0\\&\hoac{&x_1=2m\\&x_2=1-m}.} $ \\
		Phương trình đã cho có hai nghiệm $x_1,x_2$ thỏa $x_1^2+x_2^2=3$ \\
		$\Leftrightarrow\heva{&(2m)^2+m(2m)-2m^2>0\\&(1-m)^2+m(1-m)-2m^2>0\\&(2m)^2+(1-m)^2=3}\Leftrightarrow\heva{&4m^2>0\\&2m^2+m-1>0\\&5m^2-2m-2=0} $ \\
		$ \Leftrightarrow\heva{&m\neq 0\\&-1<m<\dfrac{1}{2}\\&m=\dfrac{1-\sqrt{11}}{5};m=\dfrac{1+\sqrt{11}}{5}}\Leftrightarrow m=\dfrac{1-\sqrt{11}}{5.} $ \\
		Vậy không có giá trị nguyên nào của $m$ thỏa yêu cầu đề bài.
	}
\end{ex}
\begin{ex}
	[HSG Bắc Ninh 2019]%Câu 10.
	Tìm tất cả các giá trị của tham số thực $m$ để phương trình $4\left(\log_2\sqrt{x}\right)^2-\log_{\tfrac{1}{2}}x+m=0$ có hai nghiệm phân biệt thuộc khoảng $(0;1)$. 
	\choice
	{\True $0<m<\dfrac{1}{4}$}
	{$0\leq m<\dfrac{1}{4}$}
	{$m\leq\dfrac{1}{4}$}
	{$-\dfrac{1}{4}<m<0$}
	\loigiai{
		Ta có\\
		$4\left(\log_2\sqrt{x}\right)^2-\log_{\tfrac{1}{2}}x+m=0\Leftrightarrow\left(2\log_2\sqrt{x}\right)^2+\log_2x+m=0\Leftrightarrow(\log_2x)^2+\log_2x=-m(1)$.\\
		Đặt $t=\log_2x$ với $t\in(-\infty;0)$.\\
		$(1)\Leftrightarrow t^2+t=-m$.\\
		Xét $f(t)=t^2+t$.\\
		$f'(t)=2t+1$, $f'(t)=0\Leftrightarrow t=-\dfrac{1}{2}$.\\
		Bảng biến thiên
		\begin{center}
			\begin{tikzpicture}
				\tkzTabInit[nocadre,lgt=1.2,espcl=3]{$t$/1.0,$f'(t)$/0.8,$f(t)$/2}{$-\infty$,$-\dfrac{1}{2}$,$0$}
				\tkzTabLine{,-,0,+,}
				\tkzTabVar{+/$+\infty$,-/$-\dfrac{1}{4}$,+/$0$}
			\end{tikzpicture}
		\end{center}
		Dựa vào bảng biến thiên $-\dfrac{1}{4} <-m<0\Leftrightarrow 0<m<\dfrac{1}{4}$.
	}
\end{ex}
\begin{ex}
	[THPT Đông Sơn Thanh Hóa 2019]%Câu 11.
	Tìm $m$ để phương trình: $(m-1)\log_{\tfrac{1}{2}}^2(x-2)^2+4(m-5)\log_{\tfrac{1}{2}}\dfrac{1}{x-2}+4m-4=0$ có nghiệm trên $\left[\dfrac{5}{2},4\right]$. 
	\choice
	{$m\in\mathbb{R}$}
	{\True $-3\leq m\leq\dfrac{7}{3}$}
	{$m\in\varnothing$}
	{$-3<m\leq\dfrac{7}{3}$}
	\loigiai{
		Điều kiện: $x>2$. Phương trình đã cho\\
		$\Leftrightarrow(m-1)\left[\log_{\tfrac{1}{2}}(x-2)^2\right]^2+4(m-5)\log_2(x-2)+4m-4=0 $ \\
		$\Leftrightarrow(m-1)\left[-2\log_2(x-2)\right]^2+4(m-5)\log_2(x-2)+4m-4=0 $ \\
		$ \Leftrightarrow 4(m-1)\log_2^2(x-2)+4(m-5)\log_2(x-2)+4m-4=0 $ \\
		$ \Leftrightarrow(m-1)\log_2^2(x-2)+(m-5)\log_2(x-2)+m-1=0 $. (1).\\
		Đặt $t=\log_2(x-2)$. Vì $x\in\left[\dfrac{5}{2};4\right]\Rightarrow t\in[-1; 1]$.\\
		Phương trình (1) trở thành $(m-1)t^2+(m-5)t+m-1=0$, $t\in[-1; 1]$. (2)\\
		$ \Leftrightarrow m=\dfrac{t^2+5t+1}{t^2+t+1}=f(t),t\in[-1; 1] $.\\
		Ta có $f'(t)=\dfrac{-4t^2+4}{\left(t^2+t+1\right)^2}=0\Leftrightarrow\hoac{&t=2\\&t=-2.}$ \\
		Bảng biến thiên
		{\color{red} HINH O DAY}\\
		Phương trình đã cho có nghiệm $x\in\left[\dfrac{5}{2};4\right]$ khi phương trình (2) có nghiệm $t\in[-1; 1]$.\\
		Từ bảng biến thiên suy ra $-3\leq m\leq\dfrac{7}{3}$.
	}
\end{ex}
\begin{ex}
	[Chuyên Bắc Giang 2019]%Câu 12.
	Tìm $m$ để phương trình $\log_2^2x-\log_2x^2+3=m$ có nghiệm $x\in [1;8]$. 
	\choice
	{$6\leq m\leq 9$}
	{$2\leq m\leq 3$}
	{\True $2\leq m\leq 6$}
	{$3\leq m\leq 6$}
	\loigiai{
		$\log_2^2x-\log_2x^2+3=m$ (1).\\
		Điều kiện: $x>0$ (*).\\
		pt (1) $\Leftrightarrow(\log_2x)^2-2\log_2x+3=m$.\\
		Cách 1: (Tự luận).\\
		Đặt $t=\log_2x$, với $x\in [1;8]$ thì $t\in [0;3]$.\\
		Phương trình trở thành: $t^2-2t+3=m$ (2).\\
		Để phương trình (1) có nghiệm $x\in [1;8]$ \\
		$ \Leftrightarrow $ phương trình (2) có nghiệm $t\in [0;3]$ \\
		$ \Leftrightarrow\min\limits_{[0;3]} f(t)\leq m\leq\max\limits_{[0;3]} f(t) $, trong đó $f(t)=t^2-2t+3$ \\
		$ \Leftrightarrow 2\leq m\leq 6 $. (bấm máy tính).
	}
\end{ex}
\begin{ex}
	[HSG Bắc Ninh - 2019]%Câu 13.
	Cho phương trình $\log_2^2x-2\log_2x-\sqrt{m+\log_2x}=m(*)$. Có bao nhiêu giá trị nguyên của tham số $m\in[-2019;2019]$ để phương trình (*) có nghiệm?
	\choice
	{\True $2021$}
	{$2019$}
	{$4038$}
	{$2020$}
	\loigiai{
		Điều kiện: $\heva{&x>0\\&m+\log_2x\geq 0}$.\\ $\log_2^2x-2\log_2x-\sqrt{m+\log_2x}=m\Leftrightarrow 4\log_2^2x-8\log_2x-4\sqrt{m+\log_2x}=4m$ \\
		$ \Leftrightarrow 4\log_2^2x-4\log_2x+1=4\sqrt{m+\log_2x}+4(m+\log_2x)+1 $ \\
		$\Leftrightarrow\left(2\log_2x-1\right)^2=\left(2\sqrt{m+\log_2x}+1\right)^2\Leftrightarrow\hoac{&2\sqrt{m+\log_2x}+1=2\log_2x-1\\&2\sqrt{m+\log_2x}+1=-2\log_2x+1} $ \\
		$\Leftrightarrow\hoac{&\sqrt{m+\log_2x}=\log_2x-1\\&\sqrt{m+\log_2x}=-\log_2x.} $ \\
		Trường hợp 1: $\sqrt{m+\log_2x}=-\log_2x\Leftrightarrow\heva{&\log_2x\leq 0\\&m+\log_2x=\log_2^2x}\Leftrightarrow\heva{&0<x\leq 1\\&\log_2^2x-\log_2x-m=0(1).}$ \\
		Đặt: $t=\log_2x(t\leq 0)$, phương trình $(1)$ trở thành: $t^2-t-m=0\Leftrightarrow t^2-t=m(2)$.\\
		Đặt: $g(t)=t^2-t(t\in(-\infty;0]$. Câu toán trở thành: Tìm giá trị của tham số $m$ để phương trình $(2)$ có ít nhất 1 nghiệm $t\leq 0$.\\
		Ta có $g(t)=t^2-t\Rightarrow g'(t)=2t-1<0\forall t\leq 0$.\\
		Ta có bảng biến thiên:
		\begin{center}
			\begin{tikzpicture}
				\tkzTabInit[nocadre,lgt=1.2,espcl=3]{$t$/0.8,$g'(t)$/0.8,$g(t)$/2.5}{$-\infty$,{},$0$}
				\tkzTabLine{,-,}
				\tkzTabVar{+/$+\infty$,R,-/$0$}
			\end{tikzpicture}
		\end{center}
		Dựa vào bảng biến thiên, suy ra: để phương trình $(2)$ có ít nhất 1 nghiệm $t\leq 0$ thì $m\geq 0$ (*).\\
		Trường hợp 2: $\sqrt{m+\log_2x}=\log_2x-1\Leftrightarrow\heva{&\log_2x\geq 1\\&m+\log_2x=\log_2^2x-2\log_2x+1}$ \\
		$ \Leftrightarrow\heva{&\log_2x\geq 1\\&\log_2^2x-3\log_2x+1-m=0(3).} $ \\
		Đặt: $t=\log_2x(t\geq 1)$, phương trình $(1)$ trở thành: $t^2-3t+1-m=0\Leftrightarrow m=t^2-3t+1(4)$.\\
		Đặt: $g(t)=t^2-t+1, t\in[1;+\infty)$.\\
		Ta có $g(t)=t^2-3t+1\Rightarrow g'(t)=2t-3$.\\
		$g'(t)=0\Leftrightarrow 2t-3=0\Leftrightarrow t=\dfrac{3}{2}\in[1;+\infty)$.\\
		Câu toán trở thành: Tìm giá trị của tham số $m$ để phương trình $(4)$ có ít nhất 1 nghiệm $t\geq 1$.\\
		Ta có BBT: 
		\begin{center}
			\begin{tikzpicture}
				\tkzTabInit[nocadre,lgt=1.2,espcl=3]{$t$/1.0,$g'(t)$/0.8,$g(t)$/2.5}{$1$,$\dfrac{3}{2}$,$+\infty$}
				\tkzTabLine{,-,0,+,}
				\tkzTabVar{+/$-1$,-/$-\dfrac{5}{4}$,+/$+\infty$}
			\end{tikzpicture}
		\end{center}
		Dựa vào BBT, suy ra: để phương trình $(4)$ có ít nhất 1 nghiệm $t\geq 1$ thì $m\geq-\dfrac{5}{4}$ (**).\\
		Kết hợp (*) và (**), $m\in[-2019;2019]\Rightarrow m\in\left\{-1;0;1;2;\ldots;2019\right\}$.\\
		Vậy có tất cả $2021$ giá trị của $m$ thỏa mãn yêu cầu bài toán.
	}
\end{ex}
\begin{ex}
	[Đề Tham Khảo 2017]%Câu 14.
	Hỏi có bao nhiêu giá trị $m$ nguyên trong $[-2017;2017]$ để phương trình $\log(mx)=2\log(x+1)$ có nghiệm duy nhất?
	\choice
	{$4014$}
	{\True $2018$}
	{$4015$}
	{$2017$}
	\loigiai{
		Điều kiện $x >-1, mx>0$.\\
		$\log(mx)=2\log(x+1)\Leftrightarrow mx=(x+1)^2\Leftrightarrow m=\dfrac{(x+1)^2}{x}$.\\
		Xét hàm $f(x)=\dfrac{(x+1)^2}{x}\quad\left(x >-1,x\neq 0\right)$; $f'(x)=\dfrac{x^2-1}{x^2}=0\Leftrightarrow\hoac{&x=1\\&x=-1(l).}$ \\
		Lập bảng biến thiên
		\begin{center}
			\begin{tikzpicture}
				\tkzTabInit[nocadre,lgt=1,espcl=3]{$x$/0.8,$y'$/0.8,$y$/2.5}{$-1$,$0$,$1$,$+\infty$}
				\tkzTabLine{,-,d,-,0,+,}
				\tkzTabVar{+/$0$,-D+/$-\infty$/$+\infty$,-/$4$,+/$+\infty$}
			\end{tikzpicture}
		\end{center}
		Dựa vào BBT, phương trình có nghiệm duy nhất khi và chỉ khi $\hoac{&m=4\\&m<0\cdot.}$ \\
		Vì $m\in[-2017;2017]$ và $m\in\mathbb{Z}$ nên chỉ có $2018$ giá trị $m$ nguyên thỏa yêu cầu là\\
		$m\in\left\{-2017;-2016;\ldots;-1;4\right\}$.\\
		Chú ý: Trong lời giải, ta đã có thể bỏ qua điều kiện $mx>0$ vì với phương trình $\log_af(x)=\log_ag(x)$ với $0<a\neq 1$ ta chỉ cần điều kiện $f(x)>0$.
	}
\end{ex}
\begin{ex}
	[THPT An Lão Hải Phòng 2019]%Câu 15.
	Tìm tập hợp tất cả các giá trị thực của tham số $m$ để phương trình $mx-\ln x=0$ có hai nghiệm phân biệt thuộc khoảng $(2;3)$ 
	\choice
	{$\left(\dfrac{\ln 2}{2};\dfrac{\ln 3}{3}\right)$}
	{$\left(-\infty;\dfrac{\ln 2}{2}\right)\cup\left(\dfrac{\ln 3}{3};+\infty\right)$}
	{$\left(\dfrac{\ln 2}{2};\dfrac{1}{e}\right)$}
	{\True $\left(\dfrac{\ln 3}{3};\dfrac{1}{e}\right)$}
	\loigiai{
		$mx-\ln x=0\Leftrightarrow m=\dfrac{\ln x}{x},\forall x\in(2;3)$.\\
		Đặt $f(x)=\dfrac{\ln x}{x},\forall x\in(2;3)$.\\
		$f'(x)=\dfrac{1-\ln x}{x^2}$; $f'(x)=0\Leftrightarrow x=e$.\\
		Bảng biến thiên
		\begin{center}
			\begin{tikzpicture}
				\tkzTabInit[nocadre,lgt=1.2,espcl=3]{$x$/0.8,$f'(x)$/0.8,$f(x)$/2.5}{$2$,$\mathrm{e}$,$3$}
				\tkzTabLine{,+,0,-,}
				\tkzTabVar{-/$\dfrac{\ln2}{2}$,+/$\dfrac{1}{\mathrm{e}}$,-/$\dfrac{\ln3}{3}$}
			\end{tikzpicture}
		\end{center}
		Để phương trình có hai nghiệm phân biệt thì $m\in\left(\dfrac{\ln 3}{3};\dfrac{1}{e}\right)$.
	}
\end{ex}
\begin{ex}
	[THPT Dông Sơn Thanh Hóa 2019]%Câu 16.
	Tổng tất cả các giá trị của tham số $m$ sao cho phương trình:\\
	$2^{(x-1)^2}\cdot\log_2\left(x^2-2x+3\right)=4^{|x-m|}\cdot\log_2(2|x-m|+2)$ có đúng ba nghiệm phân biệt là 
	\choice
	{$2$}
	{$\dfrac{3}{2}$}
	{$0$}
	{\True $3$}
	\loigiai{
		Tập xác định $\mathscr{D}=\mathbb{R}$.\\
		$\begin{aligned}&2^{(x-1)^2}\cdot \log_2\left(x^2-2x+3\right)=4^{|x-m|}\cdot \log_2(2|x-m|+2)\\&\Leftrightarrow 2^{(x-1)^2}\cdot \log_2\left((x-1)^2+2\right)=2^{2|x-m|}\cdot \log_2(2|x-m|+2). (*)\end{aligned}$\\
		Đặt $f(t)=2^t\log_2(t+2), t\geq 0$; $f'(t)=2^t\ln 2\cdot\log_2(t+2)+2^t\dfrac{1}{(t+2)\ln 2}>0,\forall t\geq 0$.\\
		Vậy hàm số $f(t)=2^t\log_2(t+2)$ đồng biến trên $(0;+\infty)$.\\
		Từ (*) ta có $f\left[(x-1)^2\right]=f[2|x-m|]\Leftrightarrow (x-1)^2=2|x-m|\Leftrightarrow\hoac{&2(x-m)=(x-1)^2\\&2(x-m)=-(x-1)^2}$ \\
		$ \Leftrightarrow\hoac{&g(x)=x^2-4x+1+2m=0 (a)\\&x^2=2m-1 (b).} $ \\
		Do các phương trình $(a)$ và $(b)$ là phương trình bậc hai nên để phương trình ban đầu có 3 nghiệm phân biệt ta có các trường hợp sau:\\
		Trường hợp 1: $m=\dfrac{1}{2}$, (b) chỉ có nghiệm kép bằng 0 và (a) có 2 nghiệm phân biệt khác 0 (thỏa mãn).\\
		Trường hợp 2: $m>\dfrac{1}{2}$, (b) có 2 nghiệm phân biệt $x=\pm\sqrt{2m-1}$ và (a) có 2 nghiệm phân biệt trong đó có 1 nghiệm bằng $\pm\sqrt{2m-1}$ \\
		$ \Leftrightarrow\heva{&\Delta'>0\\&g(\pm\sqrt{2m-1})=0}\Leftrightarrow\heva{&\Delta'>0\\&g(\pm\sqrt{2m-1})=0}\Leftrightarrow\heva{&m<\dfrac{3}{2}\\&m=1}\Leftrightarrow m=1 $ (thỏa mãn).\\
		Trường hợp 3: $m>\dfrac{1}{2}$, (b) có 2 nghiệm phân biệt $x=\pm\sqrt{2m-1}$ và (a) có nghiệm kép khác $\pm\sqrt{2m-1}$ \\
		$ \Leftrightarrow\heva{&\Delta'=0\\&g(\pm\sqrt{2m-1})\neq 0}\Leftrightarrow\heva{&m=\dfrac{3}{2}\\&m\neq 1}\Leftrightarrow m=\dfrac{3}{2} $ (thỏa mãn).\\
		Vậy tổng các giá trị của $m$ là $\dfrac{1}{2}+1+\dfrac{3}{2}=3$.
	}
\end{ex}
\begin{ex}
	Tìm tất cả các giá trị của $m$ để phương trình $\ln\left(m+\ln(m+\sin x)\right)=\sin x$ có nghiệm. 
	\choice
	{$\dfrac{1}{e}+1\leq m\leq e-1$}
	{\True $1\leq m\leq e-1$}
	{$1\leq m\leq\dfrac{1}{e}+1$}
	{$1\leq m<e-1$}
	\loigiai{
		Đặt $u=\ln(m+\sin x)$ ta được hệ phương trình: $\heva{&u=\ln(m+\sin x)\\&\ln(m+u)=\sin x}\Leftrightarrow\heva{&\mathrm{e}^u=m+\sin x\\&\mathrm{e}^{\sin x}=m+u.}$ \\
		Từ hệ phương trình ta suy ra: $\mathrm{e}^u+u=\mathrm{e}^{\sin x}+\sin x(*)$.\\
		Xét hàm số $f(t)=\mathrm{e}^t+t$ có $f'(t)=\mathrm{e}^t+1>0,\forall t\in\mathbb{R}$. Hàm số $f(t)$ đồng biến trên $\mathbb{R}$.\\
		$(*)\Leftrightarrow f(u)=f(\sin x)\Leftrightarrow u=\sin x$.\\
		Khi đó ta được: $\ln(m+\sin x)=\sin x\Leftrightarrow\mathrm{e}^{\sin x}-\sin x=m(**)$.\\
		Đặt $z=\sin x, z\in[-1; 1]$. Phương trình $(**)$ trở thành: $\mathrm{e}^z-z=m(**)$.\\
		Xét hàm số: $g(z)=\mathrm{e}^z-z$ trên $[-1; 1]$.\\
		Hàm số $g(z)=\mathrm{e}^z-z$ liên tục trên $[-1; 1]$ và có $\max\limits_{[-1; 1]} g(z)=g(1)=e-1,\min\limits_{[-1; 1]} g(z)=g(0)=1$.\\
		Hệ phương trình ban đầu có nghiệm $\Leftrightarrow$ phương trình $(**)$ có nghiệm $\Leftrightarrow 1\leq m\leq e-1$.}
\end{ex}
\begin{ex}
	[THPT Yên Dũng 2 - Bắc Giang 2019]%Câu 18.
	Số các giá trị nguyên của tham số $m$ để phương trình $\log_{\sqrt{2}}(x-1)=\log_2(mx-8)$ có hai nghiệm phân biệt là
	\choice
	{$5$}
	{Vô số}
	{$4$}
	{\True $3$}
	\loigiai{
		Điều kiện: $x>1$.\\
		Ta có $\log_{\sqrt{2}}(x-1)=\log_2(mx-8)\Leftrightarrow log_2(x-1)^2=\log_2(mx-8)\Leftrightarrow (x-1)^2=mx-8$ \\
		$ \Leftrightarrow x^2-2x+9=mx $. Do $x>1$ nên suy ra $\dfrac{x^2-2x+9}{x}=m$.\\
		Xét hàm số $f(x)=\dfrac{x^2-2x+9}{x}$ trên khoảng $(1;+\infty)$.\\
		$f'(x)=\dfrac{x^2-9}{x^2}$, $f'(x)=0\Leftrightarrow x=\pm 3$.\\
		Bảng biến thiên.\\
		\textbf{(VẼ HÌNH)}\\
		Nhìn vào bảng biến thiên ta thấy yêu cầu của bài toán là $4<m<8$. Do $m$ nguyên nên $m\in\{5;6;7\}$.\\
		Vậy có 3 giá trị của $m$ thỏa mãn yêu cầu bài toán.
	}
\end{ex}
\begin{ex}
	[THPT Trần Phú - 2019]%Câu 19.
	Có bao nhiêu giá trị nguyên dương của tham số $m$ để phương trình $m^2\ln\left(\dfrac{x}{e}\right)=(2-m)\ln x-4$ có nghiệm thuộc vào đoạn $\left[\dfrac{1}{e};1\right]$?
	\choice
	{\True $1$}
	{$2$}
	{$3$}
	{$4$}
	\loigiai{
		Có $m^2\ln\left(\dfrac{x}{e}\right)=(2-m)\ln x-4\Leftrightarrow m^2(\ln x-1)=(2-m)\ln x-4\Leftrightarrow\left(m^2+m-2\right)\ln x=m^2-4 (1)$.\\
		Với $m^2+m-2=0\Rightarrow m=1 (m>0)$, $(1)\Leftrightarrow 0\ln x=-3$ (Vô nghiệm) $\Rightarrow$ Loại $m=1$.\\
		Với $m\neq 1$, $(1)\Leftrightarrow\ln x=\dfrac{m-2}{m-1} (2)$.\\
		Hàm số $y=\ln x$ đồng biến trên $\left[\dfrac{1}{e};1\right]\Rightarrow\ln x\in[-1;0]$.\\
		Phương trình $(2)$ có nghiệm thuộc đoạn $\left[\dfrac{1}{e};1\right]$ khi\\
		$-1\leq\dfrac{m-2}{m-1}\leq 0\Leftrightarrow\heva{&\dfrac{m-2}{m-1}\geq-1\\&\dfrac{m-2}{m-1}\leq 0}\Leftrightarrow\heva{&\hoac{&m\geq\dfrac{3}{2}\\&m<1}\\&1<m\leq 2}\Leftrightarrow\dfrac{3}{2}\leq m\leq 2\Rightarrow m=2$.\\
		Vậy có $1$ giá trị nguyên dương của tham số $m$ thỏa yêu cầu bài toán.
	}
\end{ex}
\begin{ex}
	[THPT Trần Phú - 2019]%Câu 20.
	Có bao nhiêu giá trị của tham số $m$ để phương trình $4\log_{36}^2x-m\log_6\dfrac{x}{6}+2=0$ có hai nghiệm phân biệt $x_1$, $x_2$ thỏa mãn $x_1\cdot x_2-72\sqrt{x_1\cdot x_2}+1296\leq 0$ 
	\choice
	{\True $0$}
	{$1$}
	{$2$}
	{$3$}
	\loigiai{
		$4\log_{36}^2x-m\log_6\dfrac{x}{6}+2=0$ (Điều kiện $x>0$)\\
		$ \Leftrightarrow\log_6^2x-m\log_6x+m+2=0 $.\\
		Phương trình có hai nghiệm phân biệt khi $\Delta=m^2-4(m+2)>0\Leftrightarrow\hoac{&m<2-2\sqrt{3}\\&m>2+2\sqrt{3}.}$ \\
		$x_1\cdot x_2-72\sqrt{x_1\cdot x_2}+1296\leq 0\Leftrightarrow\sqrt{x_1\cdot x_2}=36\Leftrightarrow x_1\cdot x_2=1296$ \\
		$ \Rightarrow\log_6(x_1\cdot x_2)=4\Leftrightarrow\log_6x_1+\log_6x_2=4\Rightarrow m=4 $ (không thỏa điều kiện của $m$).
	}
\end{ex}
\begin{ex}
	[Chuyên Lê Hồng Phong Nam Định 2019]%Câu 21.
	Tập hợp các giá trị thực của tham số m để phương trình $\log_{2019}\left(4-x^2\right)+log_{\tfrac{1}{2019}}(2x+m-1)=0$ có hai nghiệm thực phân biệt là $T=(a;b)$. Tính $S=2a+b$. 
	\choice
	{$18$}
	{$8$}
	{$20$}
	{\True $16$}
	\loigiai{
		Tập xác định $\mathscr{D}=(-2;2)\cap\left(\dfrac{1-m}{2};+\infty\right)$.\\
		Khi đó, phương trình đã cho trở thành.\\
		$\log_{2019}\dfrac{4-x^2}{2x+m-1}=0\Leftrightarrow 4-x^2=2x+m-1\Leftrightarrow x^2+2x+m-5=0 (*)$.\\
		Phương trình $(*)$ có 2 nghiệm phân biệt\\
		$ \Rightarrow\Delta=1^2-1\cdot (m-5)=6-m>0\Leftrightarrow m<6 (1) $.\\
		Khi đó phương trình $(1)$ có 2 nghiệm lần lượt là $x_1=-1+\sqrt{6-m}; x_2=-1-\sqrt{6-m}$.\\
		Trường hợp 1: $\dfrac{1-m}{2}\leq-2\Leftrightarrow m\geq 5$ (2) $\Rightarrow D=(-2; 2)$.\\
		Phương trình (1) có 2 nghiệm $x_1,x_2\in \mathscr{D}\Leftrightarrow\heva{&-1+\sqrt{6-m}<2\\&-1-\sqrt{6-m} >-2}\Leftrightarrow\heva{&\sqrt{6-m}<3\\&\sqrt{6-m}<1}\Leftrightarrow m>5$ (3).\\
		Từ (1), (2) và (3) suy ra $5<m<6$.\\
		Trường hợp 2: $-2<\dfrac{1-m}{2}<2\Leftrightarrow-3<m<5$ (4)\\
		$ \Rightarrow D=\left(\dfrac{1-m}{2}; 2\right) $.\\
		Phương trình (1) có 2 nghiệm.\\
		$x_1,x_2\in \mathscr{D}\Leftrightarrow\heva{&-1+\sqrt{6-m}<2\\&-1-\sqrt{6-m}>\dfrac{1-m}{2}}\Leftrightarrow\heva{&\sqrt{6-m}<3\\&\sqrt{6-m}<\dfrac{m-3}{2}}\Leftrightarrow\heva{&m >-3\\&\hoac{&m <-3\\&m>5}}\Leftrightarrow m>5$ (5).\\
		Từ (4) và (5) suy ra $m\in\varnothing$. Vậy $5<m<6$. Suy ra $a=5, b=6\Rightarrow 2a+b=16$.
	}
\end{ex}
\begin{ex}
	[THPT Cẩm Bình 2019]%Câu 22.
	Có bao nhiêu giá trị nguyên của $m$ để phương trình $\log_3(x+3)+m\log_{\sqrt{x+3}}9=16$ có hai nghiệm thỏa mãn $-2<x_1<x_2$. 
	\choice
	{$17$}
	{$16$}
	{$14$}
	{\True $15$}
	\loigiai{
		Điều kiện xác định $x >-3$ và $x\neq-2$.\\
		Biến đổi phương trình đã cho về phương trình sau $\log_3(x+3)+4m\log_{(x+3)}3-16=0$ \\
		$ \Leftrightarrow\log_3^2(x+3)-16\log_3(x+3)+4m=0 (1) $.\\
		Đặt $\log_3(x+3)=t$ phương trình $(1)$ trở thành $t^2-16t+4m=0(2)$.\\
		Ta có $\log_3(x+3)=t\Leftrightarrow x=3^t-3$.\\
		Theo điều kiện đề bài thì $x >-2$ nên $3^t-3 >-2\Leftrightarrow t>0$.\\
		Vậy để phương trình $\log_3(x+3)+m\log_{\sqrt{x+3}}9=16$ có hai nghiệm thỏa mãn $-2<x_1<x_2$.\\
		thì phương trình $(2)$ phải có hai nghiệm $t$ dương phân biệt\\
		$ \Leftrightarrow\heva{&\Delta'>0\\&t_1+t_2=16>0\\&t_1\cdot t_2=4m>0}\Leftrightarrow\heva{&64-4m>0\\&4m>0}\Leftrightarrow 0<m<16 $. Vậy có 15 giá trị nguyên $m$ thỏa mãn.
	}
\end{ex}
\begin{ex}
	[Chuyên Hoàng Văn Thụ - Hòa Bình 2019]%Câu 23.
	Tập hợp các số thực $m$ để phương trình $\ln(3x-mx+1)=\ln\left(-x^2+4x-3\right)$ có nghiệm là nửa khoảng $[a;b)$. Tổng $a+b$ bằng
	\choice
	{$\dfrac{10}{3}$}
	{$4$}
	{$\dfrac{22}{3}$}
	{\True $7$}
	\loigiai{
		Ta có\\
		$\begin{aligned}&\ln(3x-mx+1)=\ln\left(-x^2+4x-3\right) (1)\\&\Leftrightarrow\heva{&-x^2+4x-3>0\\&3x-mx+1=-x^2+4x-3}\Leftrightarrow\heva{&1<x<3\\&x^2-x+4=mx}\\&\Leftrightarrow\heva{&1<x<3\\&\dfrac{x^2-x+4}{x}=m}\Leftrightarrow\heva{&1<x<3\\&\dfrac{x^2-x+4}{x}=m}\Leftrightarrow\heva{&1<x<3\\&x+\dfrac{4}{x}-1=m (2)}.\end{aligned}$ \\
		Xét hàm số: $f(x)=x+\dfrac{4}{x}-1; x\in(1;3)$ có $f'(x)=1-\dfrac{4}{x^2}$.\\
		$f'(x)=1-\dfrac{4}{x^2}=0\Leftrightarrow\hoac{&x=2\in(1;3)\Rightarrow f(2)=3\\&x=-2\notin(1;3).}$ \\
		Bảng biến thiên:
		\begin{center}
			\begin{tikzpicture}
				\tkzTabInit[nocadre,lgt=1.2,espcl=3]{$x$/0.8,$f'(x)$/0.8,$f(x)$/2.5}{$1$,$2$,$3$}
				\tkzTabLine{,-,0,+,}
				\tkzTabVar{+/$4$,-/$3$,+/$\dfrac{10}{3}$}
			\end{tikzpicture}
		\end{center}
		Dựa vào bảng biến thiên ta có\\
		Phương trình (1) có nghiệm\\
		$ \Leftrightarrow $ Phương trình (2) có nghiệm thuộc khoảng $(1;3)$ \\
		$ \Leftrightarrow 3\leq m<4\Leftrightarrow m\in[3;4) $.\\
		Suy ra $\heva{&a=3\\&b=4}\Rightarrow a+b=3+4=7$.
	}
\end{ex}
\begin{ex}
	[Cần Thơ 2019]%Câu 24.
	Cho phương trình $\log_2^2x-2\log_2x-4\sqrt{1-\log_2x}=m$, với $m$ là tham số thực. Số các giá trị nguyên thuộc đoạn $[-2019; 2019]$ của $m$ để phương trình đã cho có nghiệm là
	\choice
	{$2021$}
	{\True $2024$}
	{$2023$}
	{$2020$}
	\loigiai{
		Điều kiện xác định: $1-\log_2x\geq 0\Leftrightarrow\log_2x\leq 1\Leftrightarrow 0<x\leq 2$.\\
		Với điều kiện trên thì phương trình tương đương với $(1-\log_2x)^2-4\sqrt{1-\log_2x}-1=m (1)$.\\
		Đặt $t=\sqrt{1-\log_2x}$, vì $x\in(0; 2]$ nên $t\geq 0$. Khi đó, $(1)$ trở thành $t^4-4t-1=m (2)$.\\
		Để $(1)$ có nghiệm $x\in(0; 2]$ thì $(2)$ có nghiệm $t\geq 0$.\\
		Xét hàm số $f(t)=t^4-4t-1$, $t\in[0;+\infty)$.\\
		Ta có $f'(t)=4t^3-4$. Cho $f'(t)=0\Leftrightarrow t=1\in[0;+\infty)$.\\
		Ta được bảng biến thiên của $f(t)$ như sau:
		\begin{center}
			\begin{tikzpicture}
				\tkzTabInit[nocadre,lgt=1.2,espcl=3]{$t$/0.8,$f'(t)$/0.8,$f(t)$/2.5}{$0$,$1$,$+\infty$}
				\tkzTabLine{,-,0,+,}
				\tkzTabVar{+/$-1$,-/$-4$,+/$+\infty$}
			\end{tikzpicture}
		\end{center}
		Theo BBT, để $(2)$ có nghiệm $t\geq 0$ thì $m\geq-4$, mà $m\in[-2019; 2019]\cap\mathbb{Z}$ nên tập hợp các giá trị của $m$ cần tìm là $\left\{-4;-3;-2;-1;0;1;\ldots;2019\right\}$.\\
		Vậy có tất cả 2024 giá trị nguyên của $m$ thuộc đoạn $[-2019; 2019]$ để phương trình đã cho có nghiệm.
	}
\end{ex}
\begin{ex}[Nam Định - 2019]%Câu 25.
	Tìm tất cả các giá trị thực của tham số $m$ để phương trình $x\log_3(x+1)=\log_9\left[9(x+1)^{2m}\right]$ có hai nghiệm phân biệt. 
	\choice
	{$m\in(-1; 0)$}
	{$m\in(-2; 0)$}
	{\True $m\in(-1;+\infty)$}
	{$m\in[-1; 0)$}
	\loigiai{
		Cách 1.\\
		Điều kiện: $x >-1$.\\
		Ta có pt: $x\log_3(x+1)=\log_9\left[9(x+1)^{2m}\right]\Leftrightarrow x\log_3(x+1)=1+m\log_3(x+1)$ \\
		$ \Leftrightarrow(x-m)\log_3(x+1)=1 $ (1).\\
		Đặt: $\log_3(x+1)=t\Rightarrow x=3^t-1$.\\
		Ta có, Pt (1) $\Rightarrow\left(3^t-m-1\right)\cdot t=1\Rightarrow f(t)=3^t-\dfrac{1}{t}-1=m$, với $t\neq 0$.\\
		Đặt: $f(t)=3^t-\dfrac{1}{t}-1$, với $t\neq 0$ \\
		$ \Rightarrow f'(t)=3^t\cdot\ln 3+\dfrac{1}{t^2}>0, t\in(-\infty; 0),(0;+\infty) $.\\
		Suy ra, $f(t)=3^t-\dfrac{1}{t}-1$ là hàm số đồng biến trên $(-\infty; 0)$ và $(0;+\infty)$.\\
		Ta xét các giới sau:\\
		$\lim\limits_{t\to-\infty}\left(3^t-\dfrac{1}{t}-1\right)=-1$, $\lim\limits_{t\to+\infty}\left(3^t-\dfrac{1}{t}-1\right)=+\infty$.\\
		$\lim\limits_{t\to 0^+}\left(3^t-\dfrac{1}{t}-1\right)=-\infty$, $\lim\limits_{t\to 0^-}\left(3^t-\dfrac{1}{t}-1\right)=+\infty$.\\
		Ta có bảng biến thiên của hàm số $f(t)=3^t-\dfrac{1}{t}-1$, với $t\in(-\infty; 0),(0;+\infty)$. \\
		\begin{center}
			\begin{tikzpicture}
				\tkzTabInit[nocadre,lgt=1.2,espcl=3]{$t$/0.8,$f'(t)$/0.8,$f(t)$/2.5}{$-\infty$,$0$,$+\infty$}
				\tkzTabLine{,+,d,+,}
				\tkzTabVar{-/$-1$,+D-/$+\infty$/$-\infty$,+/$+\infty$}
			\end{tikzpicture}
		\end{center}
		{\color{red} HINH O DAY}\\
		Ta có, số nghiệm của Pt (1) cũng chính là số nghiệm của đồ thị hàm số (C) $f(t)=3^t-\dfrac{1}{t}-1$.\\
		và đồ thị hàm số $y=m$ (song song hoặc trùng với trục hoành).\\
		Dựa, vào đồ thị ở hình vẽ trên, để phương trình $x\log_3(x+1)=\log_9\left[9(x+1)^{2m}\right]$ có ba nghiệm khi $m\in(-1;+\infty)$.\\
		Cách 2.\\
		Điều kiện: $x >-1$.\\
		Ta có $x\log_3(x+1)=\log_9\left[9(x+1)^{2m}\right]$ (1).\\
		Nhận thấy $x=0$ không là nghiệm phương trình trên.\\
		Phương trình (1) $\Leftrightarrow(x-m)\log_3(x+1)=1\Leftrightarrow x-\dfrac{1}{\log_3(x+1)}=m$.\\
		Đặt: $f(x)=x-\dfrac{1}{\log_3(x+1)}\Rightarrow f'(x)=1+\dfrac{1}(x+1)\ln 3\cdot{\left(\log_3(x+1)\right)^2}>0,\forall x\in(-1;+\infty)$.\\
		Suy ra $f(x)=x-\dfrac{1}{\log_3(x+1)}$ là hàm số đồng biến $\forall x\in(-1;+\infty)$.\\
		Ta có bảng biến thiên của hàm số $f(x)=x-\dfrac{1}{\log_3(x+1)}$.
		\begin{center}
			\begin{tikzpicture}
				\tkzTabInit[nocadre,lgt=1.2,espcl=3]{$x$/0.8,$f'(x)$/0.8,$f(x)$/2.5}{$-\infty$,$0$,$+\infty$}
				\tkzTabLine{,+,d,+,}
				\tkzTabVar{-/$-1$,+D-/$+\infty$/$-\infty$,+/$+\infty$}
			\end{tikzpicture}
		\end{center}
		Dựa, vào bảng biến thiên ở hình vẽ trên, để phương trình $x\log_3(x+1)=\log_9\left[9(x+1)^{2m}\right]$ có ba nghiệm khi $m\in(-1;+\infty)$.
	}
\end{ex}
\begin{ex}
	[THPT Hoàng Hoa Thám - Hưng Yên 2019]%Câu 26.
	Cho $a$, $b$ là các số thực dương lớn hơn 1, thay đổi thỏa mãn $a+b=2019$ để phương trình $5\log_ax\cdot\log_bx-4\log_ax-3\log_bx-2019=0$ luôn có hai nghiệm phân biệt $x_1; x_2$. Biết giá trị lớn nhất của $\ln(x_1\cdot x_2)$ bằng $\dfrac{3}{5}\ln\left(\dfrac{m}{7}\right)+\dfrac{4}{5}\ln\left(\dfrac{n}{7}\right)$; với $m, n$ là các số nguyên dương. Tính $S=m+2n$ 
	\choice
	{\True $22209$}
	{$20190$}
	{$2019$}
	{$14133$}
	\loigiai{
		Theo bài ra ta có\\
		$\begin{aligned}&5\log_ax\cdot\log_bx-4\log_ax-3\log_bx-2019=0\\&\Leftrightarrow 5\log_ax\cdot\left(\log_ba\cdot\log_ax\right)-4\log_ax-3\left(\log_ba\cdot\log_ax\right)-2019=0\\&\Leftrightarrow 5\log_ba\cdot (\log_ax)^2-\left(4+3\log_ba\right)\log_ax-2019=0(*).\end{aligned}$\\
		Vì $a, b>1\Rightarrow\log_ba\cdot (-2019)<0\Rightarrow(*)$ luôn có hai nghiệm phân biệt $x_1; x_2$.\\
		Theo Vi-ét ta có\\
		$\begin{aligned}&\log_ax_1+\log_ax_2=\dfrac{4+3\log_ba}{5\log_ba}\Leftrightarrow\log_a(x_1\cdot x_2)=\dfrac{4+3\cdot\dfrac{\ln a}{\ln b}}{5\dfrac{\ln a}{\ln b}}\\&\Leftrightarrow\dfrac{\ln(x_1\cdot x_2)}{\ln a}=\dfrac{4\ln b+3\ln a}{5\ln a}\Leftrightarrow\ln(x_1\cdot x_2)=\dfrac{1}{5}\left(4\ln(2019-a)+3\ln a\right).\end{aligned}$\\
		Xét $f(a)=\dfrac{1}{5}\left(4\ln(2019-a)+3\ln a\right)$ với $a\in(1; 2019)$.\\
		Ta có $f'(a)=\dfrac{1}{5}\left(\dfrac{-4}{2019-a}+\dfrac{3}{a}\right)$; $f'(a)=0\Leftrightarrow a=\dfrac{6057}{7}$.\\
		Bảng biến thiên
		\begin{center}
			\begin{tikzpicture}
				\tkzTabInit[nocadre,lgt=1.2,espcl=3]{$a$/1.0,$f'(a)$/0.8,$f(a)$/2.5}{$1$,$\dfrac{6057}{7}$,$2019$}
				\tkzTabLine{,+,0,-,}
				\tkzTabVar{-/{},+/$f\left(\dfrac{6057}{7}\right)$,-/{}}
			\end{tikzpicture}
		\end{center}
		Từ bảng biến thiên ta được giá trị lớn nhất của $\ln(x_1\cdot x_2)=\dfrac{4}{5}\cdot\ln\dfrac{8076}{7}+\dfrac{3}{5}\cdot\ln\dfrac{6057}{7}$ khi $a=\dfrac{6057}{7}$.\\
		Từ đó suy ra $m=6057$; $n=8076\Rightarrow S=m+2n=6057+2\cdot 8076=22209$.
	}
\end{ex}
\begin{ex}
	[Bỉm Sơn - Thanh Hóa - 2019]%Câu 27.
	Xét các số nguyên dương $a, b$ sao cho phương trình $a\ln^2x+b\ln x+5=0$ có hai nghiệm phân biệt $x_1, x_2$ và phương trình $5\log^2x+b\log x+a=0$ có hai nghiệm phân biệt $x_3, x_4$ thỏa mãn $x_1x_2>x_3x_4$. Tìm giá trị nhỏ nhất của $S=2a+3b$ 
	\choice
	{$S_{\min} =33$}
	{\True $S_{\min} =30$}
	{$S_{\min} =17$}
	{$S_{\min} =25$}
	\loigiai{
		Điều kiện để hai phương trình $a\ln^2x+b\ln x+5=0$ và $5\log^2x+b\log x+a=0$ có hai nghiệm phân biệt là $b^2-20a>0$.\\
		Theo giả thiết ta có $\heva{&\ln x_1+\ln x_2=-\dfrac{b}{a}\\&\log x_3+\log x_4=-\dfrac{b}{5}}\Rightarrow\heva{&\ln(x_1x_2)=-\dfrac{b}{a}\\&\log(x_3x_4)=-\dfrac{b}{5}}\Rightarrow\heva{&x_1x_2=\mathrm{e}^{-\tfrac{b}{a}}\\&x_3x_4={10}^{-\tfrac{b}{5}}.}$ \\
		Mà $x_1x_2>x_3x_4\Rightarrow\mathrm{e}^{-\tfrac{b}{a}}>10^{-\tfrac{b}{5}}$ \\
		$ \Rightarrow-\dfrac{b}{a} >-\dfrac{b}{5}\ln 10 $ \\
		$ \Rightarrow a>\dfrac{5}{\ln 10}\Rightarrow a\geq 3 $.\\
		Theo điều kiện có $b^2-20a>0\Rightarrow b^2>20a\geq 60\Rightarrow b\geq 8$.\\
		Từ và suy ra $S=2a+3b\geq 30\Rightarrow S_{\min} =30\Leftrightarrow\heva{&a=3\\&b=8}$.
	}
\end{ex}
\begin{ex}
	[THPT Thuận Thành 3 - Bắc Ninh 2019]%Câu 28.
	Có bao nhiêu giá trị nguyên dương của tham số $m$ để phương trình $\log_2\left(\dfrac{\sqrt{2x^2+mx+1}}{x+2}\right)+\sqrt{2x^2+mx+1}=x+2$ có hai nghiệm phân biệt?
	\choice
	{$3$}
	{$1$}
	{\True $4$}
	{$2$}
	\loigiai{
		Điều kiện: $\heva{&2x^2+mx+1>0\\&x+2>0.}$ \\
		$\log_2\left(\dfrac{\sqrt{2x^2+mx+1}}{x+2}\right)+\sqrt{2x^2+mx+1}=x+2\Leftrightarrow\log_2\sqrt{2x^2+mx+1}+\sqrt{2x^2+mx+1}=x+2+\log_2(x+2)$.\\
		Xét hàm số $f(t)=\log_2t+t$ trên khoảng $(0;+\infty)$,\\
		có $f'(t)=\dfrac{1}{t\ln 2}+1>0,\forall t\in(0;+\infty)\Rightarrow$ hàm số $f(t)$ đồng biến trên $(0;+\infty)$.\\
		Mà $f\left(\sqrt{2x^2+mx+1}\right)=f(x+2)\Leftrightarrow\sqrt{2x^2+mx+1}=x+2\Leftrightarrow\heva{&x >-2\\&x^2+(m-4)x-3=0.}$ \\
		Do $f(x)=x^2+(m-4)x-3$ là tam thức bậc hai nên có bảng biến thiên
		\begin{center}
			\begin{tikzpicture}
				\tkzTabInit[nocadre,lgt=1.2,espcl=3]{$x$/1.0,$f'(x)$/0.8,$f(x)$/2.5}{$-\infty$,$\dfrac{4-m}{2}$,$+\infty$}
				\tkzTabLine{,-,0,+,}
				\tkzTabVar{+/$+\infty$,-/$f\left(\dfrac{4-m}{2}\right)$,+/$+\infty$}
			\end{tikzpicture}
		\end{center}
		Phương trình đã cho có hai nghiệm phân biệt khi và chỉ khi phương trình $f(x)=x^2+(m-4)x-3=0$ có hai nghiệm phân biệt lớn hơn $-2$.\\
		suy ra: $\heva{&-2<\dfrac{4-m}{2}\\&f\left(\dfrac{4-m}{2}\right)<0<9-2m}\Leftrightarrow\heva{&m<8\\&m<\dfrac{9}{2}\\&\left(\dfrac{4-m}{2}\right)^2-3<0}\Leftrightarrow m<\dfrac{9}{2.}$ \\
		Do $m\in\mathbb{N}^*\Rightarrow m\in\{1;2;3;4\}$. Vậy có 4 giá trị của $m$.
	}
\end{ex}
\begin{ex}
	[Chuyên Bắc Giang 2019]%Câu 29.
	Số các giá trị nguyên nhỏ hơn $2018$ của tham số $m$ để phương trình $\log_6(2018x+m)=\log_4(1009x)$ có nghiệm là
	\choice
	{$2018$}
	{$2017$}
	{\True $2020$}
	{$2019$}
	\loigiai{
		Đặt $\log_4(1009x)=t\Rightarrow 1009x=4^t$.\\
		Phương trình đã cho có dạng $\log_6\left({2\cdot 4}^t+m\right)=t\Leftrightarrow 2\cdot 4^t+m=6^t\Leftrightarrow m=6^t-2\cdot 4^t$.\\
		Số nghiệm của phương trình bằng số giao điểm của đồ thị hàm số $f(t)=6^t-2\cdot 4^t$ với đường thẳng $y=m$.\\
		Xét hàm số: $f(t)=6^t-2\cdot 4^t f'(t)=6^t\ln 6-2\cdot 4^t\ln 4=2^t\left(3^t\ln 6-2\cdot 2^t\ln 4\right)$.\\
		$f'(t)=0\Leftrightarrow 6^t\ln 6=2\cdot 4^t\ln 4\Leftrightarrow\left(\dfrac{3}{2}\right)^t=4\log_62\Leftrightarrow t=\log_{\tfrac{3}{2}}(4\log_62)$.\\
		$\lim\limits_{t\to+\infty} f(t)=\lim\limits_{t\to+\infty}\left(6^t-2\cdot 4^t\right)=\lim\limits_{t\to+\infty} 6^t\left(1-2\cdot\left(\dfrac{2}{3}\right)^t\right)=+\infty$.\\
		$\lim\limits_{t\to-\infty} f(t)=\lim\limits_{t\to-\infty}\left(6^t-2\cdot 4^t\right)=0$.\\
		Ta có bảng biến thiên:
		\begin{center}
			\begin{tikzpicture}
				\tkzTabInit[nocadre,lgt=1.2,espcl=3]{$t$/1.0,$f'(t)$/0.8,$f(t)$/2.5}{$-\infty$,$\log_{\frac{3}{2}}(4\log_62)$,$+\infty$}
				\tkzTabLine{,-,0,+,}
				\tkzTabVar{+/$0$,-/$f\left(\log_{\frac{3}{2}}(4\log_62)\right)$,+/$+\infty$}
			\end{tikzpicture}
		\end{center}
		Với $f\left(\log_{\tfrac{3}{2}}(4\log_62)\right)\approx-2,0136$.\\
		Từ bảng biến thiên, để phương trình có nghiệm thì $m\geq f\left(\log_{\tfrac{3}{2}}(4\log_62)\right)\approx-2,0136$.\\
		Vậy $-2\leq m<2018$. Có $2020$ số nguyên $m$.
	}
\end{ex}
\begin{ex}
	[Chuyên ĐH Vinh - Nghệ An -2020]%Câu 30.
	Có bao nhiêu số nguyên $m$ để phương trình $\log_3\left(3^x+2m\right)=\log_5\left(3^x-m^2\right)$ có nghiệm?
	\choice
	{\True $3$}
	{$4$}
	{$2$}
	{$5$}
	\loigiai{
		Đặt $\log_3\left(3^x+2m\right)=\log_5\left(3^x-m^2\right)=t\Rightarrow\heva{&3^x+2m=3^t\\&3^x-m^2=5^t}$ \\
		$ \Rightarrow 2m+m^2=3^t-5^t\Leftrightarrow m^2+2m+1=3^t-5^t+1 $ (*).\\
		Xét hàm số $f(t)=3^t-5^t+1$ với $t\in\mathbb{R}$.\\
		Ta có $f'(t)=3^t\cdot\ln 3-5^t\cdot\ln 5$.\\
		Khi đó $f'(t)=0\Leftrightarrow 3^t\cdot\ln 3-5^t\cdot\ln 5=0\Leftrightarrow\left(\dfrac{3}{5}\right)^t=\dfrac{\ln 5}{\ln 3}\Leftrightarrow t=\log_{\tfrac{3}{5}}(\log_35)=t_0$.\\
		Bảng biến thiên\\
		{\color{red} HINH O DAY}\\
		Phương trình (*) có nghiệm\\
		$ \Leftrightarrow(m+1)^2\leq f(t_0)\Leftrightarrow-\sqrt{f(t_0)}-1\leq m\leq\sqrt{f(t_0)}-1\Rightarrow-2,068\leq m\leq 0,068 $.\\
		Do $m\in\mathbb{Z}\Rightarrow m\in\{-2;-1; 0\}$.\\
		Vậy có $3$ giá trị nguyên của $m$ thỏa mãn.
	}
\end{ex}     
\begin{ex}
	[Chuyên Hưng Yên - 2020]%Câu 31.
	Tìm tất cả các giá trị thực của tham số $m$ để phương trình $\log_3x+\sqrt{\log_3x+1}-2m-1=0$ có ít nhất một nghiệm thực trong đoạn $[1;27]$. 
	\choice
	{$m\in(0;2)$}
	{\True $m\in[0;2]$}
	{$m\in[2;4]$}
	{$m\in(0;4)$}
	\loigiai{
		Đặt $t=\sqrt{\log_3x+1}$. Với $x\in[1;27]$ thì $t\in[1;2]$.\\
		Phương trình đã cho trở thành $t^2+t-2m-2=0\Leftrightarrow 2m+2=t^2+t (*)$.\\
		Xét hàm số $f(t)=t^2+t$ trên đoạn $[1;2]$.\\
		Ta có $f'(t)=2t+1>0,\forall t\in[1;2]$ nên hàm số $f(t)=t^2+t$ đồng biến trên $[1;2]$.\\
		Bảng biến thiên: 
		\begin{center}
			\begin{tikzpicture}
				\tkzTabInit[nocadre,lgt=1.2,espcl=3]{$t$/0.8,$f'(t)$/0.8,$f(t)$/2.5}{$1$,{},$2$}
				\tkzTabLine{,+,}
				\tkzTabVar{-/$2$,R,+/$6$}
			\end{tikzpicture}
		\end{center}
		Để phương trình đã cho có ít nhất một nghiệm thực trong đoạn $[1;27]$ thì phương trình $(*)$ phải có ít nhất một nghiệm thực trong đoạn $[1;2]$.\\
		Từ bảng biến thiên, suy ra $2\leq 2m+2\leq 6\Leftrightarrow 0\leq m\leq 2$.
	}
\end{ex}
\begin{ex}
	[Chuyên KHTN - 2020]%Câu 32.
	Có bao nhiêu giá trị nguyên của tham số $m$ để phương trình $\log_3^2x-m\log_9x^2+2-m=0$ có nghiệm $x\in[1;9]$. 
	\choice
	{\True $1$}
	{$5$}
	{$3$}
	{$2$}
	\loigiai{
		Điều kiện: $x>0$.\\
		Ta có $\log_3^2x-m\log_9x^2+2-m=0\Leftrightarrow\log_3^2x-m\log_3x+2-m=0$.\\
		Đặt $t=\log_3x$, với $x\in[1;9]\Rightarrow t\in[0;2]$.\\
		Phương trình đã cho trở thành: $t^2-mt+2-m=0\Leftrightarrow m=\dfrac{t^2+2}{t+1}(1)$ (Do $t\neq-1,\forall t\in[0;2]$).\\
		Xét hàm số $f(t)=\dfrac{t^2+2}{t+1}$ với $t\in[0;2]$ ta có\\
		$f'(t)=\dfrac{t^2+2t-2}{(t+1)^2}$, $f'(t)=0\Leftrightarrow t^2+2t-2=0\Leftrightarrow\hoac{&t=-1+\sqrt{3}\in[0;2]\\&t=-1-\sqrt{3}\notin[0;2].}$ \\
		Bảng biến thiên: 
		\begin{center}
			\begin{tikzpicture}
				\tkzTabInit[nocadre,lgt=1.2,espcl=3]{$t$/0.8,$f'(t)$/0.8,$f(t)$/2.5}{$0$,$-1+\sqrt{3}$,$2$}
				\tkzTabLine{,-,0,+,}
				\tkzTabVar{+/$2$,-/$-2+2\sqrt{3}$,+/$2$}
			\end{tikzpicture}
		\end{center}
		Khi đó: phương trình đã cho có nghiệm $x\in[1;9]\Leftrightarrow$ Phương trình $(1)$ có nghiệm $t\in[0;2]$ \\
		$ \Leftrightarrow-2+2\sqrt{3}\leq m\leq 2 $.\\
		Mặt khác, do $m\in\mathbb{Z}$ nên $m=2$.\\
		Vậy có một giá trị nguyên của tham số $m$ thỏa yêu cầu bài toán.
	}
\end{ex}
\begin{ex}
	[Chuyên KHTN - 2020]%Câu 33.
	Có bao nhiêu giá trị nguyên của $m$ để phương trình $\log_2(mx)=\log_{\sqrt{2}}(x+1)$ vô nghiệm?
	\choice
	{\True $4$}
	{$6$}
	{$3$}
	{$5$}
	\loigiai{
		Điều kiện $\heva{&mx>0\\&x+1>0}\Leftrightarrow\heva{&mx>0\\&x >-1.}$ \\
		Ta có $\log_2(mx)=\log_{\sqrt{2}}(x+1)\Leftrightarrow\log_2(mx)=2\log_2(x+1)$ \\
		$ \Leftrightarrow\log_2(mx)=\log_2(x+1)^2\Leftrightarrow\heva{&x+1>0\\&mx=(x+1)^2}\Leftrightarrow\heva{&x >-1\\&mx=(x+1)^2\quad(1).} $ \\
		Nhận xét với $x=0$ không là nghiệm của phương trình (1).\\
		Với $x\neq 0$ thì $(1)\Leftrightarrow m=\dfrac{(x+1)^2}{x}$.\\
		Xét hàm số $f(x)=\dfrac{(x+1)^2}{x}$ với $x\in(-1;+\infty)\setminus\{0\}$.\\
		có $f'(x)=\dfrac{x^2-1}{x^2}\Rightarrow f'(x)=0\Leftrightarrow x=\pm 1$.\\
		Bảng biến thiên\\
		{\color{red} HINH O DAY}\\
		Phương trình đã cho vô nghiệm khi và chỉ khi $0\leq m<4$. Do $m\in\mathbb{Z}$ nên $m\in\{0;1;2;3\}$.\\
		Vậy có 4 giá trị nguyên của tham số $m$ để phương trình $\log_2(mx)=\log_{\sqrt{2}}(x+1)$ vô nghiệm.
	}
\end{ex}
\begin{ex}
	[Chuyên Lương Văn Chánh - Phú Yên - 2020]%Câu 34.
	Số các giá trị nguyên nhỏ hơn 2020 của tham số $m$ để phương trình $\log_6(2020x+m)=\log_4(1010x)$ có nghiệm là
	\choice
	{$2020$}
	{$2021$}
	{$2019$}
	{\True $2022$}
	\loigiai{
		Điều kiện xác định: $\heva{&2020x+m>0\\&1010x>0}$ (*).\\
		Đặt $\log_6(2020x+m)=\log_4(1010x)=t$.\\
		Suy ra $\heva{&2020x+m=6^t\\&1010x=4^t}(1)$.\\
		Từ đó $m=6^t-2\cdot 4^t\quad(2)$.\\
		Với mỗi nghiệm $t_0$ của phương trình $(2)$ thì $x_0=\dfrac{4^{t_0}}{2010}$ là nghiệm của hệ phương trình $(1)$ đồng thời $x_0$ thỏa mãn điều kiện $(*)$. Do đó $x_0$ là nghiệm của phương trình đã cho. Từ đó, điều kiện cần và đủ để phương trình đã cho có nghiệm là phương trình $(2)$ có nghiệm.\\
		Xét hàm số $f(t)=6^t-2\cdot 4^t$ trên $\mathbb{R}$.\\
		Ta có $f'(t)=6^t\cdot\ln 6-2\cdot 4^t\cdot\ln 4$ và $f'(t)=0\Leftrightarrow t=\log_{\tfrac{3}{2}}(\log_616)\colon =\alpha$.\\
		Bảng biến thiên của hàm số $f(t)$ như sau: 
		\begin{center}
			\begin{tikzpicture}
				\tkzTabInit[nocadre,lgt=1.2,espcl=3]{$t$/0.8,$f'(t)$/0.8,$f(t)$/2.5}{$-\infty$,$\alpha$,$+\infty$}
				\tkzTabLine{,-,0,+,}
				\tkzTabVar{+/$0$,-/$-2{,}014$,+/$+\infty$}
			\end{tikzpicture}
		\end{center}
		Dựa vào bảng biến thiên, ta có phương trình $(2)$ có nghiệm khi và chỉ khi $m\geq-2(do m\in\mathbb{Z})$.\\
		Vậy tất cả các giá trị nguyên của tham số $m$ thỏa yêu cầu bài toán là các số nguyên thuộc tập hợp $\left\{-2\cdot -1,0,1,2,\ldots\cdot,2019\right\}$, có tất cả 2022 giá trị.
	}
\end{ex}
\begin{ex}
	[Chuyên Quang Trung - 2020]%Câu 35.
	Xét các số nguyên dương $a, b$ sao cho phương trình $a\ln^2x+b\ln x+5=0$ có hai nghiệm phân biệt $x_1$, $x_2$ và phương trình $5\log^2x+b\log x+a=0$ có hai nghiệm phân biệt $x_3$, $x_4$ sao cho $x_1x_2>x_3x_4$. Tìm giá trị nhỏ nhất của $S=2a+3b$. 
	\choice
	{\True $30$}
	{$25$}
	{$33$}
	{$17$}
	\loigiai{
		$a\ln^2x+b\ln x+5=0(1)$.\\
		$5\log^2x+b\log x+a=0(2)$.\\
		Điều kiện để $(1)$ có hai nghiệm phân biệt $x_1, x_2$ và $(2)$ có hai nghiệm phân biệt $x_3, x_4$ là $b^2-20a>0\Leftrightarrow b^2>20a$.\\
		Nhận xét: $x_1, x_2, x_3, x_4>0$.\\
		Do đó: $x_1x_2>x_3x_4\Leftrightarrow\ln(x_1x_2)>\ln(x_3x_4)\Leftrightarrow\ln(x_1x_2)>\dfrac{\log(x_3x_4)}{\log e}$ \\
		$ \Leftrightarrow\left(\ln x_1+\ln x_2\right)\log e>\log x_3+\log x_4 $.\\
		Mà $\ln x_1+\ln x_2=-\dfrac{b}{a};\log x_3+\log x_4=-\dfrac{b}{5}$ và $a, b$ nguyên dương.\\
		Nên $-\dfrac{b}{a}\log e >-\dfrac{b}{5}\Leftrightarrow a>5\log e$.\\
		Vì $a$ là số nguyên dương và $5\log e\approx 2,17$ nên $a\geq 3$ \\
		$ \Rightarrow 20a\geq 60\Rightarrow b^2>60\Rightarrow b>\sqrt{60} (b>0) $.\\
		Vì $b$ là số nguyên dương và $\sqrt{60}\approx 7,75$ nên $b\geq 8$.\\
		Do đó: $S=2a+3b\geq 30\Rightarrow$ Giá trị nhỏ nhất của $S$ là 30 khi $a=3$; $b=8$.
	}
\end{ex}
\begin{ex}
	[Chuyên Thái Bình - 2020]%Câu 36.
	Cho phương trình $\log_2^2x-(5m+1)\log_2x+4m^2+m=0$. Biết phương trình có 2 nghiệm phân biệt $x_1, x_2$ thỏa $x_1+x_2=165$. Giá trị của $|x_1-x_2|$ bằng
	\choice
	{$16$}
	{$119$}
	{$120$}
	{\True $159$}
	\loigiai{
		$\begin{aligned}&\log_2^2x-(5m+1)\log_2x+4m^2+m=0\\&\Leftrightarrow\hoac{&\log_2x=m\\&\log_2x=4m+1}.\end{aligned}$ \\
		Để phương trình có 2 nghiệm phân biệt khi $m\neq 4m+1\Leftrightarrow m\neq\dfrac{-1}{3}$.\\
		Khi đó phương trình có 2 nghiệm $x_1=2^m>0, x_2=2^{4m+1}=2\cdot (2^m)^4>0$.\\
		Vì $x_1+x_2=165\Leftrightarrow 2^m+2\cdot (2^m)^4=165(*)$.\\
		Xét hàm số $f(t)=2\cdot t^4+t\Rightarrow f'(t)=8t^3+1>0\forall t>0$.\\
		Mà $2^m=3$ là nghiệm của $(*)$ nên là nghiệm duy nhất. Suy ra $x_1=3, x_2=2\cdot 3^4=162$.\\
		Suy ra $|x_1-x_2|=159$.
	}
\end{ex}
\begin{ex}
	[Chuyên Thái Nguyên - 2020]%Câu 37.
	Gọi $m_0$ là giá trị thực nhỏ nhất của tham số $m$ sao cho phương trình $(m-1)\log_{\tfrac{1}{3}}^2(x-3)-(m-5)\log_{\tfrac{1}{3}}(x-3)+m-1=0$ có nghiệm thuộc $(3;6)$. Khẳng định nào sau đây là đúng?
	\choice
	{Không tồn tại $m_0$}
	{$m_0\in\left(-1;\dfrac{4}{3}\right)$}
	{$m_0\in\left(2;\dfrac{10}{3}\right)$}
	{\True $m_0\in\left(-5;\dfrac{-5}{2}\right)$}
	\loigiai{
		Đặt $t=\log_{\tfrac{1}{3}}(x-3)$.\\
		Vì $x\in(3;6)\Rightarrow t >-1$.\\
		Phương trình trở thành: $(m-1)t^2-(m-5)t+m-1=0$ (*)\\
		$ \Leftrightarrow mt^2-mt+m=t^2-5t+1 $ \\
		$ \Leftrightarrow m=\dfrac{t^2-5t+1}{t^2-t+1} $.\\
		Xét hàm số $f(t)=\dfrac{t^2-5t+1}{t^2-t+1}$.\\
		$f'(t)=\dfrac{(2t-5)\left(t^2-t+1\right)-(2t-1)\left(t^2-5t+1\right)}{\left(t^2-t+1\right)^2}=\dfrac{4t^2-4}{\left(t^2-t+1\right)^2}$.\\
		$f'(t)=0\Leftrightarrow t=\pm 1$.\\
		Bảng biến thiên: 
		\begin{center}
			\begin{tikzpicture}
				\tkzTabInit[nocadre,lgt=1.2,espcl=3]{$t$/0.8,$f'(t)$/0.8,$f(t)$/2.5}{$-1$,$1$,$+\infty$}
				\tkzTabLine{,-,0,+,}
				\tkzTabVar{+/$\dfrac{7}{3}$,-/$-3$,+/$+\infty$}
			\end{tikzpicture}
		\end{center}
		Để phương trình đã cho có nghiệm $x\in(3;6)$ thì phương trình $(*)$ có nghiệm $t >-1$ \\
		$ \Leftrightarrow m\geq-3 $.\\
		Vậy giá trị nhỏ nhất của $m$ thỏa mãn yêu cầu bài toán là $m_0=-3\in\left(-5;\dfrac{-5}{2}\right)$.
	}
\end{ex}
\begin{ex}
	[Chuyên Vĩnh Phúc - 2020]%Câu 38.
	Cho phương trình $m\ln(x+1)-x-2=0$. Biết rằng tập hợp tất cả các giá trị của tham số $m$ để phương trình đã cho có hai nghiệm $x_1$, $x_2$ thỏa mãn $0<x_1<2<4<x_2$ là khoảng $(a;+\infty)$. Khi đó $a$ thuộc khoảng nào dưới đây?
	\choice
	{\True $(3,7;3,8)$}
	{$(3,6;3,7)$}
	{$(3,8;3,9)$}
	{$(3,5;3,6)$}
	\loigiai{
		Xét trên khoảng $(0;+\infty)$ phương trình: $m\ln(x+1)-x-2=0\Leftrightarrow m=\dfrac{x+2}{\ln(x+1)}$.\\
		Đặt $f(x)=\dfrac{x+2}{\ln(x+1)}, x\in(-1;+\infty)\setminus\{0\}$.\\
		Với yêu cầu của đề bài ta xét $f(x)$ trên $2$ khoảng $(0;2)$ và $(4;+\infty)$.\\
		$f'(x)=\dfrac{\ln(x+1)-(x+2)\dfrac{1}{x+1}}{\ln^2(x+1)}$.\\
		Đặt $g(x)=\ln(x+1)-(x+2)\dfrac{1}{x+1},x\in(0;2)\cup(4;+\infty)$.\\
		$g'(x)=\dfrac{1}{x+1}+\dfrac{1}{(x+1)^2}>0,\forall x\in(0;2)\cup(4;+\infty)$.\\
		Suy ra $\heva{&g(x)<g(2)=\ln 3-\dfrac{4}{3}<0,\forall x\in(0;2)\Rightarrow f'(x)<0,\forall x\in(0;2)\\&g(x)>g(5)=\ln 5-\dfrac{6}{5}>0,\forall x\in(4;+\infty)\Rightarrow f'(x)>0,\forall x\in(4;+\infty).}$ \\
		Từ đó ta có bảng biến thiên\\
		{\color{red} HINH O DAY}\\
		Dựa vào bảng biến thiên, để phương trình đề bài có $2$ nghiệm phân biệt thỏa $0<x_1<2<4<x_2\Leftrightarrow m>\dfrac{6}{\ln 5}\left(\approx 3,728\right)$.
	}
\end{ex}
\begin{ex}[Đại Học Hà Tĩnh - 2020]%Câu 39.
	Tìm tất cả các giá trị của tham số $a$ để phương trình $\log_3x^2+a\sqrt{\log_3x^3}+a+1=0$ có nghiệm duy nhất. 
	\choice
	{Không tồn tại $a$}
	{\True $a <-1$ hoặc $a=4-2\sqrt{10}$}
	{$a<1$}
	{$a=1$}
	\loigiai{
		Điều kiện: $\heva{&x>0\\&\log_3x^3\geq 0}\Leftrightarrow\heva{&x>0\\&x\geq 1}\Leftrightarrow x\geq 1$.\\
		Khi đó phương trình $\log_3x^2+a\sqrt{\log_3x^3}+a+1=0\Leftrightarrow 2\log_3x+a\sqrt{3\log_3x}+a+1=0$ \\
		$ \Leftrightarrow 2\cdot (3\log_3x)+3a\sqrt{3\log_3x}+3a+3=0(1) $.\\
		Đặt $\sqrt{3\log_3x}=t,t\geq 0$ thì $(1)$ trở thành: $2t^2+3at+3a+3=0$.\\
		Do đó, yêu cầu bài toán trở thành: Tìm tất cả các giá trị của tham số $a$ để phương trình.\\
		$2t^2+3at+3a+3=0$ có nghiệm duy nhất thuộc nửa khoảng $[0;+\infty)$.\\
		Ta có $2t^2+3at+3a+3=0\Leftrightarrow 3a=\dfrac{-2t^2-3}{t+1},t\geq 0$.\\
		Xét hàm số: $f(t)=\dfrac{-2t^2-3}{t+1}$ trên nửa khoảng $[0;+\infty)$. Ta có\\
		$f'(t)=\dfrac{-2t^2-4t+3}{(t+1)^2}$. $f'(t)=0\Leftrightarrow-2t^2-4t+3=0\Leftrightarrow\hoac{&t=\dfrac{-2+\sqrt{10}}{2}\\&t=\dfrac{-2-\sqrt{10}}{2}}\Leftrightarrow t=\dfrac{-2+\sqrt{10}}{2.}$ \\
		$\lim\limits_{t\to+\infty} f(t)=-\infty$.\\
		Bảng biến thiên: 
		\begin{center}
			\begin{tikzpicture}
				\tkzTabInit[nocadre,lgt=1.2,espcl=3]{$t$/1.0,$f'(t)$/0.8,$f(t)$/2.5}{$0$,$\dfrac{-2+\sqrt{10}}{2}$,$+\infty$}
				\tkzTabLine{,+,0,-,}
				\tkzTabVar{-/$-3$,+/$4-2\sqrt{10}$,-/$-\infty$}
			\end{tikzpicture}
		\end{center}
		Từ bảng biến thiên suy ra phương trình có một nghiệm duy nhất khi\\
		$\hoac{&3a <-3\\&a=-4+2\sqrt{10}}\Leftrightarrow\hoac{&a <-1\\&a=-4+2\sqrt{10}.}$ \\
		Đáp số: $a <-1$ hoặc $a=-4+2\sqrt{10}$.
	}
\end{ex}
\begin{ex}
	[Sở Ninh Bình 2020]%Câu 40.
	Gọi $m_0$ là giá trị nhỏ nhất của tham số thực $m$ sao cho phương trình $(m-1)\log_{\frac{1}{2}}^2(x-2)-(m-5)\log_{\frac{1}{2}}(x-2)+m-1=0$ có nghiệm thuộc khoảng $(2;4)$. Khẳng định nào dưới đây đúng?
	\choice
	{$m_0\in\left(-1;\dfrac{4}{3}\right)$}
	{$m_0\in\left(2;\dfrac{10}{3}\right)$}
	{$m_0\in\left(4;\dfrac{16}{3}\right)$}
	{\True $m_0\in\left(-5;\dfrac{-5}{2}\right)$}
	\loigiai{
		Điều kiện: $x>2$.\\
		Đặt $t=\log_{\tfrac{1}{2}}(x-2)$, với $x\in(2;4)\Rightarrow t\in(-1;+\infty)$.\\
		Phương trình đã cho trở thành: $(m-1)t^2-(m-5)t+m-1=0$.\\
		$t^2-5t+1=mt^2-mt+m\Leftrightarrow t^2-5t+1=m\left(t^2-t+1\right)$ \\
		$ \Leftrightarrow\dfrac{t^2-5t+1}{t^2-t+1}=m(1),\left(t^2-t+1\neq 0,\forall t\right) $.\\
		Phương trình đã cho có nghiệm $\Leftrightarrow(1)$ có nghiệm $t >-1$.\\
		Xét hàm số $f(t)=\dfrac{t^2-5t+1}{t^2-t+1},(t >-1)$.\\
		Ta có $f'(t)=\dfrac{4t^2-4}{\left(t^2-t+1\right)^2}\Rightarrow f'(t)=0\Leftrightarrow 4t^2-4=0\Leftrightarrow\hoac{&x=1\Rightarrow f(1)=-3\\&x=-1\Rightarrow f(-1)=-\dfrac{7}{3}.}$ \\
		Bảng biên thiên hàm số $f(t)$:\\ 
		{\color{red} HINH O DAY}\\
		Dựa vào bảng biến thiên, phương trình $f(t)=m$ có nghiệm $t >-1$ khi và chỉ khi $-3\leq m<\dfrac{7}{3}$.\\
		Vậy giá trị nhỏ nhất của tham số thực $m$ để phương trình đã cho có nghiệm là $m_0=-3\in\left(-5;\dfrac{-5}{2}\right)$.
	}
\end{ex}
\begin{ex}
	[Sở Yên Bái - 2020]%Câu 41.
	Giả sử phương trình $\log_2^2x-(m+2)\log_2x+2m=0$ có hai nghiệm thực phân biệt $x_1,x_2$ thỏa mãn $x_1+x_2=6$. Giá trị biểu thức $|x_1-x_2|$ là
	\choice
	{$4$}
	{$3$}
	{$8$}     
	{\True $2$}
	\loigiai{
		Điều kiện $x>0$. Phương trình đã cho tương đương.\\
		$\begin{aligned}&\log_2^2x-m\log_2x-2\log_2x+2m=0\\&\Leftrightarrow (\log_2x-m)(\log_2x-2)=0\Leftrightarrow\hoac{&\log_2x=m\\&\log_2x=2}\Rightarrow\hoac{&x=2^m\\&x=4}.\end{aligned}$ \\
		Theo giả thiết $x_1+x_2=6\Rightarrow 2^m+4=6\Rightarrow m=1\Rightarrow x=2\Rightarrow|x_1-x_2|=2$.
	}
\end{ex}
\begin{ex}
	[Bỉm Sơn - Thanh Hóa - 2020]%Câu 42.
	Tìm tất cả các giá trị của $m$ để phương trình $\log_{_2}^2x-\log_2x^2+3=m$ có nghiệm $x\in[1; 8]$. 
	\choice
	{\True $2\leq m\leq 6$}
	{$3\leq m\leq 6$}
	{$6\leq m\leq 9$}
	{$2\leq m\leq 3$}
	\loigiai{
		Đặt $t=\log_2x$. Khi $x\in[1; 8]$ thì $t\in[0; 3]$. Câu toán trở thành: Tìm $m$ để phương trình $t^2-2t+3=m$ có nghiệm $t\in[0; 3]$. Xét hàm số $f(t)=t^2-2t+3$ với $t\in[0; 3]$, ta có\\
		$f'(t)=2t-2=0\Leftrightarrow t=1$; $\min\limits_{t\in[0; 3]} f(t)=f(1)=2$; $\max\limits_{t\in[0; 3]} f(t)=f(3)=6$.\\
		Đồ thị hàm số $y=f(t)=t^2-2t+3$ và đường thẳng $y=m$ sẽ cắt nhau tại điểm có hoành độ $t\in[0; 3]$ nếu như $\min\limits_{t\in[0; 3]} f(t)\leq m\leq\max\limits_{t\in[0; 3]} f(t)\Leftrightarrow 2\leq m\leq 6$.
	}
\end{ex}
\begin{ex}
	[Đô Lương 4 - Nghệ An - 2020]%Câu 43.
	Tìm các giá trị thực của tham số $m$ để phương trình $\log_3^2x-3\log_3x+2m-7=0$ có hai nghiệm thực $x_1$, $x_2$ thỏa mãn $(x_1+3)(x_2+3)=72$. 
	\choice
	{\True $m=\dfrac{9}{2}$}
	{$m=3$}
	{Không tồn tại}
	{$m=\dfrac{61}{2}$}
	\loigiai{
		Đặt $t=\log_3x$.\\
		Phương trình đã cho trở thành $t^2-3t+2m-7=0(*)$.\\
		Ứng với mỗi nghiệm $t$ của phương trình $(*)$ có một nghiệm $x$.\\
		Vậy phương trình đã cho có hai nghiệm phân biệt khi phương trình $(*)$ có hai nghiệm phân biệt\\
		$ \Leftrightarrow\Delta>0\Leftrightarrow(-3)^2-4(2m-7)>0\Leftrightarrow 9-8m+28>0\Leftrightarrow m<\dfrac{37}{8} $.\\
		Gọi $t_1$, $t_2$ là hai nghiệm phương trình $(*)$.\\
		Theo định lý Viét ta có $t_1+t_2=3\Rightarrow\log_3x_1+\log_3x_2=3\Leftrightarrow\log_3(x_1\cdot x_2)=3\Leftrightarrow x_1\cdot x_2=27$.\\
		Theo đề bài $(x_1+3)(x_2+3)=72\Leftrightarrow x_1\cdot x_2+3(x_1+x_2)+9=72\Leftrightarrow x_1+x_2=12$.\\
		Vậy ta có $\heva{&x_1+x_2=12\\&x_1\cdot x_2=27}\Leftrightarrow\heva{&x_1=9\\&x_2=3}\Rightarrow\heva{&t_1=2\\&t_2=1}\Rightarrow t_1\cdot t_2=2$.\\
		Theo định lý Viét ta có $t_1\cdot t_2=2\Leftrightarrow 2=2m-7\Leftrightarrow m=\dfrac{9}{2}$ (thỏa mãn).\\
		Kết luận: $m=\dfrac{9}{2}$ thỏa mãn yêu cầu bài toán.
	}
\end{ex}
\begin{ex}
	[Đô Lương 4 - Nghệ An - 2020]%Câu 44.
	Số các giá trị nguyên nhỏ hơn $2020$ của tham số $m$ để phương trình $\log_6(2020x+m)=\log_4(1010x)$ có nghiệm là
	\choice
	{\True $2022$}
	{$2020$}
	{$2019$}
	{$2021$}
	\loigiai{
		Ta đặt $\log_6(2020x+m)=\log_4(1010x) =t$. Khi đó\\
		$2020x+m=6^t$ và $1010x=4^t$. Ta suy ra $2.4^t+m=6^t\Leftrightarrow m=6^t-2\cdot 4^t$.\\
		Đặt $f(t)=-2\cdot 4^t+6^t$.\\
		$f'(t)=6^t\ln 6-2\cdot 4^t\cdot\ln 4$.\\
		$f'(t)=0\Rightarrow\left(\dfrac{3}{2}\right)^t=\dfrac{2\ln 4}{\ln 6}=\log_616\Leftrightarrow t=\log_{\tfrac{3}{2}}(\log_616)$.\\
		Bảng biến thiên
		\begin{center}
			\begin{tikzpicture}
				\tkzTabInit[nocadre,lgt=1.2,espcl=3]{$t$/0.8,$f'(t)$/0.8,$f(t)$/2.5}{$-\infty$,$\log_{\frac{3}{2}}(\log_{6}16)$,$+\infty$}
				\tkzTabLine{,-,0,+,}
				\tkzTabVar{+/$+\infty$,-/$f\left(\log_{\frac{3}{2}}(\log_{6}16)\right)$,+/$+\infty$}
			\end{tikzpicture}
		\end{center}
		Phương trình $f(t)=m$ có nghiệm khi và chỉ khi $m\geq f\left(\log_{\tfrac{3}{2}}(\log_616)\right)\approx-2,01$.\\
		Hơn nữa, $\heva{&m<2020\\&m\in\mathbb{Z}}$ nên suy ra $\heva{&-2\leq m\leq 2019\\&m\in\mathbb{Z}.}$ \\
		Vậy ta có $2022$ giá trị $m$ thỏa mãn.
	}
\end{ex}
\begin{ex}
	[Hậu Lộc 2 - Thanh Hóa - 2020]%Câu 45.
	Cho phương trình\\ $\left(m\mathrm{e}^x-10x-m\right)\left[\log(mx)-2\log(x+1)\right]=0$. ($m$ là tham số). Có tất cả bao nhiêu giá trị nguyên của $m$ để phương trình đã cho có ba nghiệm thực phân biệt?
	\choice
	{Vô số}
	{$10$}
	{$11$}
	{\True $5$}
	\loigiai{
		$\left(m\mathrm{e}^x-10x-m\right)\left[\log(mx)-2\log(x+1)\right]=0$ \\
		$\Leftrightarrow\heva{&mx>0&(1)\\&x+1>0&(2)\\&\hoac{&m\mathrm{e}^x-10x-m=0&(3)\\&mx=(x+1)^2&(4)}}(*) $.\\
		$m=0$ thì pt vô nghiệm.\\
		$m>0$ thì hệ $(*)\Leftrightarrow\heva{&x>0\\&\hoac{&m=\dfrac{10x}{\mathrm{e}^x-1}\\&m=\dfrac{(x+1)^2}{x}}}$. (Vì $\mathrm{e}^x-1>\mathrm{e}^{\circ}-1\Leftrightarrow\mathrm{e}^x-1>0$).\\
		Xét $f(x)=\dfrac{10x}{\mathrm{e}^x-1}$ và $g(x)=\dfrac{(x+1)^2}{x}=x+\dfrac{1}{x}+2$.\\
		$f'(x)=\dfrac{10\left(\mathrm{e}^x-1\right)-\mathrm{e}^x\cdot 10x}{\left(\mathrm{e}^x-1\right)^2}=\dfrac{10\mathrm{e}^x(1-x)-10}{\left(\mathrm{e}^x-1\right)^2}$.\\
		Xét $u(x)=10\mathrm{e}^x(1-x)-10\Rightarrow u'(x)=-10\mathrm{e}^x+10\mathrm{e}^x(1-x)=-10x\mathrm{e}^x<0\forall x\in(0;+\infty)$.\\
		Suy ra: Hàm số $u(x)$ nghịch biến trong khoảng $(0;+\infty)\Rightarrow u(x)<u(0)=0$ \\
		$ \Rightarrow f'(x)<0\forall x\in(0;+\infty)\Rightarrow f(x) $ nghịch biến trong khoảng $(0;+\infty)$.\\
		$\lim\limits_{x\to 0^+} f(x)=10,\lim\limits_{x\to+\infty} f(x)=0$ \\
		{\color{red} HINH O DAY}\\
		$g'(x)=1-\dfrac{1}{x^2}=\dfrac{x^2-1}{x^2}=0\Leftrightarrow\hoac{&x=1\\&x=-1}$. \\
		{\color{red} HINH O DAY}\\
		Suy ra phương trình có ba nghiệm thực phân biệt $\Leftrightarrow 4<m<10$. Vì $m\in\mathbb{Z}\Rightarrow m=\{5;6;7;8;9\}$.\\
		$m<0$ thì hệ $(*)\Leftrightarrow\heva{&-1<x<0\\&\hoac{&m=f(x)\\&m=g(x)}}$ \\
		Tương tự ta có $f'(x)>0\forall x\in(-1;0),\lim\limits_{x\to-1^+} f(x)=\dfrac{-10}{\mathrm{e}^{-1}-1}=\dfrac{-10e}{1-\mathrm{e}}=\dfrac{10\mathrm{e}}{\mathrm{e}-1},\lim\limits_{x\to 0^-} f(x)=10$ 
		{\color{red} HINH O DAY}\\
		$g'(x)=1-\dfrac{1}{x^2}=\dfrac{x^2-1}{x^2}=0\Leftrightarrow\hoac{&x=1\\&x=-1}$. \\
		{\color{red} HINH O DAY}\\
		Suy ra phương trình có nhiều nhất 1 nghiệm thực phân biệt, không thỏa mãn yêu cầu bài toán.\\
		Vậy có $5$ giá trị $m$.
	}
\end{ex}
\begin{ex}
	[Liên trường Nghệ An - 2020]%Câu 46.
	Cho phương trình $4^{-|x-m|}\cdot\log_{\sqrt{2}}\left(x^2-2x+3\right)+2^{2x-x^2}\cdot\log_{\frac{1}{2}}(2|x-m|+2)=0$ với $m$ là tham số. Tổng tất cả các giá trị của tham số $m$ để phương trình đã cho có ba nghiệm phân biệt là
	\choice
	{$4$}
	{$1$}
	{$2$}
	{\True $3$}
	\loigiai{
		Có $4^{-|x-m|}\cdot\log_{\sqrt{2}}\left(x^2-2x+3\right)+2^{-x^2+2x}\cdot\log_{\frac{1}{2}}(2|x-m|+2)=0$ \\
		$ \Leftrightarrow 2^{1-2|x-m|}\cdot\log_2\left[(x-1)^2+2\right]-2^{1-(x-1)^2}\cdot\log_2(2|x-m|+2)=0 $ \\
		$\Leftrightarrow\dfrac{2^{1-2|x-m|}}{\log_2(2|x-m|+2)}=\dfrac{2^{1-(x-1)^2}}{\log_2\left[(x-1)^2+2\right]} $.\\
		Xét hàm số $f(t)=\dfrac{2^{1-t}}{\log_2(t+2)}$ có $f'(t)=\dfrac{-2^{1-t}\cdot\ln(t+2)-\dfrac{2^{1-t}}{(t+2)\ln 2}}{\left[\log_2(t+2)\right]^2}<0,\forall t >-2$.\\
		Phương trình đã cho $\Leftrightarrow f\left[(x-1)^2\right]=f(2|x-m|)\Leftrightarrow(x-1)^2=2|x-m|$ \\
		$\Leftrightarrow\left(x^2-4x+1+2m\right)\left(x^2+1-2m\right)=0 $ \\
		$\Leftrightarrow\hoac{&x^2-4x+2m+1=0\\&x^2+1-2m=0}\Leftrightarrow\hoac{&\dfrac{1}{2}x^2+1=m(1)\\&-\dfrac{1}{2}x^2+2x-\dfrac{1}{2}=m(2).} $ \\
		Khi đó $ycbt\Leftrightarrow$ phương trình $(1)$ và $(2)$ có tổng cộng $3$ nghiệm thực phân biệt.\\
		Vẽ đồ thị hàm số $f(x)=\dfrac{1}{2}x^2+\dfrac{1}{2}$ và $g(x)=-\dfrac{1}{2}x^2+2x-\dfrac{1}{2}$ trên cùng một hệ trục tọa độ (tham khảo hình vẽ). 
		{\color{red} HINH O DAY}\\
		Đồ thị hàm số $f(x)$ và $g(x)$ tiếp xúc với nhau tại điểm có hoành độ $x=1$.\\
		Dựa vào đồ thị ta có $m=\dfrac{1}{2},m=1,m=\dfrac{3}{2}$ thì phương trình đã cho có $3$ nghiệm thực phân biệt.\\
		Vậy tổng các giá trị thực của $m$ thỏa ycbt là $\dfrac{1}{2}+1+\dfrac{3}{2}=3$.
	}
\end{ex}
\begin{ex}
	[THPT Nguyễn Viết Xuân - 2020]%Câu 47.
	Cho phương trình $\log_3^2(9x)-(m+5)\log_3x+3m-10=0$ (với $m$ là tham số thực). Số giá trị nguyên của tham số $m$ để phương trình đã cho có hai nghiệm phân biệt thuộc $[1;81]$ là
	\choice
	{$3$}
	{$5$}
	{\True $4$}
	{$2$}
	\loigiai{
		Ta có $\log_3^2(9x)-(m+5)\log_3x+3m-10=0\Leftrightarrow\log_3^2(x)-(m+1)\log_3x+3m-6=0,(1)$.\\
		Đặt $t=\log_3x$, khi $x\in[1;81]$ thì $t\in[0;4]$.\\
		Khi đó ta có phương trình $t^2-(m+1)t+3m-6=0\Leftrightarrow\hoac{&t=3\\&t=m-2.}$ \\
		Phương trình đã cho có hai nghiệm phân biệt thuộc $[1;81]\Leftrightarrow$ phương trình $(1)$ có hai nghiệm phân biệt $t\in[0;4]\Leftrightarrow\heva{&m-2\neq 3\\&0\leq m-2\leq 4}\Leftrightarrow\heva{&m\neq 5\\&2\leq m\leq 6.}$ \\
		Suy ra có $4$ giá trị nguyên của tham số $m$ để phương trình đã cho có hai nghiệm phân biệt thuộc $[1;81]$.\\
		Chọn đáp án C.
	}
\end{ex}
\begin{ex}
	[THPT Nguyễn Viết Xuân - 2020]%Câu 48.
	Cho $x,y$ là hai số thực dương thỏa mãn $5x+y=4$. Tổng tất cả giá trị nguyên của tham số $m$ để phương trình $\log_3\dfrac{x^2+2y+m}{x+y}+x^2-3x-y+m-1=0$ có nghiệm là
	\choice
	{$10$}
	{\True $5$}
	{$9$}
	{$2$}
	\loigiai{
		$\begin{aligned}&\log_3\dfrac{x^2+2y+m}{x+y}+x^2-3x-y+m-1=0\\&\Leftrightarrow\log_3\left(x^2+2y+m\right)+x^2+2y+m=\log_33(x+y)+3(x+y)(1)\end{aligned}$.\\
		Vì $x,y>0$ nên $x+y>0$. Xét hàm số $f(t)=\log_3t+t$ là hàm số đồng biến trên $(0;+\infty)$.\\
		Khi đó $(1)\Leftrightarrow x^2+2y+m=3x+3y\Leftrightarrow x^2-3x-y+m=0(*)$.\\
		Kết hợp với điều kiện $5x+y=4\Rightarrow y=4-5x$. Vì $x,y>0\Rightarrow 0<x<\dfrac{4}{5}$.\\
		Ta có $(*)\Leftrightarrow x^2+2x+m-4=0\Leftrightarrow m=-x^2-2x+4,\forall x\in\left(0;\dfrac{4}{5}\right)$.\\
		Hàm số $y=-x^2-2x+4$ nghịch biến trên $\left(0;\dfrac{4}{5}\right)$ (do $-1<0$) nên $\dfrac{44}{25} <-x^2-2x+4<4$.\\
		Do vậy $m\in\{2;3\}$ là các giá trị cần tìm.\\
		Vậy tổng tất cả các giá trị $m$ thỏa ycbt là 5.
	}
\end{ex}
\begin{ex}
	[Hải Hậu - Nam Định - 2020]%Câu 49.
	Biết rằng điều kiện cần và đủ của tham số $m$ để phương trình $\log_2\left(m+\sqrt{m+2^x}\right)=2x$ có nghiệm là $m\geq-\dfrac{a}{b}$ với $a, b$ là hai số nguyên dương và $b<7$. Hỏi $a+b+b^2$ bằng bao nhiêu?
	\choice
	{$31$}
	{$32$}
	{\True $21$}
	{$23$}
	\loigiai{
		$\log_2\left(m+\sqrt{m+2^x}\right)=2x\Leftrightarrow\heva{&m+2^x\geq 0\\&m+\sqrt{m+2^x}=2^{2x}\quad(*).}$ \\
		$(*)\Leftrightarrow\left(m+2^x\right)+\sqrt{m+2^x}=(2^x)^2+2^x$.\\
		Xét hàm số $f(t)=t^2+t(t\geq 0)$. Ta có $f'(t)=2t+1>0$ với mọi $t\geq 0$, suy ra hàm số luôn đồng biến với mọi $t\geq 0$.\\
		$(*)\Leftrightarrow f\left(\sqrt{m+2^x}\right)=f(2^x)\Leftrightarrow\sqrt{m+2^x}=2^x\Leftrightarrow(2^x)^2=2^x+m(**)$.\\
		Đặt $t=2^x(t>0)$, khi đó phương trình $(**)$ trở thành $t^2-t=m(***)$.\\
		Xét hàm $g(t)=t^2-t(t>0)$, ta có $g'(t)=2t-1\Rightarrow g'(t)=0\Leftrightarrow t=\dfrac{1}{2}$.\\
		Bảng biến thiên.
		\begin{center}
			\begin{tikzpicture}
				\tkzTabInit[nocadre,lgt=1,espcl=3]{$t$/1.0,$g'$/0.8,$g$/2.5}{$0$,$\dfrac{1}{2}$,$+\infty$}
				\tkzTabLine{,-,0,+,}
				\tkzTabVar{+/{},-/$-\dfrac{1}{4}$,+/$+\infty$}
			\end{tikzpicture}
		\end{center}
		Vậy để $(***)$ có nghiệm $t>0$ thì $m\geq-\dfrac{1}{4}\Rightarrow\heva{&a=1\\&b=4}\Rightarrow a+b+b^2=21$.
	}
\end{ex}
\begin{ex}
	[Lương Thế Vinh - Hà Nội - 2020]%Câu 50.
	Có bao nhiêu giá trị nguyên của tham số $m$ để phương trình $\log_2^2(4x)-m\log_{\sqrt{2}}x-2m-4=0$ có nghiệm thuộc đoạn $[1;8]$? 
	\choice
	{$1$}
	{$2$}
	{$5$}
	{\True $3$}
	\loigiai{
		Điều kiện: $x>0$.\\
		$\begin{aligned}&\log_2^2(4x)-m\log_{\sqrt{2}}x-2m-4=0\\&\Leftrightarrow(2+\log_2x)^2-2m\log_2x-2m-4=0\\&\Leftrightarrow\log_2^2x+4\log_2x=2m(\log_2x+1)(1)\\&\log_2x=t;x\in[1;8]\Rightarrow t\in[0;3]\\&\Rightarrow(1)\Leftrightarrow\dfrac{t^2+4t}{t+1}=2m\\&f(t)=\dfrac{t^2+4t}{t+1};t\in[0;3]\\&f'(t)=\dfrac{t^2+2t+4}{(t+1)^2}>0,\forall t\in[0;3].\end{aligned}$\\
		$\begin{aligned}&\Rightarrow f(0)\leq 2m\leq f(3)\\&\Leftrightarrow 0\leq m\leq\dfrac{21}{8},m\in\mathbb{Z}\Rightarrow m\in\{0,1,2\}.\end{aligned}$.
	}
\end{ex}
\begin{ex}
	[Chuyên Biên Hòa - Hà Nam - 2020]%Câu 51.
	Cho hàm số bậc ba $y=f(x)$ có đồ thị như hình vẽ. Có bao nhiêu giá trị nguyên của tham số $m\in[-5;5]$ sao cho phương trình $\log_2^3\left(f(x)+1\right)-\log_{\sqrt{2}}^2\left(f(x)+1\right)+(2m-8)\log_{\frac{1}{2}}\sqrt{f(x)+1}+2m=0$ có nghiệm $x\in(-1;1)$?
	\begin{center}
		\begin{tikzpicture}[>=stealth,scale=0.9]
			\draw[->] (-3,0)--(3,0)node[below]{$x$};
			\draw[->] (0,-3)--(0,5)node[left]{$y$};
			\draw[samples=100,domain=2.11:-2.23] plot (\x,{0.92*(\x)^3+0.17*(\x)^2-2.92*(\x)+0.83});
			\draw(0,0)node[below left]{$O$};
			\draw[dashed]
			(2,0)node[below]{$2$}
			(-2,0)node[below left]{$-2$}
			(1,0)node[above]{$1$}
			(-1,0)node[below]{$-1$}
			(0,-1)node[left]{$-1$}
			(0,3)node[above left]{$3$}
			(2,0)--(2,3)--(0,3)
			(-1,0)--(-1,3)--(0,3)
			(1,0)--(1,-1)--(0,-1)
			;
		\end{tikzpicture}
	\end{center}
	\choice      
	{\True $7$}
	{$5$}
	{$6$}
	{vô số}
	\loigiai{
		Với $x\in(-1;1)\Rightarrow-1<f(x)<3\Leftrightarrow 0<f(x)+1<4$.\\
		Đặt $t=\log_2\left(f(x)+1\right)\Rightarrow t\in(-\infty;2),\forall x\in(-1;1)$.\\
		Khi đó phương trình đã cho trở thành: $t^3-4t^2-(m-4)t+2m=0$ \\
		$\Leftrightarrow(t-2)\left(t^2-2t-m\right)=0\Leftrightarrow\hoac{&t=2\notin(-\infty;2)\\&t^2-2t-m=0}\Leftrightarrow t^2-2t-m=0\Leftrightarrow t^2-2t=m(*) $.\\
		Để phương trình đã cho có $x\in(-1;1)\Leftrightarrow$ phương trình $(*)$ có nghiệm $t\in(-\infty;2)$.\\
		Xét hàm số $f(t)=t^2-2t$ trên $(-\infty;2)$ có $f'(t)=2t-2=0\Leftrightarrow t=1\in(-\infty;2)$.\\
		Ta có bảng biến thiên của hàm số $f(t)=t^2-2t$.
		\begin{center}
			\begin{tikzpicture}
				\tkzTabInit[nocadre,lgt=1.2,espcl=3]{$t$/0.8,$f'(t)$/0.8,$f(t)$/2.5}{$-\infty$,$1$,$2$}
				\tkzTabLine{,-,0,+,}
				\tkzTabVar{+/$+\infty$,-/$-1$,+/$0$}
			\end{tikzpicture}
		\end{center}
		Từ bảng biến thiên suy ra phương trình $(*)$ có nghiệm $t\in(-\infty;2)$ khi và chỉ khi $m\geq-1$.\\
		Mà $\heva{&m\in[-5;5]\\&m\in\mathbb{Z}}\Rightarrow m\in\left\{-1;0;1;2;3;4;5\right\}$. Vậy có $7$ giá trị của $m$ thỏa mãn yêu cầu bài toán.
	}
\end{ex}
\begin{ex}
	[THPT PTNK Cơ sở 2 - TP. HCM - 2021]%Câu 52.
	Cho phương trình $\log_{2+\sqrt{3}}^2\left(\sqrt{x^2+1}+x\right)+\left(m^2-2\right)\log_{2-\sqrt{3}}\left(\sqrt{x^2+1}-x\right)-1=0$ với $m$ là tham số thực. Gọi $S$ là tập hợp tất cả các giá trị của $m$ để phương trình có 2 nghiệm phân biệt $x_1; x_2$ thỏa mãn $\dfrac{\sqrt{x_1^2+1}-x_1}{\sqrt{x_2^2+1}+x_2}=7+4\sqrt{3}$. Tích các phần tử của $S$ bằng
	\choice
	{$-4$}
	{\True $4$}
	{$0$}
	{$2$}
	\loigiai{
		Ta có $(2+\sqrt{3})(2-\sqrt{3})=1$ và $\sqrt{x^2+1}-x=\dfrac{x^2+1-x^2}{\sqrt{x^2+1}+x}=\dfrac{1}{\sqrt{x^2+1}+x}$.\\
		$\log_{2+\sqrt{3}}^2\left(\sqrt{x^2+1}+x\right)+\left(m^2-2\right)\log_{2-\sqrt{3}}\left(\sqrt{x^2+1}-x\right)-1=0$ \\
		$\Leftrightarrow\log_{2+\sqrt{3}}^2\left(\sqrt{x^2+1}+x\right)-\left(m^2-2\right)\log_{2+\sqrt{3}}\left(\dfrac{1}{\sqrt{x^2+1}+x}\right)-1=0 $ \\
		$\Leftrightarrow\log_{2+\sqrt{3}}^2\left(\sqrt{x^2+1}+x\right)+\left(m^2-2\right)\log_{2+\sqrt{3}}\left(\sqrt{x^2+1}+x\right)-1=0 (*) $.\\
		Đặt $t=\log_{2+\sqrt{3}}\left(\sqrt{x^2+1}+x\right)\Leftrightarrow\sqrt{x^2+1}+x=(2+\sqrt{3})^t$, ta có\\
		$t'=\dfrac{1}{\left(\sqrt{x^2+1}+x\right)\ln(2+\sqrt{3})}\cdot\left(\dfrac{x}{\sqrt{x^2+1}}+1\right)=\dfrac{1}{\ln(2+\sqrt{3})\sqrt{x^2+1}}>0;\forall x$.\\
		Khi đó phương trình $(*)$ đã cho trở thành: $t^2+\left(m^2-2\right)t-1=0$.\\
		Yêu cầu bài toán tương đương: Phương trình $t^2+\left(m^2-2\right)t-1=0$ có 2 nghiệm phân biệt $t_1;t_2$ thỏa mãn $\dfrac{1}{(2+\sqrt{3})^{t_1}(2+\sqrt{3})^{t_2}}=7+4\sqrt{3}\Leftrightarrow(2+\sqrt{3})^{-(t_1+t_2)}=(2+\sqrt{3})^2\Leftrightarrow t_1+t_2=-2$ \\
		$\Leftrightarrow\heva{&\Delta>0\\&2-m^2=-2}\Leftrightarrow\heva{&\left(m^2-2\right)^2+4>0\\&m=\pm 2}\Leftrightarrow m=\pm 2 $.\\
		Vậy $S=\{-2; 2\}$. Tích các phần tử của $S$ là $-4$.
	}
\end{ex}
\begin{ex}
	[THPT Nguyễn Tất Thành - Hà Nội - 2021]%Câu 53.
	Có bao nhiêu số nguyên $m$ thỏa mãn $\dfrac{\ln x}{x+1}+\dfrac{1}{x}>\dfrac{\ln x}{x-1}+\dfrac{m}{x}$, $\forall x>0$, $x\neq 1$?
	\choice
	{$2$}
	{$1$}
	{\True Vô số}
	{$0$}
	\loigiai{
		Với $x>0$, $x\neq 1$ ta có\\
		$\dfrac{\ln x}{x+1}+\dfrac{1}{x}>\dfrac{\ln x}{x-1}+\dfrac{m}{x}\Leftrightarrow\dfrac{m}{x}<\ln x\left(\dfrac{1}{x+1}-\dfrac{1}{x-1}\right)+\dfrac{1}{x}\Leftrightarrow m <-2\dfrac{x\ln x}{x^2-1}+1=f(x)$.\\
		Xét hàm số: $f(x)=-2\dfrac{x\ln x}{x^2-1}+1$ với $x>0$, $x\neq 0$.\\
		Ta có $f'(x)=2\cdot\dfrac{x^2\ln x+\ln x+1-x^2}{\left(x^2-1\right)^2}$.\\
		$f'(x)=0\Leftrightarrow x^2\ln x+\ln x+1-x^2=0 (1)$.\\
		Xét hàm số $g(x)=x^2\ln x+\ln x+1-x^2$ với $x>0$.\\
		Đạo hàm $g'(x)=2x\ln x+\dfrac{1}{x}-x$.\\
		$g''(x)=2\ln x+1-\dfrac{1}{x^2}$.\\
		$g'''(x)=\dfrac{2}{x}+\dfrac{1}{x^3}>0$, $\forall x>0\Rightarrow$ Hàm số $g''(x)$ đồng biến trên khoảng $(0;+\infty)$.\\
		Từ đó suy ra phương trình $g''(x)=0$ nếu có nghiệm thì nghiệm đó là duy nhất.\\
		Lại có $g''(1)=0$. Suy ra phương trình $g''(x)=0$ có nghiệm duy nhất $x=1$.\\
		Bảng biến thiên của hàm số $y=g'(x)$:
		\begin{center}
			\begin{tikzpicture}
				\tkzTabInit[nocadre,lgt=1.2,espcl=3]{$x$/0.8,$g''(x)$/0.8,$g'(x)$/2.5}{$0$,$1$,$+\infty$}
				\tkzTabLine{,-,0,+,}
				\tkzTabVar{+/{},-/$0$,+/{}}
			\end{tikzpicture}
		\end{center}
		Từ bảng biến thiên suy ra phương trình $g'(x)=0$ có nghiệm duy nhất $x=1$.\\
		Bảng biến thiên của hàm số $y=g(x)$: \\
		{\color{red} HINH O DAY}\\
		Từ bảng biến thiên suy ra phương trình $g(x)=0$ có nghiệm duy nhất $x=1$.\\
		Giới hạn:\\
		$\lim\limits_{x\to 1} f(x)=\lim\limits_{x\to 1}\left(1-2\cdot\dfrac{\ln x}{x-1}\cdot\dfrac{x}{x+1}\right)=1-2\lim\limits_{x\to 1}\left(\dfrac{\ln x}{x-1}\cdot\dfrac{x}{x+1}\right)=1-1=0$.\\
		Bảng biến thiên của hàm số $y=f(x)$: 
		\begin{center}
			\begin{tikzpicture}
				\tkzTabInit[nocadre,lgt=1.2,espcl=3]{$x$/0.8,$f'(x)$/0.8,$f(x)$/2.5}{$0$,$1$,$+\infty$}
				\tkzTabLine{,-,d,+,}
				\tkzTabVar{+/{},-D-/$0$/$0$,+/{}}
			\end{tikzpicture}
		\end{center}
		Bất phương trình $m<f(x)$ nghiệm đúng $\forall x>0$, $x\neq 1\Leftrightarrow m\leq 0$.\\
		Vậy có vô số các giá trị nguyên của $m$ thỏa mãn yêu cầu bài toán.
	}
\end{ex}
\begin{ex}
	[Sở Yên Bái - 2021]%Câu 54.
	Biết điều kiện cần và đủ của tham số $m$ để phương trình $\log_{\frac{1}{2}}^2(x+2)^2+4(m-5)\log_{\frac{1}{2}}\dfrac{1}{x+2}-8m+4=0$ có nghiệm thuộc $\left[-\dfrac{3}{2};6\right]$ là $m\in(-\infty; a]\cup[b;+\infty)$. Tính giá trị biểu thức $T=a+b$. 
	\choice
	{$T=-\dfrac{8}{3}$}
	{$T=-\dfrac{22}{3}$}
	{$T=\dfrac{8}{3}$}
	{\True $T=\dfrac{22}{3}$}
	\loigiai{
		Điều kiện: $x+2>0\Leftrightarrow x >-2$.\\
		Ta có $\log_{\frac{1}{2}}^2(x+2)^2+4(m-5)\log_{\frac{1}{2}}\dfrac{1}{x+2}-8m+4=0$ \\
		$ \Leftrightarrow 4\log_{\frac{1}{2}}^2(x+2)-4(m-5)\log_{\frac{1}{2}}(x+2)-8m+4=0 $ \\
		$\Leftrightarrow\log_{\frac{1}{2}}^2(x+2)-(m-5)\log_{\frac{1}{2}}(x+2)-2m+1=0(1) $.\\
		Đặt $t=\log_{\frac{1}{2}}(x+2)$, với $x\in\left[-\dfrac{3}{2};6\right]\Rightarrow t\in[-3; 1]$.\\
		Phương trình (1) trở thành $t^2-(m-5)t-2m+1=0$ \\
		$ \Leftrightarrow t^2-mt+5t-2m+1=0\Leftrightarrow m(t+2)=t^2+5t+1 $.\\
		Nhận thấy $t=-2$ không là nghiệm nên $\dfrac{t^2+5t+1}{t+2}=m(2)$.\\
		Xét hàm số $y(t)=\dfrac{t^2+5t+1}{t+2}$ trên $[-3; 1]$.\\
		Ta có $y'(t)=\dfrac{t^2+4t+9}{(t+2)^2}>0,\forall t\in[-3; 1]\setminus\{-2\}$.\\
		Bảng biến thiên: 
		\begin{center}
			\begin{tikzpicture}
				\tkzTabInit[nocadre,lgt=1.2,espcl=3]{$t$/0.8,$y'(t)$/0.8,$y(t)$/2.5}{$-3$,$-2$,$1$}
				\tkzTabLine{,+,d,+,}
				\tkzTabVar{-/$5$,+D-/$+\infty$/$-\infty$,+/$\dfrac{7}{3}$}
			\end{tikzpicture}
		\end{center}
		Để phương trình ban đầu có nghiệm thuộc $\left[-\dfrac{3}{2};6\right]\Rightarrow$ Phương trình (2) có nghiệm thuộc đoạn.\\
		$[-3; 1]$.\\
		Suy ra $m\in\left(-\infty;\dfrac{7}{3}\right]\cup[5;+\infty)$.
	}
\end{ex}
\begin{ex}
	[Chuyên Long An - 2021]%Câu 55.
	Số giá trị nguyên của tham số $m\in[-20;10]$ để bất phương trình $9\left(\log_3\sqrt[3]{x}\right)^2+\log_3x+2m\geq 0$ nghiệm đúng với mọi giá trị $x\in(3;81)$. 
	\choice
	{\True $12$}
	{$10$}
	{$11$}
	{$15$}
	\loigiai{
		Điều kiện $x>0$: $9\left(\log_3\sqrt[3]{x}\right)^2+\log_3x+2m\geq 0\Leftrightarrow\log_3^2x+\log_3x+2m\geq 0 (*)$.\\
		Đặt $t=\log_3x$, với $x\in(3;81)\Rightarrow t\in(1;4)$. $(*)$ trở thành $t^2+t+2m\geq 0\Leftrightarrow\underbrace{t^2+t}_{f(t)}\geq-2m$.\\
		Bảng biến thiên của $f(t)$: \\
		{\color{red} HINH O DAY}\\
		Vậy để bất phương trình nghiệm đúng với mọi giá trị $x\in(3;81)\Leftrightarrow-2m\leq 2\Leftrightarrow m\geq-1\Rightarrow m\in\{-1;0;\ldots;10\}$ nên có $12$ giá trị nguyên của $m$.
	}
\end{ex}
\begin{ex}
	[THPT Lương Thế Vinh – Hà Nội -2022]%Câu 56.
	Gọi $S$ là tập các số nguyên $m\in [-2022;2022]$ để phương trình $\log_2^2x-\log_{\sqrt{2}}x=m-\sqrt{m+\log_2x}$ có đúng ba nghiệm phân biệt. Số phần tử của $S$ là
	\choice
	{\True $1$}
	{$2$}
	{$2021$}
	{$2022$}
	\loigiai{
		Đặt $t=\sqrt{m+\log_2x},(t\geq 0)\Rightarrow m=t^2-\log_2x$ phương trình trở thành\\
		$\log_2^2x-2\log_2x=\left(t^2-\log_2x\right)-t
		\Leftrightarrow\log_2^2x-\log_2x-t^2+t=0
		\Leftrightarrow\hoac{&\log_2x=t\\&\log_2x=1-t.}$ \\
		(bấm máy phương trình bậc hai ần là $\log_2x$ và tham số $t=1000$.\\
		Khi đó $\hoac{&m=t^2-t\\&m=t^2+t-1}(*)$. \\
		{\color{red} HINH O DAY}\\
		ycbt $\Leftrightarrow (*)$ có 3 nghiệm phân biệt trên $[0;+\infty)\Leftrightarrow$ đường thẳng $y=m$ cắt đồ thị của hai hàm số $y=t^2-t;y=t^2+t-1$ tại ba điểm phân biệt trên $[0;+\infty)\Leftrightarrow-\dfrac{1}{4}<m\leq 0$.
	}
\end{ex}
\begin{ex}
	[THPT Yên Lạc - Vĩnh Phúc - 2022]%Câu 57.
	Có bao nhiêu giá trị nguyên của tham số $m$ để phương trình $2\log_3(2x-1)=\log_3\left(mx^2+1\right)$ có nghiệm?
	\choice
	{$1$}
	{$3$}
	{\True $7$}
	{$9$}
	\loigiai{
		Điều kiện xác định $x>\dfrac{1}{2}$.\\
		Với điều kiện trên, $2\log_3(2x-1)=\log_3\left(mx^2+1\right)\Leftrightarrow\log_3(2x-1)^2=\log_3\left(mx^2+1\right)$ \\
		$ \Leftrightarrow 4x^2-4x+1=mx^2+1 $ \\
		$ \Leftrightarrow 4-\dfrac{4}{x}=m $ (*).\\
		Xét hàm số $f(x)=4-\dfrac{4}{x}$ với $x>\dfrac{1}{2}$.\\
		Ta có $f'(x)=\dfrac{4}{x^2}>0,\forall x>\dfrac{1}{2}$.\\
		Bảng biến thiên của hàm số $f(x)=4-\dfrac{4}{x}$ \\
		{\color{red} HINH O DAY}\\
		Dựa vào bảng biến thiên, phương trình đã cho có nghiệm $\Leftrightarrow-4<m<4$.\\
		Vậy có $7$ giá trị nguyên của tham số $m$.
	}
\end{ex}
\begin{ex}
	[Sở Phú Thọ 2022]%Câu 58.
	Cho phương trình $\log_a4+\log_{\frac{1}{5}}\left(\sqrt{x^2+ax+2}+4\right)\log_a\left(x^2+ax+5\right)=0$. Gọi $S$ là tập các giá trị nguyên của tham số $a$ để phương trình có nghiệm duy nhất. Tổng các phần tử của $S$ bằng
	\choice
	{$4$}
	{$3$}
	{$0$}
	{\True $2$}
	\loigiai{
		Điều kiện: $0<a\neq 1$.\\
		Ta có phương trình: $\log_a4+\log_{\frac{1}{5}}\left(\sqrt{x^2+ax+2}+4\right)\log_a\left(x^2+ax+5\right)=0$ (1).\\
		$(1)\Leftrightarrow\log_5\left(\sqrt{x^2+ax+2}+4\right)\log_a\left(x^2+ax+5\right)=\log_a4$.\\
		Đặt $t=\sqrt{x^2+ax+2}\geq 0$ khi đó phương trình trở thành $\Leftrightarrow\log_5(t+4)\log_a\left(t^2+3\right)=\log_a4$ (2). Khi đó ta chia hai trường hợp như sau:\\
		Trường hợp 1: $0<a<1\Rightarrow\dfrac{1}{\log_4a}<0$ thì (2) thành $\Leftrightarrow\log_5(t+4)\log_4\left(t^2+3\right)=\log_a4<0$.\\
		Hàm vế trái $y=f(t)=\log_5(t+4)\log_4\left(t^2+3\right)$ có $f'(t)\geq 0,\forall t\in [0;+\infty)$ tức $f(t)$ đồng biến trên $[0;+\infty)$.\\
		Để ý $f(1)=1$ nên ta suy ra $\hoac{&\log_a4=1\\&t=1}\Leftrightarrow\hoac{&a=4\\&x^2+ax+1=0(*)}$.\\
		Phương trình có nghiệm duy nhất nên ta suy ra $a=\pm 2$. Đối chiếu với điều kiện, ta loại.\\
		Trường hợp 2: $a>1$ thì (2) ta đánh giá tương tự $\Leftrightarrow\log_5(t+4)\log_4\left(t^2+3\right)=\log_a4>0$, đánh giá tương tự trường hợp ta cunng suy ra $a=\pm 2$, đối chiếu điều kiện ta thu được $a\in\{2\}$.
	}
\end{ex}
\begin{ex}
	[THPT Hoàng Hoa Thám - Đà Nẵng 2022]%Câu 59.
	Cho phương trình $\log_{\sqrt{5}+2}\left(x^2+mx+m+7\right)+\log_{\sqrt{5}-2}x=0$. Số giá trị nguyên $m$ thuộc $[-10;9]$ để phương trình có nghiệm duy nhất là
	\choice
	{$3$}
	{\True $5$}
	{$4$}
	{$2$}
	\loigiai{
		$\log_{\sqrt{5}+2}\left(x^2+mx+m+7\right)+\log_{\sqrt{5}-2}x=0\Leftrightarrow\log_{\sqrt{5}+2}\left(x^2+mx+m+7\right)+\log_{(\sqrt{5}+2)^{-1}}x=0$ \\
		$ \Leftrightarrow\log_{\sqrt{5}+2}\left(x^2+mx+m+7\right)-\log_{(\sqrt{5}+2)}x=0\Leftrightarrow\log_{\sqrt{5}+2}\left(x^2+mx+m+7\right)=\log_{(\sqrt{5}+2)}x $ \\
		$ \Leftrightarrow\heva{&x^2+mx+m+7=x\\&x>0}\Leftrightarrow\heva{&m=\dfrac{-x^2+x-7}{x+1}\\&x>0.} $ \\
		Xét $g(x)=\dfrac{-x^2+x-7}{x+1}$ với $x>0$; $g'(x)=\dfrac{-x^2-2x+8}{(x+1)^2}$.\\
		$g'(x)=0\Leftrightarrow-x^2-2x+8=0\Leftrightarrow\hoac{&x=2\\&x=-4}$, $g(2)=-3$. \\
		{\color{red} HINH O DAY}\\
		Từ bảng biến thiên, nhận thấy phương trình $m=\dfrac{-x^2+x-7}{x+1}$ có 1 nghiệm dương khi và chỉ khi $\hoac{&m\leq-7\\&m=-3.}$ \\
		Vậy số giá trị nguyên $m$ thuộc $[-10;9]$ là $m\in\left\{-10;-9;-8;-7;-3\right\}$.
	}
\end{ex}
\begin{ex}
	[THPT Nguyễn Viết Xuân – Vĩnh Phúc – 2022]%Câu 60.
	Cho phương trình \break $\log_3\dfrac{2x-1}{27x^2-54x+9m}=3x^2-8x+m-1$. Gọi $S$ là tập tất cả các giá trị nguyên dương của $m$ để phương trình đã cho có 2 nghiệm thực phân biệt thuộc $\left(\dfrac{1}{2};+\infty\right)$. Tổng các phần tử của $S$ bằng 
	\choice
	{$4$}
	{$5$}
	{$6$}
	{\True $7$}
	\loigiai{
		Do ta xét nghiệm của phương trình thỏa $x>\dfrac{1}{2}\Rightarrow 2x-1>0$ nên $27x^2-54x+9m>0$.\\
		Đặt $\heva{&T=2x-1\\&M=27x^2-54x+9m=9\left(3x^2-6x+m\right).}$ \\
		Khi đó phương trình đã cho $\Leftrightarrow\log_3\dfrac{T}{M}=\dfrac{M}{9}-T-2\Leftrightarrow\log_3\dfrac{9T}{M}=\dfrac{M}{9}-T$ \\
		$\Leftrightarrow\log_39T+T=\log_3\left[9\left(\dfrac{M}{9}\right)\right]+\left(\dfrac{M}{9}\right)(*) $.\\
		Xét hàm $f(t)=\log_3t+\dfrac{t}{9}$ với $t>0$. Dễ dàng chứng minh $f(t)$ đồng biến trên $(0;+\infty)$.\\
		Do đó $(*)\Leftrightarrow T=\dfrac{M}{9}\Leftrightarrow 3x^2-8x+m+1=0\Leftrightarrow m=\underbrace{-3x^2+8x-1}_{g(x)},\forall x>\dfrac{1}{2}$.\\
		Ta có $g'(x)=-6x+8$. Cho $g'(x)=0\Leftrightarrow x=\dfrac{4}{3}\in\left(\dfrac{1}{2};+\infty\right)$.\\
		Lập bảng biến thiên của $g(x)$ ta có 
		\begin{center}
			\begin{tikzpicture}
				\tkzTabInit[nocadre,lgt=1.2,espcl=3]{$x$/0.8,$g(x)$/2}{$\dfrac{1}{2}$,$\dfrac{4}{3}$,$+\infty$}
				\tkzTabVar{-/$\dfrac{9}{4}$,+/$\dfrac{13}{3}$,-/$-\infty$}
			\end{tikzpicture}
		\end{center}
		Dựa vào bảng biến thiên, yêu câu bài toán $\Leftrightarrow\dfrac{9}{4}<m<\dfrac{13}{3}$.\\
		Do $m\in\mathbb{Z}\Rightarrow m\in\{3;4\}$.\\
		Vậy tổng phần tử của tập $S$ là $3+4=7$.
	}     
\end{ex}
\begin{dang}
	{Phương trình mũ chứa tham số}
\end{dang}         
\begin{ex}    
	[THPT Minh Khai - 2019]%Câu 1
	Giá trị thực của tham số $m$ để phương trình $4^x-(2 m+3) \cdot 2^x+64=$ có hai nghiệm thực $x_1, x_2$ thỏa mãn $\left(x_1+2\right)\left(x_2+2\right)=24$ thuộc khoảng nào sau đây?
	\choice
	{$\left(0 ; \dfrac{3}{2}\right)$}
	{$\left(-\dfrac{3}{2}; 0\right)$}
	{$\left(\dfrac{21}{2}; \dfrac{29}{2}\right)$}
	{\True $\left(\dfrac{11}{2}; \dfrac{19}{2}\right)$}
	\loigiai{
		Đặt $t=2^x$, điều kiện $t>0$. Phương trình ban đầu trở thành $t^2-(2 m+3) \cdot t+64=0(*)$.\\
		Để phương trình ban đầu có hai nghiệm thực $x_1$ và $x_2$ thì phương trình $(*)$ phải có hai nghiệm
		$t_1, t_2~\text{dương}~\Leftrightarrow\left\{\begin{aligned}
			&{ \Delta > 0 } \\
			&{ S > 0 } \\
			&{ P > 0 }\\
		\end{aligned} \Leftrightarrow \left\{\begin{aligned}
			&{ 4 m ^ { 2 } + 1 2 m - 2 4 7 > 0 } \\
			&{ 2 m + 3 > 0 }
		\end{aligned} \Rightarrow \left\{\begin{aligned}
			&m<-\dfrac{19}{2} \\
			&m>\dfrac{13}{2} \\
			&m>-\dfrac{3}{2}
		\end{aligned} \Rightarrow m>\dfrac{13}{2} .\right.\right.\right.$.\\
		Theo định lý Vi-ét, ta có $t_1 \cdot t_2=64 \Leftrightarrow 2^{x_1}\cdot 2^{x_2}=64 \Leftrightarrow 2^{x_1+x_2}=64 \Rightarrow x_1+x_2=6$.\\
		Ta có $\left(x_1+2\right)\left(x_2+2\right)=24 \Leftrightarrow x_1 \cdot x_2+2\left(x_1+x_2\right)+4=24 \Leftrightarrow x_1 \cdot x_2=8$.\\
		Từ $\left\{\begin{aligned}
			&x_1+x_2=6 \\ 
			&x_1 \cdot x_2=8\\
		\end{aligned}\Rightarrow\left[\begin{aligned}
			\left\{\begin{aligned}
				&x_1=2 \\ 
				&x_2=4\\
			\end{aligned}\right. \\ \left\{\begin{aligned}
				&x_1=4 \\ 
				&x_2=2\\
			\end{aligned}\right.\end{aligned}\right.\right.$.\\
		Khi đó, ta có $t_1+t_2=2^{x_1}+2^{x_2}=20=2 m+3 \Rightarrow m=\dfrac{17}{2}$.
	}
\end{ex}
\begin{ex}
	[Chuyên - Vĩnh Phúc - 2019]%Câu 2
	Tìm tât cả các giá trị thực của tham số $m$ để phương trình $e^{3 m}+e^m=2\left(x+\sqrt{1-x^2}\right)\left(1+x \sqrt{1-x^2}\right)$ có nghiệm.
	\choice
	{$\left(0 ; \dfrac{1}{e}\right)$}
	{$\left(0 ; \dfrac{1}{2}\ln 2\right)$}
	{\True $\left(-\infty ; \dfrac{1}{2}\ln 2\right]$}
	{$\left[\dfrac{1}{2}\ln 2 ;+\infty\right)$}
	\loigiai{
		Điều kiện: $x \in[-1 ; 1]$.\\
		Đặt $x+\sqrt{1-x^2}=t$. Vì $x \in[-1 ; 1] \Rightarrow t \in[-1 ; \sqrt{2}]$.\\
		Ta có $t^2=\left(x+\sqrt{1-x^2}\right)^2=1+2 x \sqrt{1-x^2}\Rightarrow x \sqrt{1-x^2}=\dfrac{t^2-1}{2}$.\\
		Phương trình đã cho trở thành: $\mathrm{e}^{3 m}+\mathrm{e}^m=t^3+t$.\\
		Xét hàm số $f(u)=u^3+u, f'(u)=3 u^2+1>0 \forall u$ do đó hàm số $f$ đồng biến trên $\mathbb{R}$.\\
		Phương trình $\mathrm{e}^{3 m}+\mathrm{e}^m=t^3+t \Leftrightarrow f\left(e^m\right)=f(t) \Leftrightarrow \mathrm{e}^m=t$.\\
		Phương trình có nghiệm khi và chỉ khi $-1 \leq \mathrm{e}^m \leq \sqrt{2}\Rightarrow 0<\mathrm{e}^m \leq \sqrt{2}\left(\right.$ do $\left.\mathrm{e}^m>0\right)$.\\
		$\Leftrightarrow m \leq \ln \sqrt{2}\Rightarrow m \in\left(-\infty ; \dfrac{1}{2}\ln 2\right]$.
	}
\end{ex}
\begin{ex}
	[Chuyên Quang Trung - Bình Phước 2019]%Câu 3
	Tìm tất cả các giá trị thực của tham số $m$ để phương trình $4^x-m-1.2^x+2=0$ có hai nghiệm $x_1, x_2$ thoả mãn $x_1+x_2=1$.
	\choice
	{$m \in \mathbb{R}$}
	{$m>1+2 \sqrt{2}; m<1-2 \sqrt{2}$}
	{$m \geq 1+2 \sqrt{2}$}
	{\True $m>1+2 \sqrt{2}$}
	\loigiai{
		Đặt $t=2^x, t>0$. Ta có phương trình $t^2-m-1 t+2=0$ (1).\\
		Phương trình đã cho có hai nghiệm $x_1, x_2$ khi phương trình (1) có hai nghiệm $t_1, t_2>0$.\\
		Khi đó $x_1+x_2=\log _2 t_1+\log _2 t_2=\log _2 t_1 \cdot t_2$.\\
		Bài toán trở thành tìm $m$ để phương trình (1) có hai nghiệm dương phân biệt $t_1, t_2$ thoả mãn\\
		$\log _2 t_1 \cdot t_2=1 \Leftrightarrow\left\{\begin{aligned}
			&\Delta=m-1^2-8>0 \\ 
			&m-1>0 \\ 
			&\log _2 t_1 t_2=1\\
		\end{aligned}\Leftrightarrow m>1+2 \sqrt{2}\right.$.
	}
\end{ex}
\begin{ex}
	[Chuyên Quang Trung - Bình Phước 2019]%Câu 4
	Cho hàm số $y=f x$ liên tục trên $\mathbb{R}$ và có đồ thị như hình vẽ\\
	\textbf{(VẼ HÌNH)}\\
	Tập hợp tất cả các giá trị thực của $m$ để phương trình $f e^{x^2}=m$ có đúng 2 nghiệm thực là
	\choice
	{$0 ; 4$}
	{$0 ; 4$}
	{\True $0 \cup 4 ;+\infty$}
	{$4 ;+\infty$}
	\loigiai{
		Đặt $t=e^{x^2}$. Ta có $x^2 \geq 0 \Rightarrow t \geq 1$, nếu $t=1$ thì $x=0$ và nếu $t>1$ thì $x=\pm \sqrt{\ln t}$.\\ 
		Phương trình $f e^{x^2}=m.1$ trở thành phương trình $f t=m.2$.\\
		Sử dụng các nhận xét ở trên và đồ thị của hàm số $y=f x$ ta có
		1 có đúng 2 nghiệm $\Leftrightarrow 2$ có đúng 1 nghiệm thuộc $1 ;+\infty$ và nghiệm này lớn hơn 1.\\
		$\Leftrightarrow\left[\begin{aligned}
			&m=0 \\
			&m>4\\
		\end{aligned}\right.$.\\
		Vậy tập hợp các giá trị nguyên của $m$ thỏa yêu cầu bài toán là $0 \cup 4 ;+\infty$.
	}
\end{ex} 
\begin{ex}%Câu 5
	[Chuyên Thái Bình - 2019]%Câu 5
	Tìm số giá trị nguyên của tham số $m \in(-10 ; 10)$ để phương trình $(\sqrt{10}+1)^{x^2}+m(\sqrt{10}-1)^{x^2}=2 \cdot 3^{x^2+1}$ có đúng hai nghiệm phân biệt.
	\choice
	{$14$}
	{\True $15$}
	{$13$}
	{$16$}
	\loigiai{
		Ta có $(\sqrt{10}+1)^{x^2}+m(\sqrt{10}-1)^{x^2}=2 \cdot 3^{x^2+1}\Leftrightarrow\left(\dfrac{\sqrt{10}+1}{3}\right)^{x^2}+m\left(\dfrac{\sqrt{10}-1}{3}\right)^{x^2}=6$.\\
		Đặt $t=\left(\dfrac{\sqrt{10}+1}{3}\right)^{x^2}, t>0$.\\
		Phương trình (1) trở thành: $t+\dfrac{m}{t}=6 \Leftrightarrow-t^2+6 t=m$ (2).\\
		Phương trình (1) có 2 nghiệm phân biệt $\Leftrightarrow$ phương trình (2) có đúng 1 nghiệm lớn hơn 1. Xét hàm số: $g(t)=-t^2+6 t$ trên khoảng $(1 ;+\infty)$.\\
		Ta có $g'(t)=-2 t+6 \Rightarrow g^{\prime}(t)=0 \Leftrightarrow t=3$.\\
		Bảng biến thiên:\\
		\textbf{(VẼ HÌNH)}\\
		Nhìn vào bảng biến thiên ta có: phương trình (2) có đúng 1 nghiệm lớn hơn $1 \Leftrightarrow m \leq 5$. Kết hợp điều kiện $m$ nguyên và $m \in(-10 ; 10) \Rightarrow m \in(-10 ; 5] \cap \mathbb{Z}\Rightarrow$ Có 15 giá trị $m$ thỏa yêu cầu đề.
	}
\end{ex}
\begin{ex}
	[Thi thử cụm Vũng Tàu - 2019]%Câu 6
	Tổng tất cả các giá trị nguyên của $m$ để phương trình $3^{x-3+\sqrt[3]{m-3 x}}+\left(x^3-9 x^2+24 x+m\right) \cdot 3^{x-3}=3^x+1$ có 3 nghiệm phân biệt.
	\choice
	{$34$}
	{\True $27$}
	{$38$}
	{$45$}
	\loigiai{
		$\begin{aligned}
			& 3^{x-3+\sqrt[3]{m-3 x}}+\left(x^3-9 x^2+24 x+m\right) \cdot 3^{x-3}=3^x+1 \\
			& \Leftrightarrow 3^{x-3+\sqrt[3]{m-3 x}}+\left[(x-3)^3+27+m-3 x\right] \cdot 3^{x-3}=3^x+1 \\
			& \Leftrightarrow 3^{\sqrt[3]{m-3 x}}+(x-3)^3+m-3 x+27=3^3+3^{3-x}(1) \\
			& a=3-x ; b=\sqrt[3]{m-3 x}\\
			& (1) \Leftrightarrow 3^b+27+b^3-a^3=27 \cdot+3^a \Leftrightarrow 3^b+b^3=3^a+a^3 \\
			& \text{Xét}f(t)=3^t+t^3 \Rightarrow f^{\prime}(t)=3^t \cdot \ln 3+3 t^2 \geq 0 \forall t \in R \\
			& \Rightarrow f(a)=f(b) \Leftrightarrow a=b \Leftrightarrow 3-x=\sqrt[3]{m-3 x}\\
			& \Leftrightarrow m=(3-x)^3+3 x=-x^3+9 x^2-24 x+27 \\
			& f(x)=-x^3+9 x^2-24 x+27 \Rightarrow f^{\prime}(x)=-3 x^2+18 x-24 \\
			& f^{\prime}(x)=0 \Leftrightarrow x=2 \vee x=4
		\end{aligned}$\\
		Dựa vào đồ thị: $7<m<11 \Rightarrow m=8,9,10$.
	}
\end{ex}   
\begin{ex}
	[Chuyên ĐH Vinh- 2019]%Câu 7
	Cho số thực $m$ và hàm số $y=f(x)$ có đồ thị như hình vẽ dưới đây.
	Phương trình $f\left(2^x+2^{-x}\right)=m$ có nhiều nhất bao nhiêu nghiệm phân biệt thuộc đoạn $[-1 ; 2]$?
	\choice
	{$2$}
	{\True $3$}
	{$4$}
	{$5$}
	\loigiai{
		Đặt $t=t(x)=2^x+2^{-x}$, với $x \in[-1 ; 2]$.\\
		Hàm số $t=t(x)$ liên tục trên $[-1 ; 2]$ và $t'(x)=2^x \cdot \ln 2-2^{-x}\cdot \ln 2, t'(x)=0 \Leftrightarrow x=0$.\\
		Bảng biến thiên:\\
		\textbf{(VẼ HÌNH)}\\
		Vậy với $x \in[-1 ; 2] \Rightarrow t \in\left[2 ; \dfrac{17}{4}\right]$.\\
		Với mỗi $t \in\left(2 ; \dfrac{5}{2}\right]$ có 2 giá trị $x$ thỏa mãn $t=2^x+2^{-x}$.\\
		Với mỗi $t \in\{2\}\cup\left(\dfrac{5}{2}; \dfrac{17}{4}\right]$ có duy nhất 1 giá trị $x$ thỏa mãn $t=2^x+2^{-x}$.\\
		Xét phương trình $f(t)=m$ với $t \in\left[2 ; \dfrac{17}{4}\right]$.\\
		Từ đồ thị trên ta thấy phương trình $f\left(2^x+2^{-x}\right)=m$ có số nghiệm nhiều nhất khi và chỉ khi phương trình $f(t)=m$ có 2 nghiệm $t_1, t_2$, trong đó có $t_1 \in\left(2 ; \dfrac{5}{2}\right], t_2 \in\left(\dfrac{5}{2}; \dfrac{17}{4}\right]$.\\ 
		Do đó phương trình $f\left(2^x+2^{-x}\right)=m$ có nhiều nhất 3 nghiệm phân biệt thuộc đoạn $[-1 ; 2]$.
	}
\end{ex}
\begin{ex}
	[THPT Thuận Thành 3 - Bắc Ninh 2019]%Câu 8
	Gọi $S$ là tổng các giá trị nguyên của tham số $m$ để phương trình $4^x+7=2^{x+3}+m^2+6 m$ có nghiệm $x \in(1 ; 3)$. Chọn đáp án đúng.
	\choice
	{$S=-35$}     
	{$S=20$}
	{$S=25$}
	{\True $S=-21$}
	\loigiai{
		Đặt $t=2^x$ thì phương trình đã cho trở thành $t^2-8 t+7=m^2+6 m$ với $2<t<8$ (vì $1<x<3$).\\
		Khi đó phương trình $4^x+7=2^{x+3}+m^2+6 m$ có nghiệm $x \in(1 ; 3)$
		$\Leftrightarrow f(t)=m^2+6 m$ có nghiệm $t \in(2 ; 8)$, với $f(t)=t^2-8 t+7$.\\
		Ta có $f'(t)=2 t-8 ; f'(t)=0 \Leftrightarrow t=4 \in(2 ; 8)$.
		Bảng biến thiên:\\
		\textbf{(VẼ HÌNH)}\\
		Dựa vào bảng biến thiên ta có: $-9 \leq m^2+6 m<7$\\
		$\Leftrightarrow\left\{\begin{aligned}
			&{m ^{2}+ 6 m + 9 \geq 0}\\
			&{m ^{2}+ 6 m - 7 < 0}\\
		\end{aligned}\Leftrightarrow \left\{\begin{aligned}
			&m \in \mathbb{R}\\
			&-7<m<1\\
		\end{aligned}\Leftrightarrow-7<m<1 .\right.\right.$.\\
		Vì $m$ nguyên nên $m \in\{-6 ;-5 ;-4 ;-3 ;-2 ;-1 ; 0\}$.\\
		Suy ra $S=-(6+5+4+3+2+1+0)=-21$.
	}
\end{ex}
\begin{ex}
	[Chuyên Bắc Giang 2019]%Câu 9
	Tập các giá trị của $m$ để phương trình $4^{1+\sqrt{1-x^2}}-(m+2) 2^{1+\sqrt{1-x^2}}+2 m+1=0$ có nghiệm là
	\choice
	{$\left(-\infty ; \dfrac{9}{2}\right)$}
	{$\left[4 ; \dfrac{9}{2}\right]$}
	{$(-\infty ; 4)$}
	{\True $[4 ;+\infty)$}
	\loigiai{
		Điều kiện: $-1 \leq x \leq 1$. Đặt $t=2^{1+\sqrt{1-x^2}};-1 \leq x \leq 1 \Rightarrow 2 \leq t \leq 4$.\\
		Phương trình trở thành: $t^2-(m+2) t+2 m+1=0 \Leftrightarrow \dfrac{t^2-2 t+1}{t-2}=m(*)$.\\
		Đặt $ f(t)=\dfrac{t^2-2 t+1}{t-2}\Rightarrow f'(t)=\dfrac{t^2-4 t+3}{(t-2)^2}$.\\
		Phương trình có nghiệm $\Leftrightarrow(*)$ có nghiệm $t \in[2 ; 4] \Leftrightarrow m \geq 4$.
	}
\end{ex}
\begin{ex}
	Cho hàm số $f(x)=3^{x-4}+(x+1) \cdot 2^{7-x}-6 x+3$, khi phương trình $f\left(7-4 \sqrt{6 x-9 x^2}\right)+3 m-1=0$ có số nghiệm nhiều nhất thì giá trị nhỏ nhất của tham số $m$ có dạng $\dfrac{a}{b}$ (trong đó $a, b \in \mathbb{N}$ và $\dfrac{a}{b}$ là phân số tối giản). Tính $T=a+b$.
	\choice
	{$T=7$}
	{$T=11$}
	{\True $T=8$}
	{$T=13$}
	\loigiai{
		Đặt $t=7-4 \sqrt{6 x-9 x^2}=7-4 \sqrt{1-(3 x-1)^2}\in[3 ; 7]$. Khi đó $f(t)=1-3 m$.\\
		Xét hàm số $f(t)=3^{t-4}+(t+1) 2^{7-t}-6 t+3$ trên đoạn $[3 ; 7]$.
		Ta có $f'(t)=3^{t-4}\ln 3+2^{7-t}-(t+1) 2^{7-t}\ln 2-6$;
		$f''(t)=3^{t-4}(\ln 3)^2-2^{7-t}\ln 2-2^{7-t}\ln 2+(t+1) 2^{7-t}(\ln 2)^2$
		$=3^{t-4}(\ln 3)^2+\underbrace{[-2+(t+1) \ln 2]}_{>0, \forall t \in[3 ; 7]}2^{7-t}\ln 2>0$.\\
		Suy ra hàm số $f'(t)$ đồng biến trên $(3 ; 7)$.\\
		Lại có $\left\{\begin{aligned}
			&f'(3)<0 \\ 
			&f'(7)>0\\
		\end{aligned}\Rightarrow f'(x)=0\right.$ có nghiệm duy nhất $t_0$ thuộc $(3 ; 7)$.\\
		Dựa vào bảng biến thiên, ta thấy phương trình $f(t)=1-3 m$ có số nghiệm nhiều nhất
		$\Leftrightarrow f\left(t_0\right)<1-3 m \leq-4 \Leftrightarrow \dfrac{5}{3}\leq m<\dfrac{1-f\left(t_0\right)}{3}$.\\
		Suy ra giá trị nhỏ nhất của $m$ là $\dfrac{5}{3}\longrightarrow\left\{\begin{aligned}
			&a=5 \\ 
			&b=3\\
		\end{aligned}\right.$ nên $a+b=8$.
	}
\end{ex}
\begin{ex}
	Có bao nhiêu giá trị nguyên của tham số $m$ để phương trình $9^{1+\sqrt{1-x^2}}-(m+3) \cdot 3^{1+\sqrt{1-x^2}}+2 m+1=0$ có nghiệm thực?
	\choice
	{$5$}
	{\True $7$}
	{Vô số}
	{$3$} 
	\loigiai{
		Ta có $0 \leq 1-x^2 \leq 1 \Leftrightarrow 1 \leq 1+\sqrt{1-x^2}\leq 2 \Leftrightarrow 3 \leq 3^{1+\sqrt{1-x^2}}\leq 9$.\\
		Đặt $t=3^{1+\sqrt{1-x^2}}$ phương trình trở thành $t^2-(m+3) \cdot t+2 m+1=0$.\\
		Phương trình $9^{1+\sqrt{1-x^2}}-(m+3) \cdot 3^{1+\sqrt{1-x^2}}+2 m+1=0$ có nghiệm thực $\Leftrightarrow$ phương trình (1) có nghiệm $t \in[3 ; 9]$.\\
		(1) $\Leftrightarrow m=\dfrac{t^2-3 t+1}{t-2}\Leftrightarrow m=t-1-\dfrac{1}{t-2}$ vì $t-2>0$.\\
		Xét $f(t)=t-1-\dfrac{1}{t-2}$ liên tục trên đoạn [3;9] có $f^{\prime}(t)=1+\dfrac{1}{(t-2)^2}>0 \quad \forall t \in[3 ; 9] \Rightarrow f(t)$ đồng biến trên đoạn $[3 ; 9]$.\\
		Có $f(3)=1 ; f(9)=\dfrac{55}{7}$ Vậy phương trình $9^{1+\sqrt{1-x^2}}-(m+3) \cdot 3^{1+\sqrt{1-x^2}}+2 m+1=0$ có nghiệm thực $\Leftrightarrow m \in\left[1 ; \dfrac{55}{7}\right]$ $m \in \mathbb{Z}\Rightarrow$ Có 7 giá trị nguyên.
	}
\end{ex}
\begin{ex}[THPT Thăng Long 2019]%Câu 12
	Cho hệ phương trình $\left\{\begin{aligned}
		&2^{x-y}-2^y+x=2 y \\ 
		&2^x+1=m^2+2.2^y \cdot \sqrt{1-y^2}1\\
	\end{aligned}\right.$ $m$ là tham số. Gọi $S$ là tập các giá trị $m$ nguyên để hệ 1 có một nghiệm duy nhất. Tập $S$ có bao nhiêu phần tử?
	\choice
	{$0$}
	{\True $1$}
	{$3$}
	{$2$}
	\loigiai{
		Điều kiện $|y| \leq 1 \Leftrightarrow y \in-1 ; 1$.\\
		Từ phương trình thứ nhất của hệ 1 ta có $2^{x-y}+x-y=2^y+y 2$.\\
		Xét hàm số $y=f t=2^t+t$ với $t \in \mathbb{R}$.\\
		Dễ thấy $y'=2^t \cdot \ln 2+1>0$ với mọi $t \in \mathbb{R}$ nên hàm số $y=f t$ đồng biến trên $\mathbb{R}$.\\
		Do đó phương trình 2 tương đương với $x-y=y \Leftrightarrow x=2 y$.
		Thay $x=2 y$ vào phương trình thứ hai của hệ 1 ta được $4^y+1=m^2+2 \cdot 2^y \cdot \sqrt{1-y^2}3$.\\
		Để hệ đã cho có nghiệm duy nhất thì phương trình 3 phải có nghiệm duy nhất $y \in-1 ; 1$.\\
		Giả sử $y_0 \in-1 ; 1$ là một nghiệm của 3 thì $4^{y_0}+1=m^2+2 \cdot 2^{y_0}\cdot \sqrt{1-y_0^2}$.\\
		Khi đó $4^{-y_0}+1=m^2+2 \cdot 2^{-y_0}\cdot \sqrt{1--y_0^2}\Leftrightarrow 4^{y_0}+1=m^2+2 \cdot 2^{y_0}\cdot \sqrt{1-y_0^2}$ nên $-y_0$ cũng là nghiệm của 3.\\ 
		Suy ra $y_0=-y_0 \Leftrightarrow y_0=0$. Thay $y=0$ vào 3 ta được $m=0$.
		Thử lại: với $m=0$ thì 3 viết thành $4^y+1=2 \cdot 2^y \cdot \sqrt{1-y^2}\Leftrightarrow 2^y+\dfrac{1}{2^y}=2 \sqrt{1-y^2}4$.
		Ta có VT $4\geq 2$, dấu bằng khi $2^y=\dfrac{1}{2^y}\Leftrightarrow y=0 ; V P 4 \leq 2$, dấu bằng khi $y=0$.
		Suy ra phương trình 4 có nghiệm duy nhất là $y=0$. Vậy $m=0$ thỏa mãn.
	}
\end{ex}
\begin{ex}%Câu 13
	Cho $a, b$ là các số thực thỏa mãn $a>0$ và $a \neq 1$, biết phương trình $a^x-\dfrac{1}{a^x}=2 \cos (b x)$ có 7 nghiệm phân biệt. Tìm số nghiệm thực phân biệt của phương trình $a^{2 x}-2 a^x(\cos b x+2)+1=0$.
	\choice
	{$28$}
	{\True $14$}
	{$0$}
	{$7$}
	\loigiai{
		Ta có
		$\begin{aligned}
			& a^{2 x}-2 a^x(\cos b x+2)+1=0 \Leftrightarrow a^x-2+\dfrac{1}{a^x}=2(\cos b x+1) \\
			&\Leftrightarrow\left(a^{\frac{x}{2}}-\dfrac{1}{a^{\frac{x}{2}}}\right)^2=4 \cos ^2 \dfrac{b x}{2}\Leftrightarrow\left[\begin{aligned}
				&{a ^{\frac{x}{2}}- \dfrac{1}{a ^{\frac{x}{2}}}= 2 \cos\dfrac{b x}{2}}\\
				&{a ^{\dfrac{x}{2}}- \dfrac{1}{a ^{\frac{x}{2}}}= - 2 \cos\dfrac{b x}{2}}
			\end{aligned}\Leftrightarrow \left[\begin{aligned}
				&a^{\frac{x}{2}}-\dfrac{1}{a^{\frac{x}{2}}}=2 \cos b\left(\dfrac{x}{2}\right) \\
				&a^{\frac{-x}{2}}-\dfrac{1}{a^{\frac{-x}{2}}}=2 \cos b\left(\dfrac{-x}{2}\right)
			\end{aligned}\right.\right.
		\end{aligned}$\\
		Nếu phương trình (1) và phương trình $(2)$ có nghiệm chung là $x_0$ thì $2 \cos \dfrac{b x_0}{2}=-2 \cos \dfrac{b x_0}{2}$ $\Rightarrow \cos \dfrac{b x_0}{2}=0 \Rightarrow a^{\frac{x_0}{2}}-\dfrac{1}{a^{\frac{x_0}{2}}}=0 \Rightarrow x_0=0 \Rightarrow \cos \dfrac{b x_0}{2}=1$ (Vô lí).\\
		Do đó phương trình (1) và phương trình (2) không có nghiệm chung.
		Mặt khác theo giả thiết phương trình (1) và phương trình (2) dều có 7 nghiệm phân biệt.\\
		Vậy phương trình đã cho có 14 nghiệm phân biệt.
	}
\end{ex}
\begin{ex}%Câu 14
	Cho hàm số bậc ba $y=f(x)$ có bảng biến thiên như sau
	Giá trị lớn nhất của $m$ để phương trình $e^{2 f^3(x)-\dfrac{13}{2}f^2(x)+7 f(x)+\dfrac{3}{2}}=m$ có nghiệm trên đoạn $0 ; 2$ là
	\choice
	{\True $\mathrm{e}^4$}
	{$\mathrm{e}^3$}
	{$\mathrm{e}^{\dfrac{15}{13}}$}
	{$\mathrm{e}^5$}
	\loigiai{
		Giả sử $f'(x)=a(x-1)(x-3)=a\left(x^2-4 x+3\right)(a \neq 0)$
		$f(x)=\displaystyle\int\limits f'(x) \mathrm{d}x=\displaystyle\int\limits a\left(x^2-4 x+3\right) \mathrm{d}x=a\left(\dfrac{x^3}{3}-2 x^2+3 x\right)+\mathrm{C}$.\\
		Do $f(1)=1, f(3)=\dfrac{15}{13}$ nên ta có hệ $\left\{\begin{aligned}
			&\dfrac{4}{3}a+\mathrm{C}=1 \\ 
			&\mathrm{C}=\dfrac{15}{13}\\
		\end{aligned}\Leftrightarrow\left\{\begin{aligned}
			&a=-\dfrac{3}{26}\\ 
			&\mathrm{C}=\dfrac{15}{13}\\
		\end{aligned}\right.\right.$.\\
		Khi đó hàm số $f(x)=-\dfrac{1}{26}x^3+\dfrac{3}{13}x^2-\dfrac{9}{26}x+\dfrac{15}{13}$.\\
		Xét hàm số $f(x)$ trên đoạn $0 ; 2$
		$f'(x)=\dfrac{-3}{26}\left(x^2-4 x+3\right), f'(x)=0.\\ \Leftrightarrow\left[\begin{aligned}
			&x=1 \\
			&x=3 \notin 0 ; 2\\
		\end{aligned}\right.$.\\
		Ta có $f(0)=\dfrac{15}{13}, f(1)=1, f(2)=\dfrac{14}{13}$ nên $1 \leq f(x) \leq \dfrac{15}{13}$
		Đặt $f(x)=t$.\\
		Xét hàm số $g(t)=2 t^3-\dfrac{13}{2}t^2+7 t+\dfrac{3}{2}$ trên đoạn $\left[1 ; \dfrac{15}{13}\right]$
		$g'(t)=6 t^2-13 t+7, g'(t)=0 \Leftrightarrow\left[\begin{aligned}
			&t=1 \\
			&t=\dfrac{7}{6}\notin\left[1 ; \dfrac{15}{13}\right]\\
		\end{aligned}\right.$.\\
		Ta có $g(1)=4, g\left(\dfrac{15}{13}\right)=\dfrac{8778}{2197}$.\\
		Suy ra GTLN của $g(x)$ trên đoạn $\left[1 ; \dfrac{15}{13}\right]$ bằng 4 . Theo yêu cầu bài toán thì $m=\mathrm{e}^4$.
	}
\end{ex}
\begin{ex}
	[Hoàng Hoa Thám - Hưng Yên 2019]%Câu 15
	Cho phương trình $(4+\sqrt{15})^x+(2 m+1)(4-\sqrt{15})^x-6=0(m$ là tham số $)$. Biết phương trình có hai nghiệm phân biệt $x_1, x_2$ thỏa mãn $x_1-2 x_2=0$. Khi đó $m$ thuộc khoảng nào sau đây?
	\choice
	{\True $(3 ; 5)$}
	{$(-1 ; 1)$}
	{$(1 ; 3)$}
	{$(-\infty ;-1)$}
	\loigiai{
		Đặt $(4+\sqrt{15})^x=t(t>0) \Rightarrow(4-\sqrt{15})^x=\dfrac{1}{t}$.\\
		Phương trình trở thành: $t+\dfrac{2 m+1}{t}-6=0 \Leftrightarrow t^2-6 t+2 m+1=0 \quad(*)$
		$x_1-2 x_2=0 \Leftrightarrow(4+\sqrt{15})^{x_1-2 x_2}=1 \Leftrightarrow \dfrac{(4+\sqrt{15})^{x_1}}{(4+\sqrt{15})^{2 x_2}}=1 \Leftrightarrow \dfrac{t_1}{t_2^2}=1$.\\
		Do đó, phương trình đã cho có hai nghiệm phân biệt $x_1, x_2$ thỏa mãn $x_1-2 x_2=0$ khi và chỉ khi phương trình $(*)$ có hai nghiệm phân biệt dương $t_1, t_2$ thỏa mãn $t_1=t_2^2$.\\
		Tức là: $\left\{\begin{aligned}
			&\Delta'>0 \\ 
			&S>0 \\ 
			&P>0 \\ 
			&t_1=t_2^2\\
		\end{aligned}\right.$.\\
		$\Delta'>0 \Leftrightarrow 8-2 m>0 \Leftrightarrow m<4$.
		$S=6>0$ luôn đúng.\\
		$P=2 m+1>0 \Leftrightarrow m>-\dfrac{1}{2}$.\\
		Theo Vi-ét: $\left\{\begin{aligned}
			&t_1+t_2=6 \\ 
			&t_1 \cdot t_2=2 m+1\\
		\end{aligned}\right.$.\\
		Từ (1) và (3) suy ra: $t_2^3=2 m+1 \Leftrightarrow t_2=\sqrt[3]{2 m+1}\Rightarrow t_1=(\sqrt[3]{2 m+1})^2$.\\
		Thay vào (2) ta được $(\sqrt[3]{2 m+1})^2+\sqrt[3]{2 m+1}-6=0 \Leftrightarrow\left[\begin{aligned}
			&\sqrt[3]{2 m+1}=2 \\ 
			&\sqrt[3]{2 m+1}=-3\\
		\end{aligned}\Rightarrow\left[\begin{aligned}
			&m=\dfrac{7}{2}(\mathrm{tm}) \\ 
			&m=-14(\mathrm{ktm})\\
		\end{aligned}\right.\right.$.\\
		Vậy $m=\dfrac{7}{2}\in(3 ; 5)$.
	}
\end{ex}
\begin{ex}
	[THPT Minh Khai 2019]%Câu 16
	Gọi $S$ là tập hợp tất cả các giá trị nguyên của tham số $m$ để phương trình $5^x+10=m \sqrt{25^x+4}$ có nghiệm duy nhất. Số tập con của $S$ là
	\choice
	{$3$}
	{$4$}
	{\True $16$}
	{$15$}
	\loigiai{
		$5^x+10=m \sqrt{25^x+4}\Leftrightarrow \dfrac{5^x+10}{\sqrt{25^x+4}}=m$.\\
		Trường hợp 1: $m \leq 0$. Phương trình (1) vô nghiệm.\\
		Trường hợp 2: $m>0$. (1) $\Leftrightarrow \dfrac{\left(5^x+10\right)^2}{25^x+4}=m^2$.\\
		Đặt $t=5^x, t>0$. Ta có $\dfrac{(t+10)^2}{t^2+4}=m^2$.\\
		Xét hàm số $f(t)=\dfrac{(t+10)^2}{t^2+4}$ trên khoảng $(0 ;+\infty)$
		$f'(t)=\dfrac{-20 t^2-192 t+80}{\left(t^2+4\right)^2}\cdot f'(t)=0 \Leftrightarrow\left[\begin{aligned}
			&t=-10(\mathrm{l}) \\
			&t=\dfrac{2}{5}(\mathrm{tm})\\
		\end{aligned}.\right.$\\
		Bảng biến thiên:\\
		\textbf{(VẼ HÌNH)}\\
		Đề phương trình (1) có đúng một nghiệm $\Leftrightarrow$ Phương trình (2) có đúng một nghiệm $t>0$.\\ $\Leftrightarrow\left[\begin{aligned}
			&m^2=26 \\ 
			&1<m^2 \leq 25\\
		\end{aligned}\right.$. Do điều kiện $\left\{\begin{aligned}
			&m>0 \\ 
			&m \in \mathbb{Z}\\
		\end{aligned}\Rightarrow m \in\{2,3,4,5\}\right.$.\\
		Vậy $S=\{2,3,4,5\}$, do đó số tập con của $S$ là $2^4=16$.
	}
\end{ex}
\begin{ex}
	[Sở Quảng Trị 2019]%Câu 17
	Tìm tập hợp tất cả các giá trị tham số $m$ để phương trình $4^{x^2-2 x+1}-m \cdot 2^{x^2-2 x+2}+3 m-2=0$ có 4 nghiệm phân biệt.
	\choice
	{$(1 ;+\infty)$}
	{$(-\infty ; 1) \cup(2 ;+\infty)$}
	{\True $(2 ;+\infty)$}
	{$[2 ;+\infty)$}
	\loigiai{
		Xét phương trình: $4^{x^2-2 x+1}-m \cdot 2^{x^2-2 x+2}+3 m-2=0$
		(1).\\
		Đặt $t=2^{x^2-2 x+1}=2^{(x-1)^2}$. Do đó, ta có $(x-1)^2=\log _2 t$.\\
		Điều kiện $(t \geq 1)$.\\
		Ta có phương trình: (1) trở thành: $t^2-2 m t+3 m-2=0$.\\
		Ta nhận thấy mỗi giá trị $t>1$ cho hai giá trị $x$ tương ứng.\\
		Như vậy phương trình (1) có 4 nghiệm phân biệt khi và chỉ khi phương trình (2) có 2 nghiệm thỏa: $1<t_1<t_2$.\\
		$(2) \Leftrightarrow(2 t-3) m=t^2-2$.\\
		Nhận xét: $t=\dfrac{3}{2}$, không là nghiệm phương trình.\\
		Xét $t \neq \dfrac{3}{2},(2) \Leftrightarrow m=\dfrac{t^2-2}{2 t-3}$.\\ 
		Xét hàm $g(t)=\dfrac{t^2-2}{2 t-3}$ trên $(1 ;+\infty) \backslash\left\{\dfrac{3}{2}\right\}$ $g'(t)=\dfrac{2 t^2-6 t+4}{(2 t-3)^2}; g'(t)=0 \Leftrightarrow\left[\begin{aligned}
			&t=1 \\ 
			&t=2\\
		\end{aligned}\right.$.\\
		\textbf{(VẼ HÌNH)}\\
		Dựa vào bảng biến thiên, ta cần $m>2$.
	}
\end{ex}
\begin{ex}%Câu 18
	Cho phương trình: $2^{x^3+x^2-2 x+m}-2^{x^2+x}+x^3-3 x+m=0$. Tập các giá trị để bất phương trình có ba nghiệm phân biệt có dạng $(a ; b)$. Tổng $a+2 b$ bằng:
	\choice
	{$1$}
	{\True $2$}
	{$-4$}
	{$0$}
	\loigiai{
		Ta có $2^{x^3+x^2-2 x+m}-2^{x^2+x}+x^3-3 x+m=0 \Leftrightarrow 2^{x^3+x^2-2 x+m}+x^3+x^2-2 x+m=2^{x^2+x}+x^2+x(*)$.\\
		Xét hàm số $f(t)=2^t+t$ trên $\mathbb{R}$.\\
		Ta có $f'(t)=2^t \ln 2+1>0, \forall t \in \mathbb{R}\Rightarrow$.\\ Hàm số $f(t)$ đồng biến trên $\mathbb{R}$.\\
		$\begin{aligned}
			&\text{Mà}~(*)\Leftrightarrow f\left(x^3+x^2-2 x+m\right)=f\left(x^2+x\right) \Leftrightarrow x^3+x^2-2 x+m=x^2+x \\
			& \Leftrightarrow x^3-3 x+m=0 \Leftrightarrow m=-x^3+3 x(* *) .
		\end{aligned}$\\
		Xét hàm số $g(x)=-x^3+3 x$ trên $\mathbb{R}$.\\
		Ta có $g'(x)=-3 x^2+3$.\\
		$g'(x)=0 \Leftrightarrow x=\pm 1$.\\
		Bảng biến thiên:\\
		\textbf{(VẼ HÌNH)}\\
		Phương trình $2^{x^3+x^2-2 x+m}-2^{x^2+x}+x^3-3 x+m=0$ có 3 nghiệm phân biệt $\Leftrightarrow$ phương trình $(* *)$ có 3 nghiệm phân biệt $\Leftrightarrow-2<m<2 \Rightarrow\left\{\begin{aligned}
			&a=-2 \\ 
			&b=2\\
		\end{aligned}\Rightarrow a+2 b=2\right.$.
	}
\end{ex}
\begin{ex}
	[Chuyên ĐH Vinh- 2019]%Câu 19
	Có bao nhiêu số nguyên $m$ để phương trình $$9.3^{2 x}-m\left(4 \sqrt[4]{x^2+2 x+1}+3 m+3\right) 3^x+1=0$$ có đúng 3 nghiệm thực phân biệt?
	\choice
	{Vô số}
	{$3$}
	{\True $1$}
	{$2$}
	\loigiai{
		Ta có $9 \cdot 3^{2 x}-m\left(4 \sqrt[4]{x^2+2 x+1}+3 m+3\right) \cdot 3^x+1=0 \Leftrightarrow 3^{2(x+1)}-m(4 \sqrt{|x+1|}+3 m+3) \cdot 3^x+1=0$.\\
		$\Leftrightarrow 3^{x+1}+\dfrac{1}{3^{x+1}}=\dfrac{1}{3}m(4 \sqrt{|x+1|}+3 m+3)$.\\
		Đặt $x+1=t$ thì PT trở thành $3^t+\dfrac{1}{3^t}=\dfrac{1}{3}m(4 \sqrt{|t|}+3 m+3)(*)$.\\
		Vậy bài toán trở thành tìm $m$ để phương trình $(*)$ có đúng 3 nghiệm thực phân biệt.\\
		Nhận thấy, nếu $t_o$ là một nghiệm của $\left(^*\right)$ thì $-t_o$ cũng là nghiệm của $\left({}^*\right)$.\\
		Suy ra, điều kiện cần để phương trình $\left({}^*\right)$ có đúng 3 nghiệm thực phân biệt là $\left({}^*\right)$ có nghiệm.\\
		$\begin{aligned}
			&t=0.\\
			&\Rightarrow 1+1=\dfrac{1}{3}m \cdot(3 m+3) \Leftrightarrow m^2+m-2=0 \Leftrightarrow\left[\begin{aligned}
				&m=-2\\
				&m=1\\
			\end{aligned}\right. \text{.}
		\end{aligned}$
		Thử lại:\\
		Với $m=-2$ thì $\left(^*\right)$ trở thành $3^t+\dfrac{1}{3^t}=\dfrac{-2}{3}(4 \sqrt{|t|}-3)$.\\
		Nhận thấy, $V T=3^t+\dfrac{1}{3^t}\stackrel{\text{Caucchy}}{\geq}2, V P \leq \dfrac{-2}{3}\cdot(-3)=2 \Rightarrow~\mathrm{PT}$ có nghiệm duy nhất $t=0$ nên $m=-2$ không thỏa mãn.\\
		Với $m=1$ thì $(*)$ trở thành $3^t+\dfrac{1}{3^t}=\dfrac{1}{3}(4 \sqrt{|t|}+6) \Leftrightarrow 3^t+\dfrac{1}{3^t}-\dfrac{2}{3}(2 \sqrt{|t|}+3)=0 \quad(* *)$.\\ 
		Xét hàm số $f(t)=3^t+\dfrac{1}{3^t}-\dfrac{2}{3}(2 \sqrt{|t|}+3)$ với $t>0$.\\
		Ta có, $f'(t)=3^t \cdot \ln 3-\dfrac{\ln 3}{3^t}-\dfrac{2}{3 \sqrt{t}}; f''(t)=3^t \ln ^2 3+\dfrac{\ln ^2 3}{3^t}+\dfrac{1}{3 \sqrt{t^3}}>0$ với mọi $t>0$. $\Rightarrow f'(t)$ đồng biến trên $(0 ;+\infty) \Rightarrow f'(t)=0$ có nhiều nhất 1 nghiệm $t>0$.\\ $\Rightarrow f(t)=0$ có nhiều nhất 2 nghiệm $t>0$.\\
		Lại có, $f(1)=0$ và $f(0)=0 \Rightarrow$ Phương trình $(* *)$ có 3 nghiệm là $t=0, t=\pm 1$. Vậy $m=1$ thỏa mãn.
	}
\end{ex}
\begin{ex}
	Có bao nhiêu giá trị nguyên của tham số $m \in-2019 ; 2019$ để phương trình $2019^x+\dfrac{2 x-1}{x+1}+\dfrac{m x-2 m-1}{x-2}=0$ có đúng 3 nghiệm thực phân biệt?
	\choice
	{$4038$}
	{$2019$}
	{\True $2017$}
	{$4039$}
	\loigiai{
		Ta có phương trình $2019^x+\dfrac{2 x-1}{x+1}+\dfrac{m x-2 m-1}{x-2}=0 \Leftrightarrow 2019^x+\dfrac{2 x-1}{x+1}+\dfrac{m(x-2)-1}{x-2}=0$.\\
		$\Leftrightarrow 2019^x+\dfrac{2 x-1}{x+1}+m-\dfrac{1}{x-2}=0 \Leftrightarrow m=\dfrac{1}{x-2}-2019^x-\dfrac{2 x-1}{x+1}$.\\
		Xét hàm số
		$y=\dfrac{1}{x-2}-2019^x-\dfrac{2 x-1}{x+1}\Rightarrow y'=-\dfrac{1}{(x-2)^2}-2019^x \ln (2019)-\dfrac{3}{(x+1)^2}<0 ; \forall x \in \mathbb{R}\backslash\{-1 ; 2\}$.\\
		Ta có bảng biến thiên\\
		\textbf{(VẼ HÌNH)}\\
		Vậy để phương trình có 3 nghiệm phân biệt thì $m \in-\infty ;-2$ mà $m \in-2019 ; 2019 ; m \in \mathbb{Z}$. Vậy ta có $2017$ số nguyên $m$ cần tìm.
	}
\end{ex}
\begin{ex}
	[Chuyên Hưng Yên - 2020]%Câu 21
	Gọi $S$ là tập hợp các giá trị của tham số $m$ sao cho hai phương trình $2 x^2+1=3^m$ và $m=3^x-2 x^2+x-1$ có nghiệm chung. Tính tổng các phần tử của $S$.
	\choice
	{$6$}
	{\True $3$}
	{$1$}
	{$\dfrac{5}{2}$}
	\loigiai{
		Vì hai phương trình đã cho có nghiệm chung nên hệ sau có nghiệm\\
		$\begin{aligned}
			& \left\{\begin{aligned}
				&{2 x ^{2}+ 1 = 3 ^{m}}\\
				&{m = 3 ^{x}- 2 x ^{2}+ x - 1}\\
			\end{aligned}\Leftrightarrow \left\{\begin{aligned}
				&m=\log _3\left(2 x^2+1\right) \\
				&m=3^x-2 x^2+x-1\\
			\end{aligned}\Rightarrow \log _3\left(2 x^2+1\right)=3^x-2 x^2+x-1\right.\right. \\
			& \Leftrightarrow \log _3\left(2 x^2+1\right)+2 x^2+1=3^x+x \Leftrightarrow 3^{\log _3\left(2 x^2+1\right)}+\log _3\left(2 x^2+1\right)=3^x+x .
		\end{aligned}$\\
		Xét hàm số $f(t)=3^t+t$ xác định trên $\mathbb{R}\Rightarrow f'(t)=3^t \cdot \ln 3+1>0$ suy ra hàm $f(t)=3^t+t$ đồng biến trên $\mathbb{R}$ suy ra $\log _3\left(2 x^2+1\right)=x \Leftrightarrow 2 x^2+1=3^x$.\\
		Xét hàm số $g(x)=2 x^2+1-3^x$ xác định và liên tục trên $\mathbb{R}$.\\
		Ta có $g'(x)=4 x-3^x \ln 3 \Rightarrow g^{\prime \prime}(x)=4-3^x \ln ^2 3 \Rightarrow g^{\prime \prime \prime}(x)=-3^x \ln ^3 3<0$.\\ 
		Suy ra hàm số $g''(x)$ nghịch biến trên $\mathbb{R}$. Do đó $g(x)=0$ có nhiều nhất là 3 nghiệm.\\
		Ta lại có $g(0)=g(1)=g(2)=0$.\\ 
		Suy ra phương trình $2 x^2+1=3^x \Leftrightarrow\left[\begin{aligned}
			&x=0 \\ 
			&x=1 \\ 
			&x=2\\
		\end{aligned}\Rightarrow\left[\begin{aligned}
			&m=0 \\
			&m=1 \\ 
			&m=2\\
		\end{aligned}\right.\right.$.
		Vậy $S=3$.
	}
\end{ex}
\begin{ex}
	[Chuyên Nguyễn Bỉnh Khiêm - Quảng Nam - 2020]%Câu 22
	Giá trị của tham số $m$ để phương trình $4^x-m \cdot 2^{x+1}+2 m=0$ có hai nghiệm $x_1, x_2$ thỏa mãn $x_1+x_2=3$ là
	\choice
	{$m=2$}
	{$m=3$}
	{\True $m=4$}
	{$m=1$}
	\loigiai{
		Đặt $t=2^x, t>0$.\\
		Phương trình trở thành $t^2-2 m t+2 m=0\left({}^*\right)$.
		Để phương trình ban đầu có 2 nghiệm thì phương trình $\left(^*\right)$ phải có 2 nghiệm dương.\\
		$\Leftrightarrow\left\{\begin{aligned}
			&{\Delta '\geq 0}\\
			&{S > 0}\\
			&{P > 0}\\
		\end{aligned}\Leftrightarrow \left\{\begin{aligned}
			&m^2-2 m \geq 0 \\
			&2 m>0 \\
			&2 m>0\\
		\end{aligned}\Leftrightarrow m \geq 2 .\right.\right.$\\
		Ta có $x_1+x_2=3 \Leftrightarrow 2^{x_1+x_2}=2^3 \Leftrightarrow 2^{x_1}\cdot 2^{x_2}=8 \Leftrightarrow t_1 \cdot t_2=8 \Leftrightarrow 2 m=8 \Leftrightarrow m=4$.\\
		Kết luận $m=4$.
	}
\end{ex}
\begin{ex}
	[Chuyên Chu Văn An - 2020]%Câu 23
	Tìm $m$ để phương trình $4^x-2^{x+1}+m=0$ có hai nghiệm trái dấu.
	\choice
	{$m<0$}
	{$m>1$}
	{$-1<m<1$}
	{\True $0<m<1$}
	\loigiai{
		Đặt $t=2^x$, điều kiện: $t>0$.\\
		Phương trình trở thành: $t^2-2 t+m=0(*)$.\\
		Phương trình $4^x-2^{x+1}+m=0$ có 2 nghiệm phân biệt trái dấu
		$\Leftrightarrow t^2-2 t+m=0$ có 2 nghiệm dương phân biệt $t_1, t_2$ thỏa mãn: $t_1<1<t_2$.\\
		Phương trình $\left(^*\right)$ có 2 nghiệm dương phân biệt $\Leftrightarrow\left\{\begin{aligned}
			&\Delta'=1-m>0 \\ 
			&t_1 \cdot t_2=m>0\\
		\end{aligned}\Leftrightarrow 0<m<1\right.$.\\
		Phương trình có một nghiệm lớn hơn 1 và một nghiệm nhỏ hơn 1 khi và chỉ khi\\
		$\begin{aligned}
			& \left(t_1-1\right)\left(t_2-1\right)<0 \\
			& \Leftrightarrow t_1 t_2-\left(t_1+t_2\right)+1<0 \\
			& \Leftrightarrow m-2+1<0 \\
			& \Leftrightarrow m<1 .
		\end{aligned}$\\
		Kết hợp các điều kiện thì ta được $0<m<1$ thỏa mãn yêu cầu bài toán.
	}
\end{ex}
\begin{ex}
	[Chuyên Phan Bội Châu - Nghệ An - 2020]%Câu 24
	Có bao nhiêu giá trị nguyên dương của tham số $m$ để phương trình $9^x-2 \cdot 6^{x+1}+(m-3) \cdot 4^x=0$ có hai nghiệm phân biệt?
	\choice
	{\True $35$}
	{$38$}
	{$34$}
	{$33$}
	\loigiai{
		Phương trình tương đương $\left(\dfrac{3}{2}\right)^{2 x}-12 \cdot\left(\dfrac{3}{2}\right)^x+(m-3)=0$.\\
		Đặt $t=\left(\dfrac{3}{2}\right)^x, t>0$.\\
		Phương trình trở thành $t^2-12 . t+(m-3)=0, t>0(*)$.\\
		Phương trình đã cho có hai nghiệm phân biệt khi phương trình (*) có hai nghiệm phân biệt dương.\\
		$\Leftrightarrow\left\{\begin{aligned}
			&{\Delta '> 0}\\
			&{P > 0}\\
			&{S > 0}\\
		\end{aligned}\Leftrightarrow \left\{\begin{aligned}
			&{3 9 - m > 0}\\
			&{m - 3 > 0}\\
			&{1 2 > 0}\\
		\end{aligned}\Leftrightarrow \left\{\begin{aligned}
			&m<39 \\
			&m>3\\
		\end{aligned}\Leftrightarrow 3<m<39 .\right.\right.\right.$.\\
		Vậy có 35 giá trị nguyên dương của tham số $m$.
	}
\end{ex}
\begin{ex}
	[ĐHQG Hà Nội - 2020]%Câu 25
	Gọi $S$ là tập hợp các số nguyên $m$ sao cho phương trình $4^x-m \cdot 2^{x+1}+3 m^2-500=0$ có 2 nghiệm phân biệt. Hỏi tập $S$ có bao nhiêu phần tử
	\choice
	{$1$}
	{$4$}
	{\True $3$}
	{$2$}
	\loigiai{
		Đặt $t=2^x~t>0$ khi đó phương trình $4^x-m \cdot 2^{x+1}+3 m^2-500=0~(1)$ trở thành: $t^2-2 m \cdot t+3 m^2-500=0 (2)$.\\
		Để 1 có 2 nghiệm phân biệt khi và chỉ khi 2 có 2 nghiệm dương phân biệt hay\\ 
		$\left\{\begin{aligned}
			&\Delta>0 \\ 
			&P>0 \\ 
			&S>0\\
		\end{aligned}\Leftrightarrow\left\{\begin{aligned}
			&m^2-3 m^2+500>0 \\ 
			&3 m^2-500>0 \\ 
			&2 m>0\\
		\end{aligned}\right.\right.$.\\ $\Leftrightarrow\left\{\begin{aligned}
			&-5 \sqrt{10}<m<5 \sqrt{10}\\ 
			&m<\dfrac{-10 \sqrt{15}}{3}\vee m>\dfrac{10 \sqrt{15}}{3}\Leftrightarrow \dfrac{10 \sqrt{15}}{3}<m<5 \sqrt{10}. \\ 
			&m>0
		\end{aligned}\right.$.\\
		Vậy tập hợp các số nguyên $m$ là $S=13 ; 14 ; 15$.
	}
\end{ex}
\begin{ex}
	[ĐHQG Hà Nội - 2020]%Câu 26
	Tìm điều kiện của tham số $a$ để phương trình sau có nghiệm: $9^{1+\sqrt{1-x^2}}-(a+2) \cdot 3^{1+\sqrt{1-x^2}}+2 a+1=0$. Hãy chọn đáp án đúng nhât?
	\choice
	{\True $4 \leq a \leq \dfrac{64}{7}$}
	{$2 \leq a \leq \dfrac{64}{9}$}
	{$3 \leq a \leq \dfrac{50}{3}$}
	{$1 \leq a \leq \dfrac{50}{3}$}
	\loigiai{
		Đặt $t=3^{1+\sqrt{1-x^2}}$ vì $0 \leq \sqrt{1-x^2}\leq 1 \Rightarrow 3 \leq t \leq 9$.\\
		Khi đó bài toán trở thành tìm điều kiện của tham số $a$ để phương trình $t^2-(a+2) \cdot t+2 a+1=0(*)$ có nghiệm trên đoạn $[3 ; 9]$.\\
		Ta có $(*) \Leftrightarrow t^2-2 t+1=a(t-2)$.\\
		Vì $t=2$ không phải nghiệm của phương trình nên $(*) \Leftrightarrow \dfrac{t^2-2 t+1}{t-2}=a$.\\
		Xét $f(t)=\dfrac{t^2-2 t+1}{t-2}\Rightarrow f'(t)=\dfrac{t^2-4 t+3}{t-2}; f'(t)=0 \Leftrightarrow\left[\begin{aligned}
			&t=1 \notin[3 ; 9] \\ 
			&t=3 \in[3 ; 9]\\
		\end{aligned}\right.$.\\
		Ta có $f(3)=4 ; f(9)=\dfrac{64}{7}\Rightarrow 4 \leq a \leq \dfrac{64}{7}$ thì phương trình bài ra có nghiệm.
	}
\end{ex}
\begin{ex}
	[ĐHQG Hà Nội - 2020]%Câu 27
	Điều kiện của $m$ để hệ bất phương trình $\left\{\begin{aligned}
		&7^{2 x+\sqrt{x+1}}-7^{2+\sqrt{x+1}}+2020 x \leq 2020 \\ &x^2-(m+2) x+2 m+3 \geq 0\\
	\end{aligned}\right.$ có nghiệm là
	\choice
	{$m \geq-3$}
	{$-2 \leq m \leq 1$}
	{$-1 \leq m \leq 2$}
	{\True $m \geq-2$}
	\loigiai{
		$7^{2 x+\sqrt{x+1}}-7^{2+\sqrt{x+1}}+2020 x \leq 2020 \Leftrightarrow 7^{2 x+\sqrt{x+1}}+1010 \cdot(2 x+\sqrt{x+1}) \leq 7^{2+\sqrt{x+1}}+1010 .(2+\sqrt{x+1})$.\\
		Hàm số $f(t)=7^t+1010 . t$ đồng biến trên $\mathbb{R}$.\\
		$(*) \Leftrightarrow f(2 x+\sqrt{x+1}) \leq f(2+\sqrt{x+1})$.\\
		Suy ra: $2 x+\sqrt{x+1}\leq 2+\sqrt{x+1}\Rightarrow-1 \leq x \leq 1$.\\
		$\begin{aligned}
			& x \in[-1 ; 1]: x^2-(m+2) x+2 m+3 \geq 0 \Leftrightarrow m \geq \dfrac{x^2-2 x+3}{x-2}. \\
			& \text{Ycbt}\Leftrightarrow \exists x \in[-1 ; 1]: m \geq \dfrac{x^2-2 x+3}{x-2}
		\end{aligned}$.\\
		Từ bảng biến thiên ta có, $(* *) \Leftrightarrow m \geq-2$.
	}
\end{ex}
\begin{ex}
	[Sở Phú Thọ - 2020]%Câu 28
	Cho phương trình $16^{x^2}-2 \cdot 4^{x^2+1}+10=m$ ( $m$ là tham số). Số giá trị nguyên của tham $m \in[-10 ; 10]$ để phương trình đã cho có đúng hai nghiệm thực phân biệt là
	\choice
	{$7$}
	{$9$}
	{\True $8$}
	{$1$}
	\loigiai{
		Đặt $t=4^{x^2}, t \geq 1$.\\
		Khi đó phương trình đã cho trở thành $t^2-8 t+10=m(1)$.\\
		Nghiệm $t=1$ cho một nghiệm $x=0$.\\
		Mỗi nghiệm $t>1$ cho hai nghiệm $x$ đối nhau.\\
		Do vậy phương trình đã cho có hai nghiệm thực phân biệt khi và chỉ khi phương trình (1) có đúng một nghiệm $t>1$, nghiệm còn lại (nếu có) phải nhỏ hơn 1.\\
		Xét hàm số $f(t)=t^2-8 t+10$.\\
		Bảng biến thiên
		\begin{center}
			\begin{tikzpicture}
				\tkzTabInit[nocadre,lgt=1.2,espcl=3]{$t$/0.8,$f'(t)$/0.8,$f(t)$/2}{$1$,$4$,$+\infty$}
				\tkzTabLine{,-,0,+,}
				\tkzTabVar{+/$3$,-/$-6$,+/$+\infty$}
			\end{tikzpicture}
		\end{center}
		Từ bảng biến thiên ta có phương trình (1) có một nghiệm lớn hơn 1 khi $\left[\begin{aligned}
			&m>3 \\ 
			&m=-6\\
		\end{aligned}\right.$.\\
		Suy ra số giá trị nguyên $m \in[-10 ; 10]$ là 8.
	}
\end{ex}
\begin{ex}
	[Sở Hà Tĩnh - 2020]%Câu 29
	Gọi $S$ là tập nghiệm của phương trình $\left(2^x-2 x\right) \sqrt{(3)^{2^x}-m}=0$ (với $m$ là tham số thực). Có tất cả bao nhiêu giá trị nguyên của $m \in[-2020 ; 2020]$ để tập hợp $S$ có hai phần tử?
	\choice
	{\True $2094$}
	{$2092$}
	{$2093$}
	{$2095$}
	\loigiai{
		Gọi $\mathscr{D}$ là tập xác định của phương trình đã cho.\\
		Nếu $m \leq 1$ thì $3^{2^x}-m>0 \forall x \in \mathbb{R}$ nên $D=\mathbb{R}$.\\
		Nếu $m>1$ thì $\mathscr{D}=\left[\log _2\left(\log _3 m\right) ;+\infty\right)$.\\
		$\left(2^x-2 x\right) \sqrt{(3)^{2^x}-m}=0 \Leftrightarrow\left[\begin{aligned}
			&2^x-2 x=0(1) \\
			&(3)^{2^x}-m=0(2)\\
		\end{aligned}.\right.$.\\
		Xét hàm số $f(x)=2^x-2 x$ có $f'(x)=2^x \ln 2-2 ; f'(x)=0 \Leftrightarrow x=\dfrac{2}{\ln 2}$ do đó phương trình $f(x)=0$ có không quá 2 nghiệm.\\
		Mặt khác $f(1)=0 ; f(2)=0$ nên $(1) \Leftrightarrow\left[\begin{aligned}
			&x=1 \\ 
			&x=2\\
		\end{aligned}\right.$.\\
		Lại có với $m>1,(2) \Leftrightarrow x=\log _2\left(\log _3 m\right)$.\\
		Nếu $m \leq 1$ thì $S=\{1 ; 2\}$ (thỏa mãn yêu cầu bài toán).\\
		Nếu $m>1$ thì $S$ có hai phần tử khi và chỉ khi $1 \leq \log _2\left(\log _3 m\right)<2 \Leftrightarrow 9 \leq m<81$.\\
		Vậy $S$ có hai phần tử khi và chỉ khi $\left[\begin{aligned}
			&m \leq 1 \\ 
			&9 \leq m<81\\
		\end{aligned}\right.(*)$.\\ 
		Số các giá trị nguyên của $m \in[-2020 ; 2020]$ thỏa mãn $(*)$ là $1+2020+1+81-9=2094$.
	}
\end{ex}
\begin{ex}
	[Sở Ninh Bình 2020]%Câu 30
	Cho hai số thực bất kỳ $a>1, b>1$. Gọi $x_1, x_2$ là hai nghiệm phương trình $a^x b^{x^2-1}=1$. Trong trường hợp biểu thức $S=\left(\dfrac{x_1 x_2}{x_1+x_2}\right)^2-6 x_1-6 x_2$ đạt giá trị nhỏ nhất, khẳng định nào dưới đây đúng?
	\choice
	{$a=b^{\sqrt[3]{3}}$}
	{$a=b^{\sqrt[3]{6}}$}
	{\True $a=b^{\sqrt[3]{\dfrac{1}{3}}}$}
	{$a=b^{\sqrt[3]{\dfrac{1}{6}}}$}
	\loigiai{
		$\begin{aligned}
			&\text{Ta có}~a^x b^{x^2-1}=1 \Leftrightarrow \dfrac{a^x b^{x^2}}{b}=1 \Leftrightarrow a^x b^{x^2}=b \Leftrightarrow \ln \left(a^x b^{x^2}\right)=\ln b \\
			& \Leftrightarrow \ln a^x+\ln b^{x^2}=\ln b \Leftrightarrow x^2 \ln b+x \ln a-\ln b=0 .
		\end{aligned}$\\
		Phương trình đã cho có hai nghiệm phân biệt $\left\{\begin{aligned}
			&\ln b \neq 0 \forall b>1 \\ 
			&\Delta=\ln ^2 a+4 \ln ^2 b \geq 0 \forall a>1, b>1
		\end{aligned}\right.$.\\
		Vậy phương trình đã cho luôn có hai nghiệm $x_1, x_2$.\\
		Theo định lý Viét ta có $\left\{\begin{aligned}
			&x_1+x_2=-\dfrac{\ln a}{\ln b}=-\log _b a \\ 
			&x_1 \cdot x_2=-\dfrac{\ln b}{\ln b}=-1\\
		\end{aligned}\right.$.\\ 
		Khi đó ta có $S=\left(\dfrac{x_1 x_2}{x_1+x_2}\right)^2-6 x_1-6 x_2=\left(\dfrac{-1}{-\log _b a}\right)^2+6 \log _b a=\dfrac{1}{\log _b^2 a}+6 \log _b a$.\\
		Do $a>1, b>1 \Rightarrow \log _b a>\log _b 1=0$.\\
		Áp dụng bất đẳng thức Cauchy ta có\\
		$S=\dfrac{1}{\log _b^2 a}+6 \log _b a=\dfrac{1}{\log _b^2 a}+3 \log _b a+3 \log _b a \geq 3 \sqrt[3]{\dfrac{1}{\log _b^2 a}\cdot 3 \log _b a \cdot 3 \log _b a}=3 \sqrt[3]{9}$.\\
		Dấu ``$=$'' xảy ra khi và chỉ khi $\dfrac{1}{\log _b^2 a}=3 \log _b a \Leftrightarrow \log _b^3 a=\dfrac{1}{3}\Leftrightarrow \log _b a=\sqrt[3]{\dfrac{1}{3}}\Leftrightarrow a=b^{\sqrt[3]{\frac{1}{3}}}$.
		Vậy khẳng định đúng là $a=b^{\sqrt[3]{\frac{1}{3}}}$.
	}
\end{ex}
\begin{ex}
	[Sở Bắc Ninh - 2020]%Câu 31
	Gọi $S$ là tập tất cả các giá trị của $m$ để phương trình $16^x-6.8^x+8.4^x-m \cdot 2^{x+1}-m^2=0$ có đúng hai nghiệm phân biệt. Khi đó $S$ có
	\choice
	{$4$ tập con}
	{Vô số tập con}
	{$8$ tập con}
	{\True $16$ tập con}
	\loigiai{
		Đặt $t=2^x,(t>0)$, phương trình đã cho trở thành $t^4-6 t^3+8 t^2-2 m t-m^2=0(*), t>0$.\\
		Phương trình đã cho có đúng hai nghiệm phân biệt khi và chỉ khi phương trình $(*)$ có đúng hai nghiệm dương phân biệt.\\
		$(*) \Leftrightarrow\left(t^4-6 t^3+9 t^2\right)-\left(t^2+2 m t+m^2\right)=0 \Leftrightarrow\left[\begin{aligned}
			&m=t^2-4 t(1) \\
			&m=-t^2+2 t(2)
		\end{aligned}.\right.$.\\
		Xét hai hàm số $f(t)=t^2-4 t ; g(t)=-t^2+2 t$ trên khoảng $(0 ;+\infty)$ có đồ thị như sau\\
		\textbf{(VẼ HÌNH)}\\
		Dựa vào đồ thị hai hàm số này ta suy ra phương trình $\left({}^*\right)$ có đúng hai nghiệm dương phân biệt khi và chỉ khi $m \in\{0 ; 1 ;-3 ;-4\}$ hay $S$ có 4 phần tử.\\
		Vậy $S$ có $2^4=16$ tập con.
	}
\end{ex}
\begin{ex}
	[Lý Nhân Tông - Bắc Ninh - 2020]%Câu 32
	Tìm tập hợp các giá trị của tham số thực $m$ để phương trình $6^x+(3-m) 2^x-m=0$ có nghiệm thuộc khoảng $(0 ; 1)$.
	\choice
	{$[3 ; 4]$}
	{$[2 ; 4]$}
	{\True $(2 ; 4)$}
	{$(3 ; 4)$}
	\loigiai{
		Phương trình $6^x+(3-m) 2^x-m=0 \Leftrightarrow m=\dfrac{6^x+3 \cdot 2^x}{1+2^x}$.\\
		Xét hàm số $f(x)=\dfrac{6^x+3 \cdot 2^x}{1+2^x}$ liên tục trên $(0 ; 1)$.\\
		Ta có $f'(x)=\dfrac{12^x \ln 3+6^x \ln 6+3 \cdot 2^x \ln 2}{\left(1+2^x\right)^2}>0, \forall x \in(0 ; 1)$.\\ 
		Suy ra hàm số $f(x)=\dfrac{6^x+3 \cdot 2^x}{1+2^x}$ đồng biến trên $(0 ; 1)$.\\
		Do đó phương trình $6^x+(3-m) 2^x-m=0$ có nghiệm thuộc khoảng $(0 ; 1)$ khi và chỉ khi $f(0)<m<f(1)$, tức là $2<m<4$.
	}
\end{ex}
\begin{ex}
	[Nguyễn Huệ - Phú Yên - 2020]%Câu 33
	Có bao nhiêu giá trị nguyên $m \in(-2019 ; 2020)$ sao cho hệ phương trình sau có nghiệm
	$\left\{\begin{aligned}
		&4+9 \cdot 3^{x^2-2 y}=\left(4+9^{x^2-2 y}\right) \cdot 7^{2 y-x^2+2}\\
		&2 x-1=\sqrt{2 y-2 x+m}\\
	\end{aligned}?\right.$
	\choice
	{\True $2017$}
	{$2021$}
	{$2019$}
	{$2020$}
	\loigiai{
		Xét phương trình: $4+9 \cdot 3^{x^2-2 y}=\left(4+9^{x^2-2 y}\right) \cdot 7^{2 y-x^2+2}$.\\
		Đặt $t=x^2-2 y$, phương trình trở thành: $4+9.3^t=\left(4+9^t\right) \cdot 7^{2-t}\Leftrightarrow 4.7^t+9.3^t \cdot 7^t=4.49+49.3^{2 t}$
		$\Leftrightarrow 4\left(7^t-7^2\right)=3^t\left(3^t .7^2-7^t .3^2\right)(*)$.\\
		Giả sử $3^t .7^2-7^t .3^2<0 \Leftrightarrow\left(\dfrac{3}{7}\right)^t<\left(\dfrac{3}{7}\right)^2 \Leftrightarrow t>2$.\\
		Nếu $t>2 \Rightarrow\left\{\begin{array}{l}V T\left(^*\right)>0 \\ V P(*)<0\end{array}\Rightarrow(*)\right.$ vô nghiệm.\\
		Nếu $t<2 \Rightarrow\left\{\begin{aligned}
			&\mathrm{VT}(*)<0 \\ 
			&\mathrm{VP}(*)>0\\
		\end{aligned}\Rightarrow(*)\right.$ vô nghiệm.\\
		Nếu $t=2 \Rightarrow V T\left({}^*\right)=V P(*) \Rightarrow\left({}^*\right)$ có nghiệm duy nhất $t=2 \Rightarrow x^2-2 y=2 \Rightarrow 2 y=x^2-2$.\\
		Ta được: $2 x-1=\sqrt{x^2-2 x-2+m}\Leftrightarrow\left\{\begin{aligned}
			&3 x^2-2 x+3=m \\ 
			&x \geq \dfrac{1}{2}\\
		\end{aligned}\right.$.\\
		Xét hàm số $f(x)=3 x^2-2 x+3$, với $x \in\left[\dfrac{1}{2};+\infty\right) \Rightarrow f'(x)=6 x-2>0, \forall x \geq \dfrac{1}{2}$, suy ra hàm số $f(x)$ đồng biến trên khoảng $\left(\dfrac{1}{2};+\infty\right) \Rightarrow f(x) \geq f\left(\dfrac{1}{2}\right)=\dfrac{11}{4}\Rightarrow(1)$ có nghiệm $x \in\left[\dfrac{1}{2};+\infty\right)$ khi $m \geq \dfrac{11}{4}\Rightarrow m \in\left[\dfrac{11}{4}; 2020\right)$.\\ Vì $m$ nguyên nên $m \in\{3 ; 4 ; 5 ; \ldots ; 2019\}$.\\
		Vậy có 2017 giá trị của $m$.
	}
\end{ex}
\begin{ex}
	[Nguyễn Trãi - Thái Bình - 2020]%Câu 34
	Tính tổng tất cả các nghiệm của phương trình $\mathrm{e}^{\sin \left(x-\frac{\pi}{4}\right)}=\tan x$ thuộc đoạn $[0 ; 50 \pi]$
	\choice
	{$\dfrac{2671 \pi}{2}$}
	{$\dfrac{1853 \pi}{2}$}
	{\True $\dfrac{2475 \pi}{2}$}
	{$\dfrac{2653 \pi}{2}$}
	\loigiai{
		Điều kiện: $\cos x \neq 0$. Nhận thấy $\mathrm{e}^{\sin \left(x-\frac{\pi}{4}\right)}>0 \forall x \in \mathbb{R} \Rightarrow \tan x>0$.\\
		Ta có $\mathrm{e}^{\sin \left(x-\frac{\pi}{4}\right)}=\tan x \Leftrightarrow \mathrm{e}^{\frac{1}{\sqrt{2}}(\sin x-\cos x)}=\dfrac{\sin x}{\cos x}\Leftrightarrow \dfrac{\mathrm{e}^{\frac{\sin x}{\sqrt{2}}}}{\mathrm{e}^{\frac{\cos x}{\sqrt{2}}}}=\dfrac{\sin x}{\cos x}\Leftrightarrow \dfrac{\mathrm{e}^{\frac{\sin x}{\sqrt{2}}}}{\sin x}=\dfrac{\mathrm{e}^{\frac{\cos x}{\sqrt{2}}}}{\cos x}\left(^*\right)$.\\
		Xét hàm số $f(t)=\dfrac{\mathrm{e}^{\frac{t}{\sqrt{2}}}}{t}, t \in(-1 ; 0) \cup(0 ; 1)$ có
		$f'(t)=\dfrac{\mathrm{e}^{\frac{t}{\sqrt{2}}}(\sqrt{2}t-2)}{2 t^2}<0, \forall t \in(-1 ; 0) \cup(0 ; 1)$
		$\Rightarrow f(t)$ nghịch biến trên khoảng $(-1 ; 0)$ và $(0 ; 1)$.\\
		Bảng biến thiên:\\
		\textbf{(VẼ HÌNH)}\\
		Từ bảng biến thiên ta thấy: $f(-1)=-\mathrm{e}^{-\frac{1}{\sqrt{2}}}<0, f(1)=\mathrm{e}^{\frac{1}{\sqrt{2}}}>0$.\\
		Do đó từ $\left(^*\right)$ ta có $f(\sin x)=f(\cos x) \Leftrightarrow \sin x=\cos x \Leftrightarrow x=\dfrac{\pi}{4}+k \pi, k \in \mathbb{Z}$.\\
		Theo giả thiết $x \in[0 ; 50 \pi] \Rightarrow 0 \leq \dfrac{\pi}{4}+k \pi \leq 50 \pi \Leftrightarrow-\dfrac{1}{4}\leq k \leq \dfrac{199}{4}(* *)$.\\
		Do $k \in \mathbb{Z}$ nên từ $(* *)$ suy ra $k \in 0 ; 1 ; \ldots ; 49$, có 50 giá trị $k$ thỏa mãn.\\
		Vậy tổng tất cả các nghiệm của phương trình trên đoạn $[0 ; 50 \pi]$ là\\
		$S=\sum\limits_{k=0}^{49}\left(\dfrac{\pi}{4}+k \pi\right)=\dfrac{2475 \pi}{2}$.
	}      
\end{ex}
\begin{ex}
	[Trần Phú - Quảng Ninh - 2020]%Câu 35
	Tìm tập hợp các giá trị của tham số $m$ để phương trình (ẩn $x): 3^{\log _2 x^2}-2(m+3) \cdot 3^{\log _2 x}+m^2+3=0$ có hai nghiệm phân biệt thỏa mãn: $x_1 x_2>2$.
	\choice
	{\True $(-1 ;+\infty) \backslash\{0\}$}
	{$(0 ;+\infty)$}
	{$\mathbb{R}\backslash[-1 ; 1]$}
	{$(-1 ;+\infty)$}
	\loigiai{
		Điều kiện xác định: $x>0$.\\
		Ta có $3^{\log _2 x^2}-2(m+3) \cdot 3^{\log _2 x}+m^2+3=0(1)$
		$\Leftrightarrow\left(3^{\log _2 x}\right)^2-2(m+3) \cdot 3^{\log _2 x}+m^2+3=0$.\\
		Đặt: $t=3^{\log _2 x}(t>0)$
		$\begin{aligned}
			&\Rightarrow \log _2 x=\log _3 t \\
			&\Leftrightarrow x=2^{\log _3 t}\\
			&\text{Khi đó:}~t^2-2(m+3) t+m^2+3=0~(2)\\
		\end{aligned}$\\
		Phương trình (1) có hai nghiệm phân biệt
		$\Leftrightarrow$ Phương trình (2) có hai nghiệm phân biệt dương\\ $t_1 ; t_2$
		$\Leftrightarrow\left\{\begin{aligned}
			&{\Delta '> 0}\\
			&{S > 0}\\
			&{P > 0}\\
		\end{aligned}\Leftrightarrow \left\{\begin{aligned}
			&{( m + 3 ) ^{2}- ( m ^{2}+ 3 ) > 0}\\
			&{2 ( m + 3 ) > 0}\\
			&{m ^{2}+ 3 > 0}\\
		\end{aligned}\Leftrightarrow \left\{\begin{aligned}
			&m>-1 \\
			&m>-3\\
		\end{aligned}\Leftrightarrow m>-1 .\right.\right.\right.$.\\
		Theo hệ thức Vi-et, ta có $\left\{\begin{aligned}
			&t_1+t_2=2(m+3) \\ 
			&t_1 \cdot t_2=m^2+3\\
		\end{aligned}\right.$.\\
		Ta có $x_1 \cdot x_2>2$
		$\begin{aligned}
			&\Leftrightarrow 2^{\log _3 t_1}\cdot 2^{\log _3 t_2}>2 \\
			&\Leftrightarrow 2^{\log _3\left(t_1 \cdot t_2\right)}>2 \\
			&\Leftrightarrow 2^{\log _3\left(m^2+3\right)}>2 \\
			&\Leftrightarrow \log _3\left(m^2+3\right)>1 \\
			&\Leftrightarrow m^2+3>3 \\
			&\Leftrightarrow m^2>0 \\
			&\Leftrightarrow m \neq 0 \\
			&\text{Vậy}~\left\{\begin{aligned}
				&m>-1 \\
				&m \neq 0\\
			\end{aligned}\right.
		\end{aligned}$.
	}
\end{ex}
\begin{ex}
	[THPT PTNK Cơ sở 2 - TP.HCM - 2021]%Câu 36
	Xét các số nguyên dương $a, b$ sao cho phương trình $a \cdot 4^x-b \cdot 2^x+50=0$ (1) có hai nghiệm phân biệt $x_1, x_2$ và phương trình $9^x-b \cdot 3^x+50 a=0$ (2) có hai nghiệm $x_3, x_4$ thỏa mãn điều kiện $x_3+x_4>x_1+x_2$. Tìm giá trị nhỏ nhất của biểu thức $S=3 a+4 b$.
	\choice
	{\True $109$}
	{$51$}
	{$49$}
	{$87$}
	\loigiai{
		Đặt $u=2^x>0$ thì phương trình (1) trở thành $a \cdot u^2-b \cdot u+50=0$ (3).\\ 
		Phương trình (1) có hai nghiệm phân biệt $x_1, x_2$ tương đương phương trình (3) có hai nghiệm phân biệt dương $u_1, u_2$, nghĩa là $\left\{\begin{aligned}
			&b^2-200 a>0 \\ 
			&\dfrac{b}{a}>0 \\ 
			&\dfrac{50}{a}>0\\
		\end{aligned}\Leftrightarrow b^2>200 a\right.$ (vì $a, b$ là số nguyên dương).\\
		Đặt $v=3^x>0$ thì phương trình (2) trở thành $v^2-b \cdot v+50 a=0$ (4).\\ 
		Phương trình (2) có hai nghiệm phân biệt $x_3, x_4$ tương đương với phương trình (4) có hai nghiệm phân biệt dương $v_3, v_4$, nghĩa là $\left\{\begin{aligned}
			&b^2-200 a>0 \\ 
			&b>0 \\ 
			&50 a>0\\
		\end{aligned}\Leftrightarrow b^2>200 a\right.$ (vì $a, b$ là số nguyên dương).\\
		Ta có $x_3+x_4>x_1+x_2 \Leftrightarrow \log _3 v_3+\log _3 v_4>\log _2 u_1+\log _2 u_2 \Leftrightarrow \log _3\left(v_3 v_4\right)>\log _2\left(u_1 u_2\right)$ $\Leftrightarrow \log _3(50 a)>\log _2\left(\dfrac{50}{a}\right) \Leftrightarrow \log _3 a+\log _2 a>\log _2 50-\log _3 50 \approx 2,08$.\\
		Mặt khác hàm số $f(a)=\log _3 a+\log _2 a(a>0)$ là hàm số tăng, $f(2) \approx 1,63$ và $f(3) \approx 2,58$ nên $a \geq 3$.\\ 
		Từ đó ta có $b^2>200 a \geq 600 \Rightarrow b \geq 25$. Vậy $\min S=3.3+4.25=109$.
	}
\end{ex}
\begin{ex}
	[Sở Lào Cai - 2021]%Câu 37
	Tìm tất cả các giá trị của tham số $a$ để phương trình $e^{e^{2 x}-a}-2 x-a=0$ có nhiều nghiệm nhất?
	\choice
	{$a \geq 0$}
	{\True $a>1$}
	{$a<e$}
	{$a \geq-1$}
	\loigiai{
		Phương trình đã cho tương đương với: $e^{e^{2 x}-a}-2 x-a=0$.\\
		Đặt $t=e^{2 x}-a$, ta có hệ: $\left\{\begin{aligned}
			&\mathrm{e}^t-2 x-a=0 \\ 
			&\mathrm{e}^{2 x}-t-a=0\\
		\end{aligned}\right.$.\\
		Suy ra: $e^t+t=e^{2 x}+2 x(*)$.\\
		Xét hàm số $f(u)=e^u+u$, ta có: $f'(u)=e^u+1>0 \quad \forall u \in \mathbb{R}$.\\
		Suy ra hàm số $f(u)$ đồng biến trên $\mathbb{R}$.\\
		Ta có $(*) \Leftrightarrow f(t)=f(2 x) \Leftrightarrow t=2 x$.\\
		Với $t=2 x$, suy ra: $a=e^{2 x}-2 x$.\\
		Xét hàm số $g(x)=e^{2 x}-2 x, g'(x)=2 e^{2 x}-2$.\\
		$g'(x)=0 \Leftrightarrow 2 e^{2 x}-2=0 \Leftrightarrow x=0$\\
		Bảng biến thiên:\\
		\textbf{(VẼ HÌNH)}\\
		Từ bảng biến thiên suy ra phương trình đã cho có nhiều nghiệm nhất bằng 2 khi $a>1$.
	}
\end{ex}
\begin{ex}
	[THPT Triệu Sơn - Thanh Hóa - 2021]%Câu 38
	Số giá trị nguyên của tham số $m$ để phương trình $(m+1) \cdot 16^x-2(2 m-3) \cdot 4^x+6 m+5=0$ có hai nghiệm trái dấu là:
	\choice
	{$4$}
	{$8$}
	{\True $2$}
	{$1$}
	\loigiai{
		Đặt $4^x=t(t>0)$. Phương trình trở thành: $(m+1) t^2-2(2 m-3) t+6 m+5=0$ $(1)$.\\
		Yêu cầu bài toán $\Leftrightarrow$ Phương trình (1) có hai nghiệm phân biệt $t_1 ; t_2$ thỏa mãn $0<t_1<1<t_2$.\\
		$\begin{aligned}
			& \Leftrightarrow\left\{\begin{aligned}
				&{m \neq - 1}\\
				&{\Delta '= - 2 m ^{2}- 2 3 m + 4 > 0}\\
				&{t _{1}t _{2}> 0}\\
				&{t _{1}+ t _{2}> 0}\\
				&{( t _{1}- 1 ) ( t _{2}- 1 ) < 0}\\
			\end{aligned}\Leftrightarrow \left\{\begin{aligned}
				&m \neq-1 \\
				&2 m^2+23 m-4<0 \\
				&\dfrac{6 m+5}{m+1}>0 \\
				&\dfrac{2(2 m-3)}{m+1}>0 \\
				&t_1 t_2-\left(t_1+t_2\right)+1<0 \\
				&2 m^2+23 m-4<0 \\
				&\dfrac{6 m+5}{m+1}>0 \\
				&\dfrac{2(2 m-3)}{m+1}>0 \\
				&\dfrac{6 m+5}{m+1}-\dfrac{2(2 m-3)}{m+1}+1<0\\
			\end{aligned}\right.\right. \\
			&\left(\dfrac{-23-\sqrt{561}}{4}<m<\dfrac{-23+\sqrt{561}}{4}\right. \\
			& \Leftrightarrow\left\{\begin{aligned}
				&{2 m ^{2}+ 2 3 m - 4 < 0}\\
				&{( 6 m + 5 ) ( m + 1 ) > 0}\\
				&{( 2 m - 3 ) ( m + 1 ) > 0}\\
				&{( 3 m + 1 2 ) ( m + 1 ) < 0}\\
			\end{aligned}\Leftrightarrow \left\{\begin{aligned}
				&{\left[\begin{aligned}
						&m<-1 \\
						&m>-\dfrac{5}{6}\\
					\end{aligned}\right.}\\
				&{\left[\begin{aligned}
						&m<-1 \\
						&m>\dfrac{3}{2}
					\end{aligned}\right.}\\
				&-4<m<-1
			\end{aligned}\right.\right. \\
			&
		\end{aligned}$\\
		Do $m$ nguyên nên $m \in\{-3 ;-2\}$.
		Vậy có 2 giá trị nguyên của tham số $m$ thỏa mãn yêu cầu bài toán.
	}
\end{ex}   
\begin{ex}
	[THPT Trần Phú - Đà Nẵng - 2021]%Câu 39
	Có bao nhiêu số nguyên dương $a$ nhỏ hơn 2021 sao cho tồn tại số nguyên $x$ thỏa mãn $2^a\left(2^{x+3}-1\right)+a^2=a\left(2^a+2^{x+3}-1\right)$?
	\choice
	{$12$}
	{$15$}
	{\True $10$}
	{$14$}
	\loigiai{
		Ta có $2^a\left(2^{x+3}-1\right)+a^2=a\left(2^a+2^{x+3}-1\right) \Leftrightarrow\left(2^a-a\right)\left(2^{x+3}-1-a\right)=0$
		$\Leftrightarrow\left[\begin{aligned}
			&2^a-a=0 \\
			&2^{x+3}-1-a=0\\
		\end{aligned}\right.$.\\
		Xét hàm số $f(a)=2^a-a$\\
		Ta có $f'(a)=2^a \ln 2-1 ; f'(a)=0 \Leftrightarrow a=\log _2\left(\dfrac{1}{\ln 2}\right)=a_0$.\\
		Bảng biến thiên\\
		\textbf{(VẼ HÌNH)}\\
		Từ bảng biến thiên, suy ra $f(a) \geq f\left(a_0\right) \approx 0,91>0$.\\
		Nên phương trình $2^a-a=0$ vô nghiệm.\\
		$\begin{aligned}
			& 2^{x+3}-1-a=0 \Leftrightarrow a=2^{x+3}-1 . \\
			& \text{Yêu cầu bài toán}~\Leftrightarrow\left\{\begin{aligned}
				&{2 ^{x + 3}- 1 > 0}\\
				&{2 ^{x + 3}- 1 < 2 0 2 1}\\
			\end{aligned}\Leftrightarrow \left\{\begin{aligned}
				&x>-3 \\
				&x<\log _2 \dfrac{1011}{4}\approx 7,9\\
			\end{aligned}\right.\right. \\
			& \text{Vì}~x \in \mathbb{Z}\Rightarrow x \in\{-2 ;-1 ; 0 ; 1 ; 2 ; 3 ; 4 ; 5 ; 6 ; 7\}.
		\end{aligned}$\\
		Ta có ứng với mỗi $x \in \mathbb{Z}$ thì có một giá trị $a \in \mathbb{Z}$ nên có 10 số nguyên $a$ thỏa mãn yêu cầu.
	}
\end{ex}
\begin{ex}
	[THPT Hương Sơn - Hà Tĩnh - 2022]%Câu 40
	Với giá trị nào của $m$ thì phương tri $9^{x+\dfrac{1}{2}}-4 m \cdot 3^x+m+2=0$ có hai nghiệm phân biệt $x_1 ; x_2$ thỏa mãn $x_1+x_2=1$?
	\choice
	{$m=\dfrac{3}{4}$}
	{$m=-\dfrac{3}{4}$}
	{\True $m=7$}
	{$m=-1$}
	\loigiai{
		Ta có $9^{x+\dfrac{1}{2}}-4 m \cdot 3^x+m+2=0 \Leftrightarrow 3 \cdot 3^{2 x}-4 m \cdot 3^x+m+2=0$ (1).\\
		Đặt $3^x=t,(t>0)$.\\
		Phương trình đã cho trở thành: $3 t^2-4 m t+m+2=0\left(^*\right)$.\\
		Để (1) có hai nghiệm phân biệt $x_1 ; x_2$ thì phương trình $\left({}^*\right)$ có 2 nghiệm $0<t_1<t_2$.\\
		$\Leftrightarrow\left\{\begin{aligned}
			&{\Delta > 0}\\
			&{t _{1}+ t _{2}> 0}\\
			&{t _{1}\cdot t _{2}> 0}\\
		\end{aligned}\Leftrightarrow \left\{\begin{aligned}
			&4 m^2-3 m-6>0 \\
			&\dfrac{4 m}{3}>0 \\
			&\dfrac{m+2}{3}>0\\
		\end{aligned}\quad \Leftrightarrow m>\dfrac{3+\sqrt{105}}{8}.\right.\right.$.\\
		Ta có $\dfrac{m+2}{3}=t_1 \cdot t_2=3^{x_1}\cdot 3^{x_2}=3^{x_1+x_2}$ mà $x_1+x_2=1$ nên ta có $3=\dfrac{m+2}{3}\Leftrightarrow m=7$ (thỏa mãn điều kiện).
	}
\end{ex}
\begin{ex}
	[Sở Hà Tĩnh 2022]%Câu 41
	Có bao nhiêu giá trị nguyên của tham số $m \in[-10 ; 10]$ để phương trình $2^{3^m}\cdot 7^{x^2-2 x}+7^{3^m}\cdot 2^{x^2-2 x}=14^{3^m}\left(7 x^2-14 x+2-7 \cdot 3^m\right)$
	có bốn nghiệm phân biệt trong đó có đúng hai nghiệm lớn hơn $-1$ ?
	\choice
	{\True $10$}      
	{$9$}
	{$11$}
	{$8$}
	\loigiai{
		Ta có\\
		$\begin{aligned}
			& 2^{3^m}\cdot 7^{x^2-2 x}+7^{3^m}\cdot 2^{x^2-2 x}=14^{3^m}\left(7 x^2-14 x+2-7 \cdot 3^m\right) \\
			& \Leftrightarrow \dfrac{7^{x^2-2 x}}{7^{3^m}}+\dfrac{2^{x^2-2 x}}{2^{3^m}}=7 x^2-14 x+2-7 \cdot 3^m \Leftrightarrow 7^{x^2-2 x-3^m}+2^{x^2-2 x-3^m}=7\left(x^2-2 x-3^m\right)+2
		\end{aligned}$\\
		Đăt $x^2-2 x-3^m=a$.\\
		Khi đó $(*)$ trở thành $7^a+2^a=7 a+2 \Leftrightarrow 7^a+2^a-7 a-2=0$.\\
		Xét hàm số $f(a)=7^a+2^a-7 a-2$.\\
		Ta có $f'(a)=7^a \ln 7+2^a \ln 2-7$.\\
		Ta có $f''(a)=7^a(\ln 7)^2+2^a(\ln 2)^2>0, \forall a \in \mathbb{R}$.\\
		Suy ra $f'(a)$ đồng biến trên $\mathbb{R}$, do đó $f'(a)=0$ có tối đa 1 nghiệm.\\
		Mà $f'(0)=\ln 7+\ln 2-7<0$ và $f'(1)=7 \ln 7+2 \ln 2-7>0$.\\
		Suy ra $f'(a)=0$ có nghiệm duy nhất $a_0 \in(0 ; 1)$.\\
		Suy ra $f(a)=0$ có tối đa 2 nghiệm.\\
		Bảng biến thiên của $y=f(a)$\\
		\textbf{(VẼ HÌNH)}\\
		Từ bảng biến thiên ta có $f(a)=0$ có đúng 2 nghiệm $a=0$ và $a=1$.\\
		Từ đó $\left[\begin{aligned}
			&a=x^2-2 x-3^m=0 \\ 
			&a=x^2-2 x-3^m=1\\
		\end{aligned}\Leftrightarrow\left[\begin{aligned}
			&3^m=x^2-2 x \\ 
			&3^m=x^2-2 x-1\\
		\end{aligned}\right.\right.(* *)$
		Để $\left(^*\right)$ có 4 nghiệm thực phân biệt trong đó có đúng hai nghiệm lớn hơn $-1$ thì (**) có 4 nghiệm thực phân biệt trong đó có đúng hai nghiệm lớn hơn $-1$ hay tương đương với đồ thị hàm số $y=3^m$ cắt đồ thị các hàm số $y=x^2-2 x$ và $y=x^2-2 x-1$ tại 4 điểm phân biệt trong đó có đúng hai điểm có hoành độ lớn hơn $-1$.\\
		\textbf{(VẼ HÌNH)}\\
		Dựa vào đồ thị ta có $3^m \geq 3 \Leftrightarrow m \geq 1$.\\
		Suy ra $m \in\{1 ; 2 ; \ldots ; 10\}$.\\
		Vậy có 10 giá trị của $m$ thỏa mãn bài toán.
	}
\end{ex}
\begin{ex}
	[THPT Kim Liên - Hà Nội - 2022]%Câu 42
	Gọi $S$ là tập nghiệm của phương trình $\left(2^x+3^x-8 x+3\right) \sqrt{(3)^{2^x}-m}=0$ (với $m$ là tham số thực). Có tất cả bao nhiêu giá trị nguyên của $m \in[-2021 ; 2021]$ để tập hợp $S$ có hai phần tử?
	\choice
	{$2095$}
	{$2092$}
	{$2093$}
	{\True $2094$}
	\loigiai{
		Xét phương trình: $2^x+3^x-8 x+3=0$.\\
		Đặt $f(x)=2^x+3^x-8 x+3$.\\
		$f'(x)=2^x \ln 2+3^x \ln 3-8 \Rightarrow f''(x)=2^x \ln ^2 2+3^x \ln ^2 3>0, \forall x \in \mathbb{R}$.\\
		Do vậy phương trình $f(x)=0$ có nhiều nhất hai nghiệm.\\
		Ta thấy $f(1)=f(2)=0 \Rightarrow f(x)=0 \Leftrightarrow\left[\begin{aligned}
			&x=1 \\ 
			&x=2\\
		\end{aligned}\right.$.\\
		Xét phương trình: $\left(2^x+3^x-8 x+3\right) \sqrt{(3)^{2^x}-m}=0 \quad(1)$.\\
		Nếu $m \leq 0 \Rightarrow(3)^{2^x}-m>0, \forall x \in \mathbb{R}\Rightarrow(1)$ có tập nghiệm $S=\{1 ; 2\}$.\\
		Khi đó $S$ có hai phần từ khi và chỉ khi $1 \leq \log _2\left(\log _3 m\right)<2 \Leftrightarrow 9 \leq m<81$.\\
		Vậy $m \in[-2021 ; 0] \cup[9 ; 80]$ có 2094 số.
	}
\end{ex}
\begin{ex}
	[THPT Yên Lạc - Vĩnh Phúc - 2022]%Câu 43
	Có bao nhiêu giá trị nguyên của tham số $m$ để phương trình $9^{1+\sqrt{1-x^2}}-(m+3) 3^{1+\sqrt{1-x^2}}+2 m+1=0$ có nghiệm thực
	\choice
	{\True $7$}
	{$5$}
	{$6$}
	{$4$}
	\loigiai{
		Điều kiện: $-1 \leq x \leq 1$.\\
		Đặt: $t=3^{1+\sqrt{1-x^2}}, t \in[3 ; 9]$.\\
		Phương trình trở thành: $t^2-(m+3) t+2 m+1=0 \Leftrightarrow m(t-2)=t^2-3 t+1 \Leftrightarrow m=\dfrac{t^2-3 t+1}{t-2}$ (1)\\
		Xét hàm số $f(t)=\dfrac{t^2-3 t+1}{t-2}, t \in[3 ; 9]$.\\
		Ta có $f'(t)=\dfrac{t^2-4 t+5}{(t-2)^2}>0, t \in[3 ; 9] \Rightarrow$ Hàm số luôn đồng biến trên đoạn $[3 ; 9]$\\
		Vậy $f(3) \leq f(t) \leq f(9)$ hay $1 \leq f(t) \leq \dfrac{55}{7}$.\\
		Phương trình (1) có nghiệm $t \in[3 ; 9] \Leftrightarrow 1 \leq m \leq \dfrac{55}{7}$.\\
		Vì $m \in \mathbb{Z}\Rightarrow m \in\{1 ; 2 ; 3 ; 4 ; 5 ; 6 ; 7\}$.
	}
\end{ex}
\begin{ex}
	[THPT Nguyễn Huệ - Huế - 2022]%Câu 44
	Phương trình $2^{x-2+\sqrt[3]{m-3 x}}+\left(x^3-6 x^2+9 x+m\right) 2^{x-2}=2^{x+1}+1$ có 3 nghiệm phân biệt khi và chỉ khi $m \in(a ; b)$. Đặt $T=b^2-a^2$ thì:
	\choice
	{\True $T=48$}
	{$T=64$}
	{$T=72$}
	{$T=36$}
	\loigiai{
		Ta có $2^{x-2+\sqrt[3]{m-3 x}}+\left(x^3-6 x^2+9 x+m\right) 2^{x-2}=2^{x+1}+1 \Leftrightarrow 2^{\sqrt[3]{m-3 x}}+(m-3 x)=2^{2-x}+(2-x)^3$ (1).\\
		Xét hàm số $f(t)=2^t+t^3 \Rightarrow f'(t)=2^t \cdot \ln 2+3 t^2>0, \forall t \in \mathbb{R}$.\\
		Hàm số $f(t)$ đồng biến trên $\mathbb{R}$
		$(1)\Leftrightarrow f(\sqrt[3]{m-3 x})=f(2-x) \Leftrightarrow \sqrt[3]{m-3 x}=2-x \Leftrightarrow m=-x^3+6 x^2-9 x+8(2)$.\\
		Xét hàm số $g(x)=-x^3+6 x^2-9 x+8$ trên $\mathbb{R}$.\\
		$g'(x)=-3 x^2+12 x-9=0 \Leftrightarrow\left[\begin{aligned}
			&x=1 \\
			&x=3\\
		\end{aligned}\right.$.\\
		\textbf{(VẼ HÌNH)}\\
		Từ bảng biến thiên, ta có (2) có 3 nghiệm phân biệt khi và chỉ khi $4<m<8 \Rightarrow m \in(4 ; 8) \Rightarrow b^2-a^2=48$.
	}
\end{ex}
\begin{ex}
	[Sở Thái Bình 2022]%Câu 45
	Cho $a, b$ là hai số thay đổi thoả mãn $a>1, b>1$ và $a+b=12$. Giả sử $x_1$, $x_2$ là hai nghiệm của phương trình $\log _a x \cdot \log _b x-\log _a x-\log _b x-1=0$. Giá trị lớn nhất của $P=x_1 x_2$ là
	\choice
	{$P_{\max}=39$}
	{\True $P_{\max}=36$}
	{$P_{\max}=32$}
	{$P_{\max}=45$}
	\loigiai{
		Đặt $t=\ln x \Rightarrow x=e^t \Rightarrow \log _a x=\log _a e^t=t \cdot \log _a e=\dfrac{t}{\log _e a}=\dfrac{t}{\ln a}$.\\
		Tương tự: $\log _b x=\dfrac{t}{\ln b}$.\\
		Phương trình $\log _a x \cdot \log _b x-\log _a x-\log _b x-1=0$ trở thành\\
		$\dfrac{t}{\ln a}\cdot \dfrac{t}{\ln b}-\dfrac{t}{\ln a}-\dfrac{t}{\ln b}-1=0 \Leftrightarrow t^2-(\ln a+\ln b) t-\ln a \cdot \ln b=0$\\
		Theo viet: $t_1+t_2=\ln a+\ln b=\ln a b \leq \ln \left(\dfrac{a+b}{2}\right)^2$\\
		(Bất đẳng thức Cosi dấu bằng xảy ra khi $a=b=6$)\\
		$\begin{aligned}
			& \Leftrightarrow t_1+t_2 \leq \ln \left(\dfrac{12}{2}\right)^2 \Leftrightarrow t_1+t_2 \leq \ln 36 . \\
			& P=x_1 x_2=e^{t_1}\cdot e^{t_2}=e^{t_1+t_2}\leq e^{\ln 36}\Leftrightarrow P \leq 36 .~\text{Nên}~P_{\text{max}}=36 .
		\end{aligned}$.
	}
\end{ex}
\begin{dang}
	{Phương trình kết hợp của mũ và logarit chứa tham số}
\end{dang}                                       
\begin{ex}
	[Mã 103 -2019]%Câu 1.
	Cho phương trình $\left(2\log_3^2x-\log_3x-1\right)\sqrt{5^x-m}=0$ ($m$ là tham số thực). Có tất cả bao nhiêu giá trị nguyên dương của $m$ để phương trình đã cho có đúng hai nghiệm phân biệt?
	\choice
	{Vô số}
	{$124$}
	{\True $123$}
	{$125$}
	\loigiai{
		Điều kiện: $\heva{&x>0\\&5^x-m\geq 0(m>0)}\Leftrightarrow\heva{&x>0\\&x\geq\log_5m.}$ \\
		$\left(2\log_3^2x-\log_3x-1\right)\sqrt{5^x-m}=0$ (1)\\
		$ \Leftrightarrow\hoac{&2\log_3^2x-\log_3x-1=0\\&5^x-m=0}\Leftrightarrow\hoac{&x=3,x=\dfrac{1}{\sqrt{3}}\\&f(x)=5^x=m.} $ \\
		Xét $f(x)=5^x$ hàm số đồng biến trên $\mathbb{R}$. 
		{\color{red} HINH O DAY}\\
		Dựa vào bảng biến thiên, để phương trình có hai nghiệm phân biệt thì\\
		$\hoac{&m=1\\&5^{\tfrac{1}{\sqrt{3}}}\leq m<125}$, $m\in\mathbb{Z}_+\Rightarrow\hoac{&0<m\leq 1\\&3\leq m\leq 124.}$ \\
		Nên có 123 giá trị $m$ thoả mãn.
	}
\end{ex}
\begin{ex}[Mã 102 - 2019]%Câu 2.
	Cho phương trình $\left(2\log_2^2x-3\log_2x-2\right)\sqrt{3^x-m}=0$ ($m$ là tham số thực). Có tất cả bao nhiêu giá trị nguyên dương của $m$ để phương trình đã cho có đúng hai nghiệm phân biệt?
	\choice
	{vô số}
	{$81$}
	{\True $79$}
	{$80$}
	\loigiai{	
		Điều kiện $\heva{&x>0\\&3^x-m\geq 0}\Leftrightarrow\heva{&x>0\\&m\leq 3^x}$ (*).\\
		Ta có $\left(2\log_2^2x-3\log_2x-2\right)\sqrt{3^x-m}=0 (1)\Leftrightarrow\hoac{&2\log_2^2x-3\log_2x-2=0(2)\\&\sqrt{3^x-m}=0(3).}$ \\
		Trong đó $(2)\Leftrightarrow\hoac{&\log_2x=2\\&\log_2x=-\dfrac{1}{2}}\Leftrightarrow\hoac{&x=4\\&x=\dfrac{1}{\sqrt{2}}}$. (4).\\
		Với $m>0$ thì $3^x=m\Leftrightarrow\log_3m=x$.\\
		Do đó, phương trình (1) có đúng hai nghiệm phân biệt khi và chỉ khi xảy ra các trường hợp sau:\\
		Trường hợp 1: (3) có nghiệm $x=\log_3m\leq 0\Leftrightarrow 0<m\leq 1$. Kết hợp điều kiện (*) và (4) ta được $m=1$ thì (1) có hai nghiệm phân biệt $x=\dfrac{1}{\sqrt{2}}$ và $x=4$.\\
		Trường hợp 2: $m>1$, khi đó (*) $\Leftrightarrow x\geq\log_3m>0$.\\
		Và do $4>\dfrac{1}{\sqrt{2}}$ nên (1) có hai nghiệm phân biệt khi và chỉ khi $\dfrac{1}{\sqrt{2}}\leq\log_3m<4\Leftrightarrow 3^{\tfrac{1}{\sqrt{2}}}\leq m<3^4$.\\
		Mà $m$ nguyên dương nên ta có $m\in\{3,4,\ldots,80\}$, có 78 giá trị của $m$.\\
		Vậy có 79 giá trị nguyên dương của $m$ để phương trình có đúng hai nghiệm phân biệt.
	}
\end{ex}
\begin{ex}[Mã 104 2019]%Câu 3.
	Cho phương trình $\left(2\log_3^2x-\log_3x-1\right)\sqrt{4^x-m}=0$ ($m$ là tham số thực). Có tất cả bao nhiêu giá trị nguyên dương của $m$ để phương trình có đúng hai nghiệm phân biệt?
	\choice
	{$64$}
	{Vô số}
	{\True $62$}
	{$63$}
	\loigiai{
		Ta có điều kiện $\heva{&x>0\\&x\geq\log_4m}$ (*) (với $m$ nguyên dương).\\
		Phương trình $\left(2\log_3^2x-\log_3x-1\right)\sqrt{4^x-m}=0(1)$ \\
		$ \Leftrightarrow\hoac{&2\log_3^2x-\log_3x-1=0(2)\\&4^x=m(3).} $ \\
		Phương trình $(2)\Leftrightarrow\hoac{&\log_3x=1\\&\log_3x=-\dfrac{1}{2}}\Leftrightarrow\hoac{&x=3\\&x=\dfrac{\sqrt{3}}{3}.}$ \\
		Phương trình $(3)\Leftrightarrow x=\log_4m$.\\
		Do $m$ nguyên dương nên ta có các trường hợp sau:\\
		Trường hợp 1: $m=1$ thì $\log_4m=0$. Do đó (*) là $x>0$.\\
		Khi đó nghiệm của phương trình $(3)$ bị loại và nhận nghiệm của phương trình $(2)$.\\
		Do đó nhận giá trị $m=1$.\\
		Trường hợp 2: $m\geq 2$ thì (*) là $x\geq\log_4m$ (vì $\log_4m\geq\dfrac{1}{2}$).\\
		Để phương trình $(1)$ có đúng hai nghiệm phân biệt\\
		$ \Leftrightarrow\dfrac{\sqrt{3}}{3}\leq\log_4m<3 $ \\
		$ \Leftrightarrow 4^{\tfrac{\sqrt{3}}{3}}\leq m<4^3 $.\\
		Suy ra $m\in\left\{3; 4; 5;\ldots; 63\right\}$.\\
		Vậy từ cả 2 trường hợp ta có: $63-3+1+1=62$ giá trị nguyên dương $m$.
	}
\end{ex}
\begin{ex}
	[Mã 101 2019]%Câu 4.
	Cho phương trình $\left(4\log_2^2x+\log_2x-5\right)\sqrt{7^x-m}=0$ ($m$ là tham số thực). Có tất cả bao nhiêu giá trị nguyên dương của $m$ để phương trình đã cho có đúng hai nghiệm phân biệt?
	\choice
	{$49$}
	{$47$}
	{Vô số}
	{\True $48$}
	\loigiai{
		Điều kiện: $\heva{&x>0\\&7^x-m\geq 0}\Leftrightarrow\heva{&x>0\\&7^x\geq m.}$ \\
		Trường hợp $m\leq 0$ thì $\left(4\log_2^2x+\log_2x-5\right)\sqrt{7^x-m}=0\Leftrightarrow 4\log_2^2x+\log_2x-5=0$ \\
		$\Leftrightarrow(\log_2x-1)\left(4\log_2x+5\right)=0\Leftrightarrow\hoac{&\log_2x=1\\&\log_2x=-\dfrac{5}{4}}\Leftrightarrow\hoac{&x=2\\&x=2^{-\tfrac{5}{4}}.} $ \\
		Trường hợp này không thỏa điều kiện $m$ nguyên dương.\\
		Trường hợp $m>0$, ta có $\heva{&x>0\\&7^x\geq m}\Leftrightarrow x\geq\log_7m$ nếu $m>1$ và $x>0$ nếu $0<m\leq 1$.\\
		Khi đó $\left(4\log_2^2x+\log_2x-5\right)\sqrt{7^x-m}=0\Leftrightarrow\hoac{&4\log_2^2x+\log_2x-5=0\\&\sqrt{7^x-m}=0}\Leftrightarrow\hoac{&x=2\\&x=2^{-\tfrac{5}{4}}\\&x=\log_7m.}$ \\
		Xét $0<m\leq 1$ thì nghiệm $x=\log_7m\leq 0$ nên trường hợp này phương trình đã cho có đúng 2 nghiệm $x=2; x=2^{-\tfrac{5}{4}}$ thỏa mãn điều kiện.\\
		Xét $m>1$, khi đó điều kiện của phương trình là $x\geq\log_7m$.\\
		Vì $2>2^{-\tfrac{5}{4}}$ nên phương trình đã cho có hai nghiệm phân biệt khi và chỉ khi $2>\log_7m\geq 2^{-\tfrac{5}{4}}$ \\
		$ \Leftrightarrow 7^{2^{-\tfrac{5}{4}}}\leq m<7^2 $.\\
		Trường hợp này $m\in\left\{3;4;5;\ldots;48\right\}$, có 46 giá trị nguyên dương của $m$.\\
		Tóm lại có 47 giá trị nguyên dương của $m$ thỏa mãn.
	}
\end{ex}
\begin{ex}
	[Mã 102 2018]%Câu 5.
	Cho phương trình $3^x+m=\log_3(x-m)$ với $m$ là tham số. Có bao nhiêu giá trị nguyên của $m\in(-15;15)$ để phương trình đã cho có nghiệm?
	\choice
	{$15$}
	{$16$}
	{$9$}
	{\True $14$}
	\loigiai{
		Ta có $3^x+m=\log_3(x-m)\Leftrightarrow 3^x+x=\log_3(x-m)+x-m\quad (*)$.\\
		Xét hàm số $f(t)=3^t+t$, với $t\in\mathbb{R}$. Có $f'(t)=3^t\ln 3+1>0,\forall t\in\mathbb{R}$ nên hàm số $f(t)$ đồng biến trên tập xác định. Mặt khác phương trình $(*)$ có dạng: $f(x)=f\left(\log_3(x-m)\right)$. Do đó ta có $f(x)=f\left(\log_3(x-m)\right)\Leftrightarrow x=\log_3(x-m)\Leftrightarrow 3^x=x-m\Leftrightarrow 3^x-x=-m$.\\
		Xét hàm số $g(x)=3^x-x$, với $x\in\mathbb{R}$. Có $g'(x)=3^x\ln 3-1$, $g'(x)=0\Leftrightarrow x=\log_3\left(\dfrac{1}{\ln 3}\right)$.\\
		Bảng biến thiên\\
		{\color{red} HINH O DAY}\\
		Từ bảng biến thiên ta thấy các giá trị của tham số để phương trình có nghiệm là $m\in\left(-\infty;-g\left(\log_3\left(\dfrac{1}{\ln 3}\right)\right)\right]$. Vậy số giá trị nguyên của $m\in(-15;15)$ để phương trình đã cho có nghiệm là $14$.
	}
\end{ex}
\begin{ex}[Mã 101 2018]%Câu 6.
	Cho phương trình $5^x+m=\log_5(x-m)$ với $m$ là tham số. Có bao nhiêu giá trị nguyên của $m\in(-20; 20)$ để phương trình đã cho có nghiệm?
	\choice
	{\True $19$}
	{$9$}
	{$21$}
	{$20$}
	\loigiai{
		Điều kiện: $x>m$.\\
		Đặt: $t=\log_5(x-m)\Rightarrow\heva{&x-m=5^t\\&5^x+m=t}\Rightarrow 5^x+x=5^t+t (1)$.\\
		Xét hàm số $f(u)=5^u+u\Rightarrow f'(u)=5^u\ln 5+1>0,\forall u\in\mathbb{R}$.\\
		Do đó: $(1)\Leftrightarrow x=t\Leftrightarrow x=5^x+m\Leftrightarrow m=x-5^x$.\\
		Xét hàm số $f(x)=x-5^x$, $x>m$.\\
		Do: $5^x>0\Rightarrow m<x$, suy ra phương trình có nghiệm luôn thỏa điều kiện.\\
		$f'(x)=1-5^x\ln 5$, $f'(x)=0\Leftrightarrow 1-5^x\ln 5=0\Leftrightarrow x=\log_5\left(\dfrac{1}{\ln 5}\right)$.\\
		Bảng biến thiên: \\
		{\color{red} HINH O DAY}\\
		Dựa vào bảng biến thiên $\Rightarrow m\leq\approx-0,917\xrightarrow{{m\in(-20; 20)}}m=\left\{-19;-18;\ldots;-1\right\}$.\\
		Vậy có $19$ giá trị nguyên của $m$ thỏa ycbt.
	}
\end{ex}
\begin{ex}
	[Mã 103 -2018]%Câu 7.
	Cho phương trình $7^x+m=\log_7(x-m)$ với $m$ là tham số. Có bao nhiêu giá trị nguyên của $m\in(-25;25)$ để phương trình đã cho có nghiệm?
	\choice
	{$9$}
	{$25$}
	{\True $24$}
	{$26$}
	\loigiai{
		Điều kiện: $x>m$.\\
		Đặt $t=\log_7(x-m)$ ta có $\heva{&7^x+m=t\\&7^t+m=x}\Rightarrow 7^x+x=7^t+t (1)$.\\
		Do hàm số $f(u)=7^u+u$ đồng biến trên $\mathbb{R}$, nên ta có $(1)\Leftrightarrow t=x$. Khi đó:\\
		$7^x+m=x\Leftrightarrow m=x-7^x$.\\
		Xét hàm số $g(x)=x-7^x\Rightarrow g'(x)=1-7^x\ln 7=0\Leftrightarrow x=-\log_7(\ln 7)$.\\
		Bảng biến thiên: \\
		{\color{red} HINH O DAY}\\
		Từ đó phương trình đã cho có nghiệm khi và chỉ khi $m\leq g\left(-\log_7(\ln 7)\right)\approx-0,856$ (cácnghiệm này đều thỏa mãn điều kiện vì $x-m=7^x>0$).\\
		Do $m$ nguyên thuộc khoảng $(-25;25)$, nên $m\in\left\{-24;-16;\ldots;-1\right\}$.
	}
\end{ex}
\begin{ex}
	Cho phương trình $5^x+m+\log_{\tfrac{1}{5}}(x-m)=0$ với $m$ là tham số. Có bao nhiêu giá trị nguyên của tham số $m\in(-20;20)$ để phương trình đã cho có nghiệm thực?
	\choice
	{$20$}
	{$21$}
	{$18$}
	{\True $19$}
	\loigiai{
		Ta có $5^x+m+\log_{\tfrac{1}{5}}(x-m)=0\Leftrightarrow 5^x=\log_5(x-m)-m=0(1)$.\\
		Điều kiện xác định: $x>m$.\\
		Đặt $t=\log_5(x-m)$, ta có $x-m=5^t$.\\
		Khi đó ta có hệ phương trình $\heva{&x-m=5^t\\&t-m=5^x}\Leftrightarrow\heva{&x-m=5^t(*)\\&5^x+x=5^t+t(2).}$ \\
		Xét hàm số $f(u)=5^u+u,u\in\mathbb{R}$..\\
		$f'(u)=5^u\ln 5+1>0,\forall u$ suy ra hàm số $f(u)=5^u+u$ đồng biến trên $\mathbb{R}$.\\
		Do đó $(2)\Leftrightarrow f(x)=f(t)\Leftrightarrow x=t$.\\
		Thay vào phương trình $(*)$ ta có $m=x-5^x(3)$.\\
		Ta có $x-m=5^x>0$, do đó phương trình $(1)$ có nghiệm $\Leftrightarrow$ phương trình $(3)$ có nghiệm $x\in\mathbb{R}$.\\
		Xét hàm số $g(x)=x-5^x,x\in\mathbb{R}$, có $g'(x)=1-5^x\ln 5, g'(x)=0\Leftrightarrow x=\log_5\left(\dfrac{1}{\ln 5}\right)$.\\
		$\lim\limits_{x\to-\infty}\left(x-5^x\right)=-\infty;\lim\limits_{x\to+\infty}\left(x-5^x\right)=-\infty$.\\
		Bảng biến thiên\\
		{\color{red} HINH O DAY}\\
		Từ bảng biến thiên ta thấy phương trình có nghiệm $\Leftrightarrow m\leq\log_5\left(\dfrac{1}{e\ln 5}\right)\simeq-0,91$.\\
		Vì $m\in(-20;20)$ và là số nguyên, suy ra $m\in\left\{-20;-19;\ldots;-1\right\}$.\\
		Vậy có $19$ giá trị của $m$.
	}
\end{ex}
\begin{ex}
	[Mã 104 2018]%Câu 9.
	Cho phương trình $2^x+m=\log_2(x-m)$ với $m$ là tham số. Có bao nhiêu giá trị nguyên của $m\in(-18;18)$ để phương trình đã cho có nghiệm?
	\choice
	{$9$}
	{$19$}
	{\True $17$}
	{$18$}
	\loigiai{
		Điều kiện: $x>m$.\\
		Đặt $t=\log_2(x-m)$ ta có $\heva{&2^x+m=t\\&2^t+m=x}\Rightarrow 2^x+x=2^t+t (1)$.\\
		Do hàm số $f(u)=2^u+u$ đồng biến trên $\mathbb{R}$, nên ta có $(1)\Leftrightarrow t=x$. Khi đó:\\
		$2^x+m=x\Leftrightarrow m=x-2^x$.\\
		Xét hàm số $g(x)=x-2^x\Rightarrow g'(x)= 1-2^x\ln 2=0\Leftrightarrow x=-\log_2(\ln 2)$.\\
		Bảng biến thiên:\\ 
		{\color{red} HINH O DAY}\\
		Từ đó phương trình đã cho có nghiệm khi và chỉ khi $m\leq g\left(-\log_2(\ln 2)\right)\approx-0,914$ (các nghiệm này đều thỏa mãn điều kiện vì $x-m=2^x>0$).\\
		Do $m$ nguyên thuộc khoảng $(-18;18)$, nên $m\in\left\{-17;-16;\ldots;-1\right\}$.
	}
\end{ex}
\begin{ex}
	[Chuyên Nguyễn Du-ĐăkLăk 2019]%Câu 10.
	Cho phương trình $5^x+m=\log_5(x-m)$. Có bao nhiêu giá trị $m$ nguyên trong khoảng $(-20;20)$ để phương trình trên có nghiệm?
	\choice
	{$15$}
	{\True $19$}
	{$14$}
	{$17$}
	\loigiai{
		Ta có phương trình $5^x+m=\log_5(x-m)$ (1) với điều kiện $x-m>0$.\\
		Đặt $\log_5(x-m)=t\Leftrightarrow x-m=5^t$ (*) thay vào phương trình (1) ta có $5^x+m=t\Leftrightarrow t-m=5^x(**)$. Từ (*) và (**) ta có hệ phương trình $\heva{&x-m=5^t\\&t-m=5^x}$. Từ hệ phương trình ta suy ra $x-t=5^t-5^x\Leftrightarrow x+5^x=t+5^t$.\\
		Xét hàm số $f(x)=x+5^x$ trên $\mathbb{R}$, ta có $f'(x)=1+5^x\cdot\ln 5>0  \forall x\in\mathbb{R}$ nên hàm số $f(x)=x+5^x$ luôn đồng biến trên $\mathbb{R}$, do đó ta có $x+5^x=t+5^t\Leftrightarrow f(x)=f(t)\Leftrightarrow x=t$ thay vào phương trình (**) ta có $x-m=5^x\Leftrightarrow x-5^x=m$. Đặt $g(x)=x-5^x$ ta có $g'(x)=1-5^x\cdot\ln 5$. Ta có $g'(x)=0\Leftrightarrow 1-5^x\cdot\ln 5=0\Leftrightarrow 5^x=\dfrac{1}{\ln 5}\Leftrightarrow x=\log_5\left(\dfrac{1}{\ln 5}\right)$.\\
		Ta có bảng biến thiên với $g\left(\log_5\left(\dfrac{1}{\ln 5}\right)\right)=\log_5\left(\dfrac{1}{\ln 5}\right)-\dfrac{1}{\ln 5} =\alpha$. \\
		{\color{red} HINH O DAY}\\
		Từ bảng biến thiên ta thấy phương trình $x-5^x=m$ có nghiệm khi $m\leq\alpha$ hay $m\leq\log_5\left(\dfrac{1}{\ln 5}\right)-\dfrac{1}{\ln 5}$. Ta suy ra có $19$ giá trị nguyên của $m$ thỏa mãn.
	}
\end{ex}
\begin{ex}
	Tổng tất cả các giá trị của tham số $m$ để phương trình $2^{x^2+4x+5-m^2}=\log_{x^2+4x+6}\left(m^2+1\right)$ có đúng 1.
	nghiệm là
	\choice
	{$-2$}
	{$1$}
	{$4$}
	{\True $0$}
	\loigiai{
		Đặt $t=x^2+4x+5$, khi đó $t\geq 1$.\\
		Thế vào phương trình đã cho ta được phương trình sau.\\
		$2^t\ln(t+1)=2^{m^2}\ln\left(m^2+1\right)\Leftrightarrow t=m^2\Leftrightarrow x^2+4x+5=m^2$.\\
		(Do hàm đặc trưng $f(u)=2^u\ln(u+1)$ có.\\
		$f'(u)=\dfrac{2^u}{u+1}+2^u\ln(u+1)\cdot\ln 2\geq 0,\forall u\geq 0\Rightarrow f(u)$ đồng biến trên $[0;+\infty)$).\\
		Vậy $2^{x^2+4x+5-m^2}=\log_{x^2+4x+6}\left(m^2+1\right)$ có đúng 1 nghiệm\\
		$ \Leftrightarrow x^2+4x+5-m^2=0 $ có đúng 1 nghiệm\\
		$ \Leftrightarrow\Delta'=m^2-1=0 $ \\
		$ \Leftrightarrow m=\pm 1\Rightarrow $ Tổng tất cả các giá trị $m$ bằng $0$.
	}
\end{ex}
\begin{ex}
	Tổng tất cả các giá trị của tham số $m$ để phương trình $3^{x^2-2x+1-2|x-m|}=\log_{x^2-2x+3}(2|x-m|+2)$ có đúng ba nghiệm phân biệt là 
	\choice
	{$2$}
	{\True $3$}
	{$1$}
	{$0$}
	\loigiai{
		Phương trình tương đương $3^{x^2-2x+3-(2|x-m|+2)}=\dfrac{\ln(2|x-m|+2)}{\ln\left(x^2-2x+3\right)}$ \\
		$ \Leftrightarrow 3^{x^2-2x+3}\cdot\ln\left(x^2-2x+3\right)=3^{2|x-m|+2}\cdot\ln(2|x-m|+2) $ (*).\\
		Xét hàm đặc trưng $f(t)=3^t\cdot\ln t, t\geq 2$ là hàm số đồng biến nên từ phương trình (*) suy ra $\Leftrightarrow x^2-2x+3=2|x-m|+2\Leftrightarrow g(x)=x^2-2x-2|x-m|+1=0$.\\
		Có $g(x)=\heva{&x^2-4x+2m+1 khi x\geq m\\&x^2-2m+1\quad\quad khi x\leq m}\Rightarrow g'(x)=\heva{&2x-4 khi x\geq m\\&2x\quad khi x\leq m.}$ \\
		và $g'(x)=0\Leftrightarrow\hoac{&x=2 khi x\geq m\\&x=0 khi x\leq m.}$ \\
		Xét các trường hợp sau:\\
		Trường hợp 1: $m\leq 0$ ta có bảng biến thiên của $g(x)$ như sau: \\
		{\color{red} HINH O DAY}\\
		Phương trình chỉ có tối đa 2 nghiệm nên không có $m$ thoả mãn.\\
		Trường hợp 2: $m\geq 2$ tương tự.\\
		Trường hợp 3: $0<m<2$, bảng biến thiên $g(x)$ như sau: \\
		{\color{red} HINH O DAY}\\
		Phương trình có 3 nghiệm khi $\hoac{&(m-1)^2=0\\&-2m+1=0>2m-3\\&-2m+1<0=2m-3}\Leftrightarrow\hoac{&m=1\\&m=\dfrac{1}{2}\\&m=\dfrac{3}{2}.}$ \\
		Cả 3 giá trị trên đều thoả mãn, nên tổng của chúng bằng 3.
	}
\end{ex}
\begin{ex}
	[Chuyên Lam Sơn - 2020]%Câu 13.
	Có bao nhiêu giá trị nguyên của tham số $a$ trên đoạn $[-10;10]$ để phương trình\\
	$\mathrm{e}^{x+a}-\mathrm{e}^x=\ln(1+x+a)-\ln(1+x)$ có nghiệm duy nhất. 
	\choice
	{$2$}
	{$10$}
	{$1$}
	{\True $20$}
	\loigiai{
		Điều kiện xác định $\heva{&x+1+a>0\\&x+1>0}$ (*).\\
		Phương trình tương đương với $\mathrm{e}^{x+a}-\mathrm{e}^x-\left(\ln(1+x+a)-\ln(1+x)\right)=0$.\\
		Đặt $f(x)=\mathrm{e}^{x+a}-\mathrm{e}^x$, $g(x)=\ln(1+x+a)-\ln(1+x)$, $Q(x)=f(x)-g(x)$.\\
		Phương trình đã cho viết lại thành $Q(x)=0$.\\
		Với $a=0$ thì $Q(x)=0$ (luôn đúng với mọi $x$ thoả mãn (*)).\\
		Với $a>0$ có (*) tương đương với $x >-1$, $f(x)$ đồng biến và $g(x)$ nghịch biến với $x >-1$.\\
		Khi đó, $Q(x)$ đồng biến với $x >-1$. (1).\\
		Ta có $\heva{&\lim\limits_{x\to(-1)^+} Q(x)=\lim\limits_{x\to(-1)^+}\left(\mathrm{e}^{x+a}-\mathrm{e}^x-\ln\dfrac{1+x+a}{1+x}\right)=\lim\limits_{x\to(-1)^+}\left[\mathrm{e}^{x+a}-\mathrm{e}^x-\ln\left(1+\dfrac{a}{1+x}\right)\right]=-\infty\\&\lim\limits_{x\to+\infty} Q(x)=\lim\limits_{x\to+\infty}\left[\mathrm{e}^x\left(\mathrm{e}^a-1\right)-\ln\left(1+\dfrac{a}{1+x}\right)\right]=+\infty}$ (2).\\
		Kết hợp (1), (2) thì phương trình $Q(x)=0$ có nghiệm duy nhất.\\
		Với $a<0$ có (*) tương đương với $x >-1-a$, $g(x)$ đồng biến và $f(x)$ nghịch biến với $x >-1-a$.\\
		Khi đó, $Q(x)$ nghịch biến với $x >-1-a$. (3).\\
		Ta có:\\
		$\heva{&\lim\limits_{x\to(-1-a)^+} Q(x)=\lim\limits_{x\to(-1-a)^+}\left(\mathrm{e}^{x+a}-\mathrm{e}^x-\ln\dfrac{1+x+a}{1+x}\right)=\lim\limits_{x\to(-1-a)^+}\left[\mathrm{e}^{x+a}-\mathrm{e}^x-\ln\left(1+\dfrac{a}{1+x}\right)\right]=+\infty\\&\lim\limits_{x\to+\infty} Q(x)=\lim\limits_{x\to+\infty}\left[\mathrm{e}^x\left(\mathrm{e}^a-1\right)-\ln\left(1+\dfrac{a}{1+x}\right)\right]=-\infty}$ (4).\\
		Kết hợp (3), (4) suy ra $Q(x)=0$ có nghiệm duy nhất.\\
		Do $a$ là số nguyên trên đoạn $[-10;10]$ nên kết hợp 3 trường hợp trên thấy có 20 giá trị của $a$.\\
		thoả mãn điều kiện của bài.
	}
\end{ex}
\begin{ex}
	[Chuyên Sơn La - 2020]%Câu 14.
	Có bao nhiêu giá trị nguyên của tham số $m$ thuộc $(-2020;2020)$ để phương trình $\mathrm{e}^x=\ln(x+2m)+2m$ có nghiệm?
	\choice
	{\True $2019$}
	{$2020$}
	{$2021$}
	{$4039$}
	\loigiai{
		Ta có $\mathrm{e}^x=\ln(x+2m)+2m\Leftrightarrow\mathrm{e}^x+x=\ln(x+2m)+x+2m\Leftrightarrow\mathrm{e}^x+x=\mathrm{e}^{\ln(x+2m)}+\ln(x+2m)$ (*).\\
		Xét hàm số $f(t)=\mathrm{e}^t+t$ với $t\in\mathbb{R}\Rightarrow f'(t)=\mathrm{e}^t+1>0,\forall t$. Suy ra hàm số $f(t)$ đồng biến trên $\mathbb{R}$.\\
		Do đó $(*)\Leftrightarrow f(x)=f\left(\ln(x+2m)\right)\Leftrightarrow x=\ln(x+2m)\Leftrightarrow x+2m=\mathrm{e}^x\Leftrightarrow 2m=\mathrm{e}^x-x$.\\
		Xét hàm số $g(x)=\mathrm{e}^x-x\Rightarrow g'(x)=\mathrm{e}^x-1\Rightarrow g'(x)=0\Leftrightarrow x=0$.\\
		Bảng biên thiên
		{\color{red} HINH O DAY}\\
		Từ bảng biên thiên suy ra phương trình có nghiệm khi và chỉ khi $2m\geq 1\Leftrightarrow m\geq\dfrac{1}{2}$.\\
		Mà $m\in\mathbb{Z},m\in(-2020;2020)$ nên $m\in\left\{1;2;3;\ldots;2019\right\}$.\\
		Vậy có 2019 giá trị nguyên của tham số $m$ thuộc $(-2020;2020)$ để phương trình $\mathrm{e}^x=\ln(x+2m)+2m$ có nghiệm.
	}
\end{ex}
\begin{ex}
	[Liên trường huyện Quảng Xương - Thanh Hóa - 2021]%Câu 15.
	Biết rằng $x,y$ là các số thực dương sao cho 3 số $u_1=8^{x+\log_2y}$, $u_2=2^{x-\log_2y}$, $u_3=5y$ theo thứ tự lập thành một cấp số cộng và một cấp số nhân. Khi đó, tích $2^x\cdot y^2$ có giá trị bằng 
	\choice
	{$10$}
	{$5$}
	{$\sqrt{5}$}
	{\True $1$}
	\loigiai{
		Điều kiện: $y>0$;\\
		Theo đề bài, ta có: $\heva{&{2\cdot 2}^{(x-\log_2y)}=8^{x+\log_2y}+5y(1)\\&2^{2(x-\log_2y)}=8^{x+\log_2y}\cdot 5y}$ \\
		$\Rightarrow\dfrac{\left(8^{x+\log_2y}+5y\right)^2}{4}=8^{x+\log_2y}\cdot 5y\Leftrightarrow\left(8^{x+\log_2y}-5y\right)^2=0\Leftrightarrow 8^{x+\log_2y}=5y $ (2).\\
		$\begin{aligned}&\Rightarrow\log_2\left(8^{x+\log_2y}\right)=\log_25y\\&\Leftrightarrow(x+\log_2y)\cdot\log_28=\log_25+\log_2y\\&\Leftrightarrow 3x+3\log_2y=\log_25+\log_2y\\&\Leftrightarrow 3x+2\log_2y=\log_25\\&\Leftrightarrow 3x=\log_2\dfrac{5}{y^2}(3)\end{aligned}$\\
		Thay $(2)$ vào $(1)$ ta được:\\
		$2.2^{(x-\log_2y)}=5y+5y\Leftrightarrow 2^{(x-\log_2y)}=5y\Leftrightarrow x-\log_2y=\log_25y\Leftrightarrow x=\log_25y^2(4)$.\\
		Từ $(3)$ và $(4)\Rightarrow\log_2\dfrac{5}{y^2}=3\cdot\log_25y^2\Rightarrow\dfrac{5}{y^2}=\left(5y^2\right)^3\Leftrightarrow y^8=\dfrac{1}{25}\Rightarrow y=\dfrac{1}{\sqrt[4]{5}}$ \\
		$ \Rightarrow x=\log_25\cdot\left(\dfrac{1}{\sqrt[4]{5}}\right)^2=\log_2\sqrt{5}\Rightarrow 2^x\cdot y^2=2^{\log_2\sqrt{5}}\cdot\left(\dfrac{1}{\sqrt[4]{5}}\right)^2=\sqrt{5}\cdot\dfrac{1}{\sqrt{5}}=1 $.
	}
\end{ex}
\begin{ex}
	[Chuyên Quang Trung - Bình Phước - 2021]%Câu 16.
	Cho phương trình $\left(\log_2^2x-\log_2\dfrac{x^3}{4}\right)\sqrt{\mathrm{e}^x-m}=0$. Gọi $S$ là tập hợp giá trị $m$ nguyên với $m\in[-10;10]$ để phương trình có đúng hai nghiệm. Tổng giá trị các phần tử của $S$ bằng
	\choice
	{\True $-28$}
	{$-3$}
	{$-27$}
	{$-12$}
	\loigiai{
		Ta có: $\left(\log_2^2x-\log_2\dfrac{x^3}{4}\right)\sqrt{\mathrm{e}^x-m}=0$.\\
		Điều kiện: $\hoac{&x>0\\&m<\mathrm{e}^x}$ \\
		$\Leftrightarrow\hoac{&\log_2^2x-\log_2\dfrac{x^3}{4}=0\\&\sqrt{\mathrm{e}^x-m}=0} $ \\
		$ \Leftrightarrow\hoac{&\log_2^2x-3\log_2x+2=0\\&\mathrm{e}^x=m} $ \\
		$\Leftrightarrow\hoac{&\log_2x=1\\&\log_2x=2\\&\mathrm{e}^x=m}\Leftrightarrow\hoac{&x=2\\&x=4\\&\mathrm{e}^x=m.} $ \\
		Ta có\\
		Trường hợp 1: $m\leq 0$.\\
		Khi đó phương trình $\mathrm{e}^x=m$ vô nghiệm.\\
		Trường hợp 2: $m>0$.\\
		$\left(\log_2^2x-\log_2\dfrac{x^3}{4}\right)\sqrt{\mathrm{e}^x-m}=0\Leftrightarrow\heva{&\hoac{&x=2\\&x=4\\&x=\ln m}\\&x\geq\ln m.}$ \\
		Để phương trình chỉ có hai nghiệm phân biệt thì: $\hoac{&2\leq\ln(m)<4\\&\ln m\leq 0}\Leftrightarrow\mathrm{e}^2\leq m<\mathrm{e}^4$.\\
		Ngoài ra khi $m=1$ thì $\mathrm{e}^x=1\Leftrightarrow x=0(l)$.\\
		Nên $m=1,8;9;10$.\\
		Vậy $m\in\left\{-10;-9;-8;-7;-6;-5;-4;-3;-2;-1;1;8;9;10\right\}\Rightarrow\sum m= -28$.\\
		Chọn $A$.
	}
\end{ex}
\begin{ex}
	[Chuyên Ngoại Ngữ Hà Nội- 2021]%Câu 17.
	Cho phương trình $\ln(x+m)-\mathrm{e}^x+m=0$, với $m$ là tham số thực. Có bao nhiêu giá trị nguyên $m\in[-2021;2021]$ để phương trình đã cho có nghiệm?
	\choice
	{$2022$}
	{$4042$}
	{$2019$}
	{\True $2021$}
	\loigiai{
		Điều kiện: $x+m>0$.\\
		Đặt $t=\ln(x+m)\Leftrightarrow x+m=\mathrm{e}^t$, ta có hệ phương trình sau: $\heva{&x+m=\mathrm{e}^t\\&t+m=\mathrm{e}^x.}$ \\
		Suy ra $x-t=\mathrm{e}^t-\mathrm{e}^x\Leftrightarrow\mathrm{e}^x+x=\mathrm{e}^t+t\quad(*)$.\\
		Xét hàm số $f(x)=\mathrm{e}^x+x$, có $f'(x)=\mathrm{e}^x+1>0,\forall x\in\mathbb{R}\Rightarrow f(x)$ luôn đồng biến trên khoảng $(-\infty;+\infty)$.\\
		Ta thấy $(*)$ có dạng $f(x)=f(t)\Leftrightarrow x=t$.\\
		Khi đó ta có phương trình $x+m=\mathrm{e}^x\Leftrightarrow m=\mathrm{e}^x-x$.\\
		Xét hàm số $g(x)=\mathrm{e}^x-x$, có $g'(x)=\mathrm{e}^x-1;\quad g'(x)=0\Leftrightarrow\mathrm{e}^x=1\Leftrightarrow x=0$.\\
		Bảng biến thiên\\
		{\color{red} HINH O DAY}\\
		Dựa vào bảng biến thiên ta thấy phương trình đã cho có nghiệm khi $m\geq 1$.\\
		Mà $\heva{&m\in[-2021;2021]\\&m\in\mathbb{Z}}$ nên ta có $m\in\left\{1;2;\ldots;2020;2021\right\}$. Tức là có 2021 số nguyên $m$ thỏa mãn đề bài.
	}
\end{ex}
\begin{ex}
	[Chuyên Lê Khiết - Quảng Ngãi - 2021]%Câu 18.
	Có bao nhiêu số nguyên dương $m$ để phương trình $\mathrm{e}^x-1=m\ln(mx+1)$ có hai nghiệm phân biệt trên đoạn $[-10;10]$?
	\choice
	{\True $2201$}
	{$2020$}
	{$2021$}
	{$2202$}
	\loigiai{
		Điều kiện $mx+1>0$.\\
		Ta có $\mathrm{e}^x-1=m\ln(mx+1)$ \\
		$ \Leftrightarrow\mathrm{e}^x+mx=mx+1+m\ln(mx+1) $ \\
		$\Leftrightarrow\mathrm{e}^x+mx=\mathrm{e}^{\ln(mx+1)}+m\ln(mx+1) $ (1).\\
		Xét hàm số $f(t)=\mathrm{e}^t+mt,t\in\mathbb{R}$.\\
		Có $f'(t)=\mathrm{e}^t+m>0,\forall t\in\mathbb{R},m>0$. Suy ra hàm $f(t)$ đồng biến trên $\mathbb{R}$.\\
		Từ (1) ta được $f(x)=f\left(\ln(mx+1)\right)\Leftrightarrow x=\ln(mx+1)\Leftrightarrow\mathrm{e}^x=mx+1$ (2).\\
		Ta thấy (2) luôn có một nghiệm $x=0\in[-10;10]$. Do đó ta cần tìm các giá trị của $m$ để (2) có đúng một nghiệm $x\neq 0,x\in[-10;10]$.\\
		Với $x\neq 0$ thì (2) $\Leftrightarrow\dfrac{\mathrm{e}^x-1}{x}=m$.\\
		Xét hàm $g(x)=\dfrac{\mathrm{e}^x-1}{x},x\in[-10;10]\setminus\{0\}$.\\
		Ta có $g'(x)=\dfrac{x\mathrm{e}^x-\mathrm{e}^x+1}{x^2}$.\\
		Đặt $h(x)=x\mathrm{e}^x-\mathrm{e}^x+1,x\in\mathbb{R}$.\\
		Có $h'(x)=x\mathrm{e}^x,h'(x)=0\Leftrightarrow x=0$.\\
		Ta thấy $\lim\limits_{x\to+\infty} h(x)=+\infty,\lim\limits_{x\to-\infty} h(x)=1,h(0)=0$.\\
		Bảng biến thiên của hàm $h(x)$ như sau.\\
		{\color{red} HINH O DAY}\\.\\
		Từ bảng biến thiên suy ra $h(x)\geq 0,\forall x\in\mathbb{R}\Rightarrow g(x)>0,\forall x\neq 0$.\\
		Ta có $\lim\limits_{x\to 0^+} g(x)=1,\lim\limits_{x\to 0^-} g(x)=1$.\\
		Bảng biến thiên của hàm $y=g(x)$ với $x\in[-10;10]\setminus\{0\}$ như sau.\\
		{\color{red} HINH O DAY}\\.\\
		Từ bảng biến thiên suy ra (2) có đúng một nghiệm $x\neq 0,x\in[-10;10]\Leftrightarrow m\in\left[\dfrac{1-\mathrm{e}^{-10}}{10},\dfrac{\mathrm{e}^{10}-1}{10}\right]\setminus\{1\}$.\\
		Do $m$ nguyên dương nên $m\in\left\{2,3,4,\ldots,2202\right\}$. Vậy có $2201$ giá trị của $m$ thỏa mãn yêu cầu bài toán.
	}
\end{ex}
\begin{ex}
	[Chuyên Lê Hồng Phong - Nam Định - 2021]%Câu 19.
	Có bao nhiêu số nguyên dương $m$ nhỏ hơn $20$ thỏa mãn phương trình $\log\left(mx+\log m^m\right)=10^x$ có đúng hai nghiệm thực $x$ phân biệt. 
	\choice
	{\True $13$}
	{$12$}
	{$10$}
	{$11$}
	\loigiai{
		Điều kiện: $\heva{&m^m>0\\&mx+\log m^m>0}\Leftrightarrow\heva{&m>0\\&x+\log m>0.}$ \\
		Đặt $t=10^x>0\Rightarrow x=\log t$.\\
		Ta có:\\
		$\log\left(mx+\log m^m\right)=10^x\Leftrightarrow\log\left(m\log t+m\log m\right)=t\Leftrightarrow m\log t+m\log m=10^t$ \\
		$ \Leftrightarrow\log t+\log m=\dfrac{10^t}{m}\Leftrightarrow\log t+\log m=\dfrac{10^t}{10^{\log m}}\Leftrightarrow\log t+\log m=10^{t-\log m} $ \\
		$ \Leftrightarrow 10^{\log t}+\log t=10^{t-\log m}+t-\log m(*) $.\\
		Xét hàm số $f(t)=10^{\log t}+t$ với $t>0$.\\
		Ta có $f'(t)=10^t\ln 10+1>0\forall t\in\mathbb{R}$. Suy ra hàm số $f(t)=10^{\log t}+t$ đồng biến với $t>0$.\\
		Vậy $(*)\Leftrightarrow\log t=t-\log m$ hay $x=10^x-\log m\Leftrightarrow\log m=10^x-x$.\\
		Xét hàm số $g(x)=10^x-x$.\\
		Ta có $g'(x)=10^x\ln 10-1$.\\
		Khi đó $g'(x)=0\Leftrightarrow 10^x\ln 10-1=0\Leftrightarrow x=-\log(\ln 10)$.\\
		Bảng biến thiên\\
		{\color{red} HINH O DAY}\\
		Vậy phương trình đã cho có hai nghiệm $\Leftrightarrow\log m>g\left(-\log(\ln 10)\right)\Leftrightarrow m>10^{g\left(-\log(\ln 10)\right)}\approx 6,3$.\\
		Do $\heva{&m\in\mathbb{Z}\\&m<20}\Rightarrow m\in\left\{7; 8; 9;\ldots; 19\right\}$. Suy ra có $13$ giá trị $m$ thỏa mãn yêu cầu bài toán.
	}
\end{ex}
\begin{ex}[Chuyên Quốc Học Huế - 2021]%Câu 20.
	Giả sử $x_0$ là nghiệm thực của phương trình $2021\cdot 2^{-\cos x}=\log_{\pi}x^{2021}+\log_x\pi^{2021}$. Khẳng định nào sau đây là đúng?
	\choice
	{$x_0\in(2\pi;4\pi)$}
	{\True $x_0\in(0;2\pi)$}
	{$x_0\in(4\pi;6\pi)$}
	{$x_0\in(-2\pi;0)$}
	\loigiai{
		Điều kiện: $x>0; x\neq 1$.\\
		Khi đó, $2021\cdot 2^{-\cos x}=\log_{\pi}x^{2021}+\log_x\pi^{2021}\Leftrightarrow 2021\cdot 2^{-\cos x}=2021\cdot\log_{\pi}x+2021\cdot\log_x\pi$ \\
		$ \Leftrightarrow 2^{-\cos x}=\log_{\pi}x+\log_x\pi(1) $.\\
		Ta có: $-1\leq-\cos x\leq 1\Rightarrow\dfrac{1}{2}\leq 2^{-\cos x}\leq 2\Rightarrow\log_{\pi}x+\log_x\pi>0\Rightarrow\heva{&\log_{\pi}x>0\\&\log_x\pi>0.}$ \\
		(do $\log_{\pi}x$ và $\log_x\pi$ luôn cùng dấu).\\
		Áp dụng bất đẳng thức Cô-si ta được: $\log_{\pi}x+\log_x\pi\geq 2\sqrt{\log_{\pi}x\cdot\log_x\pi}\Leftrightarrow\log_{\pi}x+\log_x\pi\geq 2$ \\
		$ \Rightarrow $ Phương trình (1) có: $\heva{&\mathrm{VT}\leq 2\\&\mathrm{VP}\geq 2.}$ \\
		Do đó PT(1) có nghiệm $\Leftrightarrow\heva{&\mathrm{VT}=2\\&\mathrm{VP}=2}\Leftrightarrow\heva{&-\cos x=1\\&\log_x\pi=\log_{\pi}x=1}\Leftrightarrow x=\pi (x\in (0; 2\pi))$.
	}
\end{ex}
\begin{ex}
	[Sở Bạc Liêu - 2021]%Câu 21.
	Có bao nhiêu giá trị nguyên của tham số $m$ thuộc đoạn $[-2021; 2021]$ để phương trình $2^{x-2+\sqrt[3]{m-3x}}+\left(x^3-6x^2+9x+m\right)2^{x-2}=2^{x+1}+1$ có một nghiệm duy nhất?
	\choice
	{\True $4038$}
	{$3$}
	{$2021$}
	{$4039$}
	\loigiai{
		Chia cả hai vế của phương trình cho $2^{x-2}$ ta được:\\
		$2^{\sqrt[3]{m-3x}}+\left(x^3-6x^2+9x+m\right)=8+2^{2-x}\Leftrightarrow 2^{\sqrt[3]{m-3x}}+m-3x=2^{2-x}+(2-x)^3\quad(*)$.\\
		Xét hàm đặc trưng: $f(t)=2^t+t^3$.\\
		Ta có $f'(t)=2^t\ln 2+3t^2>0,\forall t\in\mathbb{R}$ nên $f(t)$ đồng biến trên $\mathbb{R}$.\\
		$(*)\Rightarrow\sqrt[3]{m-3x}=2-x\Leftrightarrow m=-x^3+6x^2-9x+8$.\\
		Xét hàm số $f(x)=-x^3+6x^2-9x+8$ có bảng biến thiên như sau: 
		{\color{red} HINH O DAY}\\
		Phương trình có nghiệm duy nhất khi và chỉ khi: $\hoac{&m>8\\&m<4.}$ \\
		Kết hợp với yêu cầu đề bài ta có: $\hoac{&9\leq m\leq 2021\\&-2021\leq m\leq 3}$. Vậy có tất cả $4038$ số nguyên $m$ thỏa mãn.
	}
\end{ex}
\begin{ex}
	[Sở Cần Thơ - 2021]%Câu 22.
	Cho phương trình $\log_2(2x-m)=4^x+m$ với $m$ là tham số thực. Có bao nhiêu giá trị nguyên của $m\in(-27;27)$ sao cho phương trình trên có nghiệm?
	\choice
	{$10$}
	{\True $26$}
	{$1$}
	{$53$}
	\loigiai{
		Điều kiện: $2x-m>0$.\\
		Ta có $\log_2(2x-m)=4^x+m$ \\
		$ \Leftrightarrow\log_2(2x-m)+2x-m=4^x+2x\Leftrightarrow 2^{\log_2(2x-m)}+\log_2(2x-m)=2^{2x}+2x $ (1).\\
		Xét hàm số $f(t)=2^t+t,t\in\mathbb{R}$.\\
		Có $f'(t)=2^t\ln 2+1>0,\forall t\in\mathbb{R}$. Suy ra hàm $f(t)$ đồng biến trên $\mathbb{R}$.\\
		Từ (1) ta được $f\left(\log_2(2x-m)\right)=f(2x)\Leftrightarrow\log_2(2x-m)=2x$ \\
		$ \Leftrightarrow 4^x=2x-m\Leftrightarrow 4^x-2x=-m $ (2).\\
		Xét hàm $g(x)=4^x-2x$.\\
		Ta có $g'(x)=4^x\ln 4-2=4^x\cdot 2\ln 2-2$.\\
		Có $g'(x)=0\Leftrightarrow x=x_0=\log_4\dfrac{1}{\ln 2}$.\\
		Bảng biến thiên của hàm $g(x)$ như sau.\\
		{\color{red} HINH O DAY}\\.\\
		Từ bảng biến thiên suy ra phương trình (2) có nghiệm\\
		$ \Leftrightarrow-m\geq g(x_0)\Leftrightarrow m\leq-g(x_0)=-g\left(\log_4\dfrac{1}{\ln 2}\right)\approx-0,914 $.\\
		Do $m\in\mathbb{Z}$ và $m\in(-27;27)$ nên ta đươc $m\in\left\{-26;-25;\ldots;-1\right\}$.\\
		Vậy có $26$ giá trị của $m$ thỏa mãn yêu cầu bài toán.
	}
\end{ex}
\begin{ex}
	[Đại học Hồng Đức – 2022]%Câu 23.
	Tổng $S$ của tất cả các nghiệm thuộc khoảng $(0;4\pi)$ của phương trình $2022^{\sin^2x}-2022^{\cos^2x}=2\ln (\cot x$) là
	\choice
	{$S=18\pi$}
	{$S=8\pi$}
	{\True $S=7\pi$}
	{$S=16\pi$}
	\loigiai{
		Điều kiện $\cot x>0$. Ta có\\
		$\begin{aligned}&{2022}^{\sin^2x}-2022^{\cos^2x}=2\ln (\cot x)\Leftrightarrow{2022}^{\sin^2x}-2022^{\cos^2x}=\ln\left(\cos^2x\right)-\ln\left(\sin^2x\right)\\&\Leftrightarrow{2022}^{\sin^2x}+\ln\left(\sin^2x\right)={2022}^{\cos^2x}+\ln\left(\cos^2x\right).\end{aligned}$\\
		Xét hàm số $f(t)=2022^t+\ln t$ với $t>0$.\\
		$f'(t)=2022^t\cdot\ln 2022+\dfrac{1}{t}>0,\forall t>0\Rightarrow$ hàm số $f(t)$ đồng biến trên khoảng $(0;+\infty)$.\\
		Khi đó $(1)\Leftrightarrow f\left(\sin^2x\right)=f\left(\cos^2x\right)\Leftrightarrow\sin^2x=\cos^2x\Leftrightarrow\cos 2x=0\Leftrightarrow x=\dfrac{\pi}{4}+\dfrac{k\pi}{2},k\in\mathbb{Z}$.\\
		Do $\cot x>0$ nên $x=\dfrac{\pi}{4}+k\pi,k\in\mathbb{Z}$.\\
		Mà $x\in (0;4\pi)$ suy ra $x\in\left\{\dfrac{\pi}{4};\dfrac{5\pi}{4};\dfrac{9\pi}{4}\colon\dfrac{13\pi}{4}\right\}$. Suy ra $S=7\pi$.
	}
\end{ex}
\begin{ex}
	[Liên trường Hà Tĩnh – 2022]%Câu 24.
	Có bao nhiêu giá trị nguyên của $m$ để phương trình $\left[\log_2(x-1)+x-2\right]\left(4^x-2^{x+3}+m-1\right)=0$ có ba nghiệm phân biệt
	\choice
	{$2$}
	{\True $3$}
	{$4$}
	{$5$}
	\loigiai{
		$+$ Phương trình đã cho $\Leftrightarrow\hoac{&\log_2(x-1)+x-2=0(1)\\&4^x-2^{x+3}+m-1=0(2).}$ \\
		Xét hàm số $f(x)=\log_2(x-1)+x-2=0$. Ta có $f'(x)=\dfrac{1}{(x-1)\ln 2}+1>0,\forall x>1$.\\
		Lại có $f(2)=0$ suy ra phương trình (1) có đúng 1 nghiệm $x=2$.\\
		Yêu cầu bài toán $PT$ (2) phải có hai nghiệm phân biệt khác 2. Suy ra phương trình $t^2-8t-1=-m$ phải có hai nghiệm phân biệt khác 4 thỏa mãn $2<t_1<t_2$.\\
		Xét hàm số $f(t)=t^2-8t-1$ có bảng biến thiên:\\
		{\color{red} HINH O DAY}\\
		$-17 <-m <-13\Leftrightarrow 13<m<17$.
	}
\end{ex}
\begin{ex}[Sở Bạc Liêu 2022]%Câu 25.
	Cho phương trình $\left(2\log_3^2x-\log_3x-1\right)\sqrt{5^x-m}=0$ ($m$ là tham số thực). Có tất cả bao nhiêu giá trị nguyên dương của $m$ để phương trình đã cho có đúng $2$ nghiệm phân biệt?
	\choice
	{$125$}
	{\True $123$}
	{$122$}
	{$124$}
	\loigiai{
		Điều kiện: $\heva{&x>0\\&x\geq\log_5m.}$ \\
		Với $m=1$, điều kiện của phương trình là $x>0$ và phương trình trở thành.\\
		$\left(2\log_3^2x-\log_3x-1\right)\sqrt{5^x-1}=0\Leftrightarrow\hoac{&2\log_3^2x-\log_3x-1=0\\&5^x-1=0}\Leftrightarrow\hoac{&\log_3x=1\\&\log_3x=-\dfrac{1}{2}\\&x=0(l)}\Leftrightarrow\hoac{&x=3(\text{nhận})\\&x=\dfrac{1}{\sqrt{3}}(\text{nhận}).}$ \\
		Suy ra $m=1$ thỏa yêu cầu bài toán.\\
		Với $m=2$, điều kiện của phương trình là $x\geq\log_52$ và phương trình trở thành.\\
		$\left(2\log_3^2x-\log_3x-1\right)\sqrt{5^x-2}=0\Leftrightarrow\hoac{&2\log_3^2x-\log_3x-1=0\\&5^x-2=0}\Leftrightarrow\hoac{&\log_3x=1\\&\log_3x=-\dfrac{1}{2}\\&x=\log_52(n)}\Leftrightarrow\hoac{&x=3(\text{nhận})\\&x=\dfrac{1}{\sqrt{3}}(\text{nhận})\\&x=\log_52(\text{nhận}).}$ \\
		Suy ra $m=2$ không thỏa yêu cầu bài toán.\\
		Với $m\geq 3$, điều kiện phương trình là $x\geq\log_5m$, khi đó\\
		$\left(2\log_3^2x-\log_3x-1\right)\sqrt{5^x-m}=0\Leftrightarrow\hoac{&2\log_3^2x-\log_3x-1=0\\&5^x-m=0}$ \\
		$\Leftrightarrow\hoac{&\log_3x=1\\&\log_3x=-\dfrac{1}{2}\\&5^x=m}\Leftrightarrow\hoac{&x=3(\text{nhận})\\&x=\dfrac{1}{\sqrt{3}}(\text{loại})\\&5^x=m.} $ \\
		Do đó phương trình có đúng 2 nghiệm khi và chỉ khi $3>\log_5m\Leftrightarrow m<5^3=125$.\\
		Suy ra $m\in\left\{3;4;5;\ldots;124\right\}$ nên có $122$ giá trị của $m$.\\
		Vậy có tất cả $123$ giá trị của $m$.
	}
\end{ex}
\begin{ex}
	[Chuyên Lam Sơn 2022]%Câu 26.
	Có bao nhiêu số nguyên dương $m$ để phương trình $m\left(\mathrm{e}^x-1\right)\cdot\ln (mx+1)+2\mathrm{e}^x=\mathrm{e}^{2x}+1$ có 2 nghiệm phân biệt không lớn hơn 5. 
	\choice
	{$26$}
	{$27$}
	{$29$}
	{\True 28}
	\loigiai{
		Xét phương trình $m\left(\mathrm{e}^x-1\right)\cdot\ln (mx+1)+2\mathrm{e}^x=\mathrm{e}^{2x}+1(*)$ điều kiện $mx+1>0$.\\
		$(*)\Leftrightarrow\hoac{&\mathrm{e}^x-1=0\\&\mathrm{e}^x-1=m\cdot\ln (mx+1).}$ \\
		$\mathrm{e}^x-1=0\Leftrightarrow x=0\mathrm{e}^x-1=m\cdot\ln (mx+1)$,\\
		Đặt $y=\ln (mx+1)\Rightarrow\mathrm{e}^x-1=my$. Ta có hệ phương trình $\heva{&x=\ln (my+1)(1)\\&y=\ln (mx+1)(2).}$ \\
		Trừ (1) và (2) theo vế ta được: $x-y=\ln (my+1)-\ln (mx+1)$ hay $x+\ln (mx+1)=y+\ln (my+1)$ với $m>0$ thì hàm số $f(x)=x+\ln (mx+1)$ đồng biến trên tập xác định nên $x+\ln (mx+1)=y+\ln (my+1)\Leftrightarrow x=y$.\\
		Thay $x=y$ vào (1) ta được $x=\ln (mx+1)$ hay $\mathrm{e}^x=mx+1(4)$.\\
		Rõ ràng $x=0$ là 1 nghiệm của phương trình (4).\\
		Với $x\neq 0$ ta có $(4)\Leftrightarrow m=\dfrac{\mathrm{e}^x-1}{x}$.\\
		Xét hàm số $g(x)=\dfrac{\mathrm{e}^x-1}{x}$, ta có: Tập xác định $\mathscr{D}=\mathbb{R}\setminus\{0\}$ và $g'(x)=\dfrac{x\mathrm{e}^x-\mathrm{e}^x+1}{x^2} g'(x)=0\Leftrightarrow x\mathrm{e}^x-\mathrm{e}^x+1=0$.\\
		Hàm số $h(x)=x\mathrm{e}^x-\mathrm{e}^x+1$ có $h'(x)=x\mathrm{e}^x$ nên $h'(x)=0\Leftrightarrow x=0$.\\
		Ta có bảng biến thiên của $h(x)$ như sau: 
		{\color{red} HINH O DAY}\\
		Suy ra $h(x)\geq 0,\forall x$ do đó $g'(x)>0,\forall x\neq 0$.\\
		Bảng biến thiên của $g(x)$:\\
		{\color{red} HINH O DAY}\\
		Để phương trình $\mathrm{e}^x-1=\ln (mx+1)^m$ có 2 nghiệm phân biệt không lớn hơn 5 thì phương trình $m=g(x)$ có duy nhất 1 nghiệm bé hơn hoặc bằng 5. Ta có $g(5)=\dfrac{\mathrm{e}^5-1}{5}\approx 29,5$.\\
		Dựa vào bảng biến thiên của $g(x)$ ta có $\heva{&0<m\leq g(5)\\&m\neq 1}$ do $m\in\mathbb{N}^*$ nên có 28 giá trị thỏa mãn.
	}
\end{ex}
\begin{ex}
	[Chuyên Lương Văn Tụy-Ninh Bình 2022]%Câu 27.
	Cho phương trình $\left(2\log_3^2x-\log_3x-1\right)\sqrt{5^x-m}=0$ ($m$ là tham số thực). Có tất cả bao nhiêu giá trị nguyên dương của $m$ để phương trình đã cho có đúng hai nghiệm phân biệt?
	\choice
	{$125$}
	{\True $123$}
	{$122$}
	{$124$}
	\loigiai{
		Điều kiện $\heva{&x>0\\&5^x-m\geq 0}\Leftrightarrow\heva{&x>0\\&m\leq 5^x}$..\\
		$\left(2\log_3^2x-\log_3x-1\right)\sqrt{5^x-m}=0\Leftrightarrow\heva{&5^x-m\geq 0,x>0\\&\hoac{&5^x-m=0\\&2\log_3^2x-\log_3x-1=0}.}$ \\
		$\heva{&5^x-m\geq 0,x>0\\&\hoac{&5^x-m=0\\&\log_3x=\dfrac{-1}{2}\\&\log_3x=1}}\Leftrightarrow\heva{&5^x-m\geq 0,x>0\\&\hoac{&5^x-m=0\\&\log_3x=\dfrac{-1}{2}\\&\log_3x=1}.}$ \\
		$+$ Khi $m=1\Rightarrow x=\log_21=0$ vậy phương trình $\left(2\log_3^2x-\log_3x-1\right)\sqrt{5^x-m}=0$ có 2 nghiệm $\hoac{&x=3^{\tfrac{1}{\sqrt{2}}}\\&x=3.}$ \\
		$+m>1\Rightarrow x=\log_5m$ là 1 nghiệm. Để phương trình có đúng 2 nghiệm thì $\dfrac{1}{\sqrt{3}}\leq\log_5m<3\Leftrightarrow 5^{\tfrac{1}{\sqrt{3}}}\leq m<5^3\Leftrightarrow 2,53\leq m<125$.
	}
\end{ex}
\begin{ex}
	[Chuyên KHTN 2022]%Câu 28.
	Có bao nhiêu số nguyên dương $a$ để phương trình sau đây có ít nhất một nghiệm thực.\\
	$\left(a^{\log x}+1\right)^{\log a}+a^{\log x}=2x-2$ 
	\choice
	{$8$}
	{$1$}
	{$0$}
	{\True $9$}
	\loigiai{
		Điều kiện $x>0$.\\
		Vì $a^{\log x}=x^{\log a}$ nên đặt $x^{\log a}+1=t$ ta có $\heva{&x^{\log a}+1=t\\&t^{\log a}+t=2x-1}\Rightarrow x^{\log a}+2x=t^{\log a}+2t$..\\
		Vì $a\geq 1\Rightarrow\log a\geq 0\Rightarrow f(x)=x^{\log a}+2x$ đồng biến trên $(0;+\infty)$. Vì vậy $f(x)=f(t)\Rightarrow x=t$.\\
		Hay $x^{\log a}=x-1$.\\
		Nếu $a=1$ thay lại vào phương trình ta có $x^{\circ}=x-1\Leftrightarrow x=2$. Vậy $a=1$ thỏa mãn.\\
		Nếu $a>1\Rightarrow x=t=x^{\log a}+1>2$. Vậy $\log a=\dfrac{\ln (x-1)}{\ln x}<\dfrac{\ln x}{\ln x}=1$.\\
		Do vậy $a\in (0;10)$. Mà $a$ nguyên dương và lớn hơn 1 nên $a\in\{2,3,4,\ldots,9\}$.\\
		Vậy có 9 giá trị của $a$ thỏa mãn điều kiện đầu bài.
	}
\end{ex}
\begin{ex}
	[Chuyên Lê Khiết - Quảng Ngãi 2022]%Câu 29.
	Có bao nhiêu số nguyên dương $a$ sao cho tồn tại số thực $x$ thỏa mãn phương trình sau.\\
	$2021^{x^3-a^{3\log(x+1)}}\left(x^3+2020\right)=a^{3\log(x+1)}+2020$ 
	\choice
	{\True $9$}
	{$5$}
	{$8$}
	{$12$}
	\loigiai{
		Điều kiện: $x\geq-1$.\\
		Phương trình.\\
		$2021^{x^3-a^{3\log(x+1)}}\left(x^3+2020\right)=a^{3\log(x+1)}+2020$ \\
		$ \Leftrightarrow 2021^{x^3}\left(x^3+2020\right)=2021^{a^{3\log(x+1)}}\left[a^{3\log(x+1)}+2020\right] $ (1).\\
		Xét hàm số $f(t)=2021^t(t+2020), t >-1$.\\
		$f'(t)=(t+2020)2021^t\ln 2021+2021^t>0,\forall t >-1$.\\
		Hàm số $f(t)$ đồng biến trên $(-1;+\infty)$, nên từ (1) $\Rightarrow x^3=a^{3\log(x+1)}$. (*).\\
		TH1: $-1<x\leq 0$ không tồn tại $a$ thỏa mãn (*).\\
		TH2: $x>0$ (*) $\Leftrightarrow\log_ax=\log(x+1)\Leftrightarrow\log a=\dfrac{\log x}{\log(x+1)}<1\Leftrightarrow a<10$.\\
		Vì $a$ nguyên dương $\Rightarrow a\in\left\{1; 2; 3;\ldots; 9\right\}$.
	}
\end{ex}
\begin{dang}
	{Phương trình mũ – logarit chứa nhiều ẩn}
\end{dang}     
\begin{ex}
	[Đề Minh Họa 2020 Lần 1]%Câu 1.     
	Có bao nhiêu cặp số nguyên $(x; y)$ thỏa mãn $0\leq x\leq 2020$ và $\log_3(3x+3)+x=2y+9^y$?
	\choice
	{$2019$}
	{$6$}
	{$2020$}
	{\True $4$}
	\loigiai{
		Cách 1.\\
		Ta có\\ $\log_3(3x+3)+x=2y+9^y\Leftrightarrow\log_3(x+1)+x+1=2y+3^{2y}$ $(1)$.\\
		Đặt $\log_3(x+1)=t\Rightarrow x+1=3^t$.\\
		Phương trình $(1)$ trở thành: $t+3^t=2y+3^{2y} (2)$.\\
		Xét hàm số $f(u)=u+3^u$ trên $\mathbb{R}$.\\
		$f'(u)=1+3^u\ln 3>0,\forall u\in\mathbb{R}$ nên hàm số $f(u)$ đồng biến trên $\mathbb{R}$.\\
		Do đó $(2)\Leftrightarrow f(t)=f(2y)\Leftrightarrow t=2y\Rightarrow\log_3(x+1)=2y\Leftrightarrow x+1=9^y\Leftrightarrow x=9^y-1$.\\
		Vì $0\leq x\leq 2020\Rightarrow 0\leq 9^y-1\leq 2020\Leftrightarrow 1\leq 9^y\leq 2021\Leftrightarrow 0\leq y\leq\log_92021$.\\
		$\left(\log_32021\approx 3,464\right)$.\\
		Do $y\in\mathbb{Z}\Rightarrow y\in\{0; 1; 2; 3\}$, có 4 giá trị của y nên cũng có 4 giá trị của $x$.\\
		Vậy có 4 cặp số nguyên $(x; y)$.\\
		Cách 2.\\
		Ta có\\ $\log_3(3x+3)+x=2y+9^y\Leftrightarrow\log_3(x+1)+x+1=2y+3^{2y}$.\\
		Xét hàm số $f(x)=\log_3(x+1)+x+1$ với $x\in[0; 2020]$.\\
		Ta có $f'(x)=\dfrac{1}{(x+1)\ln 3}+1>0,\forall x\in x\in[0; 2020]\Rightarrow$ Hàm số $f(x)$ đồng biến trên đoạn $[0; 2020]$.\\
		Suy ra $f(0)\leq f(x)=\log_3(x+1)+x+1\leq f(2020)\Leftrightarrow 1\leq f(x)\leq\log_22021+2021$ \\
		$ \Rightarrow 1\leq 2y+9^y\leq\log_32021+2021<2028 $.\\
		Nếu $y<0\Rightarrow 2y+9^y<9^y<9^{\circ}=1\Rightarrow y\geq 0$.\\
		Khi đó $y\in\mathbb{N}\Rightarrow\left(2y+9^y\right)\in\mathbb{N}\Rightarrow 2y+9^y\leq 2027\Rightarrow 9^y\leq 2027-2y\leq 2027$ \\
		$ \Rightarrow y\leq\log_92027\approx 3,465\Rightarrow y\leq 3\Rightarrow 0\leq y\leq 3 $ \\
		$ \Rightarrow y\in\{0; 1; 2; 3\} $. Do $f(x)$ là hàm số luôn đồng biến nên với mỗi giá trị của $y$ chỉ cho 1 giá trị của $x$.\\
		$y=0\Rightarrow\log_3(x+1)+x+1=1\Leftrightarrow x=0$.\\
		$y=1\Rightarrow\log_3(x+1)+x+1=11\Leftrightarrow\log_3(x+1)+x=10\Leftrightarrow x=8$.\\
		$y=2\Rightarrow\log_3(x+1)+x+1=85\Leftrightarrow\log_3(x+1)+x=84\Leftrightarrow x=80$.\\
		$y=3\Rightarrow\log_3(x+1)+x+1=735\Leftrightarrow\log_3(x+1)+x=734\Leftrightarrow x=729$.\\
		Vậy có 4 cặp số nguyên $(x; y)$.
	}
\end{ex}
\begin{ex}
	[Đề Tham Khảo 2020 Lần 2]%Câu 2.
	Có bao nhiêu số nguyên $x$ sao cho tồn tại số thực $y$ thỏa mãn $\log_3(x+y)=\log_4\left(x^2+y^2\right)$?
	\choice
	{$3$}
	{\True $2$}
	{$1$}
	{Vô số}
	\loigiai{
		Cách 1.\\
		Đặt $t=\log_3(x+y)=\log_4\left(x^2+y^2\right)\Rightarrow\heva{&x+y=3^t\\&x^2+y^2=4^t}(1)$.\\
		Áp dụng bất đẳng thức Cauchy, ta có\\
		$9^t=(x+y)^2\leq 2\left(x^2+y^2\right)=4^t\Rightarrow\dfrac{9^t}{4^t}\leq 2\Rightarrow t\leq\log_{\frac{9}{2}}2$.\\
		Như vậy,\\
		$x^2+y^2=4^t\Rightarrow x^2\leq 4^t\leq 4^{\log_{\frac{9}{4}}2}\approx 1,89\Rightarrow x\in\{-1;0;1\}$.\\
		Trường hợp 1: $x=0\Rightarrow\heva{&y=3^t\\&y^2=4^t}\Rightarrow\heva{&t=0\\&y=1.}$ \\
		Trường hợp 2: $x=1\Rightarrow\heva{&y=3^t-1\\&y^2=4^t-1}\Rightarrow\heva{&t=0\\&y=0.}$ \\
		Trường hợp 3: $x=-1\Rightarrow\heva{&y=3^t+1\\&y^2+1=4^t\geq 1}\Rightarrow\heva{&t\geq 0\\&y=3^t+1\geq 2}\Rightarrow x^2+y^2\geq 5$ mâu thuẫn với $x^2+y^2\leq 4^{\log_{\frac{3}{2}}\sqrt{2}}$ suy ra loại $x=-1$.\\
		Vậy có hai giá trị $x\in\{0;1\}$.\\
		Cách 2.\\
		Đặt $t=\log_3(x+y)=\log_4\left(x^2+y^2\right)\Rightarrow\heva{&x+y=3^t\\&x^2+y^2=4^t}(1)$.\\
		Suy ra $x,y$ là tọa độ của điểm $M$ với $M$ thuộc đường thẳng $d\colon x+y=3^t$ và đường tròn $(C)\colon x^2+y^2=4^t$.\\
		Để tồn tại $y$ tức tồn tại $M$ nên $d,(C)$ có điểm chung, suy ra $\mathrm{d}(O,d)\leq R$ trong đó $O(0;0), R=2^t$ nên $\dfrac{|-3^t|}{\sqrt{2}}\leq 2^t\Leftrightarrow t\leq\log_{\frac{3}{2}}\sqrt{2}$.\\
		Khi đó $(1)\Rightarrow\heva{&0<x+y\leq 3^{\log_{\frac{3}{2}}\sqrt{2}}\\&x^2+y^2\leq 4^{\log_{\frac{3}{2}}\sqrt{2}}.}$ \\
		Minh họa quỹ tích điểm $M$ như hình vẽ sau
		{\color{red} HINH O DAY}\\
		Ta thấy có 3 giá trị $x\in\mathbb{Z}$ có thể thỏa mãn là $x=-1;x=0;x=1$.\\
		Thử lại:\\
		Trường hợp 1: $x=0\Rightarrow\heva{&y=3^t\\&y^2=4^t}\Rightarrow\heva{&t=0\\&y=1.}$ \\
		Trường hợp 2: $x=1\Rightarrow\heva{&y=3^t-1\\&y^2=4^t-1}\Rightarrow\heva{&t=0\\&y=0.}$ \\
		Trường hợp 3: $x=-1\Rightarrow\heva{&y=3^t+1\\&y^2+1=4^t\geq 1}\Rightarrow\heva{&t\geq 0\\&y=3^t+1\geq 2}\Rightarrow x^2+y^2\geq 5$ mâu thuẫn với $x^2+y^2\leq 4^{\log_{\frac{3}{2}}\sqrt{2}}$ suy ra loại $x=-1$.}
\end{ex}
\begin{ex}[Mã 103 - 2020 Lần 2]%Câu 3.
	Có bao nhiêu cặp số nguyên dương $(m;n)$ sao cho $m+n\leq 10$ và ứng với mỗi cặp $(m;n)$ tồn tại đúng $3$ số thực $a\in(-1;1)$ thỏa mãn $2a^m=n\ln\left(a+\sqrt{a^2+1}\right)$?
	\choice
	{$7$}
	{$8$}
	{$10$}
	{\True $9$}
	\loigiai{
		Ta có $2a^m=n\ln\left(a+\sqrt{a^2+1}\right)\Leftrightarrow\dfrac{2a^m}{n}=\ln\left(a+\sqrt{a^2+1}\right)$.\\
		Xét hai hàm số $f(x)=\ln\left(x+\sqrt{x^2+1}\right)$ và $g(x)=\dfrac{2}{n}x^m$ trên $(-1;1)$.\\
		Ta có $f'(x)=\dfrac{1}{\sqrt{x^2+1}}>0$ nên $f(x)$ luôn đồng biến và $f(-x)=\ln\left(-x+\sqrt{x^2+1}\right)=\ln\left(\dfrac{1}{x+\sqrt{x^2+1}}\right)=-\ln\left(x+\sqrt{x^2+1}\right)=-f(x)$ nên $f(x)$ là hàm số lẻ.\\
		Nếu $m$ chẵn thì $g(x)$ là hàm số chẵn và có bảng biến thiên dạng
		{\color{red} HINH O DAY}\\
		Suy ra phương trình có nhiều nhất $2$ nghiệm, do đó $m$ lẻ.\\
		Nếu $m$ lẻ thì hàm số $g(x)$ là hàm số lẻ và luôn đồng biến.\\
		Ta thấy phương trình luôn có nghiệm $x=0$. Dựa vào tính chất đối xứng của đồ thị hàm số lẻ, suy ra phương trình đã cho có đúng $3$ nghiệm trên $(-1;1)$ khi có $1$ nghiệm trên $(0;1)$, hay $f(1)>g(1)\Leftrightarrow\ln(1+\sqrt{2})<\dfrac{2}{n}\Leftrightarrow n<\dfrac{2}{\ln(1+\sqrt{2})}\approx 2,26\Rightarrow n\in\{1;2\}$.\\
		Đối chiếu điều kiện, với $n=1$ suy ra $m\in\{1;3;5;7;9\}$, có $5$ cặp số thỏa mãn.\\
		Với $n=2$ thì $m\in\{1;3;5;7\}$ có $4$ cặp số thỏa mãn.\\
		Vậy có $9$ cặp số thỏa mãn bài toán.
	}
\end{ex}
\begin{ex}
	[Mã 101 - 2020 Lần 2]%Câu 4.
	Có bao nhiêu cắp số nguyên dương $(m,n)$ sao cho $m+n\leq 14$ và ứng với mỗi cặp $(m,n)$ tồn tại đúng ba số thực $a\in(-1;1)$ thỏa mãn $2a^m=n\ln\left(a+\sqrt{a^2+1}\right)$?
	\choice
	{$14$}
	{$12$}
	{\True $11$}
	{$13$}
	\loigiai{
		Xét $f(x)=\dfrac{2}{n}\cdot x^m-\ln\left(x+\sqrt{x^2+1}\right)$ trên $(-1;1)$.\\
		Đạo hàm $f'(x)=\dfrac{2m}{n}x^{m-1}-\dfrac{1}{\sqrt{x^2+1}}=0$.\\
		Theo đề bài $f(x)=0$ có ba nghiệm nên $\dfrac{2m}{n}x^{m-1}=\dfrac{1}{\sqrt{x^2+1}}$ có ít nhất hai nghiệm.\\
		Xét đồ thị của hàm $y=x^{m-1};y=\dfrac{1}{\sqrt{x^2+1}}$, suy ra $m-1$ chẵn và $m-1>0$.\\
		Suy ra $m\in\left\{3;5; 7; 9; 11;13\right\}$. Khi đó $f'(x)=0$ có nghiệm $\heva{&x_1<0\\&x_2>0.}$ \\
		Phương trình có 3 nghiệm $\Leftrightarrow\heva{&f(1)>0\\&f(-1)<0}$ \\
		$\Leftrightarrow\heva{&\dfrac{2}{n}>\ln(\sqrt{2}+1)\\&-\dfrac{2}{n}<\ln(\sqrt{2}-1)}\Leftrightarrow n\leq 2\Rightarrow n=\{1;2\} $.\\
		$n\in\{1; 2\}$ và $m\in\left\{3;5; 7; 9; 11;13\right\}$, do $m+n\leq 14$ nên ta có 11 cặp $(m; n)$ thỏa yêu cầu bài toán.
	}
\end{ex}
\begin{ex}
	[Mã 104 - 2020 Lần 2]%Câu 5.
	Có bao nhiêu cặp số nguyên dương $(m,n)$ sao cho $m+n\leq 12$ và ứng với mỗi cặp $(m,n)$ tồn tại đúng 3 số thực $a\in (-1,1)$ thỏa mãn $2a^m=n\ln (a+\sqrt{a^2+1})$?
	\choice
	{$12$}
	{$10$}
	{$11$}
	{\True $9$}
	\loigiai{
		Ta có $2a^m=n\ln (a+\sqrt{a^2+1})\Leftrightarrow\dfrac{2}{n}a^m=\ln (a+\sqrt{a^2+1})\quad (*)$.\\
		Xét hàm $f(a)=\ln (a+\sqrt{a^2+1})$ trên $(-1,1)$ (dễ thấy hàm $f$ lẻ, đồng biến trên $R$), có BBT: 
		{\color{red} HINH O DAY}\\
		Xét hàm $g(a)=\dfrac{2}{n}\cdot a^m$ trên $(-1,1)$.\\
		Với $m$ chẵn, $g(a)$ là hàm chẵn và $g(a)\geq 0,\forall a\in \mathbb{R}$, do đó $(*)$ không thể có 3 nghiệm.\\
		Với $m$ lẻ, $g(a)$ là hàm lẻ, đồng biến trên $R$ và tiếp tuyến của đồ thị tại điểm $a=0$ là đường thẳng $y=0$.\\
		Dễ thấy $(*)$ có nghiệm $a=0\in (-1;1)$. Để $(*)$ có đúng 3 nghiệm tức là còn có 2 nghiệm nữa là $\pm a_0$ với $0<a_0<1$.\\
		Muốn vậy, thì $g(1)=\dfrac{2}{n}\cdot 1^m=\dfrac{2}{n}>f(1)=\ln (1+\sqrt{2})\Leftrightarrow n<\dfrac{2}{\ln (1+\sqrt{2})}\approx 2,26\Rightarrow n=1;n=2$.\\
		Cụ thể:\\
		$m\in\{3;5;7;9\}$ thì $n\in\{1;2\}$: Có $8$ cặp $(m,n)$.\\
		$m=11$ thì $n\in\{1\}$: Có $1$ cặp $(m,n)$.\\
		$m=1$: Đồ thị hàm số $g(a)$ là đường thẳng ($g(a)=a;g(a)=2a$) không thể cắt đồ thị hàm số $f(a)$ tại giao điểm $a_0\neq 0$ được vì tiếp tuyến của hàm số $f(a)$ tại điểm có hoành độ $a=0$ là đường thẳng $y=a$.\\
		Vậy có cả thảy $9$ cặp $(m,n)$.
	}
\end{ex}
\begin{ex}
	[Chuyên Biên Hòa - Hà Nam - 2020]%Câu 6.
	Có tất cả bao nhiêu giá trị thực của tham số $m\in[-1; 1]$ sao cho phương trình $\log_{m^2+1}\left(x^2+y^2\right)=\log_2(2x+2y-2)$ có nghiệm nguyên $(x; y)$ duy nhất?
	\choice
	{$3$}
	{\True $2$}
	{$1$}
	{$0$}
	\loigiai{
		Điều kiện: $\heva{&x^2+y^2>0\\&x+y-1>0.}$ \\
		Nhận xét: Vì $x, y$ có vai trò như nhau nên nếu phương trình có nghiệm $(x_0; y_0)$ thì $(y_0;x_0)$ cũng là một nghiệm của phương trình.\\
		Điều kiện cần: Phương trình đã cho có nghiệm duy nhất $\Leftrightarrow x_0=y_0$.\\
		Thay vào phương trình ta được $\log_{m^2+1}\left(2x_0^2\right)=\log_2(4x_0-2)$.\\
		Vì $x_0\in\mathbb{Z}\Rightarrow 4x_0-2>1$. Lại có $2x_0^2\geq 4x_0-2\Rightarrow\log_2(4x_0-2)=\log_{m^2+1}\left(2x_0^2\right)\geq\log_{m^2+1}(4x_0-2)$ \\
		$ \Rightarrow\dfrac{1}{\log_{4x_0-2}2}\geq\dfrac{1}{\log_{4x_0-2}\left(m^2+1\right)}\Rightarrow\log_{4x_0-2}\left(m^2+1\right)\geq\log_{4x_0-2}2 $ \\
		$ \Rightarrow m^2+1\geq 2\Rightarrow m^2\geq 1 $ mà $m\in[-1; 1]\Rightarrow m=\pm 1$.\\
		Điều kiện đủ: Với $m=\pm 1$ thì phương trình đã cho trở thành.\\
		$\log_2\left(x^2+y^2\right)=\log_2(2x+2y-2)\Leftrightarrow x^2+y^2=2x+2y-2\Leftrightarrow(x-1)^2+(y-1)^2=0\Leftrightarrow\heva{&x=1\\&y=1.}$ \\
		Suy ra phương trình đã cho có nghiệm duy nhất $(1; 1)$.\\
		Vậy có hai giá trị $m$ cần tìm là $m=\pm 1$.
	}
\end{ex}
\begin{ex}
	[Chuyên Lương Văn Tỵ - Ninh Bình - 2020]%Câu 7.
	Có bao nhiêu số nguyên $y$ để tồn tại số thực $x$ thỏa mãn $\log_{11}(3x+4y)=\log_4\left(x^2+y^2\right)$?
	\choice
	{$3$}
	{\True $2$}
	{$1$}
	{vô số}
	\loigiai{
		Đặt $\log_{11}(3x+4y)=\log_4\left(x^2+y^2\right)=t\Leftrightarrow\heva{&3x+4y={11}^t\\&x^2+y^2=4^t}$ (*).\\
		Hệ có nghiệm $\Leftrightarrow$ đường thẳng $\Delta\colon 3x+4y=11^t$ và đường tròn $(C)\colon x^2+y^2=4^t$ có điểm chung $\Leftrightarrow\mathrm{d}(O,\Delta)\leq R\Leftrightarrow\dfrac{11^t}{5}\leq 2^t\Leftrightarrow\left(\dfrac{11}{2}\right)^t\leq 5\Leftrightarrow t\leq\log_{\frac{11}{2}}5$.\\
		Do $x^2+y^2=4^t$ nên $|y|\leq 2^t\leq 2^{\log_{\frac{11}{2}}5}\approx 1\cdot 9239767$.\\
		Vì $y\in\mathbb{Z}$ nên $y\in\{-1;0;1\}$.\\
		Thử lại:\\
		Với $y=-1$, hệ (*) trở thành $\heva{&3x-4={11}^t\\&x^2+1=4^t}\Rightarrow\left(\dfrac{11^t+4}{3}\right)^2+1=4^t\Leftrightarrow 121^t+8\cdot 11^t+25=9\cdot 4^t$ (**).\\
		Nếu $t<0$ thì $4^t<1\Rightarrow 4^t<\left(\dfrac{11^t+4}{3}\right)^2+1$.\\
		Nếu $t\geq 0\Rightarrow\heva{&{121}^t\geq 4^t\\&{8\cdot 11}^t\geq{8\cdot 4}^t}\Rightarrow 121^t-4^t+8\left({11}^t-4^t\right)+25>0$.\\
		Vậy (**) vô nghiệm.\\
		Với $y=0$ thì hệ (*) trở thành $\heva{&3x={11}^t\\&x^2=4^t}\Rightarrow\dfrac{121^t}{9}=4^t\Leftrightarrow t=\log_{\frac{11}{2}}3\Rightarrow x=\dfrac{11^{\log_{\frac{11}{2}}3}}{3.}$ \\
		Với $y=1$ thì hệ (*) trở thành $\heva{&3x+4={11}^t\\&x^2+1=4^t}\Rightarrow\left(\dfrac{11^t-4}{3}\right)^2+1=4^t\Leftrightarrow 121^t-8\cdot 11^t+25=9\cdot 4^t$.\\
		Xét hàm số $f(t)=121^t-8\cdot 11^t+25-9\cdot 4^t$, liên tục trên $\left[\dfrac{1}{2};1\right]$ có $f\left(\dfrac{1}{2}\right)f(1)<0$ nên phương trình $f(t)=0$ luôn có nghiệm thuộc đoạn $\left[\dfrac{1}{2};1\right]$. Khi đó hiển nhiên sẽ tồn tại $x$ thỏa mãn.\\
		Vậy có 2 giá trị nguyên của $y$ thỏa mãn là $y=0, y=1$.
	}
\end{ex}
\begin{ex}
	[Chuyên Nguyễn Bỉnh Khiêm - Quảng Nam - 2020]%Câu 8.
	Có bao nhiêu cặp số thực $(x;y)$ thỏa mãn đồng thời các điều kiện $3^{\left|x^2-2x-3\right|-\log_35}=5^{-(y+4)}$ và $4|y|-|y-1|+(y+3)^2\leq 8$? 
	\choice
	{$1$}
	{$3$}
	{$4$}
	{\True $2$}
	\loigiai{
		Ta có $3^{\left|x^2-2x-3\right|-\log_35}=5^{-(y+4)}\Leftrightarrow 5^{-(y+3)}=3^{\left|x^2-2x-3\right|}\cdot (*)$.\\
		Vì $3^{\left|x^2-2x-3\right|}\geq 3^{\circ}\Rightarrow 5^{-(y+3)}\geq 1\Rightarrow y+3\leq 0\Rightarrow y\leq-3$.\\
		Với $y\leq-3$ ta có: $4|y|-|y-1|+(y+3)^2\leq 8\Leftrightarrow-4y+(y-1)+(y+3)^2\leq 8\Leftrightarrow y^2+3y\leq 0$ \\
		$ \Leftrightarrow-3\leq y\leq 0 $. Kết hợp với $y\leq-3$ suy ra $y=-3$.\\
		Thế $y=-3$ vào $(*)$ ta được: $3^{\left|x^2-2x-3\right|}=1\Leftrightarrow x^2-2x-3=0\Leftrightarrow\hoac{&x=-1\\&x=3.}$ \\
		Vậy các cặp số thực $(x;y)$ thỏa mãn là $(-1;-3);(3;-3)$.
	}
\end{ex}
\begin{ex}
	[Chuyên Bến Tre - 2020]%Câu 9.
	Giả sử $(x_0;y_0)$ là một nghiệm của phương trình $4^{x-1}+2^x\sin\left(2^{x-1}+y-1\right)+2=2^x+2\sin\left(2^{x-1}+y-1\right)$. Mệnh đề nào sau đây đúng?
	\choice
	{$x_0>7$}
	{\True $-2<x_0<4$}
	{$4<x_0<7$}
	{$-5<x_0 <-2$}
	\loigiai{
		Ta có $4^{x-1}+2^x\sin\left(2^{x-1}+y-1\right)+2=2^x+2\sin\left(2^{x-1}+y-1\right)$.\\
		$\begin{aligned}&\Leftrightarrow 4^x-4\cdot 2^x+4\cdot\left(2^x-2\right)\cdot\sin\left(2^{x-1}+y-1\right)+4+4=0\\&\Leftrightarrow\left(2^x-2\right)^2+4\left(2^x-2\right)\sin\left(2^{x-1}+y-1\right)+4\left[\sin^2\left(2^{x-1}+y-1\right)+\cos^2\left(2^{x-1}+y-1\right)\right]=0\\&\Leftrightarrow\left(2^x-2\right)^2+2\cdot\left(2^x-2\right)\cdot 2\sin\left(2^{x-1}+y-1\right)+\left[2\sin^2\left(2^{x-1}+y-1\right)\right]^2+4\cos^2\left(2^{x-1}+y-1\right)=0\\&\Leftrightarrow\left[\left(2^x-2\right)+2\sin\left(2^{x-1}+y-1\right)\right]^2+4\cos^2\left(2^{x-1}+y-1\right)=0\\&\Leftrightarrow\hoac{&2^x-2+2\sin\left(2^{x-1}+y-1\right)=0\\&\cos^2\left(2^{x-1}+y-1\right)=0}.\end{aligned}$ \\
		Vì $\cos^2\left(2^{x-1}+y-1\right)=0\Rightarrow\sin^2\left(2^{x-1}+y-1\right)=\pm 1$.\\
		$\sin^2\left(2^{x-1}+y-1\right)=1\Rightarrow 2^x=0$ (vô nghiệm).\\
		$\sin^2\left(2^{x-1}+y-1\right)=-1\Rightarrow 2^x=4\Rightarrow x=x_0=2\in(-2;4)$.
	}
\end{ex}
\begin{ex}
	[Chuyên Lào Cai - 2020]%Câu 10.
	Có bao nhiêu cặp số nguyên $(x;y)$ thỏa mãn $0\leq x\leq 4000$ và $5\left({25}^y+2y\right)=x+\log_5(x+1)^5-4$?
	\choice
	{\True $3$}
	{$2$}
	{$4$}
	{$5$}
	\loigiai{
		Đặt $\log_5(x+1)=t\Rightarrow x=5^t-1$.\\
		Phương trình trở thành:\\
		$5\left(5^{2y}+2y\right)=5^t-1+5t-4\Leftrightarrow 5^{2y}+2y=5^{t-1}+(t-1)$.\\
		Xét hàm số $f(u)=5^u+u\Rightarrow f'(u)=5^u\cdot\ln 5+1>0$ nên hàm số luôn đồng biến.\\
		Vậy để $f(2y)=f(t-1)\Leftrightarrow 2y=t-1\Leftrightarrow 2y+1=t=\log_5(x+1)$ \\
		$ \Rightarrow 0\leq 2y+1\leq\log_54001\Rightarrow 0\leq 2y+1\leq 5\Rightarrow y=\{0;1;2\} $.\\
		Với mỗi nghiệm $y$ ta tìm được một nghiệm $x$ tương ứng.
	}
\end{ex}
\begin{ex}
	[Chuyên Sơn La - 2020]%Câu 11.
	Cho $x$ là số thực dương và $y$ là số thực thỏa mãn $2^{x+\dfrac{1}{x}}=\log_2\left[14-(y-2)\sqrt{y+1}\right]$. Giá trị của biểu thức $P=x^2+y^2-xy+2020$ bằng
	\choice
	{$2022$}
	{$2020$}
	{\True $2021$}
	{$2019$}
	\loigiai{
		Áp dụng bất đẳng thức Cauchy ta có $x+\dfrac{1}{x}\geq 2\sqrt{x\cdot\dfrac{1}{x}}=2,\forall x>0\Rightarrow 2^{x+\dfrac{1}{x}}\geq 4$.\\
		Đặt $\sqrt{y+1}=t,t\geq 0$ thu được.\\
		$14-(y-2)\sqrt{y+1}=14-(t^2-3)t=-t^3+3t+14=16-(t-1)^2(t+2)\leq 16,\forall t\geq 0$.\\
		Dẫn đến $\log_2\left[14-(y-2)\sqrt{y+1}\right]\leq\log_216=4$.\\
		Như vậy hai vế bằng nhau khi dấu đẳng thức xảy ra tức là.\\
		$\heva{&t=1\\&x=\dfrac{1}{x}\\&x>0}\Rightarrow x=1;y=0\Rightarrow P=x^2+y^2-xy+2020=2021$.
	}
\end{ex}
\begin{ex}
	[Sở Hưng Yên - 2020]%Câu 12.
	Cho phương trình $log_3\left(3x^2-6x+6\right)=3^{y^2}+y^2-x^2+2x-1$. Hỏi có bao nhiêu cặp số $(x;y)$ và $0<x<2020; y\in\mathbb{N}$ thỏa mãn phương trình đã cho?
	\choice
	{$5$}
	{$6$}
	{$7$}
	{\True $4$}
	\loigiai{
		$log_3\left(3x^2-6x+6\right)=3^{y^2}+y^2-x^2+2x-1\Leftrightarrow log_33\left(x^2-2x+2\right)=3^{y^2}+y^2-x^2+2x-1$ \\
		$ \Leftrightarrow 1+log_3\left(x^2-2x+2\right)=3^{y^2}+y^2-x^2+2x-1 $ \\
		$ \Leftrightarrow log_3\left(x^2-2x+2\right)+\left(x^2-2x+2\right)=3^{y^2}+y^2 $ (1).\\
		Đặt $log_3\left(x^2-2x+2\right)=z\Rightarrow x^2-2x+2=3^z$ thì (1) trở thành:\\
		$\Leftrightarrow 3^z+z=3^{y^2}+y^2$ (2).\\
		Xét hàm số $f(t)=3^t+t\Rightarrow f'(t)=3^t\ln 3+1>0,\forall t\in\mathbb{R}$.\\
		Suy ra hàm số $f(t)$ đồng biến trên $\mathbb{R}$.\\
		(2) $\Leftrightarrow f(z)=f(y^2)\Leftrightarrow z=y^2$.\\
		Thay trở lại cách đặt ta có: $log_3\left(x^2-2x+2\right)=y^2\Leftrightarrow x^2-2x+2=3^{y^2}$.\\
		Xét hàm số: $g(x)=x^2-2x+2, x\in(0;2020)\Rightarrow g'(x)=2x-2$.\\
		$g'(x)=0\Leftrightarrow x=1$.\\
		Bảng biến thiên: 
		{\color{red} HINH O DAY}\\
		Suy ra:\\
		$\Leftrightarrow 1\leq g(x)<4076362\Leftrightarrow 1\leq 3^{y^2}<4076362\Leftrightarrow 0\leq y^2<\log_34076362$.\\
		Do $y\in\mathbb{N}\Rightarrow 0\leq y<\sqrt{\log_34076362}\approx 3,7\Rightarrow y\in\{0;1;2;3\}$ \\
		$\Rightarrow\hoac{&g(x)=1\\&g(x)=3\\&g(x)=3^4\\&g(x)=3^9.} $ \\
		Dựa vào bảng biến thiên của hàm số $g(x)$ ta thấy mỗi phương trình trên có một nghiệm $0<x<2020$.\\
		Vậy có 4 cặp số $(x;y)$ thỏa mãn đề bài.
	}
\end{ex}
\begin{ex}
	[Sở Phú Thọ - 2020]%Câu 13.
	Có bao nhiêu cặp số nguyên $(x;y)$ thỏa mãn $2\leq x\leq 2021$ và $2^y-\log_2\left(x+2^{y-1}\right)=2x-y$?
	\choice
	{$2020$}
	{$9$}
	{$2019$}
	{\True $10$}
	\loigiai{
		Đặt $\log_2\left(x+2^{y-1}\right)=t$. Suy ra $x+2^{y-1}=2^t$, $x=2^t-2^{y-1}$.\\
		Phương trình đã cho trở thành: $2^y-t=2\left(2^t-2^{y-1}\right)-y\Leftrightarrow 2\cdot 2^y+y=2\cdot 2^t+t$.\\
		Xét hàm số $g(x)=2\cdot 2^x+x$ có $g'(x)=2\cdot 2^x\ln 2+1>0,\forall x$ nên hàm số $y=g(x)$ luôn đồng biến.\\
		Khi đó $2.2^y+y=2\cdot 2^t+t\Leftrightarrow y=t$ hay $y=\log_2\left(x+2^{y-1}\right)$.\\
		Suy ra $x+2^{y-1}=2^y\Leftrightarrow x=2^y-2^{y-1}=2^{y-1}$.\\
		Mà $2\leq x\leq 2021$ nên $2\leq 2^{y-1}\leq 2021\Leftrightarrow 1\leq y-1\leq\log_22021$ hay $2\leq y\leq\left(\log_22021\right)+1$.\\
		Lại có $y$ là số nguyên nên $y\in\{2,3,\ldots,11\}$ tức 10 giá trị thỏa mãn.\\
		Xét biểu thức $x=2^{y-1}$, mỗi giá trị nguyên của $y$ cho tương ứng 1 giá trị nguyên của $x$ nên có 10 cặp số nguyên $(x,y)$ thỏa mãn yêu cầu đề bài.
	}
\end{ex}
\begin{ex}
	[Sở Bắc Ninh - 2020]%Câu 14.
	Có bao nhiêu cặp số nguyên dương $(x;y)$ thảo mãn $3^{x+y}-x^2\left(3^x-1\right)=(x+1)3^y-x^3$, với $x<2020$?
	\choice
	{$13$}
	{$15$}
	{$6$}
	{\True $7$}
	\loigiai{
		Ta có\\
		$3^{x+y}-x^2\left(3^x-1\right)=(x+1)3^y-x^3\Leftrightarrow 3^y\left(3^x-x-1\right)=x^2\left(3^x-x-1\right)\Leftrightarrow\left(3^x-x-1\right)\left(3^y-x^2\right)=0$.\\
		Ta thấy $3^x-x-1>0,\forall x>0\Rightarrow\left(3^x-x-1\right)\left(3^y-x^2\right)=0\Leftrightarrow 3^y=x^2\Leftrightarrow y=2\log_3x\Rightarrow x=3^k$.\\
		Vì $x<2020\Rightarrow 3^k<2020\Rightarrow 3^k\leq 3^6\Rightarrow k=\left\{0;1;2;3;4;5;6\right\}$.
	}
\end{ex}
\begin{ex}
	[Sở Bình Phước - 2020]%Câu 15.
	Biết $a,b$ là các số thực sao cho $x^3+y^3=a\cdot 10^{3z}+b\cdot 10^{2z},$ đồng thời $x,y,z$ là các số các số thực dương thỏa mãn $\log(x+y)=z$ và $\log\left(x^2+y^2\right)=z+1$. Giá trị của $\dfrac{1}{a^2}+\dfrac{1}{b^2}$ thuộc khoảng
	\choice
	{$(1;2)$}
	{$(2;3)$}
	{$(3;4)$}
	{\True $(4;5)$}
	\loigiai{
		Ta có: $\heva{&\log(x+y)=z\\&\log\left(x^2+y^2\right)=z+1}\Leftrightarrow\heva{&x+y={10}^z\\&x^2+y^2={10}^{z+1}={10\cdot 10}^z}\Rightarrow x^2+y^2=10(x+y)$.\\
		Khi đó $x^3+y^3=a\cdot 10^{3z}+b\cdot 10^{2z}\Leftrightarrow(x+y)\left(x^2-xy+y^2\right)=a\cdot\left({10}^z\right)^3+b\cdot\left({10}^z\right)^2$ \\
		$ \Leftrightarrow(x+y)\left(x^2-xy+y^2\right)=a\cdot (x+y)^3+b\cdot (x+y)^2\Leftrightarrow x^2-xy+y^2=a\cdot (x+y)^2+b\cdot (x+y) $ \\
		$ \Leftrightarrow x^2-xy+y^2=a\cdot\left(x^2+2xy+y^2\right)+\dfrac{b}{10}\left(x^2+y^2\right)\Leftrightarrow x^2+y^2-xy=\left(a+\dfrac{b}{10}\right)\left(x^2+y^2\right)+2a\cdot xy $.\\
		Đồng nhất hệ số ta được $\heva{&a+\dfrac{b}{10}=1\\&2a=-1}\Rightarrow\heva{&a=-\dfrac{1}{2}\\&b=15}\Rightarrow\dfrac{1}{a^2}+\dfrac{1}{b^2}=4+\dfrac{1}{225}=4,008\in(4;5)$.
	}
\end{ex}
\begin{ex}
	[Đặng Thúc Hứa - Nghệ An - 2020]%Câu 16.
	Có bao nhiêu cặp số nguyên $(x;y)$ thỏa mãn $0<y<2020$ và $3^x+3x-6=9y+\log_3y^3$. 
	\choice
	{$2020$}
	{$9$}
	{\True $7$}
	{$8$}
	\loigiai{
		Ta có: $3^x+3x-6=9y+\log_3y^3\Leftrightarrow 3^x+3(x-2)=9y+3\log_3y\Leftrightarrow 3^x+3(x-2)=3^{2+\log_3y}+3\log_3y(*)$.\\
		Xét hàm số: $f(t)=3^t+3(t-2)$.\\
		Ta có $f'(t)=3^t\cdot\ln 3+3>0,\forall t\in\mathbb{R}$. Suy ra hàm số $y=f(t)$ đồng biến trên $\mathbb{R}$.\\
		Khi đó: $(*)\Leftrightarrow f(x)=f(2+\log_3y)\Leftrightarrow x=2+\log_3y\Leftrightarrow y=3^{x-2}$.\\
		Do $0<y<2020$ và $x,y$ nguyên nên: $1\leq 3^{x-2}<2020\Leftrightarrow 2\leq x<2+\log_32020\Rightarrow x\in\left\{2;3;4;5;6;7;8\right\}$.\\
		Ứng với mỗi giá trị $x$ có một giá trị của $y$ nên có 7 cặp số $(x;y)$ nguyên thỏa mãn yêu cầu bài toán.
	}
\end{ex}
\begin{ex}
	[Đô Lương 4 - Nghệ An - 2020]%Câu 17.
	Giả sử $a,b$ là các số thực sao cho $x^3+y^3=a\cdot 10^{3z}+b\cdot 10^{2z}$ đúng với mọi các số thực dương.\\
	$x,y,z$ thỏa mãn $\log (x+y)=z$ và $\log (x^2+y^2)=z+1$. Giá trị của $a+b$ bằng
	\choice
	{$-\dfrac{25}{2}$}
	{$-\dfrac{31}{2}$}
	{$\dfrac{31}{2}$}
	{\True $\dfrac{29}{2}$}
	\loigiai{
		$\heva{&\log (x+y)=z\\&\log (x^2+y^2)=z+1}\Leftrightarrow\heva{&x+y={10}^z\\&x^2+y^2={10}^{z+1}}\Leftrightarrow\heva{&x+y={10}^z\\&(x+y)^2-2xy={10\cdot 10}^z}$ \\
		$\Leftrightarrow\heva{&x+y={10}^z\\&{10}^{2z}-2xy={10\cdot 10}^z}\Leftrightarrow\heva{&x+y={10}^z\\&xy=\dfrac{10^{2z}-10\cdot 10^z}{2}.} $ \\
		Khi đó $x^3+y^3=(x+y)^3-3xy(x+y)=10^{3z}-3\cdot\left(\dfrac{10^{2z}-10\cdot 10^z}{2}\right)\cdot 10^z$.\\
		$=\dfrac{1}{2}\left({2\cdot 10}^{3z}-3\cdot 10^{3z}+30\cdot 10^{2z}\right)=\dfrac{1}{2}\left(-10^{3z}+30\cdot 10^{2z}\right)=-\dfrac{1}{2}\cdot 10^{3z}+15\cdot 10^{2z}$.\\
		Lại có $x^3+y^3=a\cdot 10^{3z}+b\cdot 10^{2z}$.\\
		Suy ra $\heva{&a=-\dfrac{1}{2}\\&b=15}\Rightarrow a+b=\dfrac{29}{2}$.
	}
\end{ex}
\begin{ex}
	[Kim Liên - Hà Nội - 2020]%Câu 18.
	Có bao nhiêu số hữu tỉ $a$ thuộc đoạn $[-1;1]$ sao cho tồn tại số thực $b$ thỏa mãn.\\
	$\log_2\left(1-a^2-b^2+2b\right)=\dfrac{2^a}{4^a+1}+\dfrac{4^a}{2^a+1}+\dfrac{1}{2^a+4^a}-\dfrac{1}{2}$. 
	\choice
	{$0$}
	{$3$}
	{\True $1$}
	{Vô số}
	\loigiai{
		Ta có: $\dfrac{2^x}{4^x+1}+\dfrac{8^x+1}{2^x\left(1+2^x\right)}-\dfrac{1}{2}=\dfrac{2^x}{4^x+1}+\dfrac{4^x-2^x+1}{2^x}-\dfrac{1}{2}$.\\
		$=\dfrac{2^x}{4^x+1}+\dfrac{4^x+1}{2^x}-\dfrac{3}{2} =\dfrac{2^x}{4^x+1}+\dfrac{4^x+1}{4\cdot 2^x}+\dfrac{3}{4}\left(2^x+\dfrac{1}{2^x}\right)-\dfrac{3}{2}$.\\
		Áp dụng bất đẳng thức Cô si: $\dfrac{2^x}{4^x+1}+\dfrac{4^x+1}{4\cdot 2^x}\geq 1(1)$.\\
		Lại có $\dfrac{3}{4}\left(2^x+\dfrac{1}{2^x}\right)-\dfrac{3}{2}=\dfrac{3}{4}\left(2^x+\dfrac{1}{2^x}-2\right)=\dfrac{3}{4}\left(\sqrt{2^x}-\dfrac{1}{\sqrt{2^x}}\right)^2\geq 0(2)$.\\
		Từ $(1);(2)$ suy ra $\dfrac{2^x}{4^x+1}+\dfrac{4^x}{2^x+1}+\dfrac{1}{2^x+4^x}-\dfrac{1}{2}\geq 1$ \\
		$ \Rightarrow\log_2\left(1-a^2-b^2+2b\right)\geq 1\Leftrightarrow 1-a^2-b^2+2b\geq 2\Leftrightarrow a^2+b^2-2b+1\leq 0\Leftrightarrow a^2+(b-1)^2\leq 0\Leftrightarrow\heva{&a=0\\&b=1.} $ \\
		$a=0\in[-1;1]$ nên chọn phương án C.
	}
\end{ex}
\begin{ex}
	[Thanh Chương 1 - Nghệ An - 2020]%Câu 19.
	Cho các số thực $x$, $y$ thỏa mãn $x>1$, $y>1$ và $\log_3x\log_36y+2\log_3x\log_32y\left(3-\log_32xy\right)=\dfrac{9}{2}$. Giá trị của biểu thức $P=x+2y$ gần với số nào nhất trong các số sau
	\choice
	{$7$}
	{\True $8$}
	{$10$}
	{$9$}
	\loigiai{
		Đặt $a=\log_3x$, $b=\log_32y$. Do $x>1$, $y>1$ nên $a>0$, $b>\log_32$.\\
		Theo giả thiết ta có: $a(b+1)+2ab(3-a-b)=\dfrac{9}{2}\Leftrightarrow 2a^2b+a\left(2b^2-7b-1\right)+\dfrac{9}{2}=0(1)$.\\
		Coi $(1)$ là phương trình bậc hai ẩn $a$, $b$ là tham số. Để phương trình $(1)$ có nghiệm $a>0$ thì: $\heva{&\Delta\geq 0\\&-\dfrac{2b^2-7b-1}{2b}>0}\Leftrightarrow\heva{&\left(2b^2-7b-1\right)^2-36b\geq 0\\&2b^2-7b-1<0}\Leftrightarrow\heva{&4b^4-28b^3+45b^2-22b+1\geq 0\\&2b^2-7b-1<0}$ \\
		$ \Leftrightarrow\heva{&(b-1)^2\left(4b^2-20b+1\right)\geq 0\\&2b^2-7b-1<0}\Leftrightarrow\hoac{&b=1\\&\heva{&4b^2-20b+1\geq 0\\&2b^2-7b-1<0}.} $ \\
		Với $b=1\Rightarrow 2a^2-6a+\dfrac{9}{2}=0\Leftrightarrow a=\dfrac{3}{2}$. Khi đó $P=x+2y=3^{\frac{3}{2}}+3\approx 8,1$.\\
		Với $\heva{&4b^2-20b+1\geq 0\\&2b^2-7b-1<0}$: hệ vô nghiệm do $b>\log_32$.\\
		Vậy giá trị biểu thức $P=x+2y$ gần nhất với 8.
	}
\end{ex}
\begin{ex}
	[Tiên Lãng - Hải Phòng - 2020]%Câu 20.
	Có bao nhiêu cặp số nguyên dương $(x;y)$ với $x\leq 2020$ thỏa mãn.\\
	$2(3x-y)=3\left(1+9^y\right)-\log_3(2x-1)$ 
	\choice
	{$1010$}
	{$2020$}
	{\True $3$}
	{$4$}
	\loigiai{
		Đặt $\log_3(2x-1)=t\Rightarrow 2x=3^t+1$, ta được $3\left(3^t+1\right)-2y=3\left(1+3^{2y}\right)-t\Leftrightarrow 3\cdot 3^t+t=3\cdot 3^{2y}+2y$ (*).\\
		Xét hàm số $f(u)=3\cdot 3^u+u\Rightarrow f'(u)=3\cdot 3^u\ln 3+1>0,\forall u\in\mathbb{R}\Rightarrow f(u)$ đồng biến trên $\mathbb{R}$.\\
		Do đó (*) $\Leftrightarrow t=2y$, vậy nên $2x=3^{2y}+1\Leftrightarrow 9^y=2x-1$.\\
		Vì $x\leq 2020\Rightarrow 9^y\leq 4039\Leftrightarrow y\leq\log_94039$. Vì $y$ nguyên dương nên $y\in\{1;2;3\}$. Ta thấy với mỗi giá trị nguyên của $y$ thì tìm được 1 giá trị nguyên của $x$. Vậy có 3 cặp $(x;y)$ thỏa mãn.
	}
\end{ex}
\begin{ex}
	[THPT Lê Lợi - Thanh Hóa - 2021]%Câu 21.
	Cho phương trình $2\log_3\dfrac{6x}{\sqrt{2x+1}-1}+2\left(\sqrt{2x+1}-|y|\right)-3^{2|y|}+2x=0$. Với các cặp số $(x;y)$ thoả mãn phương trình trên, giá trị nhỏ nhất của $T=\dfrac{1}{3}\sqrt{2x+1}(2x+4)+2x+\dfrac{7}{3}-2\cdot 3^{2|y|}$ thuộc khoảng nào sau đây?
	\choice
	{$(-4;-2)$}
	{\True $(-11;-9,5)$}
	{$(-6;-4)$}
	{$(-9,5;-8)$}
	\loigiai{
		Điều kiện $x>0$.\\
		Ta có $2\log_3\dfrac{6x}{\sqrt{2x+1}-1}+2\left(\sqrt{2x+1}-|y|\right)-3^{2|y|}+2x=0$ \\
		$ \Leftrightarrow 2\log_3\left(3\left(\sqrt{2x+1}+1\right)\right)+2\sqrt{2x+1}-2|y|-3^{2|y|}+2x=0 $ \\
		$ \Leftrightarrow 2+2\log_3\left(\sqrt{2x+1}+1\right)+2\sqrt{2x+1}-2|y|-3^{2|y|}+2x=0 $ \\
		$ \Leftrightarrow 2\log_3\left(\sqrt{2x+1}+1\right)^2+2x+2+2\sqrt{2x+1}=2|y|+3^{2|y|} $ \\
		$ \Leftrightarrow 2\log_3\left(\sqrt{2x+1}+1\right)^2+\left(\sqrt{2x+1}+1\right)^2=\log_33^{2|y|}+3^{2|y|} $ (1).\\
		Xét hàm số $y=f(t)=\log_3t+t\Rightarrow f'(t)=\dfrac{1}{t\ln 3}+1>0,\forall t>0$.\\
		Suy ra hàm số $y=f(t)$ đồng biến trên khoảng $(0;+\infty)$.\\
		Từ $(1)\Leftrightarrow f\left(\left(\sqrt{2x+1}+1\right)^2\right)=f\left(3^{2|y|}\right)\Leftrightarrow\left(\sqrt{2x+1}+1\right)^2=3^{2|y|}$.\\
		Do đó $T=\dfrac{1}{3}\sqrt{2x+1}(2x+4)+2x+\dfrac{7}{3}-2\cdot 3^{2|y|} =\dfrac{1}{3}\sqrt{2x+1}(2x+4)+2x+\dfrac{7}{3}-2\cdot\left(\sqrt{2x+1}+1\right)^2$.\\
		Đặt $t=\sqrt{2x+1}>1\Rightarrow 2x=t^2-1$.\\
		Suy ra $T=\dfrac{1}{3}t\left(t^2-1+4\right)+t^2-1+\dfrac{7}{3}-2(t+1)^2=\dfrac{1}{3}t^3-t^2-3t-\dfrac{2}{3}$.\\
		Có $T'=t^2-2t-3\Rightarrow T'=0\Leftrightarrow\hoac{&t=-1\\&t=3.}$ \\
		Bảng biến thiên\\
		{\color{red} HINH O DAY}\\
		Từ bảng biến thiên suy ra $\min T=-\dfrac{29}{3}\approx-9,67\in(-11;-9,5)$.
	}
\end{ex}
\begin{ex}
	[THPT Quốc Oai - Hà Nội - 2021]%Câu 22.
	Có bao nhiêu số nguyên $x$ sao cho tồn tại số thực $y$ thỏa mãn $\log_3(x+y)=\log_4\left(x^2+2y^2\right)$ 
	\choice
	{$3$}
	{\True $2$}
	{Vô số}
	{$1$}
	\loigiai{
		Điều kiện $\heva{&x+y>0\\&x^2+2y^2>0.}$ \\
		Ta có $\log_3(x+y)=\log_4\left(x^2+2y^2\right)=t\Rightarrow\heva{&x+y=3^t\\&x^2+2y^2=4^t.}$ \\
		Ta có $(x+y)^2\leq\left(1+\dfrac{1}{2}\right)\left(x^2+2y^2\right)\Leftrightarrow(3^t)^2\leq\dfrac{3}{2}\cdot 4^t\Rightarrow t\leq\log_{\frac{9}{4}}\dfrac{3}{2}=\dfrac{1}{2}$.\\
		Mặt khác $x^2\leq 4^t\leq 4^{\frac{1}{2}}=2\Rightarrow-\sqrt{2}\leq x\leq\sqrt{2}$.\\
		Lại có $x\in\mathbb{Z}\Rightarrow x\in\{-1;0;1\}$.\\
		Với $x=-1\Rightarrow\heva{&y=1+3^t\\&y^2=4^t-1}\Rightarrow 4^t-1=\left(1+3^t\right)^2\Leftrightarrow 9^t+2\cdot 3^t+2-4^t=0$.\\
		Xét hàm số $f(t)=9^t+2\cdot 3^t+2-4^t$. Khi đó $\hoac{&t<0\Rightarrow 2-4^t>0\\&t\geq 0\Rightarrow 9^t-4^t>0}\Rightarrow f(t)>0,\forall t\Rightarrow f(t)=0$ vô nghiệm.\\
		Với $x=0\Rightarrow\heva{&y=3^t\\&y^2=4^t}\Rightarrow 9^t=4^t\Leftrightarrow t=0\Rightarrow y=1$ (thỏa mãn).\\
		Với $x=1\Rightarrow\heva{&y=3^t-1\\&y^2=4^t-1}\Rightarrow 4^t-1=\left(3^t-1\right)^2\Leftrightarrow 9^t-2\cdot 3^t+2-4^t=0$.\\
		Phương trình có nghiệm $t=0\Rightarrow y=0$ (thỏa mãn).\\
		Vậy có $2$ giá trị nguyên của $x$ thỏa đề là $x=0$ và $x=1$.
	}
\end{ex}
\begin{ex}
	[Trung Tâm Thanh Tường - 2021]%Câu 23.
	Có bao nhiêu cặp số $x,y$ là các số nguyên không âm thỏa mãn:\\
	$2\left(1+\sqrt{x+2y}\right)^2+\log_2(x+2y)=2\log_2\left(x^2+y^2+2xy+x\right)+2(x+y)^2+4x+4y$ 
	\choice
	{$2$}
	{\True $3$}
	{$4$}
	{$5$}
	\loigiai{
		Ta có $2\left(1+\sqrt{x+2y}\right)^2+\log_2(x+2y)=2\log_2\left(x^2+y^2+2xy+x\right)+2(x+y)^2+4x+4y$ \\
		$ \Leftrightarrow 1+2\sqrt{x+2y}+x+2y+\dfrac{1}{2}\log_2(x+2y)=\log_2\left[(x+y)^2+x\right]+(x+y)^2+2x+2y $ \\
		$ \Leftrightarrow 1+\log_2\sqrt{x+2y}+2\sqrt{x+2y}=\log_2\left[(x+y)^2+x\right]+(x+y)^2+x $ \\
		$\Leftrightarrow\log_22\sqrt{x+2y}+2\sqrt{x+2y}=\log_2\left[(x+y)^2+x\right]+\left[(x+y)^2+x\right] $ (1).\\
		Xét hàm số $g(t)=\log_2t+t$ với $t>0$, có $g'(t)=\dfrac{1}{t\cdot\ln 2}+1>0,\forall t>0$.\\
		Suy ra hàm số $y=g(t)=\log_2t+t$ đồng biến trên khoảng $(0;+\infty)$.\\
		Từ $(1)\Leftrightarrow g\left(2\sqrt{x+2y}\right)=g\left((x+y)^2+x\right)\Leftrightarrow 2\sqrt{x+2y}=(x+y)^2+x$ \\
		$ \Leftrightarrow x+2y+2\sqrt{x+2y}=(x+y)^2+2(x+y) $ (2).\\
		Xét hàm số $h(u)=u^2+u$ với $u>0$, có $h'(u)=2u+1>0,\forall u>0$.\\
		Suy ra hàm số $y=h(u)=u^2+u$ đồng biến trên khoảng $(0;+\infty)$.\\
		Từ $(2)\Leftrightarrow h\left(\sqrt{x+2y}\right)=h(x+y)\Leftrightarrow\sqrt{x+2y}=x+y\Leftrightarrow x+2y=(x+y)^2$ \\
		$ \Leftrightarrow x^2+2xy+y^2=x+2y\Leftrightarrow y^2+2(x-1)y+x^2-x=0 $ (3).\\
		Ta coi phương trình (3) là phương trình bậc hai ẩn $y$, khi đó: Phương trình (3) có nghiệm $\Leftrightarrow\Delta_y'\geq 0\Leftrightarrow(x-1)^2-\left(x^2-x\right)\geq 0\Leftrightarrow-x+1\geq 0\Leftrightarrow x\leq 1$.\\
		Do $x$ là số nguyên không âm nên $x=0$ hoặc $x=1$.\\
		Với $x=0$, suy ra $(3)\Leftrightarrow y^2-2y=0\Leftrightarrow\hoac{&y=0\\&y=2}$ (thoả mãn).\\
		Với $x=1$, suy ra $(3)\Leftrightarrow y^2=0\Leftrightarrow y=0$ (thoả mãn).\\
		Vậy có 3 cặp $(x;y)$ thoả mãn là $(0;0),(0;2),(1;0)$.
	}
\end{ex}
\begin{ex}
	Có bao nhiêu cặp số nguyên dương $(x;y)$ thỏa mãn đồng thời hai điều kiện: $1\leq x\leq 10^6$ và $\log\left(10x^2-20x+20\right)=10^{y^2}+y^2-x^2+2x-1$?
	\choice
	{$4$}
	{$2$}
	{$3$}
	{\True $1$}
	\loigiai{
		Điều kiện: $10x^2-20x+20>0$, đúng $\forall x\in\mathbb{R}$.\\
		Ta có $\log\left(10x^2-20x+20\right)=10^{y^2}+y^2-x^2+2x-1\Leftrightarrow\left(x^2-2x+1\right)+\log\left[10\left(x^2-2x+2\right)\right]=10^{y^2}+y^2\Leftrightarrow\left(x^2-2x+1\right)+\log 10+\log\left(x^2-2x+2\right)=10^{y^2}+y^2$ \\
		$ \Leftrightarrow\left(x^2-2x+2\right)+\log\left(x^2-2x+2\right)=10^{y^2}+y^2 $ \\
		$ \Leftrightarrow 10^{\log\left(x^2-2x+2\right)}+\log\left(x^2-2x+2\right)=10^{y^2}+y^2 $ (*).\\
		Xét hàm $f(t)=10^t+t$ trên $\mathbb{R}$.\\
		Ta có $f'(t)=10^t\cdot\ln 10+1>0$, $\forall t\in\mathbb{R}$. Do đó $f(t)$ đồng biến trên $\mathbb{R}$.\\
		Khi đó (*) $\Leftrightarrow f\left[\log\left(x^2-2x+2\right)\right]=f(y^2)\Leftrightarrow\log\left(x^2-2x+2\right)=y^2\Leftrightarrow x^2-2x+2=10^{y^2}\Leftrightarrow(x-1)^2+1=10^{y^2}$.\\
		Vì $1\leq x\leq 10^6$ nên $1\leq(x-1)^2+1=10^{y^2}\leq\left({10}^6-1\right)^2+1\Rightarrow 0\leq y^2\leq\log\left[\left({10}^6-1\right)^2+1\right]$.\\
		Vì $y\in\mathbb{Z}^+$ nên $y\in\{1;2;3\}$.\\
		Với $y=1\Rightarrow x^2-2x+2=10\Leftrightarrow x^2-2x-8=0\Leftrightarrow\hoac{&x=-2 (ktm)\\&x=4 (tm).}$ \\
		Với $y=2\Rightarrow x^2-2x+2=10^4\Leftrightarrow x^2-2x-9998=0$ (không có giá trị $x$ nguyên nào thỏa mãn).\\
		Với $y=3\Rightarrow x^2-2x+2=10^9\Rightarrow x^2-2x-999999998=0$ (không có giá trị $x$ nguyên nào thỏa mãn).\\
		Vậy có một cặp nguyên dương $(x; y)=(4; 1)$ thỏa mãn yêu cầu bài toán.
	}
\end{ex}
\begin{ex}
	Có bao nhiêu số nguyên $y<10$ sao cho tồn tại số nguyên $x$ thỏa mãn $5^{\sqrt{2}^y+x-2}+\sqrt2^y=5^{x^2-x-1}+(x-1)^2$?
	\choice
	{$10$}
	{$1$}
	{\True $5$}
	{Vô số}
	\loigiai{
		Ta có $$5^{\sqrt{2}^y+x-2}+\sqrt2^y=5^{x^2-x-1}+(x-1)^2\Leftrightarrow 5^{\sqrt{2}^y+x-2}+\sqrt2^y+x-1=5^{x^2-x-1}+x^2-x.$$
		Dễ thấy $f(t)=5^{t-1}+t$ đồng biến trên $\mathbb{R}$. Do đó từ phương trình trên suy ra
		$$\sqrt2^y+x-1=x^2-x\Leftrightarrow(x-1)^2=\sqrt2^y=2^{\frac{y}{2}}\Leftrightarrow x=1\pm 2^{\frac{y}{2}}.$$
		Do $x$ nguyên nên ta có $2^{\frac{y}{2}}\in\mathbb{Z}$ và $y<10$ nên $y\in\{0;2;4;6;8\}$.
	}
\end{ex}
\begin{ex}
	Có bao nhiêu cặp số nguyên dương $(x; y)$ thoả mãn $1\leq x\leq 2020$ và $2^y+y=2x+\log_2\left(x+2^{y-1}\right)$ 
	\choice
	{$2021$}
	{$10$}
	{$2020$}
	{\True $11$}
	\loigiai{
		Theo đề bài, $2^y+y=2x+\log_2\left(x+2^{y-1}\right)$ \\
		$ \Leftrightarrow 2^y+\log_2(2^y)=2x+\log_2\left(x+\dfrac{2^y}{2}\right) $ \\
		$ \Leftrightarrow 2^y+2^y+\log_2(2^y)=2x+2^y+\log_2\left(\dfrac{2x+2^y}{2}\right) $ \\
		$ \Leftrightarrow 2\cdot (2^y)+\log_2(2^y)=2\cdot\left(\dfrac{2x+2^y}{2}\right)+\log_2\left(\dfrac{2x+2^y}{2}\right) (1) $.\\
		Xét hàm số $f(t)=2t+\log_2t$, $t>0$.\\
		Vì $f'(t)=2+\dfrac{1}{t\ln 2}>0\forall t>0\Rightarrow f(t)$ đồng biến trên $(0;+\infty)$.\\
		nên $(1)\Leftrightarrow f(2^y)=f\left(\dfrac{2x+2^y}{2}\right)\Leftrightarrow 2^y=\dfrac{2x+2^y}{2}\Leftrightarrow 2\cdot 2^y=2x+2^y\Leftrightarrow 2x=2^y\Leftrightarrow x=2^{y-1}$.\\
		Do $1\leq x\leq 2020$ nên $0\leq y-1\leq\log_22020\Leftrightarrow 1\leq y\leq 11,98$.\\
		Do $y\in\mathbb{N}^*$ nên $y\in\left\{1; 2; 3;\ldots;11\right\}$, với mỗi giá trị $y$ cho ta 1 giá trị $x$ thoả đề.\\
		Vậy có $11$ cặp số nguyên $(x; y)$ thoả mãn đề bài.
	}
\end{ex}
\begin{ex}
	Có bao nhiêu số nguyên $x$ sao cho tồn tại số thực $y$ thỏa mãn $2\log_2(x+y)-\log_2(1+\sqrt{3})=\log_{\sqrt{3}}\left(x^2+y^2-1\right)$ 
	\choice
	{$1$}
	{$3$}
	{\True $2$}
	{$5$}
	\loigiai{
		Đặt: $t=2\log_2(x+y)-\log_2(1+\sqrt{3})=\log_{\sqrt{3}}\left(x^2+y^2-1\right)$.\\
		Suy ra: $\heva{&(x+y)^2=2^{t+\log_2(1+\sqrt{3})}\\&x^2+y^2-1={\sqrt{3}}^t}\Leftrightarrow\heva{&(x+y)^2=(1+\sqrt{3}){\cdot 2}^t\\&x^2+y^2=1+\sqrt{3}^t.}$ \\
		Ta có\\
		$\begin{aligned}&(x+y)^2\leq 2\left(x^2+y^2\right)\\&\Leftrightarrow(1+\sqrt{3}){\cdot 2}^t\leq 2\left(1+\sqrt{3}^t\right)\\&\Leftrightarrow\dfrac{(1+\sqrt{3})2^t}{2}\leq 1+\sqrt{3}^t\\&\Leftrightarrow\left(\dfrac{1}{2}\right)^t+\left(\dfrac{\sqrt{3}}{2}\right)^t\geq\dfrac{1+\sqrt{3}}{2}.\end{aligned}$\\
		Xét $f(t)=\left(\dfrac{1}{2}\right)^t+\left(\dfrac{\sqrt{3}}{2}\right)^t$ nghịch biến trên $\mathbb{R}$ nên.\\
		$\left(\dfrac{1}{2}\right)^t+\left(\dfrac{\sqrt{3}}{2}\right)^t\geq\dfrac{1+\sqrt{3}}{2}\Leftrightarrow f(t)\geq f(1)\Leftrightarrow t\leq 1$.\\
		Do đó $\heva{&0<x+y=2^{\frac{t+\log_2(1+\sqrt{3})}{2}}\leq\sqrt{2\log_2(1+\sqrt{3})}\\&x^2+y^2=1+\sqrt{3}^t\leq 1+\sqrt{3}}\Rightarrow x\in\{0;\pm 1\}$ (vì $x\in\mathbb{Z}$).\\
		Thử lại:\\
		Với $x=1$:\\
		$\begin{aligned}&\heva{&y=\sqrt(1+\sqrt{3})2^t-1\\&y^2=(\sqrt{3})^t}\\&\Rightarrow \left(\sqrt{(1+\sqrt{3})2^t}-1\right)^2-\sqrt{3}^t=0\\&\Rightarrow(1+\sqrt{3})2^t-2\sqrt{(1+\sqrt{3})2^t}-\sqrt{3}^t+1=0.\end{aligned}$ \\
		Ta có: $g(x)=(1+\sqrt{3})2^t-2\sqrt{(1+\sqrt{3}){\cdot 2}^t}-\sqrt3^t+1$ liên tục trên $[0;1]$ thỏa mãn $g(0)g(1)<0$ nên phương trình có nghiệm $t\in(0;1)$.\\
		Do đó với $x=1$ thì tồn tại số thực $y$ thỏa mãn $2\log_2(x+y)-\log_2(1+\sqrt{3})=\log_{\sqrt{3}}\left(x^2+y^2-1\right)$.\\
		Với $x=-1$:\\
		$\begin{aligned}&\heva{&y=\sqrt(1+\sqrt{3})2^t}+1\\&y^2={\sqrt{3}}^t\\&\Rightarrow \left(\sqrt{(1+\sqrt{3})2^t}+1\right)^2-\sqrt{3}^t=0\\&\Rightarrow(1+\sqrt{3})2^t+2\sqrt{(1+\sqrt{3}){\cdot 2}^t}-\sqrt{3}^t+1=0.\end{aligned}$ \\
		Ta có: $(1+\sqrt{3})2^t+2\sqrt{(1+\sqrt{3})\cdot 2^t}-(\sqrt{3})^t+1>0,\forall t\leq 1$ nên phương trình vô nghiệm.\\
		Do đó với $x=-1$ thì không tồn tại số thực $y$ thỏa mãn $2\log_2(x+y)-\log_2(1+\sqrt{3})=\log_{\sqrt{3}}\left(x^2+y^2-1\right)$.\\
		Với $x=0$:\\
		$\begin{aligned}&\heva{&y^2=(1+\sqrt{3})2^t\\&y^2=1+\sqrt{3}^t}\\&\Rightarrow(1+\sqrt{3})2^t={\sqrt{3}}^t+1\\&\Rightarrow(1+\sqrt{3})2^t-\sqrt{3}^t-1=0.\end{aligned}$ \\
		Ta có: $h(x)=(1+\sqrt{3})2^t-\sqrt3^t-1$ liên tục trên $[-1;0]$ thỏa mãn $h(-1)h(0)<0$ nên phương trình có nghiệm $t\in(-1;0)$.\\
		Do đó với $x=0$ thì tồn tại số thực $y$ thỏa mãn $2\log_2(x+y)-\log_2(1+\sqrt{3})=\log_{\sqrt{3}}\left(x^2+y^2-1\right)$.\\
		Vậy $x\in\{0;1\}$.
	}
\end{ex}
\begin{ex}
	Có bao nhiêu cặp số nguyên $(x;y)$ thỏa mãn $0\leq y\leq 2020$ và $\log_3\left(\dfrac{2^x-1}{y}\right)=y+1-2^x$? 
	\choice
	{$2019$}
	{\True $11$}
	{$2020$}
	{$4$}
	\loigiai{
		Từ giả thiết ta có: $\heva{&y\neq 0\\&\dfrac{2^x-1}{y}>0\\&y\geq 0}\Leftrightarrow 2^x>1\Leftrightarrow x>0$.\\
		Ta có phương trìn $\Leftrightarrow\log_3\left(2^x-1\right)+2^x-1=\log_3y+y (*)$.\\
		Xét hàm số $f(t)=\log_3t+t$ trên $(0;+\infty)$.\\
		Khi đó $f'(t)=\dfrac{1}{t\ln 3}+1>0$ do đó hàm số $f(t)=\log_3t+t$ đồng biến trên $(0;+\infty)$.\\
		(*) có dạng $f\left(2^x-1\right)=f(y)\Leftrightarrow y=2^x-1$.\\
		Vì $0\leq y\leq 2020\Leftrightarrow 0\leq 2^x-1\leq 2020\Leftrightarrow 1\leq 2^x\leq 2021\Leftrightarrow 0\leq x\leq\log_2(2021)$.\\
		$\heva{&0\leq x\leq\log_2(2021)\\&x\in\mathbb{Z}}\Leftrightarrow x\in\left\{0;1;2;3;4;5;6;7;8;9;10\right\}$. Vậy có $11$ cặp $(x;y)$ thỏa mãn.
	}
\end{ex}
\begin{ex}
	[Chuyên Hạ Long - Quảng Ninh - 2021]%Câu 29.
	Có bao nhiêu cặp số nguyên $(x;y)$ thỏa mãn.\\
	$(4xy+7y)(2x-1)\left(\mathrm{e}^{2xy}-\mathrm{e}^{4x+y+7}\right)=\left[2x(2-y)+y+7\right]\mathrm{e}^y$ 
	\choice
	{$8$}
	{$5$}
	{\True $6$}
	{$7$}
	\loigiai{
		Ta có $(4xy+7y)(2x-1)\left(\mathrm{e}^{2xy}-\mathrm{e}^{4x+y+7}\right)=\left[2x(2-y)+y+7\right]\mathrm{e}^y$ \\
		$\Leftrightarrow(4x+7)(2xy-y)\left(\mathrm{e}^{2xy-y}-\mathrm{e}^{4x+7}\right)=2x(2-y)+y+7 $ (vì $\mathrm{e}^y>0\forall y)$ \\
		$\Leftrightarrow(4x+7)(2xy-y)\cdot\left(\mathrm{e}^{2xy-y}-\mathrm{e}^{4x+7}\right)=(4x+7)-(2xy-y) $ \\
		$\Leftrightarrow(4x+7)\cdot\left[(2xy-y)\mathrm{e}^{2xy-y}-1\right]=(2xy-y)\left[(4x+7)\mathrm{e}^{4x+7}-1\right] $ \\
		$\Rightarrow\dfrac{(2xy-y)\mathrm{e}^{2xy-y}-1}{2xy-y}=\dfrac{(4x+7)\mathrm{e}^{4x+7}-1}{4x+7} $ với $x\neq-\dfrac{7}{4};x\neq\dfrac{1}{2};y\neq 0$ \\
		$\Leftrightarrow\mathrm{e}^{2xy-y}-\dfrac{1}{2xy-y}=\mathrm{e}^{4x+7}-\dfrac{1}{4x+7} $.\\
		Xét $f(t)=\mathrm{e}^t-\dfrac{1}{t}\quad(t\neq 0)$ \\
		$ \Rightarrow f'(t)=\mathrm{e}^t+\dfrac{1}{t^2}>0\forall t\neq 0 $ \\
		$ \Rightarrow f(t) $ đồng biến trên các khoảng xác định\\
		$ \Rightarrow f(2xy-y)=f(4x+7) $.\\
		Trường hợp 1: $(2xy-y)\cdot (4x+7)<0$.\\
		Giả sử $2xy-y<0$ và $4x+7>0$. Do $x,y\in\mathbb{Z}$ nên $2xy-y\leq-1$ và $4x+7\geq 1$ \\
		$ \Rightarrow\heva{&f(2xy-y)\leq f(-1)=1-\dfrac{1}{e}<1\\&f(4x+7)\geq f(1)=e-1>1} $. Do đó, $f(2xy-y)\neq f(4x+7)$.\\
		Trường hợp 2: $(2xy-y)\cdot (4x+7)>0$.\\
		$f(2xy-y)=f(4x+7)$ \\
		$ \Leftrightarrow 2xy-y=4x+7 $ \\
		$ \Rightarrow y=\dfrac{4x+7}{2x-1}=2+\dfrac{9}{2x-1} $.\\
		Theo bài, $y\in\mathbb{Z}$ nên $2x-1\in\left\{\pm 1;\pm 3;\pm 9\right\}\Rightarrow x\in\left\{-4;-1;0;1;2;5\right\}$.
	}
\end{ex}
\begin{ex}
	[Đề Minh Họa 2021]%Câu 30.
	Có bao nhiêu số nguyên $a (a\geq 2)$ sao cho tồn tại số thực $x$ thỏa mãn: $\left(a^{\log x}+2\right)^{\log a}=x-2$?
	\choice
	{\True $8$}
	{$9$}
	{$1$}
	{Vô số}
	\loigiai{
		Xét phương trình $\left(a^{\log x}+2\right)^{\log a}=x-2\Leftrightarrow\left(x^{\log a}+2\right)^{\log a}=x-2$.\\
		Vì $\left(x^{\log a}+2\right)^{\log a}>0$ nên suy ra $x>2$.\\
		Ta có: $\left(x^{\log a}+2\right)^{\log a}=x-2\Leftrightarrow\left(x^{\log a}+2\right)^{\log a}+\left(x^{\log a}+2\right)=x^{\log a}+x$.\\
		Xét hàm số $f(t)=t^{\log a}+t$ có $f'(t)=\log a\cdot t^{\log a-1}+1>0$, $\forall t>2$. Do đó $f(t)$ là hàm số đồng biến trên $(2;+\infty)$.\\
		Mà $f\left(x^{\log a}+2\right)=f(x)\Rightarrow x^{\log a}+2=x\Rightarrow x^{\log a}=x-2$.\\
		Trường hợp 1: $\log a\geq 1\Leftrightarrow a\geq 10$. 
		{\color{red} HINH O DAY}\\
		Dễ thấy hai đồ thị của hai hàm số $y=x^{\log a}$ và $y=x-2$ không có điểm chung, vậy $a\geq 10$ không thỏa mãn yêu cầu bài toán.\\
		Trường hợp 1: $\log a<1\Leftrightarrow a<10$.\\
		Dễ thấy phương trình $x^{\log a}=x-2$ luôn có nghiệm duy nhất. \\
		{\color{red} HINH O DAY}\\
		Vậy $a\in\left\{2; 3;\ldots; 9\right\}\Rightarrow$ có 8 giá trị của $a$ thỏa mãn.
	}
\end{ex}
\begin{ex}
	[Mã 101 - 2021 Lần 1]%Câu 31.
	Có bao nhiêu số nguyên $y$ sao cho tồn tại $x\in\left(\dfrac{1}{3};3\right)$ thỏa mãn $27^{3x^2+xy}=(1+xy)27^{9x}$?
	\choice
	{$27$}
	{$9$}
	{\True $11$}
	{$12$}
	\loigiai{
		Xét $f(x)=27^{3x^2-9x+xy}-(xy+1)$ và áp dụng $a^x\geq x(a-1)+1$.\\
		Suy ra: $f(x)\geq 26\left(3x^2-9x+xy\right)-xy-1=78x^2+25xy-234x-1>0,\forall y\geq 10$.\\
		Do đó $y\leq 9$.\\
		$y=0\Rightarrow 27^{3x^2-9x}=1\Rightarrow 3x^2-9x=0$: loại.\\
		$y\leq-3\Rightarrow xy <-1\Rightarrow VP<0$: loại.\\
		$y=-1,y=-2$: thỏa mãn.\\
		Xét $y>0$ có $f(3)=27^{3y}-(3y+1)\geq 0,\forall y>0$.\\
		Và $f\left(\dfrac{1}{3}\right)=3^{y-8}-\dfrac{y}{3}-1<0,\forall y\in\left\{1;2;3;\ldots;9\right\}$ \\
		$ \Rightarrow y\in\left\{-2;-1;1;2;3;4;5;6;7;8;9\right\} $.
	}
\end{ex}
\begin{ex}
	[Mã 102 - 2021 Lần 1]%Câu 32.
	Có bao nhiêu số nguyên $y$ sao cho tồn tại $x\in\left(\dfrac{1}{3}; 4\right)$ thỏa mãn $27^{3x^2+xy}=(1+xy)\cdot 27^{12x}$?
	\choice
	{$27$}
	{$15$}
	{$12$}
	{\True $14$}
	\loigiai{
		Xét $f(x)=27^{3x^2+xy-12x}-(1+xy)$.\\
		Áp dụng bất đẳng thức: $a^x\geq x(a-1)+1$, ta có $f(x)\geq 26\left(3x^2+xy-12x\right)+1-(1+xy)=78x^2+(25y-312)x>0,\forall y\geq 13$.\\
		Do đó $y\leq 12$.\\
		$y=0\Rightarrow 27^{3x^2-12x}=1\Leftrightarrow 3x^2-12x=0\Leftrightarrow\hoac{&x=0\\&x=4}$ (loại).\\
		$y\leq-3\Rightarrow xy <-1\Rightarrow VP<0$ (loại).\\
		$y=-1, y=-2$: thỏa mãn.\\
		Xét $y>0$ có $f(4)=27^{4y}-(1+4y)\geq 0,\forall y>0$ và $f\left(\dfrac{1}{3}\right)=f(x)=3^{y-11}-\dfrac{y}{3}-1<0,\forall y\in\left\{1; 2;\ldots; 12\right\}$.\\
		Do đó phương trình $f(x)=0$ có nghiệm $x\in\left(\dfrac{1}{3}; 4\right)$, $\forall y\in\left\{1; 2;\ldots; 12\right\}$,\\
		Vậy $y\in\left\{-2;-1; 0; 1; 2;\ldots;12\right\}$.
	}
\end{ex}
\begin{ex}
	[Mã 104 - 2021 Lần 1]%Câu 33.
	Có bao nhiêu số nguyên $y$ sao cho tồn tại $x\in\left(\dfrac{1}{3}; 6\right)$ thỏa mãn $27^{3x^2+xy}=(1+xy)\cdot 27^{18x}$?
	\choice
	{$19$}
	{\True $20$}
	{$18$}
	{$21$}
	\loigiai{
		Ta có $27^{3x^2+xy}=(1+xy)\cdot 27^{18x}\Leftrightarrow 27^{3x^2+xy-18x}=1+xy\Leftrightarrow 27^{3x^2+xy-18x}-1-xy=0$.\\
		Mạt khác: $27^{3x^2+xy-18x}-1-xy=(1+26)^{3x^2+xy-18x}-1-xy\geq 1+26\left(3x^2+xy-18x\right)-1-xy$.\\
		$=84x^2+25xy-468x$.\\
		Với $y\geq 19\Rightarrow 27^{3x^2+xy-18x}-1-xy\geq 78x^2+25xy-468x\geq 84x^2+7x>0\forall x\in\left(\dfrac{1}{3}; 6\right)$ (loại).\\
		Với $y\leq-3\Rightarrow 27^{3x^2+xy-18x}=1+xy\leq 1-3x<0\forall x\in\left(\dfrac{1}{3}; 6\right)$ (loại).\\
		Với $y=-2\Rightarrow x=\dfrac{1}{2}$ (thỏa mãn).\\
		Với $y=-1\Rightarrow x=1$ (thỏa mãn).\\
		Với $y=0\Rightarrow 27^{3x^2-18x}=1\Leftrightarrow 3x^2-18x=0\Leftrightarrow\hoac{&x=0\\&x=6}$ (loại).\\
		Với $y\geq 1$: Đặt $f(x)=27^{3x^2+xy-18x}-1-xy$.\\
		Do $f\left(\dfrac{1}{3}\right)=3^{y-17}-1-\dfrac{y}{3}<0\forall y\in\left\{1; 2;\ldots\cdot;18\right\}$ và $f(6)=27^{6y}-6y-1>0\forall y\in\left\{1; 2;\ldots\cdot;18\right\}$ nên phương trình $f(x)=0$ luôn có nghiệm thuộc $\left(\dfrac{1}{3}; 6\right)$.\\
		Tóm lại $y\in\left\{-2;-1; 1; 2;\ldots\cdot;18\right\}$.\\
		Vậy có $20$ số nguyên $y$ thỏa mãn yêu cầu bài toán.
	}
\end{ex}
\begin{ex}
	[Mã 103 - 2021 - Lần 1]%Câu 34.
	Có bao nhiêu số nguyên $y$ sao cho tồn tại $x\in\left(\dfrac{1}{3};5\right)$ thỏa mãn $27^{3x^2+xy}=(1+xy)27^{15x}$?
	\choice
	{\True $17$}
	{$16$}
	{$18$}
	{$15$}
	\loigiai{
		Xét $f(x)=27^{3x^2-15x+xy}-(xy+1)$ và áp dụng $a^x\geq x(a-1)+1$.\\
		Suy ra $f(x)\geq 26\left(3x^2-15x+xy\right)-xy-1=84x^2+25xy-390x-1>0,\forall y\geq 16,\forall x\in\left(\dfrac{1}{3};5\right)$.\\
		Do đó $y\leq 16$.\\
		$y=0\Rightarrow 27^{3x^2-15x}=1\Rightarrow 3x^2-15x=0$: loại.\\
		$y\leq-3\Rightarrow xy <-1\Rightarrow VP<0\colon$ loại.\\
		$y=-1,y=-2$: thỏa mãn(1).\\
		Xét $y>0$ ta có $f(4)=27^{5y}-(5y+1)\geq 0,\forall y>0$.\\
		và $f\left(\dfrac{1}{3}\right)=3^{y-14}-\dfrac{y}{3}-1<0,\forall y\in\{1;2;3;\ldots;15\}$ (2).\\
		Từ (1) và (2) $\Rightarrow y\in\left\{-2;-1;1;2;3;4;5;6;7;8;9;10;11;15\right\}$.
	}
\end{ex}
\begin{ex}
	Có bao nhiêu số nguyên dương $a$ sao cho ứng với mỗi $a$ tồn tại đúng $8$ số thực $x$ thỏa mãn $\left(x^4-4x^2-3+\log_4a\right)\left(a{\cdot 2}^{2x^4-8x^2-3}+1\right)=-3$?
	\choice
	{$1024$}
	{$1028$}
	{$1023$}
	{\True $1026$}
	\loigiai{
		Đặt $t=x^4-4x^2+\log_4a\Leftrightarrow x^4-4x^2=t-\log_4a=t-\dfrac{1}{2}\log_2a$.\\
		Phương trình trở thành $(t-3)\left(2^{2t-3}+1\right)=-3\Leftrightarrow t-3=-\dfrac{3}{2^{2t-3}+1}\Leftrightarrow g(t)=t-3+\dfrac{3}{2^{2t-3}+1}=0$ (*).\\
		Có $\Leftrightarrow g'(t)=1-\dfrac{6\cdot 2^{2t-3}\ln 2}{\left(2^{2t-3}+1\right)^2}=0$ có đúng 2 nghiệm nên (*) có tối đa 3 nghiệm. Nhận thấy $g(1)=g\left(\dfrac{3}{2}\right)=g(2)=0$ do đó (*) $\Leftrightarrow t=1$; $t=\dfrac{3}{2}$; $t=2$.\\
		Vậy $\heva{&x^4-4x^2+\log_4a=1\\&x^4-4x^2+\log_4a=\dfrac{3}{2}\\&x^4-4x^2+\log_4a=2}\Leftrightarrow\heva{&\log_4a-1=4x^2-x^4\\&\log_4a-\dfrac{3}{2}=4x^2-x^4\\&\log_4a-2=4x^2-x^4.}$ \\
		Để ý ba đường thẳng $y=\log_4a-1$; $y=\log_4a-\dfrac{3}{2}$; $y=\log_4a-2$ đôi một song song.\\
		Hàm số $g(x)=4x^2-x^4$ có bảng biến thiên như sau: \\
		{\color{red} HINH O DAY}\\
		Vậy phương trình có đúng 8 nghiệm khi và chỉ khi\\
		Trường hợp 1: (1) vô nghiệm, (2), (3) mỗi phương trình có 4 nghiệm\\
		$\Leftrightarrow\heva{&\log_4a-1>4\\&0<\log_4a-\dfrac{3}{2}<4\\&0<\log_4a-2<4}\Leftrightarrow 5<\log_4a<5,5\Leftrightarrow 1024<a<2048\Rightarrow a\in\left\{1025;\ldots;2047\right\} $.\\
		Trường hợp 2: (1) có 4 nghiệm và (2); (3) mỗi phương trình có 2 nghiệm\\
		$\Leftrightarrow\heva{&0<\log_4a-1<4\\&\log_4a-\dfrac{3}{2}<0\\&\log_4a-2<0}\Leftrightarrow 1<\log_4a<1,5\Leftrightarrow 4<a<8\Rightarrow a\in\{5;6;7\} $.\\
		Trường hợp 3: (1) có hai nghiệm, (2) có 4 nghiệm và (3) có 2 nghiệm $\Leftrightarrow\heva{&\log_4a-1=4\\&0<\log_4a-\dfrac{3}{2}<4\\&\log_4a-2<0}$ (vô nghiệm).\\
		Vậy có tất cả $2047-1025+1+3=1026$ số nguyên thỏa mãn.
	}
\end{ex}
\begin{ex}
	[Chuyên Quang Trung - Bình Phước - 2021]%Câu 36.
	Số giá trị $m$ nguyên, $m\in[-20;20]$, sao cho $\min\limits_{x\in[0,3; 1]}\left|\dfrac{\log_{0,3}x^m+16}{\log_{0,3}x+m}\right|=16$ là
	\choice
	{$5$}
	{\True $1$}
	{$20$}
	{$40$}
	\loigiai{
		Đặt $t=\log_{0,3}x$.\\
		Đặt $f(x)=\dfrac{m\log_{0,3}x+16}{\log_{0,3}x+m} (x>0)$.\\
		Khi đó: Xét $f(t)=\dfrac{mt+16}{t+m}$ trên đoạn $[0; 1]$.\\
		Từ đó $f'(t)=\dfrac{m^2-16}{(t+m)^2}$.,\\
		$f(0)=\dfrac{16}{m}$, $f(1)=\dfrac{m+16}{m+1}$ (Điều kiện $m\neq 0,-1$).\\
		Trường hợp 1: $m\in[-20;-4]\Rightarrow f'(t)=\dfrac{m^2-16}{(t+m)^2} >0,\forall t\in[0;1]$.\\
		Nên hàm số đồng biến trên khoảng $(0;1)$.\\
		Suy ra, $f(0)\leq f(t)\leq f(1)<0$ nên $\left|f(0)\right|\geq\left|f(t)\right|\geq\left|f(1)\right|>0$, $\forall t\in[0;1]$.\\
		Nên $\max\limits_{t\in[0;1]} f(t)=f(1)=\dfrac{m+16}{m+1}<0\Rightarrow\min\limits_{t\in[0;1]}\left|f(t)\right|=\left|f(1)\right|=\left|\dfrac{m+16}{m+1}\right|(m\neq-1)$.\\
		Mà $\left|\dfrac{m+16}{m+1}\right|=16\Leftrightarrow\hoac{&m=0(l)\\&m=\dfrac{-32}{17}(l).}$ \\
		Trường hợp 2: $m\in[-4;0]\Rightarrow f'(t)=\dfrac{m^2-16}{(t+m)^2} <0,\forall t\in[0;1]$.\\
		Nên hàm số nghịch biến trên đoạn $[0;1]$.\\
		Suy ra, $0>f(0)\geq f(t)\geq f(1)$ nên $\left|f(1)\right|\geq\left|f(t)\right|\geq\left|f(0)\right|>0$, $\forall t\in[0;1]$.\\
		Nên $\max\limits_{x\in[0;1]} f(x)=f(0)\Rightarrow\min\limits_{x\in[0;1]}\left|f(x)\right|=\left|f(0)\right|=\left|\dfrac{16}{m}\right|(m\neq 0)$.\\
		Mà $\left|\dfrac{16}{m}\right|=16\Leftrightarrow\hoac{&m=1(l)\\&m=-1(l).}$ \\
		Trường hợp 3: $m\in[0;4]\Rightarrow f'(t)=\dfrac{m^2-16}{(t+m)^2} <0,\forall t\in[0;1]$.\\
		Nên hàm số nghịch biến trên khoảng $(0; 1)$.\\
		Suy ra, $f(0)\geq f(t)\geq f(1)>0$ nên $\left|f(1)\right|\geq\left|f(t)\right|\geq\left|f(0)\right|>0$, $\forall t\in[0;1]$.\\
		Nên $\min\limits_{x\in[0;1]} f(t)=f(1)\Rightarrow\min\limits_{x\in[0;1]}\left|f(t)\right|=\left|f(1)\right|=\left|\dfrac{m+16}{m+1}\right|(m\neq-1)$.\\
		Mà $\left|\dfrac{m+16}{m+1}\right|=16\Leftrightarrow\hoac{&m=0(n)\\&m=\dfrac{-32}{17}(l).}$ \\
		Trường hợp 4: $m\in[4;20]\Rightarrow f'(t)=\dfrac{m^2-16}{(t+m)^2} >0,\forall t\in[0;1]$.\\
		Nên hàm số đồng biến trên khoảng $(0;1)$.\\
		Suy ra, $0<f(0)\leq f(t)\leq f(1)$ nên $0<\left|f(0)\right|\leq\left|f(t)\right|\leq\left|f(1)\right|$, $\forall t\in[0;1]$.\\
		Nên $\min\limits_{x\in[0;1]} f(t)=f(0)\Rightarrow\min\limits_{x\in[0;1]}\left|f(t)\right|=\left|f(0)\right|=\left|\dfrac{16}{m}\right|(m\neq 0)$.\\
		Mà $\left|\dfrac{16}{m}\right|=16\Leftrightarrow\hoac{&m=1(l)\\&m=-1(l).}$ \\
		Vậy tổng hợp các trường hợp: $m=1$ thì thỏa ycbt.
	}
\end{ex}
\begin{ex}
	[THPT Đồng Quan - Hà Nội - 2021]%Câu 37.
	Có bao nhiêu số nguyên $y$ nằm trong khoảng $(-2021;+\infty)$ sao cho với mỗi giá trị của $y$ tồn tại nhiều hơn hai số thực $x$ thỏa mãn $x^2+y+\left(x^2-x\right)\cdot 2020^{x+y}=\left(2x^2-x+y\right)\cdot 2020^{x-x^2}$?
	\choice
	{\True $2020$}
	{$2019$}
	{$2021$}
	{$2022$}
	\loigiai{
		Ta có $x^2+y+\left(x^2-x\right)\cdot 2020^{x+y}=\left(2x^2-x+y\right)\cdot 2020^{x-x^2}$ \\
		$ \Leftrightarrow\left(x^2+y\right)\cdot 2020^{x^2-x}+\left(x^2-x\right)\cdot 2020^{x^2+y}=2x^2-x+y $ \\
		$\Leftrightarrow\left(x^2+y\right)\cdot\left({2020}^{x^2-x}-1\right)+\left(x^2-x\right)\cdot\left({2020}^{x^2+y}-1\right)=0 (1) $.\\
		Nếu $\left(x^2+y\right)\left(x^2-x\right)\neq 0$ thì $(1)\Leftrightarrow\dfrac{2020^{x^2-x}-1}{x^2-x}+\dfrac{2020^{x^2+y}-1}{x^2+y}=0 (2)$.\\
		Dễ thấy vế trái của $(2)$ luôn dương nên suy ra $(1)$ không xảy ra.\\
		Do đó $(1)\Leftrightarrow\hoac{&x^2-x=0\\&x^2+y=0}\Leftrightarrow\hoac{&x=0\\&x=1\\&x^2=-y.}$ \\
		Với $y=-1$ thì có ba giá trị $x$ thỏa mãn đề bài.\\
		Yêu cầu bài toán $\Leftrightarrow -y>0\Leftrightarrow y<0$.\\
		Do $y$ nguyên nằm trong khoảng $(-2021;+\infty)$ nên $y\in\left\{-2020;-2019;\ldots;-1\right\}$.\\
		Vậy có $2020$ số nguyên $y$ thỏa mãn yêu cầu bài toán.
	}
\end{ex}
\begin{ex}
	[THPT Trần Phú - Đà Nẵng - 2021]%Câu 38.
	Có bao nhiêu số nguyên dương $a$ nhỏ hơn $2021$ sao cho tồn tại số nguyên $x$ thỏa mãn $2^a\left(2^{x+3}-1\right)+a^2=a\left(2^a+2^{x+3}-1\right)$?
	\choice
	{$12$}
	{$15$}
	{\True $10$}
	{$14$}
	\loigiai{
		Ta có $2^a\left(2^{x+3}-1\right)+a^2=a\left(2^a+2^{x+3}-1\right)\Leftrightarrow\left(2^a-a\right)\left(2^{x+3}-1-a\right)=0$ \\
		$ \Leftrightarrow\hoac{&2^a-a=0\\&2^{x+3}-1-a=0.} $ \\
		Xét hàm số $f(a)=2^a-a$.\\
		Ta có $f'(a)=2^a\ln 2-1; f'(a)=0\Leftrightarrow a=\log_2\left(\dfrac{1}{\ln 2}\right)=a_0$.\\
		Bảng biến thiên
		{\color{red} HINH O DAY}\\
		Từ bảng biến thiên, suy ra $f(a)\geq f(a_0)\approx 0,91>0$.\\
		Nên phương trình $2^a-a=0$ vô nghiệm.\\
		$2^{x+3}-1-a=0\Leftrightarrow a=2^{x+3}-1$.\\
		Yêu cầu bài toán $\Leftrightarrow\heva{&2^{x+3}-1>0\\&2^{x+3}-1<2021}\Leftrightarrow\heva{&x >-3\\&x<\log_2\dfrac{1011}{4}\approx 7,9.}$ \\
		Vì $x\in\mathbb{Z}\Rightarrow x\in\left\{-2;-1;0;1;2;3;4;5;6;7\right\}$.\\
		Ta có ứng với mỗi $x\in\mathbb{Z}$ thì có một giá trị $a\in\mathbb{Z}$ nên có $10$ số nguyên $a$ thỏa mãn yêu cầu.
	}
\end{ex}
\begin{ex}
	Có bao nhiêu số nguyên $m(m\geq 2)$ sao cho tồn tại số thực $x$ thỏa mãn $\left(m^{\ln x}+4\right)^{\ln m}+4=x$? 
	\choice
	{$8$}
	{$9$}
	{\True $1$}
	{Vô số}
	\loigiai{
		Điều kiện: $x>0$.\\
		Đặt $y=m^{\ln x}+4>0$ thế vào phương trình ta có $y^{\ln m}+4=x\Leftrightarrow x=4+m^{\ln y}$ vì $m^{\ln y}=y^{\ln m}$.\\
		Khi đó ta có hệ phương trình: $\heva{&y=m^{\ln x}+4\quad(1)\\&x=m^{\ln y}+4\quad(2).}$ \\
		Xét hàm số $f(t)=m^t+4\Rightarrow f'(t)=\ln m\cdot m^t>0$ (Do $m\geq 2$). Nên hàm số $f(t)$ đồng biến trên $\mathbb{R}$.\\
		Khi đó: $x=y$.\\
		Từ (2): $x=m^{\ln x}+4\Leftrightarrow x^{\ln m}=x-4\Leftrightarrow\ln\left(x^{\ln m}\right)=\ln(x-4)\Leftrightarrow\ln m\cdot\ln x=\ln(x-4)\Leftrightarrow\ln m=\dfrac{\ln(x-4)}{\ln x}$.\\
		Do $x>0$ nên $x-4<x\Rightarrow\ln(x-4)<\ln x\Rightarrow\dfrac{\ln(x-4)}{\ln x}<1$.\\
		Nên $\ln m<1\Leftrightarrow m<e$ hay $m\in\{2\}$.
	}
\end{ex}
\begin{ex}
	Có bao nhiêu giá trị nguyên của $m$ với $m>1$ sao cho tồn tại số thực $x$ thỏa mãn: $\left(m^{\log_5x}+3\right)^{\log_5m}=x-3(1)$. 
	\choice
	{$4$}
	{\True $3$}
	{$5$}
	{$8$}
	\loigiai{
		Điều kiện: $x>0$.\\
		Đặt $m^{\log_5x}+3=u$ thay vào phương trình $(1)$ ta được: $u^{\log_5m}=x-3\Leftrightarrow x=u^{\log_5m}+3$.\\
		Vì $u^{\log_5m}=m^{\log_5u}$. Từ đó ta có hệ Phương trình $\heva{&u=m^{\log_5x}+3\\&x=u^{\log_5m}+3.}$ \\
		Xét hàm đặc trưng $f(t)=m^t+3$ trên $\mathbb{R}$.\\
		Do $m>1$. Suy ra hàm số $f(t)$ đồng biến trên $\mathbb{R}$.\\
		Do đó, $f(\log_5x)=f(\log_5u)\Leftrightarrow x=u$.\\
		Vì thế, ta đưa về xét phương trình: $x=m^{\log_5x}+3\Leftrightarrow x=x^{\log_5m}+3\Leftrightarrow x-3=x^{\log_5m}$ \\
		$\Leftrightarrow\log_5(x-3)=\log_5\left(x^{\log_5m}\right)\Leftrightarrow\log_5(x-3)=\log_5x\cdot\log_5m\Leftrightarrow\log_5m=\dfrac{\log_5(x-3)}{\log_5x} $.\\
		Do $x>0$ nên $x-3<x$ nên $\log_5m=\dfrac{\log_5(x-3)}{\log_5x}<1\Leftrightarrow m<5$.\\
		Suy ra $\heva{&m\in\mathbb{Z}\\&1<m<5}\Rightarrow m\in\{2,3,4\}$.\\
		Vậy, có $3$ giá trị tham số $m$ thỏa mãn.
	}
\end{ex}
\begin{ex}
	Cho các số thực $x, y, z$ thỏa mãn $\log_3\left(2x^2+y^2\right)=\log_7\left(x^3+2y^3\right)=\log z$. Có bao giá trị nguyên của $z$ để có đúng hai cặp $(x,y)$ thỏa mãn đẳng thức trên. 
	\choice
	{$2$}
	{\True $211$}
	{$99$}
	{$4$}
	\loigiai{
		Ta có $\log_3\left(2x^2+y^2\right)=\log_7\left(x^3+2y^3\right)=\log z=t \Leftrightarrow\heva{&2x^2+y^2=3^t(1)\\&x^3+2y^3=7^t(2)\\&z={10}^t(3).}$ \\
		Nếu $y=0 (2)\Rightarrow x=7^{\frac{t}{3}}$ thay vào $(1)$ ta được $2.7^{\frac{2t}{3}}=3^t\Leftrightarrow t=\log_{\frac{3}{\sqrt[3]{49}}}2$ do đó $z=10^{\log_{\frac{3}{\sqrt[3]{49}}}2}$.\\
		Nếu $y\neq 0$.\\
		Từ $(1)\& (2)$ suy ra $\heva{&\left(2x^2+y^2\right)^3={27}^t\\&\left(x^3+2y^3\right)^2={49}^t}
		\Rightarrow\dfrac{\left(x^3+2y^3\right)^2}{\left(2x^2+y^2\right)^3}=\left(\dfrac{49}{27}\right)^t
		\Leftrightarrow\dfrac{\left(\left(\frac{x}{y}\right)^3+2\right)^2}{\left(2\left(\frac{x}{y}\right)^2+1^2\right)^3}=\left(\dfrac{49}{27}\right)^t,(*)$.\\
		Đặt $\dfrac{x}{y}=u,u\neq-\sqrt[3]{2}$. Xét $f(u)=\dfrac{\left(u^3+2\right)^2}{\left(2u^2+1\right)^3}\Rightarrow f'(u)=\dfrac{6u\left(u^3+2\right)(u-4)}{\left(2u^2+1\right)^4}=0 \Leftrightarrow\hoac{&u=0\\&u=-\sqrt[3]{2}\\&u=4.}$ \\
		Ta có bảng biến thiên
		{\color{red} HINH O DAY}\\
		Nhận xét với mỗi giá trị $u$ tương ứng với duy nhất 1 cặp $(x,y)$ thỏa mãn bài toán do đó.\\
		Yêu cầu bài toán tương đương $\hoac{&\dfrac{1}{8}\leq\left(\dfrac{49}{27}\right)^t<4\\&0<\left(\dfrac{49}{27}\right)^t<\dfrac{4}{33}}
		\Leftrightarrow\hoac{&{10}^{\log_{\frac{49}{27}}\left(\frac{1}{8}\right)}\leq z<{10}^{\log_{\frac{49}{27}}4}\\&0<z<{10}^{\log_{\frac{49}{27}}\left(\frac{4}{33}\right)}.}$ \\
		Vì $z$ là số nguyên nên có $211$ giá trị thỏa mãn.
	}
\end{ex}
\begin{ex}
	Có bao nhiêu số tự nhiên $a$ sao cho tồn tại số thực $x$ thoả $2021^{x^3-a^{3\log(x+1)}}\left(x^3+2020\right)=a^{3\log(x+1)}+2020$ 
	\choice
	{\True $9$}
	{$8$}
	{$5$}
	{$12$}
	\loigiai{
		Xét phương trình: $2021^{x^3-a^{3\log(x+1)}}=\dfrac{a^{3\log(x+1)}+2020}{x^3+2020}$, điều kiện: $x >-1$,\\
		$\begin{aligned}&\Leftrightarrow x^3-a^{3\log(x+1)}=\log_{2021}\left(a^{3\log(x+1)}+2020\right)-\log_{2021}\left(x^3+2020\right)\\&\Leftrightarrow x^3+\log_{2021}\left(x^3+2020\right)=a^{3\log(x+1)}+\log_{2021}\left(a^{3\log(x+1)}+2020\right)(*).\end{aligned}$\\
		Xét hàm số $f(t)=t^3+\log_{2021}\left(t^3+2020\right)$, trên $(0;+\infty)$.\\
		$f'(t)=3t^2+\dfrac{3t^2}{\left(t^3+2020\right)\ln 2021}>0,\forall t>0$ nên hàm số $f(t)$ đồng biến trên $(0;+\infty)$.\\
		Do đó $(*)$ trở thành: $x=a^{\log(x+1)}\Leftrightarrow x=(x+1)^{\log a}\Leftrightarrow\log x=\log a\cdot\log (x+1)$ \\
		$ \Leftrightarrow\log a=\dfrac{\log x}{\log(x+1)}<1,\forall x >-1 $ nên $a<10\Rightarrow a\in\left\{1,2,3,4,5,6,7,8,9\right\}$.
	}
\end{ex}
\begin{ex}
	Cho các số thực $x, y, z$ thỏa mãn $\log_3\left(2x^2+y^2\right)=\log_7\left(x^3+2y^3\right)=\log z$. Có bao giá trị nguyên của $z$ để có đúng hai cặp $(x,y)$ thỏa mãn đẳng thức trên. 
	\choice
	{$2$}
	{\True $211$}
	{$99$}
	{$4$}
	\loigiai{
		Ta có $\log_3\left(2x^2+y^2\right)=\log_7\left(x^3+2y^3\right)=\log z=t\Leftrightarrow\heva{&2x^2+y^2=3^t(1)\\&x^3+2y^3=7^t(2)\\&z={10}^t(3).}$ \\
		Nếu $y=0 (2)\Rightarrow x=7^{\frac{t}{3}}$ thay vào $(1)$ ta được $2.7^{\frac{2t}{3}}=3^t\Leftrightarrow t=\log_{\frac{3}{\sqrt[3]{49}}}2$ do đó $z=10^{\log_{\frac{3}{\sqrt[3]{49}}}2}$.\\
		Nếu $y\neq 0$.\\
		Từ $(1)\&(2)$ suy ra $\heva{&\left(2x^2+y^2\right)^3={27}^t\\&\left(x^3+2y^3\right)^2={49}^t}
		\Rightarrow\dfrac{\left(x^3+2y^3\right)^2}{\left(2x^2+y^2\right)^3}=\left(\dfrac{49}{27}\right)^t
		\Leftrightarrow\dfrac{\left(\left(\dfrac{x}{y}\right)^3+2\right)^2}{\left(2\left(\dfrac{x}{y}\right)^2+1^2\right)^3}=\left(\dfrac{49}{27}\right)^t,(*)$.\\
		Đặt $\dfrac{x}{y}=u,u\neq-\sqrt[3]{2}$. Xét $f(u)=\dfrac{\left(u^3+2\right)^2}{\left(2u^2+1\right)^3}
		\Rightarrow f'(u)=\dfrac{6u\left(u^3+2\right)(u-4)}{\left(2u^2+1\right)^4}=0
		\Leftrightarrow\hoac{&u=0\\&u=-\sqrt[3]{2}\\&u=4.}$ \\
		Ta có bảng biến thiên\\
		{\color{red} HINH O DAY}\\
		Nhận xét với mỗi giá trị $u$ tương ứng với duy nhất 1 cặp $(x,y)$ thỏa mãn bài toán do đó.\\
		Yêu cầu bài toán tương đương $\hoac{&\dfrac{1}{8}\leq\left(\dfrac{49}{27}\right)^t<4\\&0<\left(\dfrac{49}{27}\right)^t<\dfrac{4}{33}}\Leftrightarrow\hoac{&{10}^{\log_{\frac{49}{27}}\left(\dfrac{1}{8}\right)}\leq z<{10}^{\log_{\frac{49}{27}}4}\\&0<z<{10}^{\log_{\frac{49}{27}}\left(\dfrac{4}{33}\right)}.}$ \\
		Vì $z$ là số nguyên nên có $211$ giá trị thỏa mãn.
	}
\end{ex}
\begin{ex}[Chuyên Long An - 2021]%Câu 44.
	Cho $x,y$ là các số thực thoả mãn $\log_3(3x+6)+x-2y=3\cdot 9^y$. Biết $5\leq x\leq 2021$, tìm số cặp $x,y$ nguyên thoả mãn đẳng thức trên. 
	\choice
	{$5$}
	{\True $2$}
	{$4$}
	{$3$}
	\loigiai{
		Ta có $\log_3(3x+6)+x-2y=3\cdot 9^y\Leftrightarrow\log_3(x+2)+x+2=3^{2y+1}+2y+1$ \\
		$\Leftrightarrow\log_3(x+2)+x+2=3^{2y+1}+\log_33^{2y+1}\quad(1) $.\\
		Xét hàm số $f(t)=\log_3(t)+t, t>0$.\\
		Ta có: $f'(t)=\dfrac{1}{t\cdot\ln 3}+1>0,\forall t>0\Rightarrow f(t)$ là hàm số đồng biến trên tập xác định.\\
		Từ (1) suy ra $x+2=3^{2y+1}\Leftrightarrow x=3^{2y+1}-2$.\\
		Do $5\leq x\leq 2021$ nên $5\leq 3^{2y+1}-2\leq 2021\Leftrightarrow\dfrac{\log_37-1}{2}\leq y\leq\dfrac{\log_32023-1}{2}$.\\
		Mà $y$ nguyên nên $\Rightarrow\hoac{&y=1\\&y=2}\Rightarrow\hoac{&x=25\\&x=241}$.
	}
\end{ex}
\begin{ex}
	[Chuyên Hùng Vương - Gia Lai - 2021]%Câu 45.
	Có bao nhiêu số nguyên $a(a\geq 2)$ sao cho tồn tại số thực $x$ thỏa mãn $\ln\left(a^{\log x^4}+4a^{\log x^2}+4\right)=\dfrac{\ln(x-2)}{\log a}$?
	\choice
	{\True $2$}
	{$3$}
	{$1$}
	{$9$}
	\loigiai{
		Ta có:\\
		$\ln\left(a^{\log x^4}+4a^{\log x^2}+4\right)=\dfrac{\ln(x-2)}{\log a}\Leftrightarrow\ln\left(a^{4\log x}+4a^{2\log x}+4\right)=\dfrac{\ln(x-2)}{\log a}\Leftrightarrow 2\ln\left(a^{2\log x}+2\right)=\dfrac{\ln(x-2)}{\log a}$.\\
		Đặt $a^{2\log x}+2=t\Rightarrow\log a\cdot 2\log x=\log(t-2)\Rightarrow\log a=\dfrac{\log(t-2)}{2\log x}\Rightarrow\ln t\cdot\ln(t-2)=\ln x\cdot\ln(x-2)$.\\
		Xét hàm $f(u)=\ln u\cdot\ln(u-2)$ \\
		$ \Rightarrow f'(u)=\dfrac{\ln(u-2)}{u}+\dfrac{\ln u}{u-2}>0 $.\\
		Do $t-2=a^{2\log x}\geq 2^{2\log 2}>1\Rightarrow u=x\Rightarrow a^{2\log x}=x-2\Leftrightarrow x-x^{2\log a}=2\Rightarrow x>x^{2\log a}\Leftrightarrow 2\log a<1$ \\
		$ \Leftrightarrow\log a<\dfrac{1}{2}\Leftrightarrow a<\sqrt{10}\Rightarrow a\in\{2;3\} $.
	}
\end{ex}
\begin{ex}
	[Liên Trường Nghệ An – 2021]%Câu 46.
	Có bao nhiêu cặp $(x;y)$ thỏa mãn $10^{\frac{10}{x+y}}=\left(x+y+\dfrac{1}{x}+\dfrac{1}{y}\right)10^{\frac{1}{xy}}$ và $x\in\mathbb{N}^*, y>0$. 
	\choice
	{\True $14$}
	{$7$}
	{$21$}
	{$10$}
	\loigiai{
		$10^{\frac{10}{x+y}}=\left(x+y+\dfrac{1}{x}+\dfrac{1}{y}\right)\cdot 10^{\frac{1}{xy}}\Leftrightarrow 10^{\frac{10}{x+y}}=\dfrac{(x+y)(xy+1)}{xy}\cdot 10^{\frac{1}{xy}}\Leftrightarrow\dfrac{10}{x+y}\cdot 10^{\frac{10}{x+y}}=\left(1+\dfrac{1}{xy}\right)\cdot 10^{\frac{1}{xy}+1}$.\\
		Xét hàm số $f(t)=t\cdot 10^t$ trên khoảng $(0;+\infty)$.\\
		$f'(t)=10^t+t\cdot 10^t\ln 10>0,\forall t>0$ nên hàm số $f(t)=t\cdot 10^t$ đồng biến trên khoảng $(0;+\infty)$.\\
		Do đó: $\dfrac{10}{x+y}\cdot 10^{\frac{10}{x+y}}=\left(1+\dfrac{1}{xy}\right)\cdot 10^{\frac{1}{xy}+1}\Leftrightarrow\dfrac{10}{x+y}=1+\dfrac{1}{xy}$ \\
		$ \Leftrightarrow 10=(x+y)\left(1+\dfrac{1}{xy}\right)\Leftrightarrow 10-\left(y+\dfrac{1}{y}\right)=x+\dfrac{1}{x} $.\\
		Vì $y+\dfrac{1}{y}\geq 2,\forall y>0$ nên $x+\dfrac{1}{x}\leq 8\Leftrightarrow x^2-8x+1\leq 0\Leftrightarrow 4-\sqrt{15}\leq x\leq 4+\sqrt{15}$, $x\in\mathbb{N}^*$ \\
		$ \Rightarrow x\in A=\left\{1;2;3;4;5;6;7\right\} $.\\
		Với mỗi số $a>0$ phương trình $y+\dfrac{1}{y}=a\Leftrightarrow y^2-ay+1=0 (*)$ có $\heva{&\Delta>0\\&S=a>0\\&P=1>0}$ \\
		$ \Rightarrow $ Phương trình (*) luôn có hai nghiệm $y>0$. Vậy có $14$ cặp $(x;y)$ thỏa mãn yêu cầu bài toán.
	}
\end{ex}
\begin{ex}
	[Chuyên Vinh - 2021]%Câu 47.
	Có bao nhiêu số thực của $y$ để với mỗi $y$ tồn tại đúng 2 giá trị thực của $x$ sao cho $\ln\left(4x^2\right)=xy+y$?
	\choice
	{$1$}
	{Vô số}
	{$2$}
	{\True $3$}
	\loigiai{
		Điều kiện: $x\neq 0$.\\
		Xét $x=-1$, khi đó: $\ln\left(4x^2\right)=xy+y\Rightarrow\ln(4)=0$ (vô lí).\\
		Suy ra: $x\neq-1$.\\
		Xét: $\ln\left(4x^2\right)=xy+y\Leftrightarrow y=\dfrac{\ln\left(4x^2\right)}{x+1}$.\\
		Để với mỗi $y$ tồn tại đúng 2 giá trị thực của $x$ khi đồ thị hàm số $y=\dfrac{\ln\left(4x^2\right)}{x+1}$ và đường thẳng song song với trục hoành là $y=m$ cắt nhau tại 2 điểm phân biệt.\\
		Xét hàm số $y=\dfrac{\ln\left(4x^2\right)}{x+1}$, $y'=\dfrac{\dfrac{2}{x}+2-\ln\left(4x^2\right)}{(x+1)^2}=\dfrac{g(x)}{(x+1)^2}$.\\
		Xét: $g'(x)=\dfrac{-2}{x^2}-\dfrac{2}{x}=\dfrac{-2(x+1)}{x^2}, g'(x)>0\Leftrightarrow x >-1$, khi đó $g(x)$ đồng biến trên $(-\infty;-1)$.\\
		Xét: $\lim\limits_{x\to(-1)^-} g(x)=-\ln(4)<0$, suy ra: $g(x)<0,\forall x <-1$.\\
		Xét: $\lim\limits_{x\to(-1)^+} g(x)=-\ln(4)<0$, suy ra: $g(x)<0,\forall x\in(-1; 0)$.\\
		Trên $(0;+\infty)$, $g(x)$ nghịch biến, liên tục và $g(3)\cdot g(4)<0\Rightarrow\exists a\in(3; 4)\colon g(a)=0$.\\
		Xét: $\lim\limits_{x\to-\infty} y=0$, $\lim\limits_{x\to+\infty} y=0$, $\lim\limits_{x\to 0^+} y=-\infty$, $\lim\limits_{x\to 0^-} y=-\infty$, $\lim\limits_{x\to(-1)^-} y=-\infty$, $\lim\limits_{x\to(-1)^+} y=+\infty$.\\
		Bảng biến thiên của hàm số $y=\dfrac{\ln\left(4x^2\right)}{x+1}$: \\
		{\color{red} HINH O DAY}\\
		Vậy: Để với mỗi $y$ tồn tại đúng 2 giá trị thực của $x$ khi $y=y(a)$ hoặc $y=0$. Khi đó tồn tại 2 giá trị của $y$ thỏa yêu cầu đề bài.
	}
\end{ex}
\begin{ex}
	[Chuyên Thái Bình - 2021]%Câu 48.
	Có bao nhiêu số nguyên $a (a\geq 3)$ sao cho tồn tại số thực $x$ thỏa mãn $\left(a^{\log_{2021}x}+3\right)^{\log_{2021}a}=x-3$?
	\choice
	{$2019$}
	{\True $2018$}
	{$2020$}
	{$2003$}
	\loigiai{
		Điều kiện có nghiệm là $x>3$.\\
		Ta có $\left(a^{\log_{2021}x}+3\right)^{\log_{2021}a}=x-3\Leftrightarrow\left(x^{\log_{2021}a}+3\right)^{\log_{2021}a}+\left(x^{\log_{2021}a}+3\right)=x^{\log_{2021}a}+x (1)$.\\
		Xét hàm số $f(t)=t^{\log_{2021}a}+t$ trên $(0;+\infty)$.\\
		$f'(t)=t^{\log_{2021}a-1}\cdot\log_{2021}a+1>0\forall t>0 v\mu a\geq 3\Rightarrow$ Hàm số $f(t)$ đồng biến trên $(0;+\infty)$.\\
		Do đó $(1)\Leftrightarrow f\left(x^{\log_{2021}a}+3\right)=f(x)\Leftrightarrow x^{\log_{2021}a}+3=x\Leftrightarrow x^{\log_{2021}a}=x-3$ \\
		$ \Leftrightarrow\log_{2021}x^{\log_{2021}a}=\log_{2021}(x-3)\Leftrightarrow\log_{2021}a\log_{2021}x=\log_{2021}(x-3) $ \\
		$ \Leftrightarrow\log_{2021}a=\dfrac{\log_{2021}(x-3)}{\log_{2021}x} $.\\
		Mà $\log_{2021}(x-3)<\log_{2021}x\forall x>3$ và $\log_{2021}x>0\forall x>3$ nên $\dfrac{\log_{2021}(x-3)}{\log_{2021}x}<1\forall x>3$.\\
		Lại có: Hàm số $g(x)=\dfrac{\log_{2021}(x-3)}{\log_{2021}x}$ liên tục trên $(3;+\infty)$.\\
		và $\lim\limits_{x\to+\infty} g(x)=1$; $\lim\limits_{x\to 3^+} g(x)=-\infty$.\\
		Do đó tập giá trị của hàm số $g(x)=\dfrac{\log_{2021}(x-3)}{\log_{2021}x}$ là $(-\infty; 1)$.\\
		Khi đó phương trình đã cho có nghiệm thì $\log_{2021}a<1\Leftrightarrow a<2021$.\\
		Lại có $a$ nguyên và $a\geq 3$ nên $a\in\left\{3; 4;\ldots; 2020\right\}$.\\
		Vậy có $2018$ số nguyên $a (a\geq 3)$ thỏa mãn yêu cầu bài toán.
	}
\end{ex}
\begin{ex}
	[Chuyên Thái Bình - 2021]%Câu 49.
	Tìm số giá trị nguyên của tham số thực $m$ để tồn tại các số thực $x$; $y$ thỏa mãn $\mathrm{e}^{x^2+y^2-m}+\mathrm{e}^{x+y+xy-m}=x^2+y^2+x+y+xy-2m+2$: 
	\choice
	{$6$}
	{\True $9$}
	{$8$}
	{$7$}
	\loigiai{
		Ta có $\mathrm{e}^{x^2+y^2-m}+\mathrm{e}^{x+y+xy-m}=x^2+y^2+x+y+xy-2m+2$ \\
		$ \Leftrightarrow\mathrm{e}^{x^2+y^2-m}-\left(x^2+y^2-m\right)-1+\mathrm{e}^{x+y+xy-m}-(x+y+xy-m)-1=0 (1) $.\\
		Xét hàm số $f(t)=\mathrm{e}^t-t-1$ trên $\mathbb{R}$.\\
		$f'(t)=\mathrm{e}^t-1$.\\
		$f'(t)=0\Leftrightarrow\mathrm{e}^t-1=0\Leftrightarrow t=0$.\\
		Bảng biến thiên: 
		{\color{red} HINH O DAY}\\
		Từ bảng biến thiên suy ra $f(t)\geq 0\forall t\in\mathbb{R}$.\\
		Do đó $(1)\Leftrightarrow f\left(x^2+y^2-m\right)+f(x+y+xy-m)=0$ \\
		$\Leftrightarrow\heva{&x^2+y^2-m=0\\&x+y+xy-m=0}\Leftrightarrow\heva{&(x+y)^2-2xy-m=0\\&x+y+xy-m=0} (2) $.\\
		Đặt $\heva{&x+y=S\\&xy=P}\left(S^2\geq 4P\right)$.\\
		$(2)$ trở thành $\heva{&S^2-2P-m=0\\&S+P-m=0}\Leftrightarrow\heva{&S^2+2S=3m\\&S+P=m} (3)$.\\
		Xét phương trình $S^2+2S=3m\Leftrightarrow S^2+2S-3m=0 (4)$.\\
		Hệ $(3)$ có nghiệm $\Leftrightarrow (4)$ có nghiệm $\Leftrightarrow\Delta'=1+3m\geq 0\Leftrightarrow m\geq-\dfrac{1}{3}$.\\
		Do $S^2\geq 4P$ nên $P\leq\dfrac{1}{4}S^2$.\\
		Ta có: $S^2+2S=3m=3S+3P\leq 3S+\dfrac{3}{4}S^2\Rightarrow\dfrac{1}{4}S^2-S\leq 0\Rightarrow 0\leq S\leq 4$ \\
		$ \Rightarrow 0\leq S^2+2S\leq 24\Rightarrow 0\leq 3m\leq 24\Rightarrow 0\leq m\leq 8 $.\\
		Do $m$ nguyên nên $m\in\left\{0; 1;\ldots; 8\right\}$.\\
		Vậy có $9$ giá trị nguyên của tham số thực $m$ thỏa mãn yêu cầu bài toán.
	}
\end{ex}
\begin{ex}
	[Sở Đồng Tháp 2021]%Câu 50.
	Có bao nhiêu cặp số nguyên $(x; y)$ thỏa mãn $0\leq y\leq 2020$ và $\log_2(4y+4)-x=1+2^x-y$? 
	\choice
	{\True $11$}
	{$12$}
	{$2021$}
	{$10$}
	\loigiai{
		Điều kiện: $4y+4>0\Leftrightarrow y >-1$.\\
		Ta có $\log_2(4y+4)-x=1+2^x-y$.\\
		$\begin{aligned}&\Leftrightarrow\log_2[4(y+1)]-x=1+2^x-y\\&\Leftrightarrow 2+\log_2(y+1)-x=1+2^x-y\\&\Leftrightarrow\log_2(y+1)+y+1=x+2^x\\&\Leftrightarrow\log_2(y+1)+2^{\log_2(y+1)}=x+2^x.\end{aligned}$\\
		Xét hàm đặc trưng $f(t)=t+2^t$ ta có $f'(t)=1+2^t\ln 2>0$ \\
		$ \Rightarrow $ Hàm số $y=f(t)$ đồng biến trên R, do đó ta có $\log_2(y+1)=x\Leftrightarrow y+1=2^x$.\\
		Theo đề bài ta có $0\leq y\leq 2020$.\\
		$\begin{aligned}&\Leftrightarrow 0\leq 2^x-1\leq 2020\\&\Leftrightarrow 1\leq 2^x\leq 2021\\&\Leftrightarrow 0\leq x\leq\log_22021\approx 10,98\end{aligned}$\\
		Mà $y\in \mathbb{Z}\Rightarrow y\in\left\{0, 1; 2; 3; 4; 5; 6; 7; 8; 9; 10\right\}$.\\
		Ứng với mỗi giá trị của y cho 1 giá trị x tương ứng.\\
		Vậy có 11 cặp số nguyên $(x; y)$ thỏa yêu cầu bài toán.
	}
\end{ex}
\begin{ex}
	[Chuyên Lê Hồng Phong - 2021]%Câu 51.
	Có bao nhiêu số nguyên $y$ thuộc đoạn $[-2021; 2021]$ sao cho ứng với mỗi $y$ tồn tại số thực $x$ thỏa mãn $\log_2\left(y-5+\sqrt{y+2^x-5}\right)=2x$?
	\choice
	{\True $2017$}
	{$2016$}
	{$4041$}
	{$2021$}
	\loigiai{
		Ta có\\
		$\log_2\left(y-5+\sqrt{y+2^x-5}\right)=2x\Leftrightarrow y-5+\sqrt{y+2^x-5}=2^{2x}\Leftrightarrow y+2^x-5+\sqrt{y+2^x-5}=2^{2x}+2^x(1)$.\\
		Xét hàm số $g(t)=t^2+t$, $\forall t>0$.\\
		Có $g'(t)=2t+1>0,\forall t>0$. Do đó, hàm số $g(t)$ đồng biến trên khoảng $(0;+\infty)$.\\
		$(1)\Leftrightarrow g\left(\sqrt{y+2^x-5}\right)=g(2^x)\Leftrightarrow y+2^x-5=2^{2x}\Leftrightarrow y=2^{2x}-2^x+5$.\\
		Xét hàm số $f(x)=2^{2x}-2^x+5$.\\
		Có $f'(x)=2\cdot\ln 2\cdot 2^{2x}-\ln 2\cdot 2^x$, $f'(x)=0\Leftrightarrow x=-1$.\\
		Ta có BBT
		{\color{red} HINH O DAY}\\
		Vì $y\in[-2021; 2021]$, $y\in\mathbb{Z}$ nên $\dfrac{19}{4}\leq y\leq 2021$.\\
		Vậy có $2017$ giá trị cần tìm của $y$.
	}
\end{ex}
\begin{ex}
	[Sở Phú Thọ - 2021]%Câu 52.
	Có bao nhiêu cặp số nguyên $(x;y)$ thỏa mãn $0<y\leq 2021$ và $3^x+3x-6=9y+\log_3y^3$?
	\choice
	{$2021$}
	{\True $7$}
	{$9$}
	{$2020$}
	\loigiai{
		$3^{x-1}+(x-1)=3y+\log_3(3y)$.\\
		Đặt $3^{x-1}=u\Rightarrow x-1=\log_3u,(u>0)$, suy ra: $u+\log_3u=3y+\log_3(3y)$. $(*)$.\\
		Xét hàm số $f(t)=t+\log_3t$ trên $(0;+\infty)$.\\
		Ta có: $f'(t)=1+\dfrac{1}{t\ln 3}>0$, $\forall t>0$ nên từ $(*)$ suy ra:\\
		$(*)\Leftrightarrow f(u)=f(3y)\Leftrightarrow u=3y$.\\
		Khi đó ta có: $3y=3^{x-1}\Leftrightarrow y=3^{x-2} (**)$.\\
		Theo giả thiết: $\heva{&y\in\mathbb{Z}\\&0<y\leq 2021}\Rightarrow 1\leq y\leq 2021$, suy ra:\\
		$\heva{&x\in\mathbb{Z}\\&1\leq 3^{x-2}\leq 2021}\Leftrightarrow\heva{&x\in\mathbb{Z}\\&0\leq x-2\leq\log_32021\approx 6,928}$ \\
		$ \Leftrightarrow\heva{&x\in\mathbb{Z}\\&0\leq x-2\leq 6}\Leftrightarrow\heva{&x\in\mathbb{Z}\\&2\leq x\leq 8}\Rightarrow x\in\left\{2;3;4;5;6;7;8\right\} $ (có $7$ số).\\
		Từ $(**)$ ta có, ứng với mỗi giá trị của $x$, cho duy nhất một giá trị của $y$ nên có $7$ cặp.
	}
\end{ex}
\begin{ex}
	[THPT Nghi Sơn - Thanh Hóa - 2021]%Câu 53.
	Có bao nhiêu cặp số $(x;y)$ thuộc đoạn $[1;2020]$ thỏa mãn $y$ là số nguyên và $x+\ln x=y+\mathrm{e}^y$?
	\choice
	{$6$}
	{\True $7$}
	{$2021$}
	{$2020$}
	\loigiai{
		Ta có $x+\ln x=y+\mathrm{e}^y\Leftrightarrow x+\ln x=\mathrm{e}^y+\ln\mathrm{e}^y$.\\
		Xét hàm số $f(t)=t+\ln t$ trên $(0;+\infty)$, ta có $f'(t)=1+\dfrac{1}{t}>0,\forall t>0\Rightarrow f(t)$ đồng biến trên $(0;+\infty)\Rightarrow f(x)=f(\mathrm{e}^y)\Rightarrow x=\mathrm{e}^y$.\\
		Do $x\in[1;2020]$ nên $1\leq\mathrm{e}^y\leq 2020\Leftrightarrow 0\leq y\leq\ln 2020$.\\
		Mà $y\in\mathbb{Z}$ và $y\in[1;2020]$ nên $y\in\left\{1;2;3;4;5;6;7\right\}$.\\
		Vậy có $7$ cặp số $(x;y)$ thỏa mãn bài toán.
	}
\end{ex}
\begin{ex}
	[THPT Cẩm Bình - Hà Tĩnh - 2021]%Câu 54.
	Có bao nhiêu cặp số nguyên $(x;y)$ thoả mãn $2\leq x\leq 2021$ và $2^y-\log_2\left(x+2^{y-1}\right)=2x-y$ 
	\choice
	{$2020$}
	{\True $10$}
	{$9$}
	{$2021$}
	\loigiai{
		Đặt $\log_2\left(x+2^{y-1}\right)=t\Leftrightarrow x+2^{y-1}=2^t\Leftrightarrow x+\dfrac{2^y}{2}=2^t\Leftrightarrow 2x+2^y=2\cdot 2^t\Leftrightarrow 2x=2\cdot 2^t-2^y$.\\
		Phương trình đã cho trở thành: $2^y-t=2\cdot 2^t-2^y-y\Leftrightarrow 2\cdot 2^y+y=2\cdot 2^t+t(*)$.\\
		Xét hàm số: $f(u)=2\cdot 2^u+u$. Tập xác định $\mathscr{D}=\mathbb{R}$.\\
		Ta có: $f'(u)=2\cdot 2^u\cdot\ln 2+1>0\forall u$. Suy ra hàm số $f(u)=2\cdot 2^u+u$ đồng biến trên $\mathbb{R}$.\\
		Do đó phương trình $(*)$ có một nghiệm duy nhất:\\
		$y=t\Leftrightarrow 2x=2\cdot 2^y-2^y\Leftrightarrow 2x=2^y\Leftrightarrow y=\log_2(2x)$.\\
		Do $2\leq x\leq 2021\Leftrightarrow\log_24\leq y\leq\log_24042\Leftrightarrow 2\leq y\leq\log_24042\approx 11,98$.\\
		Lại có $y$ là số nguyên nên: $y\in\left\{2;3\cdot;\ldots;11\right\}$.\\
		Với mỗi giá trị của $y$ ta tìm được duy nhất một giá trị của x. Vậy có $10$ cặp số nguyên $(x;y)$ để thoả mãn điều kiện bài toán.
	}
\end{ex}
\begin{ex}
	[Chuyên Vinh – 2022]%Câu 55.
	Có bao nhiêu số nguyên $a$ sao cho ứng với mỗi $a$, tồn tại số thực $b\geq a$ thỏa mãn $4^a=2^b+b$ và đoạn $[a;b]$ chứa không quá 5 số nguyên?
	\choice
	{$5$}
	{$10$}
	{$6$}
	{\True $11$}
	\loigiai{
		Do đoạn $[a;b]$ chứa không quá 5 số nguyên nên ta có điều kiện đủ là $a\leq b<a+5$. Khi đó ta có\\
		Phương trình ban đầu tương đương với: $2^b+b-4^a=0$.\\
		Xét hàm số $y=f_a(b)=2^b+b-4^a$ có $f_a'(b)=2^b\ln 2+1>0,\forall b\in\mathbb{R}$ tức hàm $f_a(b)$ đồng biến trên $\mathbb{R}$.\\
		Kết hợp điều kiện cần ban đầu ta suy ra hàm số $f_a(b)$ đồng biến trên $[a;a+5)$.\\
		Như vậy điều kiện tồn tại nghiệm là $\heva{&f_a(a+5)>0\\&f_a(a)\leq 0}\Leftrightarrow\heva{&2^{a+5}+a+5>4^a\\&2^a+a\leq 4^a}$..\\
		Trường hợp 1: Nếu $a+5<0\Leftrightarrow a\leq-6$ thì $2^{a+5}+a+5\leq 2^{\circ}-1=0<4$ (loại).\\
		Trường hợp 2: Nếu $a+5\geq 0\Leftrightarrow a\geq-5$ thì khi đó $2^{a+5}+a+5>4^a\Leftrightarrow 4^a<2^{a+6}\Leftrightarrow 2^{2a}<2^{a+6}\Leftrightarrow a<6$ Đối chiếu với điều kiện ta suy ra $-5\leq a\leq 5$.\\
		Đến đây với mọi $a\in [-5;5]$ thì bất phương trình $2^{a+5}+a+5>4^a$ luôn xảy ra vì $4^a\leq 2^{a+5}\leq 2^{a+5}+a+5$ (không có dấu bằng xảy ra).\\
		Xét bất phương trình còn lại: $2^a+a\leq 4^a$ ta thấy cũng luôn đúng với mọi $a\in [-5;5]$.\\
		Vậy $a\in [-5;5]$ thì thỏa mãn yêu cầu đề bài tức có 11 giá trị nguyên $a$.
	}
\end{ex}
\begin{ex}
	[Cụm Trường Nghệ An - 2022]%Câu 56.
	Có bao nhiêu cặp số nguyên $(x; y)$ thỏa mãn đẳng thức sau:\\
	$\log_{2022}\left[\left(x^4-2x^2+2023\right)^{y^2+2022}\right]=2y+2021$. 
	\choice
	{$3$}
	{$1$}
	{$0$}
	{\True $2$}
	\loigiai{
		Ta có $x^4-2x^2+2023 =\left(x^2-1\right)^2+2022\geq 2022$ với $\forall x\in\mathbb{R}$ \\
		$ \Rightarrow\log_{2022}\left(x^4-2x^2+2023\right)\geq 1 $ với $\forall x\in\mathbb{R}$.\\
		Lại có $\heva{&y^2+2022\geq 2022\\&y^2+2022=y^2+1+2021\geq 2y+2021}$ với $\forall y\in\mathbb{R}$ \\
		$\Rightarrow\log_{2022}\left[\left(x^4-2x^2+2023\right)^{y^2+2022}\right]=\left(y^2+2022\right)\log_{2022}\left(x^4-2x^2+2023\right)\geq 2y+2021\forall x, y\in\mathbb{R} $.\\
		Do đó $\log_{2022}\left[\left(x^4-2x^2+2023\right)^{y^2+2022}\right]=2y+2021\Leftrightarrow\heva{&x^2-1=0\\&y=1}\Leftrightarrow\hoac{&\heva{&x=1\\&y=1}\\&\heva{&x=-1\\&y=1}.}$ \\
		Có hai cặp số nguyên $(x; y)$ thỏa mãn đề bài.
	}
\end{ex}
\begin{ex}
	[Đại học Hồng Đức – 2022]%Câu 57.
	Cho $x$ là số nguyên dương và $y$ là số thực. Có tất cả bao nhiêu cặp số $(x;y)$ thỏa mãn $\ln (1+x+2y)=2y+3x-10$? 
	\choice
	{$10$}
	{Vô số}
	{$11$}
	{\True $9$}
	\loigiai{
		Điều kiện: $1+x+2y>0\Leftrightarrow y >-\dfrac{x+1}{2}$.\\
		Ta luôn chứng minh được $\mathrm{e}^x\geq x+1,\forall x\in\mathbb{R}$.\\
		Xét hàm số $y=g(x)=\mathrm{e}^x-x-1\Rightarrow g'(x)=\mathrm{e}^x-1=0\Leftrightarrow x=0$.\\
		Bảng biến thiên: 
		{\color{red} HINH O DAY}\\
		Suy ra $g(x)\geq 0\forall x\in\mathbb{R}\Leftrightarrow\mathrm{e}^x\geq x+1\forall x\in\mathbb{R}$.\\
		Ta có $\ln (1+x+2y)=2y+3x-10\Leftrightarrow 1+x+2y=\mathrm{e}^{2y+3x-10}\geq (2y+3x-10)+1\Leftrightarrow x\leq 5$.\\
		Do $x\in\mathbb{N}^*$, nên $x\in\{1;2;3;4;5\}$.\\
		Lại có: $\ln (1+x+2y)=2y+3x-10\Leftrightarrow\ln (1+x+2y)-2y-3x+10=0\Leftrightarrow f(y)=0$.\\
		Xét hàm số $f(y)=\ln (1+x+2y)-2y-3x+10$ trên khoảng $\left(-\dfrac{x+1}{2};+\infty\right)$.\\
		Suy ra $f'(y)=\dfrac{2}{1+x+2y}-2;f'(y)=0\Leftrightarrow y=-\dfrac{x}{2}\in\left(-\dfrac{x+1}{2};+\infty\right)$.\\
		Bảng biến thiên của hàm số $f(y)=\ln (1+x+2y)-2y-3x+10$. 
		{\color{red} HINH O DAY}\\
		Dựa vào bảng biến thiên ta nhận thấy:\\
		Với một giá trị $x\in\{1;2;3;4\}$, phương trình $\ln (1+x+2y)-2y-3x+10=0$ theo ẩn $y$ có 2 nghiệm phân biệt. Với $x=5$ phương trình $\ln (1+x+2y)-2y-3x+10$ theo ẩn $y$ có 1 nghiệm.\\
		Vậy có 9 nghiệm $(x;y)$ thỏa mãn bài toán.
	}
\end{ex}
\begin{ex}
	[THPT Hương Sơn - Hà Tĩnh - 2022]%Câu 58.
	Trong mặt phẳng với hệ tọa độ $Oxy$, điểm $M(x; y)$ biểu diễn nghiệm của phương trình $\log_3(9x+18)+x=y+3^y$. Có bao nhiêu điểm $M$ có tọa độ nguyên thuộc hình tròn tâm $O$ bán kính $R=7$?
	\choice
	{$7$}
	{\True $2$}
	{$3$}
	{$49$}
	\loigiai{
		Điều kiện: $x >-2$.\\
		Ta có $\log_3(9x+18)+x=y+3^y\Leftrightarrow\log_3(x+2)+x+2=y+3^y$.\\
		Đặt $t=\log_3(x+2)\Leftrightarrow x+2=3^t$ ta được $3^t+t=3^y+y$.\\
		Xét hàm số $f(u)=3^u+u$ thì $f(u)$ là hàm đồng biến, vậy $y=t=\log_3(x+2)$.\\
		Do $M$ có tọa độ nguyên và nằm trên hình tròn tâm $O$ bán kính bằng $R=7$ nên: $x^2+y^2\leq 49$.\\
		$M$ có tọa độ nguyên nên $\heva{&-2<x\leq 7\\&x\in\mathbb{Z}\\&\log_3(x+2)\in\mathbb{Z}.}$ \\
		Trường hợp 1: $x=-1\Rightarrow y=0$ (thỏa mãn).\\
		Trường hợp 2: $x=1\Rightarrow y=1$ (thỏa mãn).\\
		Trường hợp 3: $x=7\Rightarrow y=2$ (loại).\\
		Vậy có 2 điểm thỏa mãn yêu cầu là $(-1;0),(1;1)$.
	}
\end{ex}
\begin{ex}
	[THPT Lê Thánh Tông - HCM - 2022]%Câu 59.
	Có tất cả bao nhiêu số $b$ nguyên dương sao cho tồn tại đúng hai số thực $a$ thỏa mãn đẳng thức $b.2^{a^2-6a-1}+b^2\cdot 2^{2a^2-12a-1}-3=7\log_2\left(a^2-6a+\log_2b\right)$?
	\choice
	{$1024$}
	{\True $1023$}
	{$2047$}
	{$2048$}
	\loigiai{
		Đặt $\heva{&x=a^2-6a\\&y=\log_2b\geq 0\Rightarrow b=2^y\\&x+y>0.}$ \\
		Ta được $2^y\cdot 2^{x-1}+2^{2y}\cdot 2^{2x-1}-3=7\log_2(x+y)\Leftrightarrow 2^{x+y}+2^{2(x+y)}-14\log_2(x+y)-6=0$.\\
		Xét hàm số $f(t)=2^{2t}+2^t-14\log_2t-6, t>0$.\\
		$f'(t)=4^t\ln 4+2^t\ln 2-\dfrac{14}{t\cdot\ln 2}\Rightarrow f''(t)=4^t\ln^24+2^t\ln^22+\dfrac{14}{t^2\ln 2}>0\forall t>0$ \\
		$ \Rightarrow f(t)=0 $ có nhiều nhất hai nghiệm dương. Ta thấy $f(1)=f(2)=0$.\\
		Do vậy, ta được $f(t)=0\Leftrightarrow\hoac{&t=1\\&t=2}\Rightarrow\hoac{&x+y=1\\&x+y=2}\Rightarrow\hoac{&\log_2b=1+6a-a^2\\&\log_2b=2+6a-a^2}$. \\
		{\color{red} HINH O DAY}\\
		Từ đồ thị, ta có: $10<\log_2b<11\Leftrightarrow 2^{10}<b<2^{11}$.\\
		Vậy có $1023$ số nguyên dương $b$ thỏa mãn.
	}
\end{ex}
\begin{ex}
	[Sở Bắc Giang 2022]%Câu 60.
	Có bao nhiêu số nguyên $x$ sao cho tồn tại số thực $y$ thỏa mãn $2\log_3(x+y+1)=\log_2\left(x^2+2x+2y^2+1\right)$?
	\choice
	{$4$}
	{\True $2$}
	{$3$}
	{$1$}
	\loigiai{
		Đặt $X=x+1$. Khi đó, ta có $2\log_3(X+y)=\log_2\left(X^2+2y^2\right)\Leftrightarrow\log_3(X+y)=\log_4\left(X^2+2y^2\right)$.\\
		Đặt $\log_3(X+y)=\log_4\left(X^2+2y^2\right)=t\Rightarrow\heva{&X+y=3^t\\&X^2+2y^2=4^t}$ \\
		$ \Rightarrow 4^t=X^2+\dfrac{y^2}{\dfrac{1}{2}}\geq\dfrac{(X+y)^2}{1+\dfrac{1}{2}}=\dfrac{3^{2t}}{\dfrac{3}{2}}=\dfrac{2}{3}\cdot 9^t\Rightarrow 3\cdot 4^t\geq 2\cdot 9^t\Rightarrow\left(\dfrac{4}{9}\right)^t\geq\dfrac{2}{3}\Rightarrow t\leq\dfrac{1}{2} $..\\
		Suy ra $\heva{&0<X+y\leq\sqrt{3}\\&0<X^2+2y^2\leq 2.}$ \\
		Ta có: $X^2+2y^2\leq 2\Rightarrow 0\leq X^2\leq 2\Rightarrow-\sqrt{2}\leq X\leq\sqrt{2}\Rightarrow X\in\{-1;0;1\}$ do $X$ nguyên.\\
		Với $X=0$, ta có\\
		$\heva{&y=3^t\\&2y^2=4^t}\Rightarrow 2\cdot 9^t=4^t\Rightarrow\left(\dfrac{4}{9}\right)^t=2\Leftrightarrow t=\log_{\frac{4}{9}}2\Rightarrow y=3^{\log_{\frac{4}{9}}2}$.\\
		Với $X=1$, ta có $\heva{&y=3^t-1\\&2y^2=4^t-1}\Rightarrow 2\cdot\left(3^t-1\right)^2=4^t-1\cdot (^*)$.\\
		Ta thấy $t=0$ là nghiệm của $(*)\Rightarrow$ Phương trình đã cho có nghiệm $y=0$.\\
		Với $X=-1$, ta có $\heva{&y=3^t+1\\&2y^2=4^t-1.}$ \\
		Vì $y=3^t+1\Rightarrow y>1$.\\
		Mặt khác, ta có:\\
		$X^2+2y^2\leq 2\Rightarrow 0\leq 2y^2\leq 2\Rightarrow 2y^2\leq 2\Rightarrow y^2\leq 1\Rightarrow|y|\leq 1$.\\
		Do vậy $y>1$ là không thỏa mãn nên $X=-1$ không thỏa mãn Vậy $X\in\{0;1\}$ hay $x\in\{-1;0\}$ thì tồn tại số thực $y$ thỏa mãn $2\log_3(x+y+1)=\log_2\left(x^2+2x+2y^2+1\right)$.
	}
\end{ex}
\begin{ex}
	[Sở Hà Tĩnh 2022]%Câu 61.
	Có bao nhiêu cặp số $(x;y)$ (trong đó $x,y$ nguyên dương thuộc đoạn $[0;2022]$ thỏa mãn điều kiện $2^x-\log_2\left(y^2+615\right)=y^2-x+615$ 
	\choice
	{\True $1$}
	{$3$}
	{$4$}
	{$2$}
	\loigiai{
		$2^x+x=\left(y^2+615\right)+\log_2\left(y^2+615\right)$.\\
		Đặt $t=\log_2\left(y^2+615\right)\Leftrightarrow y^2+615=2^t$. Khi đó ta có $2^x+x=2^t+t$.\\
		Xét hàm số $f(x)=2^x+x,x\in\mathbb{R}$, ta có $f'(x)=2^x\ln 2+1>0,\forall x\in\mathbb{R}$ nên hàm số $f(x)$ đồng biến trên $\mathbb{R}$.\\
		Do vậy $2^x+x=2^t+t\Leftrightarrow x=t\Leftrightarrow x=\log_2\left(y^2+615\right)\Leftrightarrow y^2=2^x-615$.\\
		Theo bài ra $0\leq y\leq 2022\Leftrightarrow 0\leq y^2\leq 2022^2\Leftrightarrow 0\leq 2^x-615\leq 2022^2\Leftrightarrow\log_2615<x<\log_2\left({2022}^2+615\right)$ \\
		$ \Leftrightarrow 9,3\leq x\leq 21,9 $.\\
		Khi $x$ chạy từ $10$ tới $21$, ta thấy $x=12\Rightarrow y=59$. Vậy có $1$ cặp số $(x;y)$ thỏa mãn yêu cầu bài toán.
	}
\end{ex}
\begin{ex}
	[Sở Vĩnh Phúc 2022]%Câu 62.
	Có bao nhiêu cặp số nguyên dương $(x,y)$ thoả mãn $3^{4x^2-1}\log\left(4x^2+4x+2\right)=3^{y-2x-4}\log(2x+y-1)$ đồng thời $x,y\leq 2021$ 
	\choice
	{$15$}
	{$28$}
	{\True $22$}
	{$35$}
	\loigiai{
		$3^{4x^2-1}\log\left(4x^2+4x+2\right)=3^{y-2x-4}\log(2x+y-1)$ \\
		$ \Leftrightarrow 3^{4x^2+4x+2}\log\left(4x^2+4x+2\right)=3^{2x+y-1}\log(2x+y-1)(1) $.\\
		Xét $f(t)=3^t\cdot\log t$ ($t\geq 2$).\\
		$f'(t)=3^t\ln 3\cdot\log t+\dfrac{3^t}{t\cdot\ln 10}=3^t\left(\ln 3\cdot\log t+\dfrac{1}{t\cdot\ln 10}\right)>0,\forall t\geq 2$.\\
		Do đó $(1)\Leftrightarrow 4x^2+4x+2=2x+y-1\Leftrightarrow y=4x^2+2x+3$.\\
		$y\leq 2021\Leftrightarrow 4x^2+2x+3\leq 2021\Leftrightarrow\dfrac{-1-3\sqrt{897}}{4}\leq x\leq\dfrac{-1+3\sqrt{897}}{4}$.\\
		Vì $x\in\mathbb{Z}^+\cap(0;2022)$ nên $x\in\left\{1;2;3;\ldots;21;22\right\}$.\\
		Vì mỗi giá trị của $x$ có đúng một giá trị của $y$ nên có 22 cặp số nguyên dương $(x,y)$ thoả mãn.
	}
\end{ex}
\begin{ex}
	[THPT Lương Tài 2 - Bắc Ninh - 2022]%Câu 63.
	Có tất cả bao nhiêu số nguyên dương $y$ sao cho tồn tại số thực $x\in(1;8)$ thỏa mãn:\\
	$(x-1)\left(2\mathrm{e}^x-y^2\right)=y\left(\mathrm{e}^x-x^2\right)$?
	\choice
	{$11$}
	{$14$}
	{$12$}
	{\True $13$}
	\loigiai{
		Xét $f(x)=(x-1)\left(2\mathrm{e}^x-y^2\right)-y\left(\mathrm{e}^x-x^2\right)$ trên $(1;8)$ với $y$ là tham số.\\
		Ta có\\
		$f'(x)=2x\mathrm{e}^x-y\mathrm{e}^x-y^2+2yx=\left(\mathrm{e}^x+y\right)(2x-y)=0\Rightarrow x=\dfrac{y}{2}$.\\
		Ta thấy $f(1)=-y(e-1)<0$ do $y$ nguyên dương; $f(8)=7\left(2\mathrm{e}^8-y^2\right)-y\left(\mathrm{e}^8-64\right)=-7y^2-\left(\mathrm{e}^8-64\right)y+14\mathrm{e}^8$.\\
		Trường hợp 1: Khi $\dfrac{y}{2}\leq 1\Leftrightarrow y\leq 2\Rightarrow f'(x)>0,\forall x\in(1;8)$. Lập bảng biến thiên cho $f(x)$, từ yêu cầu bài toán $\Rightarrow f(8)>0\Rightarrow y<13,85\Rightarrow y\in\{1;2\}$.\\
		Trường hợp 2: Khi $\dfrac{y}{2}\geq 8\Leftrightarrow y\geq 16\Rightarrow f'(x)<0,\forall x\in(1;8)\Rightarrow f(8)<f(1)<0$ suy ra phương trình vô nghiệm trên $(1;8)$.\\
		Trường hợp 3: Khi $1<\dfrac{y}{2}<8\Leftrightarrow 2<y<16\Rightarrow x_{CT}=\dfrac{y}{2}$. Lập bảng biến thiên cho $f(x)$, từ yêu cầu bài toán $\Rightarrow f(8)>0\Rightarrow y<13,85\Rightarrow y\in\left\{3;4;5;\ldots 13\right\}$.\\
		Như vậy có tất cả 13 giá trị $y$ thỏa mãn.
	}
\end{ex}
\begin{ex}
	[Chuyên Hà Tĩnh 2022]%Câu 64.
	Có bao nhiêu cặp số nguyên dương $(x;y)$ thỏa mãn: $\log_7\dfrac{2x^3-3x^2+1}{6xy+1+2x+3y}=14x+3y-7\left(x^2+1\right)$ đồng thời $1<x<2022$?
	\choice
	{\True $1347$}
	{$1348$}
	{$674$}
	{$673$}
	\loigiai{
		Ta có $\log_7\dfrac{2x^3-3x^2+1}{6xy+1+2x+3y}=14x+3y-7\left(x^2+1\right)$ \\
		$ \Leftrightarrow\log_7\dfrac{(x-1)^2(2x+1)}{(3y+1)(2x+1)}=14x+3y-7\left(x^2+1\right) $ \\
		$ \Leftrightarrow\log_7(x-1)^2-\log_7(3y+1)=3y-7(x-1)^2 $ \\
		$ \Leftrightarrow\log_7(x-1)^2+7(x-1)^2=\log_7(3y+1)+3y $ \\
		$ \Leftrightarrow\log_77(x-1)^2+7(x-1)^2=\log_7(3y+1)+3y+1 $ (*).\\
		Xét hàm số $f(t)=\log_7t+t$. Ta có $f'(t)=\dfrac{1}{t\ln 7}+1>0$, $\forall t>0$ nên $f(t)$ đồng biến trên khoảng $(0;+\infty)$.\\
		Do đó, (*) $\Leftrightarrow 7(x-1)^2=3y+1\Rightarrow y=\dfrac{7(x-1)^2-1}{3}$.\\
		$y$ nguyên dương khi $7(x-1)^2-1\vdots 3\Rightarrow x$ chia hết cho $3$ hoặc $x$ chia $3$ dư 2.
		Do đó $x\in\left\{2;3;\ldots;2021\right\}\setminus\left\{4;7;10;\ldots\cdot;2020\right\}$. Vậy có $2020-673 =1347$ cặp $(x;y)$ thỏa mãn.}
\end{ex}
\begin{ex}
	[Chuyên Lê Hồng Phong - Nam Định 2022]%Câu 65.
	Có bao nhiêu cặp số nguyên $(x;y)$ thỏa mãn $0<y\leq 2022$ và $3^x+3x-6=9y+\log_3y^3$. 
	\choice
	{$2022$}
	{$9$}
	{$8$}
	{\True $7$}
	\loigiai{
		Đặt $t=\log_3y\Rightarrow y=3^t$.\\
		Phương trình trở thành.\\
		$\begin{aligned}&3^x+3x-6={9\cdot 3}^t+3t\\&\Leftrightarrow 3^x+3x=3^{t+2}+3(t+2).\end{aligned}$\\
		Xét hàm số $f(a)=3^a+3a\Rightarrow f'(a)=3^a\ln 3+3>0,\forall a\in\mathbb{R}$, tức là $f(a)$ đồng biến trên $\mathbb{R}$.\\
		Suy ra $x=t+2\Rightarrow x=\log_3y+2$.\\
		Để $x\in\mathbb{Z}$ thì $\log_3y\in\mathbb{Z}$.\\
		Mà $0<y\leq 2022\Rightarrow y\in\left\{1;3;3^2;3^3;\ldots\cdot;3^6\right\}$.\\
		Suy ra có 7 cặp số nguyên $(x;y)$ thỏa mãn yêu cầu bài toán.
	}
\end{ex}
\begin{ex}
	[Chuyên Thái Bình 2022]%Câu 66.
	Có bao nhiêu bộ $(x;y)$ với $x$, $y$ là hai số tự nhiên thỏa mãn $11+10^x+6^x=(\sqrt{3})^{y!}$?
	\choice
	{$0$}
	{\True $1$}
	{$2$}
	{vô số}
	\loigiai{
		Với $x=0$, khi đó $VT=13$, dễ thấy không tồn tại $y\in\mathbb{N}$ để $(\sqrt{3})^{y!}=13$.\\
		Với $x\geq 1$, khi đó $VT$ là số có chữ số tận cùng là $7$.\\
		Do $(\sqrt{3})^{y!}=11+10^x+6^x>11$ nên $y\geq 3$.\\
		Nếu $y=3$, khi đó $VP=27=VT$ với $x=1$. Vậy $(1;3)$ là một cặp nghiệm.\\
		Nếu $y\geq 4$, khi đó $\dfrac{y!}{2}$ có dạng $4k$, $(k\in \mathbb{N}*)$, suy ra $VP$ có chữ số tận cùng bằng $1$ (mâu thuẫn). Vậy chỉ có đúng 1 cặp nghiệm thỏa mãn.
	}
\end{ex}
\begin{ex}
	[Cụm trường Nam Định 2022]%Câu 67.
	Có bao nhiêu cặp số nguyên $(x;y)$ thỏa mãn đồng thời $1\leq x\leq 2022$ và $384\cdot 128^{x^2-2x}-6\cdot 8^y+6=3y-7x^2+14x$? 
	\choice
	{$674$}
	{\True $1348$}
	{$1346$}
	{$2022$}
	\loigiai{
		Ta có $384\cdot 128^{x^2-2x}-6\cdot 8^y+6=3y-7x^2+14x\Leftrightarrow 3\cdot 2^{7(x-1)^2}+7(x-1)^2=3\cdot 2^{3y+1}+3y+1$.\\
		Xét hàm số $f(t)=3\cdot 2^t+t, t\geq 0$ có $f'(t)=3\cdot 2^t\cdot\ln 2+1>0$ nên hàm số đồng biến trên $[0;+\infty)$.\\
		Do đó: $7(x-1)^2=3y+1\Leftrightarrow 7\left(x^2-2x\right)+6=3y$.\\
		Vì $x,y\in\mathbb{Z}$ nên $\left(x^2-2x\right)\vdots 3\Rightarrow\hoac{&x-2\vdots 3\\&x\vdots 3}\Rightarrow\hoac{&x=3n+2\\&x=3n}, n\in\mathbb{Z.}$ \\
		mà $1\leq x\leq 2022\Rightarrow\hoac{&1\leq 3n+2\leq 2022\\&1\leq 3n\leq 2022}\Rightarrow\hoac{&\dfrac{-1}{3}\leq n\leq\dfrac{2020}{3}\\&\dfrac{1}{3}\leq n\leq 674}\Rightarrow\hoac{&0\leq n\leq 673\\&1\leq n\leq 674.}$ \\
		hay có $1348$ số nguyên $n$. Mỗi giá trị của $n$ cho chúng ta một cặp số nguyên $(x;y)$ thỏa mãn điều kiện của bài toán.\\
		Vậy có $1348$ cặp số nguyên $(x;y)$ thỏa mãn điều kiện của bài toán.
	}
\end{ex}
\begin{ex}
	[Liên trường Quảng Nam 2022]%Câu 68.
	Có bao nhiêu cặp số nguyên $(x;y)$ thoả mãn $x\in[2;4374)$ và $2.3^y-\log_3\left(x+3^{y-1}\right)=3x-y$?
	\choice
	{\True $7$}
	{$8$}
	{$6$}
	{$9$}
	\loigiai{
		Đặt $t=3^{y-1}\Rightarrow 3^y=3t\Rightarrow y=\log_3(3t)$.\\
		Phương trình đã cho tương đương với: $6t-\log_3(x+t)=3x-\log_3(3t)$ \\
		$ \Leftrightarrow 6t+\log_3(3t)=3x+\log_3(x+t) $ \\
		$ \Leftrightarrow 9t+\log_3(3t)=3(x+t)+\log_3(x+t)(*) $.\\
		Xét hàm số: $f(u)=3u+\log_3u$ liên tục trên $(0;+\infty)$, có $f'(u)=3+\dfrac{1}{u\ln 3}>0,\forall u\in(0;+\infty)$.\\
		Suy ra: hàm số $y=f(u)$ đồng biến trên $(0;+\infty) (1)$.\\
		$(*)\Leftrightarrow f(3t)=f(x+t)(2)$.\\
		Từ $(1)$ và $(2)\Leftrightarrow 3t=x+t\Leftrightarrow 2t=x\Leftrightarrow 2\cdot 3^{y-1}=x$.\\
		Ta có: $2\leq x<4374$, suy ra: $2\leq 2\cdot 3^{y-1}<4374\Leftrightarrow\log_32\leq\log_32+y-1<\log_34374$ \\
		$ \Leftrightarrow 1\leq y<1+\log_34374-\log_32\Leftrightarrow 1\leq y<8 $. Vì $y$ nguyên, nên $y\in\left\{1;2;3;4;5;6;7\right\}$.\\
		Dễ thấy, ứng mỗi giá trị của $y$, ta tìm được một giá trị nguyên của $x$.\\
		Vậy có 7 cặp số nguyên $(x;y)$ thoả mãn yêu cầu bài toán.
	}
\end{ex}
\begin{ex}
	[Sở Hưng Yên 2022]%Câu 69.
	Giả sử $(x;y)$ là cặp số nguyên thỏa mãn đồng thời $8\leq x\leq 2022$ và $2^y-\log_2\left(x+2^{y-1}\right)=2x-y$. Tổng các giá trị của $y$ bằng
	\choice
	{$60$}
	{$63$}
	{$2022$}
	{\True $49$}
	\loigiai{
		Ta có $2^y-\log_2\left(x+2^{y-1}\right)=2x-y\Leftrightarrow 2\cdot 2^y+y+1=2x+2^y+\log_2\left(2x+2^y\right)$ \\
		$ \Leftrightarrow 2^{y+1}+\log_22^{y+1}=2x+2^y+\log_2\left(2x+2^y\right)(*) $.\\
		Xét hàm số $f(t)=t+\log_2(t), t>0$ là hàm số đồng biến trên $(0;+\infty)$ nên $(*)\Leftrightarrow 2^{y+1}=2x+2^y\Leftrightarrow x=2^{y-1}$.\\
		Ta có $8\leq x\leq 2022\Rightarrow 8\leq 2^{y-1}\leq 2022\Leftrightarrow 4\leq y\leq\log_22024\Rightarrow 4\leq y\leq 10$.\\
		Suy ra $y\in\left\{4;5;6;7;8;9;10\right\}$.
	}
\end{ex}
\begin{ex}
	[Sở Nam Định 2022]%Câu 70.
	Có bao nhiêu số thực $x,y$ thỏa mãn hệ thức $2^{2|y|-x^2}=\log_{2|y|+1}x$. Hỏi có tất cả bao nhiêu giá trị nguyên của $m\in[-2022;2022]$ để tồn tại duy nhất một số thực $x$ thỏa mãn hệ thức $4y^2=10x^2+mx+1$?
	\choice
	{$2036$}
	{$2033$}
	{$2034$}
	{\True $2035$}
	\loigiai{
		Vì $2^{2|y|-x^2}>0$ nên $\log_{2|y|+1}x>0$. Mà $2|y|+1>1$ nên $\log_{2|y|+1}x>0\Rightarrow x>1$.\\
		$2^{2|y|-x^2}=\log_{2|y|+1}x\Leftrightarrow 2\cdot 2^{2|y|-x^2}=2\log_{2|y|+1}x\Leftrightarrow\dfrac{2^{2|y|+1}}{2^{x^2}}=\dfrac{\log_2(x^2)}{\log_2(2|y|+1)}$ \\
		$ \Leftrightarrow 2^{2|y|+1}\cdot\log_2(2|y|+1)=2^{x^2}\cdot\log_2(x^2)(*) $.\\
		Đặt $u=2|y|+1,v=x^2(u>1,v>1)$. Khi đó $(*)\Leftrightarrow 2^u\cdot\log_2u=2^v\cdot\log_2v$.\\
		Xét hàm đặc trưng $f(t)=2^t\cdot\log_2t(t>1)\Rightarrow f'(t)=2^t\ln 2\cdot\log_2t+\dfrac{2^t}{t\ln 2}>0,\forall t>1$.\\
		Vậy hàm số $f(t)=2^t\cdot\log_2t$ đồng biến trên khoảng $(1;+\infty)$.\\
		Ta có $2^u\cdot\log_2u=2^v\cdot\log_2v\Leftrightarrow f(u)=f(v)\Leftrightarrow u=v\Leftrightarrow 2|y|+1=x^2$.\\
		Khi đó: $4y^2=10x^2+mx+1\Leftrightarrow x^4-2x^2+1=10x^2+mx+1\Leftrightarrow x^3-12x=m\Leftrightarrow g(x)=m$. 
		{\color{red} HINH O DAY}\\
		Phương trình $g(x)=m$ có nghiệm duy nhất $x>1\Leftrightarrow\hoac{&m\geq g(1)\\&m=g(2)}\Leftrightarrow\hoac{&m\geq-11\\&m=-16.}$ \\
		Mà $m\in[-2022;2022]$ nên $m\in\left\{-16;-11;-10;-9;\ldots;2022\right\}$.\\
		Vậy có 2035 giá trị nguyên của $m\in[-2022;2022]$.
	}
\end{ex}
\begin{ex}
	[THPT Đồng Lộc - Hà Tĩnh 2022]%Câu 71.
	Có bao nhiêu giá trị nguyên dương $a$ sao cho ứng với mỗi $a$, tồn tại ít nhất số thực $b$ thỏa mãn $a^{\log_58}+2^{\log_5(5a)}=\left(b+\sqrt{4-b^2}\right)\left(6+2b\sqrt{4-b^2}\right)$?
	\choice
	{\True $11$}
	{$10$}
	{$9$}
	{$2002$}
	\loigiai{
		PT đã cho tương đương với.\\
		$8^{\log_5a}+2\cdot 2^{\log_5a}=\left(b+\sqrt{4-b^2}\right)\left(4+2b\sqrt{4-b^2}+2\right)$ \\
		$ \Leftrightarrow\left(2^{\log_5a}\right)^3+2\cdot 2^{\log_5a}=\left(b+\sqrt{4-b^2}\right)\left[\left(b+\sqrt{4-b^2}\right)^2+2\right] $ \\
		$ \Leftrightarrow\left(2^{\log_5a}\right)^3+2\cdot 2^{\log_5a}=\left(b+\sqrt{4-b^2}\right)^3+2\left(b+\sqrt{4-b^2}\right) $.\\
		Xét hàm số $f(t)=t^3+2t\Rightarrow f'(t)=3t^2+2>0,\forall t$ nên hàm số đồng biến trên $\mathbb{R}$.\\
		Do đó PT đã cho tương đương với: $2^{\log_5a}=b+\sqrt{4-b^2}(1)$.\\
		Xét hàm số $g(b)=b+\sqrt{4-b^2}$, $b\in[-2;2]\Rightarrow g'(b)=1-\dfrac{b}{\sqrt{4-b^2}}\Rightarrow g'(b)=1-\dfrac{b}{\sqrt{4-b^2}}=0\Leftrightarrow b=\sqrt{2}$.\\
		Ta có bảng biến thiên như sau: \\
		{\color{red} HINH O DAY}\\
		Phương trình (1) có nghiệm $b$ khi và chỉ khi $-2\leq 2^{\log_5a}\leq 2\sqrt{2}\Leftrightarrow 2^{\log_5a}\leq 2\sqrt{2}\Leftrightarrow\log_5a\leq\dfrac{3}{2}\Leftrightarrow a\leq 5^{\frac{3}{2}}$.\\
		Mà $a$ là số nguyên dương nên $a\in\{1;2;\ldots;11\}$.
	}
\end{ex}
\begin{ex}
	[Sở Lai Châu 2022]%Câu 72.
	Tìm số giá trị nguyên của tham số thực $m$ để tồn tại các số thực $x,y$ thỏa mãn $\mathrm{e}^{x^2+y^2-m}+\mathrm{e}^{x+y+xy-m}=x^2+y^2+x+y+xy-2m+2$. 
	\choice
	{$6$}
	{\True $9$}
	{$8$}
	{$7$}
	\loigiai{
		Xét hàm số $f(t)=\mathrm{e}^t-t-1,\forall t\in\mathbb{R}\cdot f'(t)=\mathrm{e}^t-1$ và $f'(t)=0\Leftrightarrow t=0$.\\
		Ta thấy $f'(t)$ đổi dấu từ ``$-$'' sang ``$+$'' khi qua $t=0$ nên $f(t)\geq f(0)=0,\forall t\in\mathbb{R}$.\\
		Do đó: $\heva{&\mathrm{e}^{x^2+y^2-m}-\left(x^2+y^2-m\right)-1\geq 0,\forall x,y\in\mathbb{R}\\&\mathrm{e}^{x+y+xy-m}-(x+y+xy-m)-1\geq 0,\forall x,y\in\mathbb{R}}$. Dấu "=" xảy ra khi và chỉ khi $\heva{&x^2+y^2=m\\&x+y+xy=m.}$ \\
		Hay $\mathrm{e}^{x^2+y^2-m}+\mathrm{e}^{x+y+xy-m}=x^2+y^2+x+y+xy-2m+2\Leftrightarrow\heva{&x^2+y^2=m\\&x+y+xy=m.}$ \\
		Đặt $S=x+y,P=x\cdot y$, ta có: $\heva{&S^2-2P=m\\&S+P=m}\Rightarrow S^2-S-3P=0$. Vì $S^2\geq 4P\Rightarrow S\in [0;4]$.\\
		Lấy (1)+2. (2) vế theo vế ta được: $S^2+2S=3m(3)$.\\
		Xét hàm số $f(S)=S^2+2S,S\in [0;4]$, có $f'(S)=2S+2>0,\forall S\in [0;4]$.\\
		Yêu cầu bài toán $\Leftrightarrow (3)$ có nghiệm\\
		$ \Leftrightarrow f(0)\leq 3m\leq f(4)\Leftrightarrow 0\leq m\leq 8 $. Vậy, có 9 giá trị nguyên của tham số $m$ thỏa mãn.
	}
\end{ex}
\begin{ex}
	[THPT Nguyễn Viết Xuân – Vĩnh Phúc – 2022]%Câu 73.
	Có bao nhiêu giá trị nguyên y sao cho tồn tại số thực $x$ thỏa mãn $\log_4\left(\sqrt{x^2+3^y}-x\right)\cdot\log_3\left(\sqrt{x^2+3^y}+x\right)=y^2-7y$? 
	\choice
	{$8$}
	{\True $9$}
	{$11$}
	{$10$}
	\loigiai{
		Nhận xét: $\left(\sqrt{x^2+3^y}-x\right)\left(\sqrt{x^2+3^y}+x\right)=3^y$.\\
		Suy ra $\log_3\left(\sqrt{x^2+3^y}-x\right)+\log_3\left(\sqrt{x^2+3^y}+x\right)=y$ (1).\\
		Phương trình đã cho $\log_43\cdot\log_3\left(\sqrt{x^2+3^y}-x\right)\cdot\log_3\left(\sqrt{x^2+3^y}+x\right)=y^2-7y$.\\
		Suy ra $\log_3\left(\sqrt{x^2+3^y}-x\right)\cdot\log_3\left(\sqrt{x^2+3^y}+x\right)=\left(y^2-7y\right)\cdot\log_34$ (2).\\
		(1), (2) theo yêu cầu bài toán ta cần $y^2-4\cdot\log_34\left(y^2-7y\right)\geq 0$ \\
		$ \Leftrightarrow\left(1-4\log_34\right)\cdot y^2+(28\log_34)y\geq 0\Leftrightarrow 0\leq y\leq\underbrace{\dfrac{28\log_34}{4\log_34-1}}_{\approx 8\cdot 73} $.\\
		Do $y\in\mathbb{Z}\Rightarrow y\in\{0;1;2;\ldots;8\}$.
	}
\end{ex}
\begin{ex}
	[Chuyên Hùng Vương – Gia Lai 2022]%Câu 74.
	Có bao nhiêu cặp số nguyên $(x;y)$ thỏa mãn đồng thời các điều kiện $2\leq x\leq 2022$, $1\leq y\leq 2022$ và $\log_2\sqrt[4]{\dfrac{y+3}{2x+1}}+4^x=2^{y+2}$? 
	\choice
	{\True $1011$}
	{$1010$}
	{$1009$}
	{$1012$}
	\loigiai{
		Ta có $\dfrac{1}{4}\log_2\dfrac{y+3}{2x+1}+4^x=2^{y+2}\Leftrightarrow\log_2\dfrac{y+3}{2x+1}+2^{2x+2}=2^{y+4}\Leftrightarrow\log_2(y+3)-2^{y+4}=\log_2(2x+1)-2^{2x+2}(1)$.\\
		Xét hàm số: $f(t)=\log_2t-2^{t+1}\forall t\in [4;+\infty)$.\\
		$f'(t)=\dfrac{1-t\cdot 2^{t+1}\ln 2\cdot\ln 2}{t\cdot\ln 2}<0\forall t\in [4;+\infty)$.\\
		Suy ra: (1) $\Leftrightarrow y+3=2x+1\Leftrightarrow y=2x-2$.\\
		Do $1\leq y\leq 2022\Rightarrow\dfrac{3}{2}\leq x\leq 1012\Rightarrow x\in\{2;3;\ldots,1012\}$.\\
		Do đó: $(x;y)\in\{(2;2);(3;4);\ldots;(1012;2022)\}$ có 1011 cặp thỏa mãn yêu cầu bài toán.
	}
\end{ex}
\begin{ex}
	[THPT Cò Nòi - Sơn La 2022]%Câu 75.
	Có bao nhiêu số nguyên $a$ sao cho ứng với mỗi $a$, tồn tại ít nhất $7$ số nguyên $b\in(0;10)$ thỏa mãn $\log_5\left(b^2+16\right)+\log_3\left(b\sqrt{13-a}\right)-\log_7(a-3)\geq 5$?
	\choice
	{$9$}
	{$8$}
	{$11$}
	{\True $1$}
	\loigiai{
		Điều kiện: $\heva{&b>0\\&3<a<13.}$ \\
		Ta có $\log_5\left(b^2+16\right)+\log_3\left(b\sqrt{13-a}\right)-\log_7(a-3)\geq 5$ \\
		$\Leftrightarrow\log_5\left(b^2+16\right)+\log_3b+\log_3\sqrt{13-a}-\log_7(a-3)-5\geq 0 $.\\
		Đặt $f(b)=\log_5\left(b^2+16\right)+\log_3b+\log_3\sqrt{13-a}-\log_7(a-3)-5$, điều kiện $b>0$.\\
		Bất phương trình trở thành $f(b)\geq 0$.\\
		$f'(b)=\dfrac{2b}{\left(b^2+16\right)\ln 5}+\dfrac{1}{b\ln 3}$ do $b>0$ nên $f'(b)>0\Rightarrow$ Hàm số $f(b)$ đồng biến trên $(0;10)$ suy ra $f(1)<f(2)<f(3) <\cdots <f(9)$.\\
		Do đó để có ít nhất $7$ giá trị $b$ nguyên thuộc $(0;10)$ thì $f(3)\geq 0$ \\
		$ \Leftrightarrow\log_3\sqrt{13-a}-\log_7(a-3)-2\geq 0(*) $.\\
		Đặt $g(a)=\log_3\sqrt{13-a}-\log_7(a-3)-2, a\in(3;13)$.\\
		Bất phương trình $(*)$ trở thành $g(a)\geq 0$.\\
		$g'(a)=\dfrac{-1}{2(13-a)\ln 3}-\dfrac{1}{(a-3)\ln 7}<0,\forall a\in(3;13)$ nên hàm số $g(a)$ nghịch biến trên $(3;13)$.\\
		Mặt khác $g(4)=0$ bất phương trình $(*)$ trở thành $g(a)\geq g(4)$, $g(a)$ nghịch biến nên $a\leq 4$ mà $a\in(3;13)$, $a$ nguyên nên $a=4$.\\
		Vậy có duy nhất một giá trị nguyên $a=4$ thỏa mãn bài toán.
	}
\end{ex}
\begin{ex}
	[THPT Trần Quốc Tuấn - Quảng Ngãi - 2022]%Câu 76.
	Có bao nhiêu cặp số nguyên $(x; y)$ thỏa mãn $0<y\leq 2022$ và $\log_2\left(\dfrac{3^x-1}{y}\right)=y+1-3^x$?
	\choice
	{$7$}
	{\True $6$}
	{$2022$}
	{$2021$}
	\loigiai{
		Điều kiện: $3^x>1\Leftrightarrow x>0$.\\
		$\log_2\left(\dfrac{3^x-1}{y}\right)=y+1-3^x\Leftrightarrow\log_2\left(3^x-1\right)+3^x-1=\log_2y+y$ (*).\\
		Xét hàm số $f(t)=\log_2t+t$ với $t>0$.\\
		$f'(t)=\dfrac{1}{t\ln 2}+1>0,\forall t>0\Rightarrow f(t)$ đồng biến trên khoảng $(0;+\infty)$ suy ra\\
		$(*)\Leftrightarrow 3^x-1=y$, vì $0<y\leq 2022\Leftrightarrow 0<3^x-1\leq 2022\Leftrightarrow 0<x\leq\log_32023$.\\
		Mà $x$ nguyên nên $x\in\left\{1; 2; 3; 4; 5; 6\right\}$. Vì ứng với mỗi gia trị $x=x_0$ có duy nhất 1 giá trị $y=3^{x_o}-1$ do đó có 6 cặp cặp số nguyên $(x; y)$ thỏa mãn.
	}
\end{ex}
\begin{ex}
	[THPT Ngũ Hành Sơn - Đà Nẵng 2022]%Câu 77.
	Cho $0\leq x\leq 2020$ và $\log_2(2x+2)+x-3y=8^y$. Có bao nhiêu cặp số $(x;y)$ nguyên thỏa mãn các điều kiện trên?
	\choice
	{$1$}
	{$2019$}
	{\True $4$}
	{$2018$}
	\loigiai{
		Đặt $\log_2(2x+2)=u\Rightarrow x=2^{u-1}-1$.\\
		Khi đó phương trình đã cho trở thành: $u+2^{u-1}-1-3y=8^y\Leftrightarrow 2^{u-1}+(u-1)=2^{3y}+3y$.\\
		Xét hàm đặc trưng $f(t)=2^t+t,$ có $f'(t)=2^t\ln 2+1>0,\forall t\Rightarrow f(t)$ là hàm số đồng biến trên $\mathbb{R}$.\\
		Ta có $2^{u-1}+(u-1)=2^{3y}+3y\Leftrightarrow f(u-1)=f(3y)\Leftrightarrow u-1=3y\Leftrightarrow x=2^{3y}-1$.\\
		Theo bài ra: $0\leq 2^{3y}-1\leq 2020\Leftrightarrow 0\leq y\leq\dfrac{\log_22021}{3}\approx 3,7$.\\
		Vậy có 4 cặp $(x;y)$ thỏa mãn là $(0;0),(7;1),(63;2),(511; 3)$.
	}
\end{ex}
\begin{ex}
	[THPT Lương Bắc Bằng - Thanh Hóa 2022]%Câu 78.
	Có bao nhiêu giá trị nguyên dương của tham số $m$ để tồn tại các số thực $x,y$ thoả mãn $\mathrm{e}^{3x+5y-10}-\mathrm{e}^{x+3y-9}=1-2x-2y$ và $\log_5^2(3x+2y+4)-(m+6)\log_5(x+5)+m^2+9=0$?
	\choice
	{\True $4$}
	{$3$}
	{$6$}
	{$5$}
	\loigiai{
		Ta có $\mathrm{e}^{3x+5y-10}-\mathrm{e}^{x+3y-9}=1-2x-2y\Leftrightarrow\mathrm{e}^{3x+5y-10}+3x+5y-10=\mathrm{e}^{x+3y-9}+x+3y-9$ \\
		$ \Leftrightarrow f(3x+5y-10)=f(x+3y-9)(1) $ (với $f(t)=\mathrm{e}^t+t$).\\
		Xét $f(t)=\mathrm{e}^t+t$ \\
		$ \Rightarrow f'(t)=\mathrm{e}^t+1>0\forall t $ \\
		$ \Rightarrow(1)\Leftrightarrow 3x+5y-10=x+3y-9\Leftrightarrow 2x+2y-1=0 $.\\
		Khi đó $\log_5^2(3x+2y+4)-(m+6)\log_5(x+5)+m^2+9=0$ \\
		$ \Leftrightarrow\log_5^2(x+5)-(m+6)\log_5(x+5)+m^2+9=0(*) $.\\
		Đặt $t=\log_5(x+5)$ \\
		$ \Rightarrow(*)\Leftrightarrow t^2-(m+6)t+m^2+9=0(2) $.\\
		Phương trình $(*)$ có nghiệm $\Leftrightarrow$ phương trình $(2)$ có nghiệm\\
		$ \Leftrightarrow(m+6)^2-4\left(m^2+9\right)\geq 0\Leftrightarrow-3m^2+12m\geq 0\Leftrightarrow 0\leq m\leq 4 $.\\
		Vì $m\in\mathbb{Z}^+$ nên $m\in\{1;2;3;4\}$.
	}
\end{ex}
\begin{ex}
	[THPT Nguyễn Cảnh Quân - Nghệ An 2022]%Câu 79.
	Có bao nhiêu số nguyên $y$ sao cho tồn tại $x\in\left(\dfrac{1}{3};3\right)$ thỏa mãn $27^{3x^2+xy}=(1+xy)27^{9x}$? 
	\choice
	{$10$}
	{$12$}
	{\True $11$}
	{$9$}
	\loigiai{
		Viết lại phương trình thành $27^{3x^2-9x}=(1+xy)27^{-xy}$.\\
		Ta có $(1)\Leftrightarrow 3x^2-9x=\log_{27}(1+xy)-xy\Leftrightarrow 3x^2-9x-1=\log_{27}t-t$, với $t=1+xy>0$.\\
		Xét hàm số $f(x)=3x^2-9x-1$. Ta có $-\dfrac{31}{4}\leq f(x) <-1\forall x\in\left(\dfrac{1}{3};3\right)$.\\
		Xét hàm số $g(t)=\log_{27}t-t, t>0$. $g'(t)=\dfrac{1}{t\ln 27}-1$; $g'(t)=0\Leftrightarrow t=\dfrac{1}{\ln 27}$ 
		{\color{red} HINH O DAY}\\
		Ta có $-\dfrac{31}{4}\leq f(x) <-1$, $\forall x\in\left(\dfrac{1}{3};3\right)$ suy ra $-\dfrac{31}{4}\leq g(t) <-1\Leftrightarrow\hoac{&t\in\left(\approx 8,{07\cdot 10}^{-12};\approx 0,04\right)\\&t\in\left(1;\approx 8,4\right).}$ \\
		hay $\hoac{&\approx 8,{07\cdot 10}^{-12}<1+xy<\approx 0,04\\&1<1+xy<\approx 8,4}\Leftrightarrow\hoac{&\approx\dfrac{-1+8,{07\cdot 10}^{-12}}{x}<y<\approx\dfrac{-1+0,04}{x}\\&0<y<\approx\dfrac{7,4}{x}}$ \\
		$ \Rightarrow\hoac{&-3<y <-\dfrac{1}{3}\\&0<y\leq 22} $, ($x\in\left(\dfrac{1}{3};3\right)$, $y$ nguyên).\\
		Nhận thấy $y=-2;y=-1$ thỏa mãn đề.\\
		Với $0<y\leq 22$, ta có $(1)\Leftrightarrow 3x^2-9x-1-\log_{27}(1+xy)+(1+xy)= 0$.\\
		Nhập hàm, thay các giá trị nguyên của y, kiểm tra nghiệm $x\in\left(\dfrac{1}{3};3\right)$ dẫn đến chọn $1\leq y\leq 9$.\\
		Vậy $y\in\left\{-2;-1;1;2;\ldots;9\right\}$ nên có $11$ giá trị nguyên của $y$ thỏa mãn đề.
	}
\end{ex}
\begin{ex}
	[Sở Lạng Sơn 2022]%Câu 80.
	Có bao nhiêu cặp số $(x;y)$ nguyên thoả mãn $0\leq x\leq 2022$ và $\log_2(2x+2)+x-3y=8^y$?
	\choice
	{$2018$}
	{$1$}
	{\True $4$}
	{$2019$}
	\loigiai{
		Đặt $\log_2(2x+2)=t\Rightarrow 2x+2=2^t$.\\
		Khi đó ta được phương trình $t+\dfrac{2^t-2}{2}-3y=8^y\Leftrightarrow t+2^{t-1}-1-3y=2^{3y}\Leftrightarrow 2^{t-1}+t-1=2^{3y}+3y(*)$.\\
		Xét hàm sô $f(a)=2^a+a$ trên $\mathbb{R}$ có $f(a)=2^a\ln 2+1>0\forall a\in\mathbb{R}$.\\
		Khi đó hàm số đồng biến trên $\mathbb{R}$ nên phương trình $(*)\Leftrightarrow f(t-1)=f(3y)\Leftrightarrow t-1=3y\Leftrightarrow t=3y+1$.\\
		hay $\log_2(2x+2)=3y+1$.\\
		Do $0\leq x\leq 2022$ nên $\log_22-1\leq 3y\leq\log_24046-1\Leftrightarrow\dfrac{\log_22-1}{3}\leq y\leq\dfrac{\log_24046-1}{3}$.\\
		Và $y\in\mathbb{Z}\Rightarrow y\in\{0;1;2;3\}$.
	}
\end{ex}
\begin{ex}
	[Sở Thái Bình 2022]%Câu 81.
	Có tất cả bao nhiêu cặp số nguyên $x$ và $y$ sao cho đẳng thức sau thỏa mãn $\log_{2021}\left(4^x-2^{x+1}+2022\right)^{y^2+101}=20y+1$. 
	\choice
	{$2$}
	{$3$}
	{\True $1$}
	{$0$}
	\loigiai{
		Nhận xét: $y^2+1>0,\forall y\in\mathbb{R}; 4^x-2^{x+1}+2022=\left(2^x-1\right)^2+2021\geq 2021,\forall x\in\mathbb{R}$.\\
		Ta có $\log_{2021}\left(4^x-2^{x+1}+2022\right)^{y^2+101}=20y+1\Leftrightarrow\log_{2021}\left[\left(2^x-1\right)^2+2021\right]=\dfrac{20y+1}{y^2+101}$.\\
		Mà $\log_{2021}\left[\left(2^x-1\right)^2+2021\right]\geq 1\Leftrightarrow\dfrac{20y+1}{y^2+101}\geq 1\Leftrightarrow(y-10)^2\leq 0\Leftrightarrow y=10\Rightarrow 2^x-1=0\Leftrightarrow x=0$.\\
		Vậy có $1$ cặp số $(x;y)=(0;10)$ thỏa mãn.
	}
\end{ex}
\begin{dang}
 	{Một số bài toán hàm ẩn liên quan đến logarit}
\end{dang}
\begin{ex}
	[THPT Thạch Thành 1 - Thanh Hóa - 2021]%Câu 1.
	Cho hàm số $y=f(x)$ và $f(x)>0,\forall x\in\mathbb{R}$. Biết hàm số $y=f'(x)$ có bảng biến thiên như hình vẽ và $f\left(\dfrac{1}{2}\right)=\dfrac{137}{16}$. 
	{\color{red} HINH O DAY}\\
	Có bao nhiêu giá trị nguyên của $m\in[-2020; 2020]$ để hàm số $g(x)=\mathrm{e}^{-x^2+4mx-5}\cdot f(x)$ đồng biến trên $\left(-1;\dfrac{1}{2}\right)$. 
	\choice
	{$2019$}
	{\True $2020$}
	{$4040$}
	{$4041$}
	\loigiai{
		Ta có $g'(x)=(-2x+4m)\cdot\mathrm{e}^{-x^2+4mx-5}\cdot f(x)+\mathrm{e}^{-x^2+4mx-5}\cdot f'(x)$ \\
		$ \Leftrightarrow g'(x)=\left[(-2x+4m)\cdot f(x)+f'(x)\right]\cdot\mathrm{e}^{-x^2+4mx-5} $.\\
		Yêu cầu bài toán $\Leftrightarrow g'(x)\geq 0,\forall x\in\left(-1;\dfrac{1}{2}\right)$ và $g'(x)=0$ chỉ xảy ra tại một số hữu hạn điểm thuộc $\left(-1;\dfrac{1}{2}\right)$ \\
		$ \Leftrightarrow(-2x+4m)\cdot f(x)+f'(x)\geq 0,\forall x\in\left(-1;\dfrac{1}{2}\right) $ (vì $\mathrm{e}^{-x^2+4mx-5}>0$)\\
		$ \Leftrightarrow-2x+4m\geq-\dfrac{f'(x)}{f(x)},\forall x\in\left(-1;\dfrac{1}{2}\right) $, (vì $f(x)>0,\forall x\in\mathbb{R}$)\\
		$ \Leftrightarrow 4m\geq 2x-\dfrac{f'(x)}{f(x)},\forall x\in\left(-1;\dfrac{1}{2}\right) (*) $.\\
		Xét $h(x)=2x-\dfrac{f'(x)}{f(x)},\forall x\in\left(-1;\dfrac{1}{2}\right)$. Ta có $h'(x)=2-\dfrac{f''(x)\cdot f(x)-[f'(x)]^2}{f^2(x)}$.\\
		Mà $\heva{&f''(x)<0\\&f(x)>0},\forall x\in\left(-1;\dfrac{1}{2}\right)\Rightarrow\dfrac{f''(x)\cdot f(x)-[f'(x)]^2}{f^2(x)}<0,\forall x\in\left(-1;\dfrac{1}{2}\right)$.\\
		Từ đó suy ra $h'(x)>0,\forall x\in\left(-1;\dfrac{1}{2}\right)$. Vậy hàm số $h(x)$ đồng biến trên $\left(-1;\dfrac{1}{2}\right)$.\\
		Bảng biến thiên\\
		{\color{red} HINH O DAY}\\
		Vậy điều kiện $(*)\Leftrightarrow 4m\geq h\left(\dfrac{1}{2}\right)\Leftrightarrow 4m\geq 2\cdot\left(\dfrac{1}{2}\right)-\dfrac{f'\left(\dfrac{1}{2}\right)}{f\left(\dfrac{1}{2}\right)}\Leftrightarrow 4m\geq\dfrac{225}{137}\Leftrightarrow m\geq\dfrac{225}{548}$.\\
		Lại có $\heva{&m\in\mathbb{Z}\\&m\in[-2020; 2020]}\Rightarrow m\in\left\{1; 2; 3;\ldots; 2020\right\}$.\\
		Vậy có $2020$ giá trị nguyên của $m$ thỏa mãn yêu cầu bài toán.
	}
\end{ex}
\begin{ex}
	[Chuyên Lê Thánh Tông - Quảng Nam - 2021]%Câu 2.
	Cho hàm số $y=f(x)$ có đồ thị $y=f'(x-1)$ như hình vẽ. \\
	{\color{red} HINH O DAY}\\
	Khi đó hàm số $y=\mathrm{e}^{f(x)+2x}$ đạt cực đại tại điểm $x_0$. Mệnh đề nào dưới đây là đúng?
	\choice
	{$x_0\in(-1;0)$}
	{\True $x_0\in(-4;-2)$}
	{$x_0\in(0;1)$}
	{$x_0\in(-2;-1)$}
	\loigiai{
		$y=\mathrm{e}^{f(x)+2x}\Rightarrow y'=\left(f'(x)+2\right)\mathrm{e}^{f(x)+2x}$.\\
		Do đó.\\
		$y'=0\Leftrightarrow f'(x)+2=0\Leftrightarrow f'(x)=-2$.\\
		Dựa vào đồ thị ta suy ra\\
		$f'(x)=-2\Leftrightarrow f'(x+1-1)=-2\Leftrightarrow\hoac{&x+1=0\\&x+1=\alpha <-1}\Leftrightarrow\hoac{&x=-1\\&x=\alpha-1 <-2}$.
	}
\end{ex}
\begin{ex}
	[Sở Thái Nguyên - 2021]%Câu 3.
	Cho hàm số $y=f(x)$ thỏa mãn $f(x)>0,\forall x\in(1;4)$ và $f'(4)=0$. Hàm số $f'(x)$ có bảng biến thiên như sau: 
	{\color{red} HINH O DAY}\\
	Có bao nhiêu giá trị nguyên của tham số $m$ thuộc đoạn $[-2019;2019]$ để hàm số $y=\mathrm{e}^{-x^2+mx+1}f(x)$ đồng biến trên khoảng $(1;4)$?
	\choice
	{$2010$}
	{$2008$}
	{\True $2012$}
	{$2007$}
	\loigiai{
		Ta có $y=\mathrm{e}^{-x^2+mx+1}f(x)\Rightarrow y'=(-2x+m)\mathrm{e}^{-x^2+mx+1}f(x)+\mathrm{e}^{-x^2+mx+1}f'(x)$.\\
		Ycbt $\Leftrightarrow(-2x+m)\mathrm{e}^{-x^2+mx+1}f(x)+\mathrm{e}^{-x^2+mx+1}f'(x)\geq 0$ với $x\in(1;4)$ \\
		$ \Leftrightarrow(-2x+m)f(x)+f'(x)\geq 0 $ với $x\in(1;4)$ \\
		$ \Leftrightarrow m\geq 2x-\dfrac{f'(x)}{f(x)} $ với $x\in(1;4)$.\\
		Xét hàm số $g(x)=2x-\dfrac{f'(x)}{f(x)}$ với $x\in(1;4)$:\\
		$\Rightarrow g'(x)=2-\dfrac{f''(x)\cdot f(x)-[f'(x)]^2}{f^2(x)}>0$ với $x\in(1;4)$.\\
		(Vì với $x\in(1;4)$ thì $f''(x)<0; f'(x)>0; f(x)>0$).\\
		Suy ra, với $x\in(1;4)$ thì $g(x)<g(4)=2\cdot 4-\dfrac{f'(4)}{f(4)}=8$.\\
		Suy ra $m\geq 8$.\\
		Vậy có $2019-8+1=2012$ giá trị $m$ thỏa mãn bài ra.
	}
\end{ex}
\begin{ex}
	[THPT Cẩm Bình - Hà Tĩnh - 2021]%Câu 4.
	Cho hàm số $y=f(x)$ có đồ thị hàm số $y=f'(x)$ như hình vẽ\\
	{\color{red} HINH O DAY}\\
	Hàm số $y=g(x)=f(\mathrm{e}^x-2)-2021$ nghịch biến trên khoảng nào dưới đây?
	\choice
	{\True $\left(-1;\dfrac{3}{2}\right)$}
	{$\left(\dfrac{3}{2}; 2\right)$}
	{$(-1; 2)$}
	{$(0;+\infty)$}
	\loigiai{
		{\color{red} HINH O DAY}\\
		Ta có $g'(x)=\mathrm{e}^xf'\left(\mathrm{e}^x-2\right)$.\\
		$g'(x)\leq 0\Leftrightarrow f'\left(\mathrm{e}^x-2\right)\leq 0\Leftrightarrow\mathrm{e}^x-2\leq 3\Leftrightarrow x\leq\ln 5$.\\
		Vậy hàm số $y=g(x)=f(\mathrm{e}^x-2)-2021$ nghịch biến trên khoảng $(-\infty;\ln 5)$.\\
		Do $\left(-1;\dfrac{3}{2}\right)\subset(-\infty;\ln 5)$ nên hàm số đã cho nghịch biến trên khoảng $\left(-1;\dfrac{3}{2}\right)$.
	}
\end{ex}
\begin{ex}
	[Chuyên Nguyễn Trãi - Hải Dương - 2021]%Câu 5.
	Cho hàm số $f(x)$. Hàm số $y=f'(x)$ có bảng biến thiên như sau:\\ 
	{\color{red} HINH O DAY}\\
	Bất phương trình $f(x)\leq\mathrm{e}^{x^2}+m$ đúng với mọi $x\in(-1;1)$ khi và chỉ khi
	\choice
	{$m>f(-1)-e$}
	{\True $m\geq f(0)-1$}
	{$m>f(0)-1$}
	{$m\geq f(-1)-e$}
	\loigiai{
		$f(x)\leq\mathrm{e}^{x^2}+m,\forall x\in(-1;1)\Leftrightarrow f(x)-\mathrm{e}^{x^2}\leq m,\forall x\in(-1;1)$.\\
		Xét $g(x)=f(x)-\mathrm{e}^{x^2}$ trên $(-1;1)$.\\
		Lập bảng biến thiên hàm số $y=f(x)$ trên $[-1;1]$\\ 
		{\color{red} HINH O DAY}\\
		Ta có $\max\limits_{(-1;1)} f(x)=f(0)$.\\
		Khi $x\in(-1;1)\Rightarrow x^2\in[0;1)\Rightarrow\mathrm{e}^{x^2}\in[1;e)\Rightarrow\max\limits_{(-1;1)}\left(-\mathrm{e}^{x^2}\right)=-1$.\\
		Suy ra $\max\limits_{(-1;1)} g(x)=g(0)=f(0)-1$.\\
		Vậy $m\geq f(x)-\mathrm{e}^{x^2},\forall x\in(-1;1)\Leftrightarrow m\geq f(0)-1$.
	}
\end{ex}
\begin{ex}
	[THPT Hậu Lộc 4 - Thanh Hóa - 2021]%Câu 6.
	Cho hàm số $y=f(x)$ có bảng biến thiên như sau\\
	{\color{red} HINH O DAY}\\
	Số nghiệm thuộc đoạn $[-2020; 1]$ của phương trình $f(ln x)=4$ là
	\choice
	{$2020$}
	{$2021$}
	{$4$}
	{\True $3$}
	\loigiai{
		Đặt $t=\ln x, x\in(0,1]\Rightarrow t\in(-\infty;0]$ và một nghiệm $t$ thì cho một nghiệm $x$.\\
		Phương trình tương đương $f(t)=4\Leftrightarrow\hoac{&t=a (a <-2)\\&t=b (-2<b <-1)\\&t=c (-1<c<0)\\&t=d (d>2)(L).}$ \\
		Vậy phương trình có $3$ nghiệm thuộc đoạn $[-2020; 1]$.
	}
\end{ex}
\begin{ex}
	[THPT Nguyễn Trung Thiên - Hà Tĩnh - 2021]%Câu 7.
	Cho hàm số $y=f(x)$ thỏa mãn $2020^{f(x)}=x+\sqrt{x^2+2020},\forall x\in\mathbb{R}$. Có bao nhiêu số nguyên $m$ thỏa mãn $f(\log m)<f\left(\log_m2020\right)$?
	\choice
	{$66$}
	{$63$}
	{$65$}
	{\True $64$}
	\loigiai{
		Vì $x+\sqrt{x^2+2020}>x+|x|\geq 0\Rightarrow x+\sqrt{x^2+2020}>0,\forall x\in\mathbb{R}$.\\
		Từ giả thiết $2020^{f(x)}=x+\sqrt{x^2+2020}\Leftrightarrow f(x)=\log_{2020}\left(x+\sqrt{x^2+2020}\right)$.\\
		Ta có $f'(x)=\dfrac{1+\dfrac{x}{\sqrt{x^2+2020}}}{\left(x+\sqrt{x^2+2020}\right)\ln 2020}=\dfrac{x+\sqrt{x^2+2020}}{\left(x+\sqrt{x^2+2020}\right)\ln 2020\sqrt{x^2+2020}}>0,\forall x\in\mathbb{R}$.\\
		Suy ra hàm số $f(x)$ luôn đồng biến trên $\mathbb{R}$.\\
		Mà với $\heva{&m>0\\&m\neq 1}$ thì $f(\log m)<f\left(\log_m2020\right)\Leftrightarrow\log m<\log_m2020$ \\
		$ \Leftrightarrow\dfrac{\log^2m-\log 2020}{\log m}<0\Leftrightarrow\hoac{&0<\log m<\sqrt{\log 2020}\\&\log m <-\sqrt{\log 2020}}\Leftrightarrow\hoac{&1<m<{10}^{\sqrt{\log 2020}}\approx 65,78\\&m<{10}^{-\sqrt{\log 2020}}\approx 0,02.} $ \\
		Kết hợp với $\heva{&m>0\\&m\neq 1}$ và $m\in\mathbb{Z}$ nên $m\in\left\{2;3;\ldots;65\right\}$.\\
		Vậy có tất cả $64$ giá trị nguyên $m$ thỏa mãn yêu cầu bài toán.
	}
\end{ex}
\begin{ex}
	[THPT Nguyễn Huy Hiệu - Quảng Nam - 2021]%Câu 8.
	Cho hàm số $f(x)=2020^x-2020^{-x}$. Tìm giá trị nguyên lớn nhất của tham số $m$ để phương trình $f(\log_2x-m)+f\left(\log_2^3x\right)=0$ có nghiệm $x\in(1;16)$ 
	\choice
	{$68$}
	{$65$}
	{\True $67$}
	{$69$}
	\loigiai{
		Xét hàm số $f(x)=2020^x-2020^{-x}$.\\
		Tập xác định: $\mathscr{D}=\mathbb{R}$.\\
		Ta có: $\forall x\in \mathscr{D}\Rightarrow-x\in \mathscr{D}$; $f(-x)=2020^{-x}-2020^x=-\left({2020}^x-2020^{-x}\right)=-f(x)$.\\
		Vậy hàm số $f(x)=2020^x-2020^{-x}$ là hàm số lẻ.\\
		Lại có\\
		$f'(x)=2020^x\cdot\ln 2020-2020^{-x}\cdot\ln 2020\cdot (-x)'=2020^x\cdot\ln 2020+2020^{-x}\cdot\ln 2020>0\forall x\in \mathscr{D}$.\\
		Do đó hàm số $f(x)=2020^x-2020^{-x}$ luôn đồng biến trên $\mathbb{R}$.\\
		Theo đề bài ta có\\
		$\begin{aligned}&f(\log_2x-m)+f\left(\log_2^3x\right)=0\\&\Leftrightarrow f(\log_2x-m)=-f\left(\log_2^3x\right)\end{aligned}$ \\
		$ \Leftrightarrow f(\log_2x-m)=f\left(-\log_2^3x\right) $ (Do $f(x)$ là hàm số lẻ).\\
		Mặt khác hàm số $f(x)$ luôn đồng biến trên $\mathbb{R}$ nên phương trình có nghiệm duy nhất:\\
		$\log_2x-m=-\log_2^3x\Leftrightarrow m=\log_2^3x+\log_2x$.\\
		Đặt $\log_2x=t$. Với $x\in(1;16)\Rightarrow t\in(0;4)$.\\
		Yêu cầu bài toán trở thành, tìm m để phương trình:\\
		$m=t^3+t$ có nghiệm $t\in(0;4)$.\\
		Xét hàm số $f(t)=t^3+t$ trên khoảng $(0;4)$.\\
		Ta có: $f'(t)=3t^2+1>0\forall t$ nên hàm số $f(t)$ đồng biến trên $(0;4)$.\\
		Bảng biến thiên: \\
		{\color{red} HINH O DAY}\\
		Từ bảng biến thiên ta thấy, để phương trình có nghiệm trên khoảng $(0;4)$ thì: $0<m<68$.\\
		Vậy giá trị nguyên lớn nhất của tham số m để phương trình $f(\log_2x-m)+f\left(\log_2^3x\right)=0$ có nghiệm $x\in(1;16)$ là $m=67$.
	}
\end{ex}
\begin{ex}
	[THPT Nguyễn Thị Minh Khai - Hà Tĩnh - 2021]%Câu 9.
	Cho hàm số $y=f(x)$. Hàm số $y=f'(x)$ có bảng biến thiên như sau\\
	{\color{red} HINH O DAY}\\
	Bất phương trình $f(x)<m+\ln x$ đúng với mọi $x\in(1;3)$ khi và chỉ khi
	\choice
	{$m>f(3)-\ln 3$}
	{$m>f(1)$}
	{\True $m\geq f(1)$}
	{$m\geq f(3)-\ln 3$}
	\loigiai{
		Ta có $f(x)-\ln x<m$.\\
		Xét hàm số $g(x)=f(x)-\ln x$.\\
		Ta có $g'(x)=f'(x)-\dfrac{1}{x}$. Dựa vào bảng biến thiên của hàm số $f'(x)$ ta thấy $f'(x)<0$ với mọi $x\in(1;3)$ và $\dfrac{1}{x}>0$ với mọi $x\in(1;3)$. Do đó $g'(x)<0$ với mọi $x\in(1;3)$.\\
		Khi đó $m>g(x)$ với mọi $x\in(1;3)$ tương đương với $m\geq g(1)=f(1)-\ln 1=f(1)$.
	}
\end{ex}
\begin{ex}
	[Nam Định - 2021]%Câu 10.
	Cho hàm số $f(x)$ xác định, liên tục trên $\mathbb{R}$ và $f'(x)$ có bảng xét dấu như sau\\
	{\color{red} HINH O DAY}\\
	Số điểm cực trị của hàm số $f\left(\mathrm{e}^{x^2-|x|-2}\right)$ là
	\choice
	{$2$}
	{$3$}
	{$4$}
	{\True $5$}
	\loigiai{
		Xét hàm số $g(x)=f\left(\mathrm{e}^{x^2-x-2}\right); g'(x)=(2x-1)\mathrm{e}^{x^2-x-2}f'\left(\mathrm{e}^{x^2-x-2}\right);$.\\
		$g'(x)=0\Leftrightarrow\hoac{&2x-1=0\\&f'\left(\mathrm{e}^{x^2-x-2}\right)=0.}$ \\
		Với $2x-1=0\Leftrightarrow x=\dfrac{1}{2}$.\\
		Với $f'\left(\mathrm{e}^{x^2-x-2}\right)=0\Leftrightarrow\hoac{&\mathrm{e}^{x^2-x-2}=-2(VN)\\&\mathrm{e}^{x^2-x-2}=0(VN)\\&\mathrm{e}^{x^2-x-2}=1}\Leftrightarrow x^2-x-2=0\Leftrightarrow\hoac{&x=-1\\&x=2.}$ \\
		Suy ra phương trình $g'(x)=0$ có 3 nghiệm phân biệt nên hàm số $g(x)$ có 3 điểm cực trị trong đó có 2 điểm cực trị có hoành độ dương. Vì vậy hàm số $g(|x|)=f\left(\mathrm{e}^{x^2-|x|-2}\right)$ có 5 điểm cực trị.
	}
\end{ex}
\begin{ex}
	[THPT Hoàng Hoa Thám - Đà Nẵng - 2021]%Câu 11.
	Cho hàm số $f(x)$ có bảng xét dấu của $f'(x)$ như sau:\\
	\textbf{(VẼ HÌNH)}\\
	Hàm số $y=f\left(2-\mathrm{e}^x\right)-\dfrac{1}{3}\mathrm{e}^{3x}+3\mathrm{e}^{2x}-5\mathrm{e}^x+1$ đồng biến trên khoảng nào dưới đây?
	\choice
	{\True $\left(0;\dfrac{3}{2}\right)$}
	{$(1;3)$}
	{$(-3;0)$}
	{$(-4;-3)$}
	\loigiai{
		$y'=-\mathrm{e}^x\cdot f'\left(2-\mathrm{e}^x\right)-\mathrm{e}^{3x}+6\mathrm{e}^{2x}-5\mathrm{e}^x$.\\
		Đặt $t=2-\mathrm{e}^x\Rightarrow\mathrm{e}^x=2-t$.\\
		$y'=(t-2)\cdot f'(t)+t^3-7t+6=(t-2)\cdot f'(t)+(t-2)\left(t^2+2t-3\right)$.\\
		$\begin{aligned}&y'=0\Leftrightarrow(t-2)\cdot f'(t)+(t-2)\left(t^2+2t-3\right)=0\\&\Leftrightarrow(t-2)\left[f'(t)+t^2+2t-3\right]=0\end{aligned}$ \\
		$\Leftrightarrow\hoac{&t=2\\&f'(t)=-t^2-2t+3}\Leftrightarrow\hoac{&t=2\\&t=-3\\&t=1.} $ \\
		Bảng xét dấu:\\
		\textbf{(VẼ HÌNH)}\\
		$y'\geq 0\Leftrightarrow\hoac{&t\geq 2\\&-3\leq t\leq 1}\Rightarrow\hoac{&2-\mathrm{e}^x\geq 2\\&\heva{&2-\mathrm{e}^x\geq-3\\&2-\mathrm{e}^x\leq 1}}\Leftrightarrow\heva{&2-\mathrm{e}^x\geq-3\\&2-\mathrm{e}^x\leq 1}\Leftrightarrow\heva{&\mathrm{e}^x\leq 5\\&\mathrm{e}^x\geq 1}\Leftrightarrow\heva{&x\leq\ln 5\\&x\geq 0.}$ \\
		Vậy hàm số đồng biến trên $(0;\ln 5)$.\\
		Cách 2.\\
		$y'=-\mathrm{e}^x\cdot f'\left(2-\mathrm{e}^x\right)-\mathrm{e}^{3x}+6\mathrm{e}^{2x}-5\mathrm{e}^x$.\\
		Đặt $t=2-\mathrm{e}^x\Rightarrow\mathrm{e}^x=2-t\left(t<2\Rightarrow x<\ln 2\right)$.\\
		$y'=(t-2)\cdot f'(t)+t^3-7t+6=(t-2)\cdot f'(t)+(t-2)\left(t^2+2t-3\right)$.\\
		$\begin{aligned}&y'=0\Leftrightarrow(t-2)\cdot f'(t)+(t-2)\left(t^2+2t-3\right)=0\\&\Leftrightarrow(t-2)\left[f'(t)+t^2+2t-3\right]=0\end{aligned}$ \\
		$ \Leftrightarrow\hoac{&t=2\\&f'(t)=-t^2-2t+3}\Leftrightarrow\hoac{&t=2(ptvn)\\&t=-3\\&t=1}\Rightarrow\hoac{&x=0\\&x=\ln 5.} $ \\
		Bảng xét dấu: $y'=-\mathrm{e}^x\cdot f'\left(2-\mathrm{e}^x\right)-\mathrm{e}^{3x}+6\mathrm{e}^{2x}-5\mathrm{e}^x$.\\
		Lưu ý: $t=0\Leftrightarrow\mathrm{e}^x=2\Leftrightarrow x=\ln 2$.\\
		\textbf{(VẼ HÌNH)}\\
		Và kết luận.
	}
\end{ex}
\begin{ex}
	[Chuyên Bắc Ninh - 2021]%Câu 12.
	Cho hàm số $y=f(x)$. Biết hàm số $y=f'(x)$ có đồ thị như hình vẽ bên dưới. Số điểm cực trị của hàm số $y=2021^{f(x)}+2020^{f(x)}$ là\\
	{\color{red} HINH O DAY}
	\choice
	{$2$}
	{$5$}
	{\True $3$}
	{$4$}
	\loigiai{
		Ta có\\
		$y=2021^{f(x)}+2020^{f(x)}\Rightarrow y'=f'(x)\left({2021}^{f(x)}\cdot\ln 2021+2020^{f(x)}\cdot\ln 2020\right)$.\\
		Do đó các điểm cực trị của hàm số cần tính là số điểm cực trị của hàm số $y=f(x)$.\\
		Dựa vào đồ thị đã cho, ta kết luận có 3 điểm cực trị.
	}
\end{ex}
\begin{ex}
	[THPT Quãng Xương 1-Thanh Hóa - 2021]%Câu 13.
	Cho hàm số $f(x)$ có bảng biến thiên như hình vẽ bên. Bất phương trình $f(\mathrm{e}^x)<m\left(3\mathrm{e}^x+2019\right)$ có nghiệm đúng với mọi $x\in(0; 1)$ khi và chỉ khi
	{\color{red} HINH O DAY}
	\choice
	{$m>\dfrac{-4}{1011}$}
	{$m\geq\dfrac{-2}{1011}$}
	{\True $m\geq\dfrac{f(e)}{3e+2019}$}
	{$m>\dfrac{f(e)}{3e+2019}$}
	\loigiai{
		Đặt $\mathrm{e}^x=t$.\\
		$x\in(0; 1)\Leftrightarrow t\in(1; e)$.\\
		$f(\mathrm{e}^x)<m\left(3\mathrm{e}^x+2019\right)\Leftrightarrow\dfrac{f(\mathrm{e}^x)}{3\mathrm{e}^x+2019}<m\Leftrightarrow\dfrac{f(t)}{3t+2019}<m$.\\
		Yêu cầu bài toàn tương đương với tìm $m$ để $m>\dfrac{f(t)}{3t+2019}$ với mọi $t\in(1; e)$.\\
		Xét $y=\dfrac{f(t)}{3t+2019}\Rightarrow y'=\dfrac{f'(t)\cdot (3t+2019)-3\cdot f(t)}{(3t+2019)^2}$.\\
		Từ bảng biến thiên ta có: $f'(t)>0,\forall t\in(1; e)$, $f(t)<0,\forall t\in(1; e)\Rightarrow y'>0,\forall t\in(1; e)$.\\
		Để $m>y(t),\forall t\in(1; e)\Rightarrow m\geq y(e)=\dfrac{f(e)}{3e+2019}$.
	}
\end{ex}
\begin{ex}
	[Sở Đồng Tháp 2021]%Câu 14.
	Cho hàm số $y=f(x)$ có đồ thị như hình bên. Có tất cả bao nhiêu giá trị nguyên dương của tham số m để phương trình $9^{f(x)}+9m=m\cdot 3^{f(x)}+3^{f(x)+2}$ có đúng 5 nghiệm thực phân biệt. 
	{\color{red} HINH O DAY}
	\choice
	{\True $8$}
	{$9$}
	{$7$}
	{$10$}
	\loigiai{
		Để PT có đúng 5 nghiệm thực phân biệt $\Leftrightarrow(2)$ có đúng 3 nghiệm thực khác $-2; 1$.\\
		$\begin{aligned}&\Leftrightarrow-2<\log_3m<1\\&\Leftrightarrow\dfrac{1}{9}<m<9\\&m\in \mathbb{Z}\Rightarrow m\in\left\{1,2,3,4,5,6,7,8\right\}.\end{aligned}$.
	}
\end{ex}
\begin{ex}[Sở Ninh Bình - 2021]%Câu 15.
	Cho hàm số bậc ba $y=f(x)$ có đồ thị như hình vẽ. \\
	{\color{red} HINH O DAY}\\
	Có bao nhiêu giá trị nguyên của tham số $m\in[-5; 5]$ để phương trình.\\
	$\log_2^3(f(x)+1)-\log_{\sqrt{2}}^2(f(x)+1)+(2m-8)log_{\tfrac{1}{2}}\sqrt{f(x)+1}+2m=0$ có nghiệm $x\in (-1; 1)$? 
	\choice
	{\True $7$}
	{$5$}
	{vô số}
	{$6$}
	\loigiai{
		Phương trình $\log_2^3(f(x)+1)-\log_{\sqrt{2}}^2(f(x)+1)+(2m-8)log_{\tfrac{1}{2}}\sqrt{f(x)+1}+2m=0$ \\
		$\Leftrightarrow\log_2^3(f(x)+1)-4\log_2^2(f(x)+1)+(m-4)log_2(f(x)+1)+2m=0 (1) $.\\
		Đặt $t=\log_2^(f(x)+1)$, Do $x\in (-1; 1)\Rightarrow f(x)\in (-1; 3)$ nên $t\in (-\infty; 2)$.\\
		Phương trình (1) trở thành $t^3-4t^2-(m-4)t+2m=0\Leftrightarrow (t-2)(t^2-2t-m)=0\Leftrightarrow\hoac{&t=2(\text{loại})\\&t^2-2t=m.}$ \\
		Do đó yêu cầu bài toán trở thành có bao nhiêu giá trị nguyên của tham số $m\in[-5; 5]$ để phương trình $t^2-2t=m$ có nghiệm trên khoảng $(-\infty; 2)$. Xét hàm số $g(t)=t^2-2t$ trên khoảng $(-\infty; 2)$, ta có bảng biến thiên
		{\color{red} HINH O DAY}\\
		Dựa vào bảng biến thiên ta được $m\geq-1$ thì phương trình $t^2-2t=m$ có nghiệm trên khoảng.\\
		$(-\infty; 2)$. Do đó có 7 giá trị nguyên của $m$ thỏa mãn yêu cầu.
	}
\end{ex}
\begin{ex}
	[Chuyên Biên Hòa - 2021]%Câu 16.
	Cho hàm số bậc ba $y=f(x)$ có đồ thị như hình vẽ: \\
	{\color{red} HINH O DAY}\\
	Có bao nhiêu giá trị nguyên của tham số $m$ sao cho phương trình $8^{f(x)-2}-3\cdot 4^{f(x)-2}+(m+3)\cdot 2^{f(x)-1}-4-2m=0$ có nghiệm $x\in(-1;0)$?
	\choice
	{$3$}
	{$2$}
	{$1$}
	{\True $0$}
	\loigiai{
		Đặt $t=2^{f(x)-2}$. Với $x\in(-1;0)\Rightarrow f(x)\in(0;2)\Rightarrow t\in\left(\dfrac{1}{4};1\right)$.\\
		Phương trình trở thành: $t^3-3t^2+2(m+3)t-4-2m=0$.\\
		$(t-1)\left(t^2-2t+4+2m\right)=0\Leftrightarrow\hoac{&t=1\notin\left(\dfrac{1}{4};1\right)\\&m=\dfrac{-t^2+2t-4}{2}=g(t)(1).}$ \\
		Để phương trình ban đầu có nghiệm $x\in(-1;0)$ thì phương trình $(1)$ có nghiệm $t\in\left(\dfrac{1}{4};1\right)$.\\
		Ta có: $g'(t)=-t+1$, $\forall t\in\left(\dfrac{1}{4};1\right)$.\\
		$g'(t)=0\Leftrightarrow t=1(l)$ 
		{\color{red} HINH O DAY}\\
		Suy ra $g(t)\in\left(-\dfrac{57}{32};-\dfrac{3}{2}\right)$ hay $m\in\left(-\dfrac{57}{32};-\dfrac{3}{2}\right)$.\\
		Vậy không có giá trị nguyên nào của $m$ thỏa mãn.
	}
\end{ex}
\begin{ex}
	[Chuyên Lương Văn Chánh - Phú Yên - 2021]%Câu 17.
	Cho hàm số $f(x)=x^3+2x-5^m$. Có bao nhiêu giá trị nguyên của $m$ thuộc đoạn $[-6;6]$ để bất phương trình $f\left(f(x)\right)\geq x$ đúng với mọi $x$ thuộc khoảng $(2;6)$. 
	\choice
	{$5$}
	{$11$}
	{\True $8$}
	{$6$}
	\loigiai{
		Ta có $f\left(f(x)\right)\geq x\Leftrightarrow f\left(f(x)\right)+f(x)\geq f(x)+x$.\\
		Xét hàm số $h(t)=f(t)+t$ trên tập số thực $\mathbb{R}$.\\
		Ta có $h'(t)=3t^2+3>0,\forall t\in\mathbb{R}$. Suy ra hàm số $h(t)=f(t)+t$ luôn đồng biến trên tập $\mathbb{R}$.\\
		Suy ra $f\left(f(x)\right)+f(x)\geq f(x)+x\Leftrightarrow f(x)\geq x\Leftrightarrow x^3+x\geq 5^m$.\\
		Để $f\left(f(x)\right)\geq x\Leftrightarrow x^3+x\geq 5^m$ đúng với mọi $x$ thuộc khoảng $(2;6)$ \\
		$ \Leftrightarrow\min\limits_{[2;6]}\left(x^3+x\right)\geq 5^m\Leftrightarrow 2^3+2\geq 5^m\Leftrightarrow m\leq\log_510 $.\\
		Do $m$ là số nguyên thuộc đoạn $[-6;6]$ nên $m\in\left\{-6;-5;-4;-3;-2;-1;0;1\right\}$.\\
		Vậy có $8$ giá trị nguyên của $m$ cần tìm.
	}
\end{ex}
\begin{ex}[Sở Hà Tĩnh - 2021]%Câu 18.
	Cho hàm số $y=f(x)$. Hàm số $y=f'(x)$ có bảng biến thiên như hình vẽ.\\
	{\color{red} HINH O DAY}\\
	Bất phương trình $f(x)<\cos^2x+3m$ đúng với mọi $x\in\left(0;\dfrac{\pi}{2}\right)$ khi và chỉ khi
	\choice
	{$m\geq\dfrac{1}{3}f\left(\dfrac{\pi}{4}\right)-\dfrac{1}{6}$}
	{$m>\dfrac{1}{3}f\left(\dfrac{\pi}{2}\right)$}
	{$m\geq\dfrac{1}{3}[f(0)-1]$}
	{\True $m\geq\dfrac{1}{3}f\left(\dfrac{\pi}{2}\right)$}
	\loigiai{
		Từ bảng biến thiên của hàm số $y=f'(x)$ ta thấy $f'(x)>1,\forall x\in\left(0;\dfrac{\pi}{2}\right)$.\\
		Xét bất phương trình $f(x)<\cos^2x+3m$ với mọi $x\in\left(0;\dfrac{\pi}{2}\right)$ \\
		$ \Leftrightarrow 3m>f(x)-\cos^2x $ với mọi $x\in\left(0;\dfrac{\pi}{2}\right)$.\\
		Xét hàm số $g(x)=f(x)-\cos^2x$ với $x\in\left(0;\dfrac{\pi}{2}\right)$.\\
		Ta có $g'(x)=f'(x)+2\cos x\sin x=f'(x)+\sin 2x>0,\forall x\in\left(0;\dfrac{\pi}{2}\right)$.\\
		Bảng biến thiên của hàm số $g(x)=f(x)-\cos^2x$ với $x\in\left(0;\dfrac{\pi}{2}\right)$ 
		{\color{red} HINH O DAY}\\
		Khi đó: $3m>g(x)$ với mọi $x\in\left(0;\dfrac{\pi}{2}\right)\Leftrightarrow 3m\geq g\left(\dfrac{\pi}{2}\right)\Leftrightarrow 3m\geq f\left(\dfrac{\pi}{2}\right)-\cos^2\left(\dfrac{\pi}{2}\right)$ \\
		$ \Leftrightarrow 3m\geq f\left(\dfrac{\pi}{2}\right)-0\Leftrightarrow m\geq\dfrac{1}{3}f\left(\dfrac{\pi}{2}\right) $.
	}
\end{ex}
\begin{ex}[Sở Hà Tĩnh - 2021]%Câu 19.
	Cho hàm số bậc 4 có đồ thị như hình vẽ. Có bao nhiêu giá trị nguyên của tham số $m$ và $m\in[-2021; 2021]$ để phương trình $\log\dfrac{f(x)}{mx^2}+x[f(x)-mx]=mx^3-f(x)$ có hai nghiệm dương phân biệt?\\
	{\color{red} HINH O DAY}\\
	\choice
	{$2022$}
	{$2020$}
	{$2019$}
	{\True $2021$}
	\loigiai{
		Điều kiện: $\dfrac{f(x)}{mx^2}>0\Rightarrow\heva{&f(x)>0\\&x\neq 0}\Leftrightarrow x\notin\{-1; 0; 1\}$.\\
		Xét: $\log\dfrac{f(x)}{mx^2}+x[f(x)-mx]=mx^3-f(x)$ \\
		$ \Leftrightarrow\log f(x)+\log(x+1)+xf(x)+f(x)=\log mx^2+\log(x+1)+mx^3+mx^2 $ \\
		$\Leftrightarrow\log[(x+1)f(x)]+(x+1)f(x)=\log\left[mx^2(x+1)\right]+mx^2(x+1) $.\\
		Điều kiện bổ sung: $x+1>0\Leftrightarrow x >-1$.\\
		Xét hàm số $g(t)=\log t+t$ trên $(0;+\infty)$, khi đó: $g'(t)=\dfrac{1}{t\ln 10}+1>0,\forall t\in(0;+\infty)$.\\
		Suy ra: $g$ là hàm tăng trên $(0;+\infty)$.\\
		Khi đó: $(x+1)f(x)=mx^2(x+1)\Leftrightarrow f(x)=mx^2$.\\
		Dựa vào đồ thị, để hàm số $y=f(x)$ và $y=mx^2$ cắt nhau có 2 điểm có hoành độ dương thì $m>0$.\\
		Kết hợp với đề bài: $m\in\mathbb{Z}$ và $m\in[-2021; 2021]$, ta được 2021 giá trị của $m$ thỏa yêu cầu đề bài.
	}
\end{ex}
\begin{ex}
	[THPT Ngô Quyền - Quảng Ninh - 2021]%Câu 20.
	Cho hàm số $y=f(x)$ là hàm số chẵn và xác định trên $\mathbb{R}$, sao cho $f(0)\neq 0$ và phương trình $5^x-5^{-x}=f(x)$ có đúng $5$ nghiệm phân biệt. Khi đó số nghiệm của phương trình $5^x+5^{-x}=f^2\left(\dfrac{x}{2}\right)+2$ là
	\choice
	{$5$}
	{$15$}
	{\True $10$}
	{$20$}
	\loigiai{
		Ta có $5^x+5^{-x}=f^2\left(\dfrac{x}{2}\right)+2\Leftrightarrow f^2\left(\dfrac{x}{2}\right)=5^x+5^{-x}-2=\left(5^{\tfrac{x}{2}}-5^{-\tfrac{x}{2}}\right)^2$ \\
		$ \Leftrightarrow\hoac{&f\left(\dfrac{x}{2}\right)=5^{\tfrac{x}{2}}-5^{-\tfrac{x}{2}}\\&f\left(\dfrac{x}{2}\right)=5^{-\tfrac{x}{2}}-5^{\tfrac{x}{2}}}\Leftrightarrow\hoac{&f(t)=5^t-5^{-t}\\&f(t)=5^{-t}-5^t} $ (với $t=\dfrac{x}{2}$).\\
		Do $f(x)$ là hàm số chẵn và xác định trên $\mathbb{R}$ nên $f(x)=f(-x),\forall x\in\mathbb{R}$.\\
		Khi đó từ phương trình $5^x-5^{-x}=f(x)$, thay $x$ bởi $-x$ ta được $f(-x)=f(x)=5^{-x}-5^x$.\\
		Vì phương trình $5^x-5^{-x}=f(x)$ có đúng $5$ nghiệm phân biệt nên phương trình $f(x)=5^{-x}-5^x$ cũng có đúng $5$ nghiệm phân biệt.\\
		Suy ra phương trình $f(t)=5^t-5^{-t}$ có $5$ nghiệm phân biệt $t_1,t_2,\ldots,t_5$ và phương trình $f(t)=5^{-t}-5^t$ cũng có $5$ nghiệm phân biệt $t_6,t_7,\ldots,t_{10}$ (*).\\
		Giả sử phương trình $5^x-5^{-x}=f(x)$ và $5^{-x}-5^x=f(x)$ có nghiệm chung $x=x_0$.\\
		Khi đó $\heva{&f(x_0)=5^{x_0}-5^{-x_0}(1)\\&f(x_0)=5^{-x_0}-5^{x_0}(2).}$ \\
		Lấy $(1)-(2)$ ta được $2\left(5^{x_0}-5^{-x_0}\right)=0\Leftrightarrow 5^{x_0}=5^{-x_0}\Leftrightarrow x_0=0$.\\
		Lấy $(1)+(2)$ ta được $2f(x_0)=0\Leftrightarrow f(x_0)=0$.\\
		Suy ra $x_0=0$ là nghiệm của phương trình $f(x)=0$ hay $f(0)=0$ (mâu thuẫn với giả thiết).\\
		Suy ra hai phương trình $f(t)=5^t-5^{-t}$ và $f(t)=5^{-t}-5^t$ không có nghiệm chung (**).\\
		Từ (*) và (**) ta suy ra phương trình $5^x+5^{-x}=f^2\left(\dfrac{x}{2}\right)+2$ có tổng cộng $10$ nghiệm phân biệt.
	}
\end{ex}
\begin{ex}
	[THPT Nguyễn Huệ - Phú Yên - 2021]%Câu 21.
	Cho hàm số $f(x)=\mathrm{e}^x-\mathrm{e}^{-x}+2020x$. Có bao nhiêu số nguyên dương $m$ sao cho ứng với mỗi $m$ có đúng $10$ số nguyên dương $x$ thỏa mãn bất phương trình $f(mx+1)+f(2x-2021)<0$?
	\choice
	{$2018$}
	{$19$}
	{\True $18$}
	{$2019$}
	\loigiai{
		Ta có: $f(-x)=\mathrm{e}^{-x}-\mathrm{e}^x-2020x=-f(x),\forall x\in\mathbb{R}\Rightarrow$ Hàm số $f(x)$ là hàm số lẻ.\\
		Lại có: $f'(x)=\mathrm{e}^x+\mathrm{e}^{-x}+2020>0,\forall x\in\mathbb{R}\Rightarrow$ Hàm số $f(x)$ đồng biến trên $\mathbb{R}$.\\
		Khi đó: $f(mx+1)+f(2x-2021)<0\Leftrightarrow f(mx+1)<f(2021-2x)$ \\
		$ \Leftrightarrow mx+1<2021-2x\Leftrightarrow (m+2)x<2020\Leftrightarrow x<\dfrac{2020}{m+2} $ (do $m>0$).\\
		Yêu cầu bài toán $\Leftrightarrow 10<\dfrac{2020}{m+2}\leq 11\Leftrightarrow\dfrac{1998}{11}\leq m<200$.\\
		Do $m$ nguyên dương nên $m\in\left\{182; 183;\ldots; 199\right\}$.\\
		Vậy có $18$ số nguyên dương $m$ thỏa mãn yêu cầu bài toán.
	}
\end{ex}
\begin{ex}
	[THPT Thanh Chương 1 - Nghệ An - 2021]%Câu 22.
	Cho hàm số $y=f(x)$ có bảng biên thiên như sau: \\
	{\color{red} HINH O DAY}\\
	Có tất cả bao nhiêu giá trị nguyên của tham số $m$ để phương trình $4^{f(x)-m}+3^{f(x)-m}-5f(x)+5m-2=0$ có nghiệm?
	\choice
	{$5$}
	{$3$}
	{$6$}
	{\True $4$}
	\loigiai{
		Xét hàm số $g(t)=4^t+3^t-5t-2$ trên $\mathbb{R}$.\\
		$g'(t)=4^t\cdot\ln 4+3^t\cdot\ln 3-5$; $g''(t)=4^t\cdot\ln^24+3^t\cdot\ln^23>0\forall t\in\mathbb{R}$ \\
		$ \Rightarrow $ Phương trình $g(t)=0$ có tối đa $2$ nghiệm.\\
		Mà $g(0)=g(1)=0$.\\
		Do đó phương trình $4^{f(x)-m}+3^{f(x)-m}-5f(x)+5m-2=0\Leftrightarrow g\left(f(x)-m\right)=0$ \\
		$\Leftrightarrow\hoac{&f(x)-m=0\\&f(x)-m=1}\Leftrightarrow\hoac{&f(x)=m\\&f(x)=m+1.} $ \\
		Yêu cầu bài toán $\Leftrightarrow\hoac{&-1\leq m\leq 1\\&-1\leq m+1\leq 1}\Leftrightarrow -2\leq m\leq 1$.\\
		Do $m$ nguyên nên $m\in\{-2;-1;0;1\}$.\\
		Vậy có $4$ giá trị nguyên của tham số $m$ thỏa mãn yêu cầu bài toán.
	}
\end{ex}
\begin{ex}
	[THPT Yên Định 2 - Thanh Hóa - 2021]%Câu 23.
	Cho hàm số $f(x)$ liên tục trên $\mathbb{R}$ và có đồ thị $f'(x)$ như hình vẽ bên. \\
	{\color{red} HINH O DAY}\\
	Bất phương trình $\log_5[f(x)+m+2]+f(x)>4-m$ đúng với mọi $x\in(-1;4)$ khi và chỉ khi
	\choice
	{\True $m\geq 3-f(4)$}
	{$m\geq 3-f(1)$}
	{$m<4-f(-1)$}
	{$m\geq 4-f(-1)$}
	\loigiai{
		Điều kiện $f(x)+m+2>0$.\\
		Đặt $t=\log_5[f(x)+m+2]\Rightarrow f(x)+m+2=5^t$.\\
		Bất phương trình đã cho trở thành $t+5^t>6$.\\
		Xét hàm $g(t)=t+5^t$.\\
		$g'(t)=1+5^t\cdot\ln 5>0,\forall t$ do đó $g(t)$ là hàm đồng biến.\\
		Mà $g(1)=6$ nên $t+5^t>6\Leftrightarrow t>1$.\\
		Bất phương trình $\log_5[f(x)+m+2]+f(x)>4-m$ đúng với mọi $x\in(-1;4)$ khi và chỉ khi $\heva{&f(x)+m+2>0\\&\log_5[f(x)+m+2]>1},\forall x\in(-1;4)\Leftrightarrow\heva{&f(x) >-m-2\\&f(x) >-m+3},\forall x\in(-1;4)\Leftrightarrow f(x) >-m+3,\forall x\in(-1;4)$.\\
		Xét hàm $f(x)$ trên $(-1;4)$ 
		{\color{red} HINH O DAY}\\
		Quan sát đồ thị của hàm số $f'(x)$ ta có\\
		$\displaystyle\int\limits_{-1}^1 f'(x)\mathrm{\,d}x <-\displaystyle\int\limits_1^4 f'(x)\mathrm{\,d}x\Leftrightarrow f(1)-f(-1)<f(1)-f(4)\Leftrightarrow f(-1)>f(4)$.\\
		Dựa vào bảng biến thiên của hàm $f(x)$ trên $(-1;4)$ và dựa vào nhận xét $f(-1)>f(4)$ ta có $f(x) >-m+3,\forall x\in(-1;4)$ khi $f(4)\geq-m+3\Leftrightarrow m\geq 3-f(4)$.
	}
\end{ex}
\begin{ex}
	[THPT Yên Định 2 - Thanh Hóa - 2021]%Câu 24.
	Cho hàm số $y=f(x)$ có nhận giá trị dương và có đạo hàm cấp một không âm trên $(0;+\infty)$ đồng thời thoả mãn: $\dfrac{3}{x^2}f(x)f'(x)[xf(x)]'+\dfrac{1}{x^3}\ln (1+\dfrac{xf'(x)}{f(x)})+[f'(x)]^3=0$. Giá trị của $P=2019+2020f'(2021)$ là
	\choice
	{$2020$}
	{\True $2019$}
	{$2021$}
	{$0$}
	\loigiai{
		$\dfrac{3}{x^2}f(x)f'(x)[xf(x)]'+\dfrac{1}{x^3}\ln (1+\dfrac{xf'(x)}{f(x)})+[f'(x)]^3=0$ \\
		$\Leftrightarrow\dfrac{3}{x^2}f(x)f'(x)\left[f(x)+xf'(x)\right]+\dfrac{1}{x^3}\ln (1+\dfrac{xf'(x)}{f(x)})+[f'(x)]^3=0 $.\\
		Do:\\
		$f(x)>0,f'(x)\geq 0\forall x>0$.\\
		$+)f(x)+xf'(x)\geq f(x)\Rightarrow f(x)+xf'(x)>0$.\\
		Nên ta có: $\dfrac{3}{x^2}f(x)f'(x)\left[f(x)+xf'(x)\right]\geq 0$.\\
		$\ln (1+\dfrac{xf'(x)}{f(x)})\geq\ln 1\Rightarrow\ln (1+\dfrac{xf'(x)}{f(x)})\geq 0$.\\
		$[f'(x)]^3\geq 0$.\\
		Suy ra: $\dfrac{3}{x^2}f(x)f'(x)\left[f(x)+xf'(x)\right]+\dfrac{1}{x^3}\ln (1+\dfrac{xf'(x)}{f(x)})+[f'(x)]^3\geq 0\forall x>0$.\\
		Dấu bằng xảy ra $\Leftrightarrow f'(x)=0\forall x>0\Rightarrow f'(2021)=0$.\\
		Do đó: $P=2019+2020f'(2021)=2019$.
	}
\end{ex}
\begin{ex}
	[Chuyên Nguyễn Bỉnh Khiêm - Quảng Nam - 2021]%Câu 25.
	Cho hàm số $f(x)$ liên tục trên $\mathbb{R}$ và có đồ thị $f'(x)$ như hình vẽ bên. 
	{\color{red} HINH O DAY}\\
	Bất phương trình $\log_5[f(x)+m+2]+f(x)>4-m$ đúng với mọi $x\in(-1; 4)$ khi và chỉ khi
	\choice
	{$m\geq 3-f(1)$}
	{\True $m\geq 3-f(4)$}
	{$m\geq 4-f(-1)$}
	{$m<4-f(-1)$}
	\loigiai{
		Ta có, bất phương trình $\log_5[f(x)+m+2]+f(x)>4-m$ \\
		$ \Leftrightarrow\log_5[f(x)+m+2]+f(x)+m+2>6 $ \\
		$ \Leftrightarrow\log_5[f(x)+m+2]+f(x)+m+2>\log_5(5)+5 $.\\
		Đặt: $t=f(x)+m+2,(t>0)$. Ta được bất phương trình $\log_5(t)+t>\log_5(5)+5$ (1).\\
		Xét hàm số $f(t)=\log_5(t)+t,(t>0)$.\\
		Ta có $f'(t)=\dfrac{1}{t\ln 5}+1>0,\forall t\in(0;+\infty)$ \\
		$ \Rightarrow f(t)=\log_5(t)+t $ là hàm số đồng biến trên khoảng $(0;+\infty)$.\\
		Từ (1) ta có $f(t)>f(5)\Leftrightarrow t>5$.\\
		Vậy, bất phương trình $\log_5[f(x)+m+2]+f(x)>4-m$ đúng với mọi $x\in(-1; 4)$ khi và chỉ khi $f(x)+m+2>5,\forall x\in(-1;4)\Leftrightarrow m>3-f(x),\forall x\in(-1;4)$.\\
		Dựa vào đồ thị $f'(x)$ ta có:\\
		$\displaystyle\int\limits_{-1}^4 f'(x)\mathrm{\,d}x<0\Leftrightarrow f(4)-f(-1)<0\Leftrightarrow f(4)<f(-1)\Rightarrow 3-f(4)>3-f(1)$.\\
		Mặt khác, dựa vào đồ thị hàm số $f'(x)$, ta có BBT vủa hàm số $f(x)$ trên $(-1; 4)$ như sau:\\
		{\color{red} HINH O DAY}\\
		Vậy, hàm số $3-f(x)$ có bảng biến thiên như sau:\\
		{\color{red} HINH O DAY}\\
		Vậy, $m>3-f(x),\forall x\in(-1;4)\Leftrightarrow m\geq 3-f(4)$.\\
		Do đó bất phương trình $\log_5[f(x)+m+2]+f(x)>4-m$ đúng với mọi $x\in(-1; 4)$ khi và chỉ khi $m\geq 3-f(4)$.
	}
\end{ex}
\begin{ex}
	[Chuyên Nguyễn Bỉnh Khiêm - Quảng Nam - 2021]%Câu 26.
	Cho hàm số $y=f(x)$ liên tục và có đạo hàm trên $R$. Hàm số $y=f'(x)$ có bảng xét dấu như bảng bên cạnh. \\
	{\color{red} HINH O DAY}\\
	Bất phương trình $f(x)>\mathrm{e}^{\cos x}+m$ có nghiệm $x\in\left(0;\dfrac{\pi}{2}\right)$ khi và chỉ khi
	\choice
	{$m>f(0)-e$}
	{\True $m<f(\dfrac{\pi}{2})-1$}
	{$m\leq f(\dfrac{\pi}{2})-1$}
	{$m\leq f(0)-e$}
	\loigiai{
		Có $f(x)>\mathrm{e}^{\cos x}+m\Leftrightarrow f(x)-\mathrm{e}^{\cos x}>m$.\\
		Xét hàm $g(x)=f(x)-\mathrm{e}^{\cos x}$ trên $(0;\dfrac{\pi}{2})$.\\
		Có $g'(x)=f'(x)-(-\sin x\cdot\mathrm{e}^{\cos x})=f'(x)+\sin x\cdot\mathrm{e}^{\cos x}>0,\forall x\in (0;\dfrac{\pi}{2})$.\\
		Suy ra $g(x)$ đồng biến trên $(0;\dfrac{\pi}{2})$, nên để bất phương trình $g(x)>m$ có nghiệm $x\in\left(0;\dfrac{\pi}{2}\right)$ khi và chỉ khi $m<g(\dfrac{\pi}{2})$ hay $m<f(\dfrac{\pi}{2})-1$.
	}
\end{ex}
\begin{ex}
	[Sở Quảng Bình - 2021]%Câu 27.
	Cho hàm số $y=f(x)$ liên tục trên $\mathbb{R}$ và hàm số $y=f'(x)$ có bảng biến thiên như sau:\\
	{\color{red} HINH O DAY}\\
	Điều kiện cần và đủ để $3^{f(x)+m}+4^{f(x)+m}-5f(x)\leq 2+5m,\forall x\in(-1; 2)$ là
	\choice
	{$-f(2)<m<1-f(-1)$}
	{$-f(-1)<m<1-f(2)$}
	{$-f(-1)\leq m\leq 1-f(2)$}
	{\True $-f(2)\leq m\leq 1-f(-1)$}
	\loigiai{
		Xét hàm số $y=g(t)=3^t+4^t-5t-2$ trên tập $\mathbb{R}$.\\
		Ta có $y'=g'(t)=3^t\cdot\ln 3+4^t\cdot\ln 4-5$.\\
		$y''=g''(t)=3^t\cdot\ln^23+4^t\cdot\ln^24>0,\forall t\in\mathbb{R}$.\\
		Suy ra phương trình $y'=g'(t)=3^t\cdot\ln 3+4^t\cdot\ln 4-5=0$ có không quá một nghiệm trên $\mathbb{R}$.\\
		Do đó phương trình $g(t)=3^t+4^t-5t-2=0$ có không quá hai nghiệm trên $\mathbb{R}$.\\
		Lại có $g(0)=0; g(1)=0$.\\
		Suy ra phương trình $3^t+4^t-5t-2=0\Leftrightarrow\hoac{&t=0\\&t=1.}$ \\
		Lập bảng xét dấu của biểu thức $g(t)=3^t+4^t-5t-2$ trên tập $\mathbb{R}$ ta suy ra\\
		$g(t)=3^t+4^t-5t-2\leq 0\Leftrightarrow 0\leq t\leq 1$.\\
		Bất phương trình $3^{f(x)+m}+4^{f(x)+m}-5f(x)\leq 2+5m\Leftrightarrow 3^{f(x)+m}+4^{f(x)+m}-5[f(x)+m]-2\leq 0$ \\
		$ \Leftrightarrow g\left(f(x)+m\right)\leq 0\Leftrightarrow 0\leq f(x)+m\leq 1 $.\\
		Yêu cầu bài toán $\Leftrightarrow 0\leq f(x)+m\leq 1,\forall x\in(-1; 2)\Leftrightarrow-f(x)\leq m\leq 1-f(x),\forall x\in(-1; 2)$.\\
		Dựa vào bảng biến thiên của hàm số $y=f'(x)$ ta thấy $f'(x)\leq 0,\forall x\in(-1; 2)\Rightarrow-f'(x)\geq 0,\forall x\in(-1; 2)$; ($f'(x)=0\Leftrightarrow x=0$).\\
		Từ đó $-f(x)\leq m\leq 1-f(x),\forall x\in(-1; 2)\Leftrightarrow-f(2)\leq m\leq 1-f(-1)$.
	}
\end{ex}
\begin{ex}
	[Sở Quảng Bình - 2021]%Câu 28.
	Cho hàm số $y=f(x)$ liên tục trên $\mathbb{R}$ và có đồ thị như hình vẽ bên.\\ 
	\textbf{(VẼ HÌNH)}\\
	Hỏi trong khoảng $(0; 2020)$ phương trình $f\left({2020}^{cos2x}\right)=f(\tan x)$ có bao nhiêu nghiệm?
	{\color{red} HINH O DAY}
	\choice
	{$321$}
	{\True $643$}
	{$642$}
	{$322$}
	\loigiai{
		Ta thấy đồ thị có chiều hướng đi lên từ trái qua phải nên hàm số đồng biến trên $\mathbb{R}$. Do đó $f\left({2020}^{cos2x}\right)=f(\tan x)\Leftrightarrow 2020^{cos2x}=\tan x\Leftrightarrow 2020^{cos^2x-\sin^2x}=\dfrac{\sin x}{\cos x}\Leftrightarrow \cos x\cdot 2020^{cos^2x}=\sin x\cdot 2020^{\sin^2x}(*)$ Xét hàm $g(t)=t\cdot 2020^{t^2}$ có $g'(t)=2020^{t^2}+2t^2\cdot 2020^{t^2}\ln 2020>0;\forall t$.\\
		$(*)\Leftrightarrow g(\cos x)=g(\sin x)\Leftrightarrow\cos x=\sin x\Leftrightarrow\tan x=1\Leftrightarrow x=\dfrac{\pi}{4}+k\pi$.\\
		$x\in(0; 2020)\Rightarrow 0<\dfrac{\pi}{4}+k\pi<2020\Leftrightarrow-\dfrac{1}{4}<k<642,74$; mà $k\in\mathbb{Z}$ nên $k\in\left\{0; 1; 2;\ldots; 642\right\}$.
	}
\end{ex}
\begin{ex}
	[THCS - THPT Nguyễn Khuyến 2019]%Câu 29.
	Cho hàm số $y=f(x)$ liên tục trên $\mathbb{R}$, có bảng biến thiên như hình vẽ và có đạo hàm cấp hai $f''(x)>0,\forall x\in\mathbb{R}$. 
	{\color{red} HINH O DAY}\\
	Gọi $a,b,c,n$ là các số thực và biểu thức: $P=-\left(\mathrm{e}^{f(a)}+\mathrm{e}^{f(b)}+\mathrm{e}^{f(c)}\right)+\dfrac{3}{2}\left[f\left(\dfrac{a+b+c}{3}\right)+1\right]^2$. Khẳng định đúng với mọi $a,b,c,n\in\mathbb{R}$ là
	\choice
	{$0<P<3$}
	{$7-3e\leq P\leq 0$}
	{$P\geq 3$}
	{\True $P<7-3e$}
	\loigiai{
		Ta có $\mathrm{e}^{f(a)}+\mathrm{e}^{f(b)}+\mathrm{e}^{f(c)}\geq 3\sqrt[3]{\mathrm{e}^{f(a)+f(b)+f(c)}}$.\\
		Mặt khác do $f''(x)>0,\forall x\in\mathbb{R}$ nên $f(x)$ là hàm lồi, áp dụng bất đẳng thức lồi ta có\\
		$f(a)+f(b)+f(c)\geq 3f\left(\dfrac{a+b+c}{3}\right)$.\\
		Do đó $\mathrm{e}^{f(a)}+\mathrm{e}^{f(b)}+\mathrm{e}^{f(c)}\geq 3\sqrt[3]{\mathrm{e}^{3f\left(\dfrac{a+b+c}{3}\right)}}=3\mathrm{e}^{f\left(\dfrac{a+b+c}{3}\right)}$.\\
		Suy ra $P\leq-3\mathrm{e}^{f\left(\dfrac{a+b+c}{3}\right)}+\dfrac{3}{2}\left[f\left(\dfrac{a+b+c}{3}\right)+1\right]^2$. Đặt $t=f\left(\dfrac{a+b+c}{3}\right),t\geq n^2\geq 0$.\\
		Ta có $P\leq g(t)$ với $g(t)=-3\mathrm{e}^t+\dfrac{3}{2}(t+1)^2$.\\
		Nên $g'(t)$ là hàm nghịch biến trên $[0;+\infty)$. $\Rightarrow g'(t)\leq g(0)=0,\forall t\in[0;+\infty)\Rightarrow g(t)\leq g(0)$.\\
		Do đó $P\leq g(0)=\dfrac{-3}{2}<7-3e$.
	}
\end{ex}
\begin{ex}
	[Sở Thanh Hóa 2022]%Câu 30.
	Cho hàm số $f(x)=2^x-2^{-x}+2022x^3$. Biết rằng tồn tại só thực $m$ sao cho bất phương trình $f\left(4^x-mx+37m\right)+f\left((x-m-37)\cdot 2^x\right)\geq 0$ nghiệm đúng với mọi $x\in\mathbb{R}$. Khi đó $m$ thuộc khoảng nào dưới đây?
	\choice
	{\True $(30;50)$}
	{$(10;30)$}
	{$(50;70)$}
	{$(-10;10)$}
	\loigiai{
		Có $f'(x)=2^x\ln 2+2^{-x}\ln 2+6066x^2>0,\forall x;f(-x)=2^{-x}-2^x-2022x^3=-\left(2^x-2^{-x}+2022x^3\right)=-f(x)$.\\
		$f\left(4^x-mx+37m\right)+f\left((x-m-37){\cdot 2}^x\right)\geq 0\Leftrightarrow f\left(4^x-mx+37m\right)\geq f\left(-(x-m-37){\cdot 2}^x\right)$.\\
		$4^x-mx+37m\geq-(x-m-37)\cdot 2^x\Leftrightarrow 4^x-m\cdot 2^x-m(x-37)+(x-37)\cdot 2^x\geq 0$ \\
		$ \Leftrightarrow 2^x\left(2^x-m\right)+(x-37)\left(2^x-m\right)\geq 0\Leftrightarrow\left(2^x-m\right)\left(2^x+x-37\right)\geq 0 $.\\
		ycbt $\Leftrightarrow g(x)=\left(2^x-m\right)\left(2^x+x-37\right)\geq 0,\forall x$. Do $2^x+x-37=0$ có nghiệm $x=5$ do đó điều kiện cần là $2^x-m=0$ có nghiệm $x=5\Leftrightarrow m=2^5=32$.\\
		Điều kiện đủ thử lại $m=32\Rightarrow g(x)=\left(2^x-32\right)\left(2^x+x-37\right)\geq 0,\forall x$ (thoả mãn).
	}
\end{ex}
\begin{ex}
	[Sở Vĩnh Phúc 2022]%Câu 31.
	Cho hàm số bậc ba $y=f(x)$ có đồ thị như hình vẽ
	{\color{red} HINH O DAY}\\
	Số nghiệm của phương trình $\log_2^3\left(f(x)+1\right)-\log_{\sqrt{2}}^2\left(f(x)+1\right)-2\log_{\tfrac{1}{2}}\sqrt{f(x)+1}+6=0$ là
	\choice
	{$7$}
	{$5$}
	{\True $6$}
	{$8$}
	\loigiai{
		Xét phương trình $\log_2^3\left(f(x)+1\right)-\log_{\sqrt{2}}^2\left(f(x)+1\right)-2\log_{\tfrac{1}{2}}\sqrt{f(x)+1}+6=0$ (1).\\
		Điều kiện: $f(x) >-1$.\\
		Ta có $(1)\Leftrightarrow\log_2^3\left(f(x)+1\right)-4\log_2^2\left(f(x)+1\right)+\log_2\left(f(x)+1\right)+6=0$.\\
		Đặt $t=\log_2\left(f(x)+1\right)$, ta có $t^3-4t^2+t+6=0\Rightarrow\hoac{&t=-1\\&t=3\\&t=2}\Leftrightarrow\hoac{&\log_2\left(f(x)+1\right)=-1\\&\log_2\left(f(x)+1\right)=3\\&\log_2\left(f(x)+1\right)=2}$ \\
		$\Leftrightarrow\hoac{&f(x)=-\dfrac{1}{2}\quad(2)\\&f(x)=7(3)\\&f(x)=3(4)} $ 
		{\color{red} HINH O DAY}\\
		Dễ thấy số nghiệm phương trình $(2)$ là số giao điểm của đồ thị hàm số $y=f(x)$ và $y=-\dfrac{1}{2}$. Suy ra phương trình (2) có 3 nghiệm phân biệt.\\
		Tương tự phương trình $(3)$ có nghiệm duy nhất và phương trình $(4)$ có 2 nghiệm phân biệt.\\
		Hiển nhiên các nghiệm này phân biệt. Do dó phương trình $(1)$ có 6 nghiệm phân biệt.
	}
\end{ex}
\begin{ex}
	[THPT Trần Phú – Hà Tĩnh – 2022]%Câu 32.
	Cho hàm số $f(x)=-x^3+\dfrac{13}{2}x^2-12x-\mathrm{e}^x-2022$. Bất phương trình ẩn m sau đây $f\left[\log_{0,5}\left(\log_2(2m+1)\right)-2021\right]<f[f(0)]$ có bao nhiêu nghiệm nguyên?
	\choice
	{$14$}
	{$10$}
	{$11$}
	{\True $7$}
	\loigiai{
		Ta có $f'(x)=-3x^2+13x-12-\mathrm{e}^x<0,\forall x$ do đó:\\
		$\begin{aligned}&f\left[\log_{0,5}\left(\log_2(2m+1)\right)-2021\right]<f[f(0)]\\&\Leftrightarrow\log_{0,5}\left(\log_2(2m+1)\right)-2021>f(0)=-2023\\&\Leftrightarrow\log_{0,5}\left(\log_2(2m+1)\right) >-2\Leftrightarrow 0<\log_2(2m+1)<(0,5)^{-2}=4\\&\Leftrightarrow 1<2m+1<16\Leftrightarrow 0<m<\dfrac{15}{2}\Rightarrow m\in\{1,\ldots,7\}.\end{aligned}$.
	}
\end{ex}
\begin{ex}
	[THPT Yên Phong 1 - Bắc Ninh - 2022]%Câu 33.
	Cho hàm số $y=f(x)$. Hàm số $y=f'(x)$ có bảng biến thiên như sau
	{\color{red} HINH O DAY}\\
	Bất phương trình $2f(x)<\mathrm{e}^{\cos x}+m$ đúng $\forall x\in\left(0;\dfrac{\pi}{2}\right)$ khi và chỉ khi
	\choice
	{$m\geq 2f(0)-\mathrm{e}$}
	{\True $m\geq 2f\left(\dfrac{\pi}{2}\right)-1$}
	{$m>2f\left(\dfrac{\pi}{2}\right)-1$}
	{$m>2f(0)-\mathrm{e}$}
	\loigiai{
		Có $2f(x)<\mathrm{e}^{\cos x}+m\Leftrightarrow 2f(x)-\mathrm{e}^{\cos x}<m$.\\
		Xét hàm $g(x)=2f(x)-\mathrm{e}^{\cos x}$, có $g'(x)=2f'(x)+\sin x\cdot\mathrm{e}^{\cos x}>0,\forall x\in\left(0;\dfrac{\pi}{2}\right)$.\\
		Suy ra $g(x)=2f(x)-\mathrm{e}^{\cos x}$ đồng biến trên $\left[0;\dfrac{\pi}{2}\right]$.\\
		Vậy bất phương trình đã cho nghiệm đúng với mọi $x\in\left(0;\dfrac{\pi}{2}\right)$. $\Leftrightarrow m\geq\max\limits_{\left[0;\dfrac{\pi}{2}\right]} g(x)=g\left(\dfrac{\pi}{2}\right)=2f\left(\dfrac{\pi}{2}\right)-\mathrm{e}^{\cos\dfrac{\pi}{2}}=2f\left(\dfrac{\pi}{2}\right)-1$.
	}
\end{ex}
\begin{ex}
	[THPT Quảng Xương 1 – Thanh Hóa 2022]%Câu 34.
	Có bao nhiêu cặp số nguyên $(a;b)$ thỏa mãn $0<a,b\leq 20$ sao cho đồ thị của hai hàm số $y=\dfrac{1}{a^x}+\dfrac{1}{b}$ và $y=\dfrac{1}{b^x}+\dfrac{1}{a}$ cắt nhau tại đúng hai điểm phân biệt?
	\choice
	{\True $340$}
	{$342$}
	{$361$}
	{$324$}
	\loigiai{
		Đầu tiên ta nhận thấy $a>1,b>1$.\\
		Dễ dàng nhận ra nếu $a=b$ thì hai đồ thị lúc này trùng nhau nên có vô số điểm chung, suy ra loại.\\
		Do đó $a\neq b$, tiếp đến, vì vai trò của $a$ và $b$ như nhau nên ta chỉ cần tìm cặp số nguyên $(a;b)$ với $a>b>1$ sao cho phương trình $\dfrac{1}{a^x}+\dfrac{1}{b}=\dfrac{1}{b^x}+\dfrac{1}{a}\Leftrightarrow\dfrac{1}{a^x}+\dfrac{1}{b}-\dfrac{1}{b^x}-\dfrac{1}{a}=0$ có 2 nghiệm phân biệt.\\
		Xét hàm số $y=f(x)=\dfrac{1}{a^x}+\dfrac{1}{b}-\dfrac{1}{b^x}-\dfrac{1}{a}$ có $f'(x)=-\dfrac{\ln a}{a^x}+\dfrac{\ln b}{b^x}$ và $f(1)=0$.\\
		Giải phương trình $f'(x)=-\dfrac{\ln a}{a^x}+\dfrac{\ln b}{b^x}=0\Leftrightarrow x=x_0=\log_{\tfrac{b}{a}}\left(\dfrac{\ln b}{\ln a}\right),f'(x)>0$.\\
		Cùng với $f'(x)<0$ khi $x>x_0$ và $f'(x)>0$ khi $x<x_0$. Nên ta xét 2 trường hợp sau:\\
		Trường hợp 1: $x_0=\log_{\tfrac{b}{a}}\left(\dfrac{\ln b}{\ln a}\right)=1\Leftrightarrow\dfrac{\ln a}{a}=\dfrac{\ln b}{b}\Leftrightarrow (a;b)=(4;2)$.\\
		Xét hàm số $y=g(t)=\dfrac{\ln t}{t}$ ta có: $\dfrac{\ln 3}{3}>\dfrac{\ln 2}{2}=\dfrac{2\ln 2}{4}=\dfrac{\ln 4}{4}>\dfrac{\ln 5}{5}>\ldots\cdot >\dfrac{\ln 20}{20}$.\\
		Khi đó $f(x)\geq f(x_0)=f(1)=0$ tức $f(x)$ có đúng 1 nghiệm $x_0=1$.\\
		Trường hợp 2: $x_0=\log_{\tfrac{b}{a}}\left(\dfrac{\ln b}{\ln a}\right)\neq 1$. Khi đó vẽ bảng biến thiên ta thấy $f(x)=0$ luôn có hai nghiệm phân.\\
		biệt. Với mỗi $b=k\in\{2;3;\ldots;19\}$ thì $a\in\{k+1;\ldots$.; 20 $\}$ tức có $20-k$ cách chọn giá trị $a$.\\
		Suy ra có $\sum\limits_{k=2}^{19} (100-k)=171$ cặp $(a;b)$ với $a>b>1$ và loại đi cặp $(4;2)$ ta có 170 cặp.\\
		Xét tương tự với $b>a>1$ ta cũng có 170 cặp. Vậy có tất cà 340 cặp số thỏa mãn.
	}
\end{ex}
\begin{ex}
	[Sở Hải Dương 2022]%Câu 35.
	Cho hàm số $f(x)=\log_3\left(\sqrt{4x^2+1}+2x\right)+3x^{2021}$. Có tất cả bao nhiêu giá trị nguyên của $m$ thuộc đoạn $[-2021;2021]$ để bất phương trình $f\left(x^2+1\right)+f(-2mx)\geq 0$ nghiệm đúng với mọi $x\in(0;+\infty)$. 
	\choice
	{\True $2023$}
	{$4020$}
	{$4022$}
	{$2021$}
	\loigiai{
		Tập xác định: $\mathscr{D}=x\in\mathbb{R}$.\\
		Ta có $f'(x)=\dfrac{2\left(\sqrt{4x^2+1}+2x\right)}{\sqrt{4x^2+1}\left(\sqrt{4x^2+1}+2x\right)\ln 3}+6063x^{2020}>0\Rightarrow f(x)$ đồng biến trên $\mathbb{R}$.\\
		Ta thấy:\\
		$f(-x)=\log_3\left(\sqrt{4(-x)^2+1}+2(-x)\right)+3(-x)^{2021}=\log_3\left(\sqrt{4x^2+1}+2x\right)^{-1}-3x^{2021}=-f(x)$.\\
		Vậy $f(x)$ là hàm số lẻ. Khi đó:\\
		$f\left(x^2+1\right)\geq-f(-2mx)\Leftrightarrow f\left(x^2+1\right)\geq f(2mx)\Leftrightarrow x^2+1\geq 2mx\Leftrightarrow x+\dfrac{1}{x}\geq 2m,\forall x>0$.\\
		Xét $g(x)=x+\dfrac{1}{x},\forall x>0\Rightarrow g'(x)=1-\dfrac{1}{x^2}=0\Leftrightarrow\hoac{&x=-1(\mathrm{L})\\&x=1(\mathrm{N}).}$ \\
		Ta có bảng biến thiên của hàm số $y=g(x)$: \\
		{\color{red} HINH O DAY}\\
		Theo yêu cầu bài toán thì $2m\leq 2\Leftrightarrow m\leq 1$.\\
		Vì $m\in[-2021;2021]\Rightarrow$ số giá trị của $m$ bằng $\left(1-(-2021)\right)+1=2023$.
	}
\end{ex}
\begin{ex}
	[Cụm trường Nam Định 2022]%Câu 36.
	Cho đồ thị hàm số $f(x)=ax^3+bx^2+cx+d$ có hai điểm cực trị là $A(0;3)$ và $B(2;-1)$. Số nghiệm thực của phương trình $4^{f\left(f(x)\right)}-2^{f(x)+f\left(f(x)\right)}+3\cdot 2^{f\left(f(x)\right)}=3\cdot 2^{f(x)}$ là
	\choice
	{\True $7$}
	{$6$}
	{$3$}
	{$9$}
	\loigiai{
		Đồ thị hàm số $f(x)=ax^3+bx^2+cx+d$ có hai điểm cực trị là $A(0;3)$ và $B(2;-1)$ nên ta có hệ phương trình $\heva{&d=3\\&8a+4b+2c+d=-1\\&c=0\\&3\cdot 4\cdot a+2\cdot b\cdot 2+c=0}\Leftrightarrow\heva{&a=1\\&b=-3\\&c=0\\&d=3.}$ \\
		Suy ra $f(x)=x^3-3x^2+3$.\\
		Xét phương trình $4^{f\left(f(x)\right)}-2^{f(x)+f\left(f(x)\right)}+3\cdot 2^{f\left(f(x)\right)}=3\cdot 2^{f(x)}$.\\
		Đặt $t=f(x)$, phương trình đã cho trở thành.\\
		$4^{f(t)}-2^{t+f(t)}+3\cdot 2^{f(t)}=3\cdot 2^t$ \\
		$ \Leftrightarrow 2^{f(t)}\left(2^{f(t)}-2^t\right)+3\left(2^{f(t)}-2^t\right)=0 $ \\
		$\Leftrightarrow\left(2^{f(t)}-2^t\right)\left(2^{f(t)}+3\right)=0\Leftrightarrow\hoac{&2^{f(t)}-2^t=0\\&2^{f(t)}+3=0.} $ \\
		$2^{f(t)}>0\Rightarrow 2^{f(t)}+3>0$ nên $2^{f(t)}+3=0$ (vô nghiệm).\\
		$2^{f(t)}-2^t=0\Leftrightarrow 2^{f(t)}=2^t\Leftrightarrow f(t)=t$ \\
		$ \Rightarrow t^3-3t^2+3=t\Leftrightarrow t^2(t-3)-(t-3)=0\Leftrightarrow(t-3)\left(t^2-1\right)=0\Leftrightarrow\hoac{&t=3\\&t=1\\&t=-1} $ \\
		$\Rightarrow\hoac{&x^3-3x^2+3=3\\&x^3-3x^2+3=1\\&x^3-3x^2+3=-1} $ \\
		{\color{red} HINH O DAY}\\
		Dựa vào đồ thị hàm số, ta thấy.\\
		Đường thẳng $y=3$ cắt đồ thị hàm số $f(x)=x^3-3x^2+3$ tại 2 điểm phân biệt.\\
		Đường thẳng $y=1$ cắt đồ thị hàm số $f(x)=x^3-3x^2+3$ tại 3 điểm phân biệt.\\
		Đường thẳng $y=-1$ cắt đồ thị hàm số $f(x)=x^3-3x^2+3$ tại 2 điểm phân biệt.\\
		Và 7 điểm này không có 2 điểm nào trùng nhau.\\
		Vậy phương trình có 7 nghiệm phân biệt.
	}
\end{ex}
\begin{ex}
	[Sở Sơn La 2022]%Câu 37.
	Cho hàm số $f(x)=\mathrm{e}^{\sqrt{x^2+1}}\left(\mathrm{e}^x-\mathrm{e}^{-x}\right)$. Có bao nhiêu số nguyên dương m thỏa mãn bất phương trình $f(m-7)+f\left(\dfrac{12}{m+1}\right)<0$?
	\choice
	{Vô số}
	{$4$}
	{\True $3$}
	{$5$}
	\loigiai{
		Hàm số $f(x)=\mathrm{e}^{\sqrt{x^2+1}}\left(\mathrm{e}^x-\mathrm{e}^{-x}\right)$ xác định $\forall x\in\mathbb{R}$.\\
		Khi đó với $-x\in\mathbb{R}$, ta có $f(-x)=\mathrm{e}^{\sqrt{x^2+1}}\left(\mathrm{e}^{-x}-\mathrm{e}^x\right)=-f(x)$.\\
		Suy ra $f(x)$ là hàm số lẻ. (1).\\
		Mặt khác $f(x)=\mathrm{e}^{\sqrt{x^2+1+x}}-\mathrm{e}^{\sqrt{x^2+1}-x}\Rightarrow f'(x)=\left(\dfrac{x}{\sqrt{x^2+1}}+1\right)\mathrm{e}^{\sqrt{x^2+1+x}}-\left(\dfrac{x}{\sqrt{x^2+1}}-1\right)\mathrm{e}^{\sqrt{x^2+1}-x}$.\\
		$=\left(\dfrac{x+\sqrt{x^2+1}}{\sqrt{x^2+1}}\right)\mathrm{e}^{\sqrt{x^2+1}+x}+\left(\dfrac{\sqrt{x^2+1}-x}{\sqrt{x^2+1}}\right)\mathrm{e}^{\sqrt{x^2+1}-x}>0,\forall x\in\mathbb{R}$..\\
		Do đó hàm số $f(x)$ đồng biến trên $\mathbb{R}$. (2).\\
		Ta có $f(m-7)+f\left(\dfrac{12}{m+1}\right)<0\Leftrightarrow f(m-7) <-f\left(\dfrac{12}{m+1}\right)$.\\
		Theo $(1)$ suy ra $\Leftrightarrow f(m-7)<f\left(-\dfrac{12}{m+1}\right)$.\\
		Theo (2) ta được $m-7 <-\dfrac{12}{m+1}\Leftrightarrow\dfrac{m^2-6m+5}{m+1}<0\Leftrightarrow\hoac{&1<m<5\\&m <-1.}$ \\
		Vi $m\in\mathbb{Z}^+$ nên $m\in\{2;3;4\}$.
	}
\end{ex}
\begin{ex}
	[Sở Nam Định 2022]%Câu 38.
	Cho hàm số $f(x)$ là hàm số đa thức bậc năm. Biết hàm số $y=f'(x)$ có đồ thị như hình vẽ bên dưới. Có bao nhiêu giá trị nguyên của tham số $m$ để hàm số $g(x)=2021^{f\left(x^3-3x^2+m\right)}+2022$ có $8$ điểm cực trị?\\
	{\color{red} HINH O DAY}
	\choice
	{$4$}
	{$3$}
	{$2$}
	{\True $1$}
	\loigiai{
		Từ đồ thị hàm số bậc bốn $y=f'(x)$ suy ra $f'(x)=ax(x-1)^2(x-2)$ với $a>0$.\\
		Ta có $g'(x)=\left(3x^2-6x\right)\cdot f'\left(x^3-3x^2+m\right)\cdot 2021^{f\left(x^3-3x^2+m\right)}\cdot\ln 2021$.\\
		$g'(x)=0\Leftrightarrow\hoac{&x=0\\&x=2\\&x^3-3x^2+m=0\\&x^3-3x^2+m=1\\&x^3-3x^2+m=2}\Leftrightarrow\hoac{&x=0\\&x=2\\&x^3-3x^2=-m(1)\\&x^3-3x^2+m=1(2)\\&x^3-3x^2=2-m(3)}$ trong đó các nghiệm của phương trình $(2)$ đều là nghiệm bội chẵn.\\
		Xét hàm số $y=x^3-3x^2$, ta có: $y'=3x^2-6x$, $y'=0\Leftrightarrow\hoac{&x=0\\&x=2.}$ \\
		Ta có bảng biến thiên: \\
		{\color{red} HINH O DAY}\\
		Để hàm số $g(x)$ có $8$ điểm cực trị thì các phương trình $(1),(3)$ đều phải có $3$ nghiệm phân biệt khác $0$ và $2\Leftrightarrow\heva{&-4 <-m<0\\&-4<2-m<0}\Leftrightarrow\heva{&0<m<4\\&2<m<6}\Leftrightarrow 2<m<4\Rightarrow m=3$.
	}
\end{ex}       
\Closesolutionfile{ans}
\indapan{10}{ans/CD19/Muc_9_10}